% r03 The torus comparison
\section{The torus comparison}

The rotation enters the rotation part through a single statement, the torus comparison, posed on King's tori for the two lattices $\Lk$ and $R^{-1}\Lk$. This section states it, reduces it by an exact symmetry to one-sided bounds at one lattice spacing, and proves it along the lines of \cite{K2}.

\subsection{The torus comparison}

Throughout, for a lattice $\Lambda$ of spacing $\eps_k$ descending to a King torus $\DJ$ (\Cref{rot:def:frames,rot:def:tori}), $Z^a_{\Lambda,\DJ}(\zeta)\defeq e^{-\sum_jw_ja_k\sS_\Lambda(h_j)}\langle\Theta(\zeta)\rangle_{\Lambda,\DJ}$ is the functional of the periodic state in the scheme $a$ (\Cref{def:s02:sources}), $\Zun_{\Lambda,\DJ}(s)$ the unnormalised functional of \Cref{rot:def:Zun}, and $\langle\cdot\rangle_{\fc,\Lambda}$ the walled state of the configuration $\fc=(J',w)$.

\begin{theorem}[\Kz: the torus comparison, curvature sector]\label{rot:hyp:K}
For every $R\in\Rgen$, every walled configuration $\fc=(J',w)$ of the class of \Cref{rot:rem:class} (the conditions of \Cref{rot:thm:CT}) and every real smooth $D_{J'}$-periodic $2$-form $G$, the limits $\lim_k\langle e^{F(G)}\rangle_{\fc,\Lk}$ and $\lim_k\langle e^{F(G)}\rangle_{\fc,R^{-1}\Lk}$ exist and are equal.
\end{theorem}

The proof is given at the end of the next subsection, which reduces the statement to one-sided bounds at a single lattice spacing; the bounds themselves are proved in \cref{ssec:rot:Kproof}, and \Cref{rot:rem:behindK} lists the inputs of that proof from the literature. With scalar sources the comparison is a property of a scheme $a$, which is not proved here (it is listed among the open problems, in the final section) and is used only by the optional extension of the Wick leg (\Cref{thm:s23:FormA}):

\begin{definition}[\Ka: the torus comparison with scalar sources]\label{rot:hyp:Ka}
For every $R\in\Rgen$, every King torus $\DJ$, every finite family of real smooth $\DJ$-periodic $2$-forms $G_j$ and functions $h_j$, and every $(z,w)\in\R^{2N}$, the limits
\[
\lim_{k\to\infty}Z^a_{\Lk,\DJ}\Big(\sum_jz_jG_j,\sum_jw_jh_j\Big)\qquad\text{and}\qquad
\lim_{k\to\infty}Z^a_{R^{-1}\Lk,\DJ}\Big(\sum_jz_jG_j,\sum_jw_jh_j\Big)
\]
exist in $(0,\infty)$ and are equal.
\end{definition}

\Kz\ is stated on the sub-class of walled configurations described in \Cref{rot:rem:class}, which is the one on which the results that consume it use it. Here $R^{-1}\Lk$ carries the frame $(R^{-1}e_\alpha)$ and descends to $\DJ$ by \Cref{rot:lem:geom}(b); both lattices have spacing $\eps_k$ and carry the same $a_k$, $\dm_{\eps_k}$ and $\sW_{\eps_k}$. Since $\Rgen=\Rgen^{-1}$ (\Cref{rot:lem:geom}(a)), the statement for all $R\in\Rgen$ is the same as the statement for the lattices $R\Lk$ used in \cite{K2}. Read with a single scalar function on all plaquettes, as printed in \cite[(2.12), Thm.~2.4]{K2}, the comparison is false already at $e=0$ (\Cref{rot:prop:scalar}); the $2$-form reading, in which each lattice samples the source in its own frame, is consistent with King's proof, which treats the source components separately (\Cref{rot:rem:frame}).

\begin{lemma}\label{rot:lem:K}
\begin{enumerate}[label=\textup{(\alph*)},nosep]
\item Under \Ka\ the convergence holds locally uniformly on $\C^{2N}$, the two limits are the same entire function, and for every walled configuration $\fc$ and every $(f,g)$ supported in one fundamental cell with $w=0$ on $\supp g$, $\lim_k\Tw{a}{\fc,\Lk}(f,g)$ and $\lim_k\Tw{a}{\fc,R^{-1}\Lk}(f,g)$ exist and are equal. Under \Kz\ the same holds, for the configurations of its class, for $\langle e^{zF(G)}\rangle_{\fc,\Lk}$ and $\langle e^{zF(G)}\rangle_{\fc,R^{-1}\Lk}$ at complex $z$.
\item \Ka\ implies \Kz, and \Ka\ implies $(\mathrm K)_{a'}$ if $a'_k-a_k$ converges.
\item If the statement of \Ka, respectively \Kz, holds for $R=R_{12}(\thz)$, it holds for every $R\in\Rgen$.
\item Under \Ka, for every King torus $\DJ$ and every $0\le q\in C^\infty_c$ supported in one fundamental cell, $C_{\DJ}(q)\defeq\sup_kZ^a_{\Lk,\DJ}(0,q)<\infty$.
\end{enumerate}
\end{lemma}

\begin{proof}
(a) As in the proof of \Cref{rot:thm:full}(i): at fixed $k$ the functionals are entire; by \Cref{lem:s02:alg}(d), whose proof applies verbatim to periodic states, $|Z^a(\zeta)|\le Z^a(\Re\zeta)$ and $\ln Z^a$ is convex in the real coefficients, so the convergence on $\R^{2N}$ gives local boundedness of the sequence on $\C^{2N}$; Montel's theorem gives convergent subsequences, whose limits agree on $\R^{2N}$, a set of uniqueness for entire functions on $\C^{2N}$ \cite{Hor90}; hence both sequences converge locally uniformly, to the same entire function. The walled ratios are quotients of such limits: with the family $(f;g,w)$ and the coefficients $(\mathrm i;1,-1)$ in the numerator and $(0;0,-1)$ in the denominator, $\Tw{a}{\fc,\Lambda}(f,g)=Z^a_{\Lambda,D_{J'}}(\mathrm if,g-w)/Z^a_{\Lambda,D_{J'}}(0,-w)$ (\Cref{rot:def:tori}), and the denominators have positive limits. For \Kz\ at complex $z$ use $|\langle e^{zF(G)}\rangle|\le\langle e^{(\Re z)F(G)}\rangle$ and that \Kz\ for all $G$ includes all real multiples $tG$; the same argument applies.
(b) $\langle e^{F(G)}\rangle_{\fc,\Lambda}=Z^a_{\Lambda,D}(G,-w)/Z^a_{\Lambda,D}(0,-w)$, the wall being the scalar source $-w$ (\Cref{rot:def:tori}; the scheme multiplier cancels). For the second claim, the factors $e^{-\sum_jw_j(a'_k-a_k)\sS_\Lambda(h_j)}$ have the same limit on both lattices, since $\sS_{\Lk}(h)$ and $\sS_{R^{-1}\Lk}(h)$ are Riemann sums of the smooth periodic function $h$ over lattices descending to $\DJ$ and both tend to $\int_{\GamJ}h$.
(c) Let $\sigma_2=\diag(1,-1,1)$. It preserves $\Lk$ and $\DJ$, and changes the frame only by a sign, which does not change the periodic model (\Cref{lem:s02:pass}); moreover $\sigma_2R_{12}(\thz)^{-1}\sigma_2=R_{12}(\thz)$ (\Cref{rot:lem:geom}(a)). By \Cref{lem:s02:pass} with $O=\sigma_2$ (for walled configurations \Cref{rot:lem:walledalg}(b), the wall becoming $\sigma_2\cdot w$) the statement for $(R_{12}(\thz),\text{sources }s)$ is equivalent to the statement for $(R_{12}(-\thz),\sigma_2\cdot s)$; as $s$ ranges over all sources, so does $\sigma_2\cdot s$. For the other planes, $R_{\mu\nu}(\pm\thz)=\Pi R_{12}(\pm\thz)\Pi^{-1}$ for a permutation matrix $\Pi$ mapping $e_1,e_2$ to $e_\mu,e_\nu$; apply \Cref{lem:s02:pass} with $O=\Pi$, which preserves $\Lk$ and $\DJ$ and permutes frames.
(d) A convergent sequence with a finite limit is bounded.
\end{proof}

\subsection{Reduction to one-sided bounds at one lattice spacing}

King's proof of his Theorem~2.4 compares the lattice of spacing $\eps$ with the rotated lattice of spacing $5\eps$, and obtains the reverse inequality by ``switching around the argument'' \cite[p.~337]{K2}. We compare the two lattices at the \emph{same} spacing, and obtain the reverse inequality from an exact symmetry.

\begin{lemma}[the involution]\label{rot:lem:inv}
Let $R=R_{\mu\nu}(\pm\thz)\in\Rgen$, let $\sigma\defeq\sigma_\nu$ be the reflection $x_\nu\mapsto-x_\nu$, and put $S\defeq R^{-1}\sigma$. Then $S^2=1$, $S\Lk=R^{-1}\Lk$ and $S(R^{-1}\Lk)=\Lk$, and $S$ maps $\DJ=\R^3/P_J$, $P_J\defeq(50MJ\Z)^3$, onto the torus $R^{-1}\DJ\defeq\R^3/R^{-1}P_J$. Both lattices descend to both tori, and for all real smooth $\DJ$-periodic sources $s=(G,h)$,
\begin{gather*}
\Zun_{R^{-1}\Lk,\DJ}(s)=\Zun_{\Lk,R^{-1}\DJ}(S\cdot s),\qquad
\Zun_{\Lk,\DJ}(s)=\Zun_{R^{-1}\Lk,R^{-1}\DJ}(S\cdot s),\\
\sS_{R^{-1}\Lk,\DJ}(h)=\sS_{\Lk,R^{-1}\DJ}(S\cdot h).
\end{gather*}
\end{lemma}

\begin{proof}
By \Cref{rot:lem:geom}(a), $\sigma R\sigma=R^{-1}$, so $S^2=R^{-1}(\sigma R^{-1}\sigma)=R^{-1}R=1$; $\sigma\Lk=\Lk$ gives $S\Lk=R^{-1}\Lk$, and then $S(R^{-1}\Lk)=S^2\Lk=\Lk$. Also $\sigma P_J=P_J$, so $SP_J=R^{-1}P_J$. The lattice $R^{-1}\Lk$ is invariant under $R^{-1}P_J$ trivially, and $\Lk$ is invariant under $R^{-1}P_J$ because $R^{-1}p\in10MJ\,\Z^3$ for $p\in P_J$ (\Cref{rot:lem:geom}(b)). The map $S$ carries the frame $(R^{-1}e_\alpha)$ to $(SR^{-1}e_\alpha)=(\sigma e_\alpha)$ and the frame $(e_\alpha)$ to $(R^{-1}\sigma e_\alpha)$; these differ from the frames of $\Lk$ and $R^{-1}\Lk$ only by signs, which do not change the periodic models (\Cref{lem:s02:pass}). \Cref{lem:s02:pass} with $O=S$, applied to the unnormalised integrals (the bijection of configurations preserves $d\nu$), gives the identities.
\end{proof}

\begin{definition}[\hypF, \hypV\ and \hypCc: one-sided bounds at one spacing]\label{rot:hyp:F}
Let $R\in\Rgen$ and $J\in\N$. Write $T_1\defeq\DJ$, $T_2\defeq R^{-1}\DJ$. The following are properties of the pair $(R,J)$.
\hypF: for $i=1,2$ there are real numbers $v_k(T_i)$ and $c_k(T_i)$, $k\ge k_0(J)$, independent of the fields and the sources, with $|c_k(T_i)|\le C\eps_k^{-q}$ for some $q<\infty$ and $v_k(T_i)$ bounded, such that for every real smooth $T_i$-periodic source $s=(G,h)$,
\[
\Zun_{R^{-1}\Lk,T_i}(s)\le\big(1+\delta_k(T_i,s)\big)\exp\!\big(-v_k(T_i)-c_k(T_i)\,\sS_{\Lk,T_i}(h)\big)\,\Zun_{\Lk,T_i}(s),
\qquad\lim_{k\to\infty}\delta_k(T_i,s)=0 .
\]
\hypV: $v_k(T_1)+v_k(T_2)\to0$. \hypCc: $c_k(T_1)-c_k(T_2)\to0$.
\hypFz\ is \hypF\ restricted to sources with $h=0$ and to the sources $(G,\pm h_0)$, where $h_0=-w$ is a wall; it involves no Wick scheme.
\end{definition}

On each torus, \hypF\ bounds the functional of the rotated lattice from above by that of the axis lattice. The rotated lattice is the one whose flow is King's standard block-spin flow in its own frame, and the axis lattice is the one run through King's hybrid flow; this is the direction of \cite[(4.7)--(4.9)]{K2}. Since smooth periodic functions have rapidly convergent Riemann sums, $|\sS_{\Lambda,T}(h)-\int_\Gamma h|\le C_N(h)\eps_k^N$ for every $N$ and every lattice $\Lambda$ of spacing $\eps_k$ descending to $T$, $\Gamma$ a fundamental cell; with $N>q$ the error $c_k(\sS_{\Lambda,T}(h)-\int h)$ tends to $0$.

\begin{theorem}[closing at one spacing]\label{rot:thm:sandwich}
Let $R\in\Rgen$, $J\in\N$, and assume \hypF. Write $v=v_k(T_1)$, $v'=v_k(T_2)$, $c=c_k(T_1)$, $c'=c_k(T_2)$, and $\rho_k(s)\defeq\Zun_{R^{-1}\Lk,\DJ}(s)/\Zun_{\Lk,\DJ}(s)$.
\begin{enumerate}[label=\textup{(\alph*)},nosep]
\item For every real smooth $\DJ$-periodic source $s=(G,h)$,
$e^{v'+c'\int h}(1+o(1))\le\rho_k(s)\le e^{-v-c\int h}(1+o(1))$ as $k\to\infty$.
\item $v+v'\le o(1)$. If \hypV\ holds, then $c+c'\to0$.
\item If \hypV\ holds, then $\rho_k(s)=e^{-v-c\int h}(1+o(1))$ for every $s$. Consequently, for sources $s_0,s_1$ with $\int h_0=\int h_1$,
\[
\frac{\Zun_{R^{-1}\Lk,\DJ}(s_1)}{\Zun_{R^{-1}\Lk,\DJ}(s_0)}=\frac{\Zun_{\Lk,\DJ}(s_1)}{\Zun_{\Lk,\DJ}(s_0)}\,(1+o(1)).
\]
\item If \hypV\ and \hypCc\ hold, then $c,c'\to0$, and the identity in (c) holds for all $s_0,s_1$.
\end{enumerate}
No convergence of either family and no Wick scheme is assumed.
\end{theorem}

\begin{proof}
(a) The upper bound is \hypF\ on $T_1$, with $c\,\sS_{\Lk,\DJ}(h)=c\int h+o(1)$. Apply \hypF\ on $T_2$ to the $T_2$-periodic source $S\cdot s$ and use \Cref{rot:lem:inv}: $\Zun_{R^{-1}\Lk,T_2}(S\cdot s)=\Zun_{\Lk,\DJ}(s)$, $\Zun_{\Lk,T_2}(S\cdot s)=\Zun_{R^{-1}\Lk,\DJ}(s)$ and $\sS_{\Lk,T_2}(S\cdot h)=\sS_{R^{-1}\Lk,\DJ}(h)=\int h+O(\eps_k^N)$. This is the lower bound.
(b) The interval in (a) is non-empty, so $(v+v')+(c+c')\int h\le o(1)$ for every $s$. Apply this to $(G,h)$ and to $(G,-h)$ and add: $v+v'\le o(1)$. The same two inequalities give $|c+c'|\,|\int h|\le o(1)-(v+v')$; under \hypV\ the right side tends to $0$, and choosing $h$ with $\int h\ne0$ gives $c+c'\to0$.
(c) Under \hypV\ and (b), the ratio of the two ends of the interval in (a) tends to $1$. Dividing the statement for $s_1$ by the one for $s_0$, the prefactors cancel when $\int h_1=\int h_0$.
(d) Now $c+c'\to0$ and $c-c'\to0$.
\end{proof}

The proof uses \hypF\ only for the sources in question and their images under $h\mapsto-h$. In particular, under \hypFz\ and \hypV\ the conclusions (a)--(c) hold for the sources $(G,\pm h_0)$ (where $\int h_0\ne0$ unless $w=0$) and, with no condition on $c,c'$, for the sources with $h=0$.

\begin{corollary}[walled curvature ratios]\label{rot:cor:walled}
Assume \hypFz\ and \hypV\ for $R$ and $J'$. Then for every walled configuration $\fc=(J',w)$ and every real smooth periodic $2$-form $G$, $\langle e^{zF(G)}\rangle_{\fc,R^{-1}\Lk}-\langle e^{zF(G)}\rangle_{\fc,\Lk}\to0$ locally uniformly in $z\in\C$. If moreover $\lim_k\langle e^{tF(G)}\rangle_{\fc,\Lk}$ exists for real $t$, then \Kz\ holds for $\fc$ and $G$.
\end{corollary}

\begin{proof}
With $h_0=-w$, $\langle e^{tF(G)}\rangle_{\fc,\Lambda}=\Zun_\Lambda(tG,h_0)/\Zun_\Lambda(0,h_0)$ (the Wick constants cancel; \Cref{rot:lem:walledalg}(c)), a ratio of two sources with the same $\int h$. By \Cref{rot:thm:sandwich}(c), $u_k(t)\defeq\langle e^{tF(G)}\rangle_{\fc,R^{-1}\Lk}/\langle e^{tF(G)}\rangle_{\fc,\Lk}\to1$ for real $t$. The walled states are BFS measures (\Cref{rot:def:tori}), so by \Cref{lem:s02:GD} and the symmetry $(A,\vartheta)\mapsto(-A,-\vartheta)$, $1\le\langle e^{tF(G)}\rangle\le e^{t^2\sigma_\Lambda/2}$ for real $t$, where $\sigma_\Lambda=\ip v{Cv}$ is the variance of $F_\Lambda(G)=\ip vA$ under the free periodic photon on the lattice; since the periodic photon form is at least $\lmu^2$ times the $\ell^2$ norm, $\sigma_\Lambda\le\lmu^{-2}\norm{\delta_\eps G}^2$, $\delta_\eps$ the lattice codifferential, and this is bounded uniformly in $k$ (a Riemann sum for $\norm{\dstar G}^2_{L^2(\Gamma_{J'})}$). Hence the difference of the two functions tends to $0$ on $\R$, and it is bounded on every strip $|\Re z|\le r$ uniformly in $k$, because $|\langle e^{zF}\rangle|\le\langle e^{(\Re z)F}\rangle$. Vitali's theorem \cite{Hor90} gives locally uniform convergence to $0$ on $\C$. The last statement follows by the same argument.
\end{proof}

\begin{corollary}[the full comparison]\label{rot:cor:Kfull}
Let $a$ be a scheme with $|a_k|\le C(1+k)$. Assume \hypF, \hypV\ and \hypCc\ for $R=R_{12}(\thz)$ and every $J$, and that for every King torus and every finite family of real smooth periodic sources the limit $\lim_kZ^a_{\Lk,\DJ}$ exists in $(0,\infty)$ for real coefficients. Then \Ka\ holds.
\end{corollary}

\begin{proof}
$Z^a_{\Lambda,\DJ}(s)=e^{-a_k\sS_\Lambda(h)}\Zun_\Lambda(s)/\Zun_\Lambda(0)$, and the multipliers of the two lattices differ by $e^{-a_k(\sS_{R^{-1}\Lk}(h)-\sS_{\Lk}(h))}=1+O(k\eps_k^N)$. \Cref{rot:thm:sandwich}(d) with $s_0=0$ gives $Z^a_{R^{-1}\Lk,\DJ}(s)/Z^a_{\Lk,\DJ}(s)\to1$ for every real source, which with the assumed convergence gives \Ka\ for $R_{12}(\thz)$; \Cref{rot:lem:K}(c) gives the other generators.
\end{proof}

\begin{remark}[what this subsection supplies]\label{rot:rem:closing}
\Cref{rot:lem:inv}, \Cref{rot:thm:sandwich} and its corollaries are a reduction of the torus comparison to the one-sided bounds \hypF, \hypV, \hypCc; the bounds themselves come from the renormalisation-group analysis. In the curvature sector, what \Cref{rot:cor:walled} needs is \hypFz\ and \hypV\ together with the existence of the axis-lattice limits, and these are what the proof of \Kz\ along the lines of \cite{K2} produces (\Cref{rot:rem:behindK}). In that proof King's torus bound \cite[(4.2)]{K2} is of no use, because at a non-zero wall it is vacuous; its place is taken by the large-field clause of King's upper bound \cite[Thm.~3.1, (3.4)]{K1} at blocking factor $5$, and at a mass-increasing wall by a polynomial lower bound for the walled partition function. \Cref{rot:cor:Kfull} concerns scalar sources and is not used in the curvature sector (\Cref{rot:thm:main}).
\end{remark}

\begin{proof}[Proof of \Cref{rot:hyp:K}]
By \Cref{rot:lem:K}(c) it suffices to treat $R=R_{12}(\thz)$. Fix a walled configuration $\fc=(J',w)$ of the class of \Cref{rot:rem:class} and a real smooth periodic $2$-form $G$. By \Cref{rot:cor:walled} the statement follows from \hypFz\ and \hypV\ for $R$ and $J'$ together with the existence of $\lim_k\langle e^{tF(G)}\rangle_{\fc,\Lk}$ for real $t$. The existence of the axis-lattice limit is King's torus theorem \cite[Thm.~2.1]{K1} with the wall as scalar source, read on the normalised ratio, in which the vacuum and the $h$-linear constants cancel (\Cref{rot:thm:H6W}). The bounds \hypFz\ on the two tori $T_1,T_2$, in total form, with \hypV, come from the comparison, at equal spacing, of the effective actions of King's standard block-spin flow on the rotated lattice and of the hybrid flow on the axis lattice; they are proved in \cref{ssec:rot:Kproof}, and the proof is completed there (\Cref{rot:thm:CT} and the paragraph after \Cref{rot:prop:Konly}).
\end{proof}

\begin{remark}[the inputs of the proof of \Kz]\label{rot:rem:behindK}
\Kz\ is the corrected form of King's torus theorem \cite[Thm.~2.4]{K2}: $2$-form sources in each lattice's own frame, equal spacing, the free-boundary counterterm, walls allowed. The proof of \cref{ssec:rot:Kproof} follows King's lines. Passive rotations and \Cref{rot:lem:inv} reduce the comparison to the two lattices on $T_1,T_2$. The effective actions of two block-spin flows, the standard flow of $R\Lk$ in its own frame and a hybrid flow of $\Lk$ whose first step averages over rotated blocks, are compared at equal spacing on King's small-field domain, in total form (field-independent constants included). Inserting the comparison into King's upper bound \cite[Thm.~3.1]{K1} for the rotated lattice and into the hybrid analogue of \cite[Prop.~4.1]{K2} for the axis lattice gives a bound of type \hypF\ on each torus, the large-field term absorbed by a polynomial lower bound on the partition function; \Cref{rot:thm:sandwich} closes, with \hypV\ automatic because the comparison is in total form. The inputs: Ba{\l}aban's renormalisation bound \cite[Prop.~1]{Ba3}, with one element weakened by a small degree, for the renormalised families of \cite[\S3.5]{K1} with lines from both flows; decay and regularity of the block-spin operators, \cite[Thm.~3.3]{K1}, \cite{Ba4}, and of the hybrid kernels; King's upper bound, also on the torus skew in the rotated frame; the hybrid lower bound; the polynomial lower bound \cite[(1.14)]{Ba1} at zero wall and at a mass-increasing wall; the existence of the walled values on the axis lattice, \cite[Thm.~2.1]{K1} with the wall as scalar source, read on the normalised ratio; Reisz's theorems \cite{Rei88,Rei88b} in a corrected three-dimensional form; the local limit law \cite{LL10}; the bridge between King's torus counterterm \cite[(3.33)]{K1} and the free-boundary one \cite[p.~330]{K2}. The estimates of Ba{\l}aban's papers that the proof shares with the proof of \EC\ (\cref{sec:k01}), stated in those papers without a printed argument, are among them \cite[(2.33), (2.39)--(2.42), (3.20)--(3.21), (3.25), (3.34)--(3.36)]{Ba1} and the frame proof of its (1.14), \cite[Prop.~2.6, Lemma~2.7, (3.5), (3.7)--(3.8)]{Ba2}, \cite[Prop.~1 with (1.33)]{Ba3}, \cite[Lemma~2.4, (1.8), (1.15), (1.16), (1.20)]{Ba4}, \cite[Thm.~3.1, Props.~3.6--3.9, Lemma~4.5]{K1} with the treatment of the sources there (pp.~665, 677), and the sentence of \cite{K2} that Ba{\l}aban's results apply to both models; the mass clause of \cite[Prop.~3.1]{Ba2} is cited under it as in \Cref{rem:s14:BALm}. Statements false as printed are replaced: \cite[Thm.~3.5(i), (4.21), (4.36)--(4.37)]{K1} ((4.33) itself holds, \Cref{rot:prop:433}; (3.78) fails on certain vacuum graphs whose contribution is a removed constant and is not used), the count (3.45) and the first inequality of (3.46) of \cite{Ba2}, the gap value in \cite[(2.27)]{Ba4}, and the naturalness clause of Reisz's theorem.
\end{remark}

\begin{remark}[the content of \Kz]\label{rot:rem:Kstatus}
\Kz\ is \cite[Thm.~2.4]{K2} with the corrections of \Cref{rot:prop:scalar} and \Cref{rot:rem:log}, for arbitrary smooth periodic sources, walls included, and with free-boundary counterterms. By \Cref{rot:cor:walled,rot:cor:Kfull} its two-lattice content is \hypF\ and \hypV, plus \hypCc\ for scalar sources, and its one-lattice content is the existence of the torus limits on the axis lattice. The rotation enters the whole argument of the rotation part only through \Cref{rot:hyp:K}.
\end{remark}

\begin{remark}[what the torus comparison alone does not give]\label{rot:rem:shape}
Given \Kz, the walled values of a configuration $\fc$ on $\DJ$ agree on the two lattices, and by passive covariance (\Cref{rot:lem:walledalg}(b)) the walled value of $R\cdot\fc$ on the rotated torus $R\,\DJ$ for the axis lattice equals the walled value of $\fc$ on $\DJ$ for the rotated lattice. Rotation invariance of a limit along tori would therefore follow from \Kz\ if the limit did not depend on the shape and orientation of the tori, which is the content of King's expected identity \cite[(2.24)]{K2} between torus and cube limits. That identity is a uniqueness statement for the infinite-volume limit, and we do not use it. The walls of the next section replace it by a one-sided statement about one lattice, \TIz\ (\Cref{rot:hyp:TIGHT}), which is proved from the heavy-shell estimates (\Cref{rot:thm:shell}, \Cref{rot:thm:REStoTIGHT}).
\end{remark}

% r03_fill_K1: proof of the torus comparison (K_0), part I: the statement at the strength proved, the two block-spin flows.
% Input from sec/r03_rot_torus.tex after the proof of \Cref{rot:hyp:K}. Parts II-VI follow as sec/r03_fill_K2a.tex ... K6.tex.
\providecommand{\Xk}{\ensuremath{X_k}}
\providecommand{\Yk}{\ensuremath{Y_k}}
\providecommand{\hs}{\ensuremath{\hat s}}
\providecommand{\vr}{\ensuremath{\varrho}}
\providecommand{\Bdel}{\ensuremath{(\mathrm B_\Delta)}}
\providecommand{\ipp}[1]{\ensuremath{\langle #1\rangle}}

\subsection{Proof of the torus comparison}\label{ssec:rot:Kproof}

This subsection proves the one-sided bounds \hypFz\ and \hypV\ of \Cref{rot:hyp:F}, in the total form in which the vacuum constants are absent and the source-linear constants tend to zero, for the walled configurations of the class of \Cref{rot:rem:class}; together with the existence of the axis-lattice limits (\Cref{rot:thm:H6W}) and \Cref{rot:cor:walled} this completes the proof of \Cref{rot:hyp:K}. Throughout, $R\in\Rgen$ is fixed, $5R$ and $5R^{-1}$ are integer matrices (\Cref{rot:lem:geom}(a)), and we write
\[
\Xk\defeq\Lk=\eps_k\Z^3\quad\text{(the axis lattice)},\qquad\Yk\defeq R^{-1}\Lk\quad\text{(the rotated lattice)},\qquad
Y_{k-j}\defeq R^{-1}(5^j\eps_k\Z^3),
\]
so that $Y_{k-j}\subset\Xk\cap\Yk$ for $j\ge1$; the frame of $\Xk$ is $(e_\alpha)$ and that of $\Yk$ is $(R^{-1}e_\alpha)$. The tori are $T_1=\DJ=\R^3/P_J$, $P_J=50MJ\,\Z^3$, and $T_2=R^{-1}\DJ=\R^3/R^{-1}P_J$ (\Cref{rot:hyp:F}); the involution is $S=R^{-1}\sigma_\nu$ of \Cref{rot:lem:inv}. In this subsection $r$ denotes the number of block-spin steps, never the modulus $r_x=|\phi(x)|$, which does not occur.

\begin{theorem}[the torus comparison at one configuration]\label{rot:thm:CT}
Let the couplings be in King's range and $\fc=(J,\psi,n)$ a walled configuration with $J\in5^{n_*}\N$, $J\ge J_0(\delta_r)$, $n>0$, the wall condition $n\sS_{\Xk}(\psi)>\Delta_*|T_1|/2$ ($\Delta_*$ the constant of \Cref{rot:prop:CT}(b)) and block alignment (the source box and the torus unions of blocks of the final scale); here $n_*$ is the number of large-field scales of the flow of $\Yk$, fixed by the couplings, and $J_0(\delta_r)$ the non-wrapping threshold of \Cref{rot:lem:TC} at the coarsest spacing $\delta_r$ of the flow (a condition void for large $k$, since $\delta_r\to0$). Then for every real smooth $P_J$-periodic $2$-form $G$:
\begin{enumerate}[label=\textup{(\roman*)},leftmargin=*]
\item $\langle e^{zF(G)}\rangle_{\fc,\Yk}-\langle e^{zF(G)}\rangle_{\fc,\Xk}\to0$ as $k\to\infty$, locally uniformly in $z\in\C$; in particular all moments $\langle F(G)^m\rangle$ differ by $o(1)$;
\item both families converge, locally uniformly on $\C$, to the same entire function $\Wv_\fc(zG)$, and all moments converge to the same limits;
\item the same holds with $R$ replaced by any $R'\in\Rgen$.
\end{enumerate}
At $e=0$ the statement holds for every walled configuration with no condition: $\langle e^{zF(G)}\rangle_{\fc,\Lambda}=e^{z^2q_\Lambda(G)/2}$ exactly (\Cref{rot:lem:CM}), and $q_{\Xk}(G),q_{\Yk}(G)\to q(G)$ (\Cref{rot:lem:q}).
\end{theorem}

The three conditions on $\fc$ are those of \Cref{rot:rem:class}, and the configuration produced by \Cref{rot:thm:REStoTIGHT} satisfies them (its wall height may be raised at will, \Cref{rot:rem:Eproof}(iii), and $J$, the source box and the torus are chosen in \Cref{rot:thm:shell}). The proof is assembled in \Cref{rot:prop:Kforward} and the paragraph after it, from the comparison of the effective actions of two block-spin flows at equal spacing (\Cref{rot:thm:Asrc}), King's upper bound for the rotated lattice (\Cref{rot:prop:H3}), the hybrid lower bound (\Cref{rot:prop:H4}) and a polynomial lower bound at the wall (\Cref{rot:thm:FS}), closed by \Cref{rot:thm:sandwich}.

\begin{lemma}[admissibility]\label{rot:lem:adm}
\textup{(i)} $P_J\subset M\Z^3\cap R^{-1}(M\Z^3)$. \textup{(ii)} $SP_J=R^{-1}P_J\subset M\Z^3\cap R^{-1}(M\Z^3)$. \textup{(iii)} For $0\le j\le k$, $P_J\subset5^j\eps_kM\Z^3\cap R^{-1}(5^j\eps_kM\Z^3)$, and the same for $R^{-1}P_J$: every block lattice of both flows below descends to both tori.
\end{lemma}
\begin{proof}
(i) $50J\,R^{\mathrm T}=10J\cdot5R^{\mathrm T}$ is integral, so $RP_J=R^{-\mathrm T}P_J\subset10MJ\,\Z^3$, and $R^{-1}(M\Z^3)\supset P_J$. (ii) $\sigma_\nu\Z^3=\Z^3$ gives $SP_J=R^{-1}P_J$, and $R^{-1}P_J\subset10MJ\,\Z^3$ as in (i). (iii) $5^{k-j}\in\Z$.
\end{proof}

\subsubsection*{The two block-spin flows}

\begin{definition}[the standard flow of $\Yk$ and the hybrid flow of $\Xk$]\label{rot:def:flows}
Fix $k$ and $r\le k$; $\delta_r\defeq5^r\eps_k$ is the physical spacing of the coarse lattice $Y_{k-r}$. In the unit scale of level $r$ write $\eta\defeq5^{1-r}$ and $\eta'\defeq\eta/5=5^{-r}$. Blocks are centred at every level (\Cref{rot:conv:C}).
\begin{itemize}[leftmargin=*]
\item \emph{Standard flow on $\Yk$} (flow $Y$): Ba{\l}aban's and King's block-spin transformation \cite[(2.1)--(2.11)]{Ba1}, \cite[(2.10)--(2.17)]{K1} in the frame of $\Yk$: centred cubes of $125$ points, $Y_{k-j+1}\to Y_{k-j}$, with contours made of fine $\Yk$-bonds in the axial order $f_1,f_2,f_3$ of the frame inside each block.
\item \emph{Hybrid flow on $\Xk$} (flow $X$): the first step averages over the centred tiles $y+\eps B_0$, $y\in Y_{k-1}$, with
\[
B_0\defeq\Z^3\cap R^{-1}\big(-\tfrac52,\tfrac52\big)^3
\]
($125$ sites, a complete residue system of $\Z^3/5R^{-1}\Z^3$, \Cref{rot:lem:tile}), onto $Y_{k-1}$; this is \cite[(4.3)--(4.4)]{K2}, with the photon average read frame-covariantly, $(\hat QA)_\alpha(y)=\frac1{125}\sum_{x\in B(y)}f_\alpha\cdot A(x)$, $f_\alpha=R^{-1}e_\alpha$ (\Cref{rot:rem:frameav}). Steps $2,\dots,r$ are the standard flow's steps $Y_{k-1}\to Y_{k-r}$ with the same blocks, except that every rotated bond of $Y_{k-1}$ occurring in a contour is replaced by a fixed path of fine $\Xk$-bonds (\cite[p.~336]{K2}: ``succeeding steps of the renormalization transformation in the usual way, that is without rotating the blocks'').
\item \emph{Contours, rotation-paired.} Let $\Gamma^{\rm st}_{z,w}$ be King's coordinate staircase from $z$ to $w$ \cite[(2.12)]{K1} and $\Gamma^{(r)}_{y,x}$ his composite contour \cite[(2.11)]{K1}, which the standard flow uses in $Y$'s frame. The legs of a $Y$-frame staircase between two centres of $Y_{k-1}$-blocks are concatenations of the $\eta$-bonds $(u,u+5\eps f_\alpha)$, $u\in Y_{k-1}$. The hybrid composite contour from the level-$r$ block centre $y$ to $x'\in\Xk$ is the same staircase with every $\eta$-bond $(u,u+5\eps f_\alpha)$ replaced by the \emph{fixed} fine axis staircase $u+\eps\gamma_\alpha$ (and the reversed bond by its reverse), where, for $R=R_{12}(\thz)$ and $f_\alpha=R^{-1}e_\alpha$ with $5f_1=(4,3,0)$, $5f_2=(-3,4,0)$,
\[
\gamma_1:\ e_1,e_2,e_1,e_2,e_1,e_2,e_1\ \ (\text{to }(4,3,0)),\qquad
\gamma_2:\ e_2,-e_1,e_2,-e_1,e_2,-e_1,e_2\ \ (\text{to }(-3,4,0)),\qquad
\gamma_3:\ e_3^{\,5},
\]
followed, inside the first-step tile, by a fine axis path of at most $5$ bonds (\Cref{rot:lem:tile}(ii)), chosen once as a function of the displacement only and the same in every tile, so that the block and contour system of the flow is invariant under $Y_{k-1}$, in particular under the periods of both tori; this is the first-step contour of \cite[(4.4)]{K2}. (For the other generators conjugate by the permutation and sign matrices of \Cref{rot:lem:K}(c).) The gauge field along a composite contour is read by King's convention \cite[(2.19)]{K1} (\Cref{rot:conv:C}(ii)).
\end{itemize}
Both flows use King's coefficients $a_j=a(1-5^{-2})/(1-5^{-2j})$ \cite[p.~653]{K1}, take $r$ steps and land on the same coarse lattice $Y_{k-r}/P$ with the same unit-scaled fields $(A_r,\phi_r)$ and the same small-field indicator $\chi_r$ \cite[(3.2)]{K1}. The effective actions are written so that
\[
\Zun_{\Xk,T}(s)=\int(dA_r)(d\phi_r)\,e^{-\hs^{(r)}_T(A_r,\phi_r;s)},\qquad
\Zun_{\Yk,T}(s)=\int(dA_r)(d\phi_r)\,e^{-S^{(r)}_T(A_r,\phi_r;s)}
\]
\emph{exactly}, every $c$-number kept inside: this is \cite[(3.14), (3.90)]{K1} with $E_0$ inserted from \cite[(3.89)]{K1}.
\end{definition}

\begin{lemma}[the paired contours]\label{rot:lem:pair}
Every point of $u+\eps\gamma_\alpha$ lies within $0.6\eps$ of the segment $[u,u+5\eps f_\alpha]$, and the two bound a planar region of area at most $5\eps^2$. Each level-$i$ piece of the hybrid composite contour lies inside its level-$(i+1)$ block and the in-tile piece inside its tile, so the contour lies inside the level-$r$ block; in the unit scale of level $r$ its length is at most $2.1$ (the standard flow's at most $1.5$); the Hausdorff distance between the level-$i$ pieces of the two flows is at most $\eta'$ for $i\ge1$ and at most $2\sqrt{13}\,\eta'$ in the first-step cell; and the closed polygon formed by the two composite contours and the matching segment $[x',\tau x']$ (\Cref{rot:def:match}) bounds a piecewise planar surface of area at most $1.5\,\eta'+33\,\eta'^2$.
\end{lemma}
\begin{proof}
A finite check. The staircase $\gamma_1$ alternates $e_1$ and $e_2$ steps from $0$ to $(4,3,0)$ along the segment of slope $3/4$; its vertices $(1,0),(1,1),(2,1),(2,2),(3,2),(3,3),(4,3)$ have distances $\le0.6$ from the segment (the maximum, $3/5$, at $(1,0)$ and $(3,3)$), and the region between them is a union of triangles of total area $\le5$; $\gamma_2$ is $\gamma_1$ rotated by a right angle, and $\gamma_3$ is the segment itself. The centres of the blocks are at distance $\ge\frac52\eps$ from the block faces, against the deviation $0.6\eps$, which gives the containments; the lengths are $7$ fine bonds per $\eta$-bond for $\alpha=1,2$ and $5$ for $\alpha=3$, against $5$ for the rotated bond, plus the in-tile piece of at most $5$ bonds, which gives $2.1$ against $1.5$ per unit length of staircase in the unit scale; the Hausdorff distances and the area are the sums over the pieces of the two levels, the first-step cell contributing $\max_{B_0}|x|\le\sqrt{13}$ (\Cref{rot:lem:tile}(ii)).
\end{proof}

\begin{convention}[centred blocks; the contour convention]\label{rot:conv:C}
\textup{(i)} King's blocks are corner-rooted; both flows here use blocks centred at the sites of the next coarse lattice. Since $5$ is odd, centred cubes and centred tiles are inversion-symmetric. King's blocks are translates of these, and the printed estimates use only that blocks are translates of one cube and that contours stay inside blocks with bounded length (tiling, nesting, site in block with bounded diameter, all blocks translates of one cube); the four properties hold for centred blocks at block factor $5$ by exact enumeration, the fluctuation-integral estimates of \cite{Ba1,Ba2,Ba3,Ba4} being insensitive to the root of the block once the central symmetry of the fixed staircases is given. With centred blocks the inversion about any site of the coarsest lattice is an exact symmetry of both flows at every scale, which removes every first-moment term from both effective actions. \textup{(ii)} The gauge field along a composite contour is read in King's convention \cite[(2.19)]{K1}: if $A=\sum_{j<r}A'^{(j),\eps}+A^{(r),\eps}$ is the decomposition of the fine field into the lifts of the integrated fluctuations \cite[(2.18)]{K1}, then $A(\Gamma^{(r)}_{y,x})=\sum_{j=0}^{r-1}A'^{(j),\eps}(\Gamma^{(j+1)}_{x_{j+1},x})+A^{(r),\eps}(\Gamma^{(r)}_{y,x})$: the stage-$j$ fluctuation is integrated only along the part of the contour inside the $(j+1)$-block. This is what iterating the one-step transformations produces, and it is a definition in all three prints, \cite[(2.19)]{K1}, \cite[(3.34)]{Ba1}, \cite[(1.3)]{Ba3}. Under the other reading (one fine field along the whole composite contour) a staircase and the straight rotated bond with the same displacement carry different order-$e^2$ self-energies, by $0.0265\,e^2$ per unit length, an extensive amount; under \cite[(2.19)]{K1} the mismatch is $O(e^2\eps)$ per contour and its total vanishes.
\end{convention}

\begin{remark}[the photon average is frame-covariant]\label{rot:rem:frameav}
\cite[(4.3)]{K2} prints $(\hat QA_\mu)(y)=5^{-d}\sum_{x\in B(y)}A_\mu(x)$, which puts the axis component $A_\mu$ on the rotated bond direction. For a constant field $B$ it produces the coarse field $R^{-1}B$ while the fine contours see $B$, so King's constant-background step \cite[(3.43)--(3.46)]{K1} is inconsistent with it. With the frame-covariant average of \Cref{rot:def:flows}, $\hat Q^*\hat Q=\hat P\otimes\mathbf 1$ because $\sum_\alpha f_\alpha f_\alpha^{\mathrm T}=\mathbf 1$, the hybrid photon covariances are the scalar ones times the identity, and a constant background passes through every step unchanged. This repair is used throughout.
\end{remark}

\begin{lemma}[the tile average is a block-spin average]\label{rot:lem:tile}
Let $B(y)\defeq y+\eps B_0$ for $y\in Y_{k-1}$.
\begin{enumerate}[label=\textup{(\roman*)},leftmargin=*]
\item $B_0$ is a complete residue system of $\Z^3/5R^{-1}\Z^3$ with $|B_0|=125$; the tiles $B(y)$, $y\in Y_{k-1}$, partition $\Xk$, and for every period lattice $P$ of \Cref{rot:lem:adm} the tiles $B(y)$, $y\in Y_{k-1}/P$, partition $\Xk/P$.
\item $B_0=-B_0$, $\max_{x\in B_0}|x|^2=13$, $|x|_\infty\le3$ on $B_0$, and every $x\in B_0$ is joined to $0$ by a nearest-neighbour path inside $B_0$, monotone in each coordinate, of length $|x|_1\le5$.
\item With $\hat Q$ and $Q(A)$ of \Cref{rot:def:flows} and the pairings $\ip uv_\Lambda=\sum_xs^3u(x)\cdot v(x)$ of the two lattices, $\hat Q\hat Q^*=\mathbf 1$ on coarse bond fields and $Q(A)Q(A)^*=\mathbf 1$ on coarse scalar fields, for every bond field $A$. Hence $\hat Q^*\hat Q$ and $Q(A)^*Q(A)$ are orthogonal projections, the Gaussian block-spin kernels $\exp[-\frac a2\|A_1-\hat QA\|^2]$, $\exp[-\frac a2\|\phi_1-Q(A)\phi\|^2]$ have the normalisations of \cite[(2.10)--(2.14)]{K1}, and the first hybrid step is King's renormalisation transformation with $B(y)$ in place of the cube.
\item The composite of the first step with the $j-1$ standard steps that follow it is the equal-weight average over the composite blocks $B^{(j)}(y)$, $y\in Y_{k-j}$ (unions of $125^{j-1}$ tiles), which partition $\Xk$ and have diameter $\le C5^j\eps$; the composite Gaussian kernel has King's coefficient $a_j$, as for the standard flow.
\end{enumerate}
\end{lemma}
\begin{proof}
(i) The open cube $R^{-1}(-\frac52,\frac52)^3$ is a fundamental domain of the lattice $5R^{-1}\Z^3$, and no point of $\Z^3$ lies on its boundary: for $R=R_{12}(\thz)$ the boundary planes are $4x_1+3x_2=\pm\frac{25}2$, $-3x_1+4x_2=\pm\frac{25}2$, $x_3=\pm\frac52$, none of which has an integer solution. A fundamental domain whose boundary carries no lattice point contains exactly one point of every coset, so $B_0$ is a complete residue system and $|B_0|=|\det5R^{-1}|=125$. Hence every $x\in\Z^3$ is uniquely $x=5R^{-1}m+b$ with $m\in\Z^3$, $b\in B_0$, i.e.\ $\eps x=y+\eps b$ with $y=5\eps R^{-1}m\in Y_{k-1}$ unique: the tiles partition $\Xk$. For admissible $P$, $P\subset Y_{k-1}$ (\Cref{rot:lem:adm}(iii) with $j=1$), so $P$-translates of tiles are tiles and the partition descends. (ii) The cube is inversion-symmetric, so $B_0=-B_0$. The rest is a finite check over the $125$ points: $|x|^2\le13$ with equality at $(\pm3,0,\pm2)$, $(0,\pm3,\pm2)$; $|x|_\infty\le3$; for each $x$ a coordinate-monotone path from $0$ stays in $B_0$, the longest having $|x|_1=5$ steps. (iii) For a coarse bond field $\psi$ on $Y_{k-1}$ the adjoint of $\hat Q$ with respect to the pairings is $(\hat Q^*\psi)(x)=\sum_\alpha\psi_\alpha(y(x))f_\alpha$, $y(x)$ the centre of the tile containing $x$, because $(5\eps)^3/125=\eps^3$; then $(\hat Q\hat Q^*\psi)_\alpha(y)=\frac1{125}\sum_{x\in B(y)}f_\alpha\cdot\sum_\beta\psi_\beta(y)f_\beta=\psi_\alpha(y)$. For the scalar field, $(Q(A)\phi)(y)=\frac1{125}\sum_{x\in B(y)}U(A(\Gamma_{yx}))\phi(x)$ with $|U|=1$ has adjoint $(Q(A)^*\chi)(x)=\overline{U(A(\Gamma_{y(x)x}))}\chi(y(x))$, and $Q(A)Q(A)^*=\mathbf 1$. An operator with $QQ^*=\mathbf 1$ has $Q^*Q$ an orthogonal projection, and the Gaussian integrals $\int dA_1\exp[-\frac a2\|A_1-\hat QA\|^2]$, $\int d\phi_1\exp[-\frac a2\|\phi_1-Q(A)\phi\|^2]$ are constants independent of $A$ and $\phi$, which is all that \cite[(2.10)--(2.14)]{K1} uses of the block. (iv) The composition of equal-weight averages over successive partitions is the equal-weight average over the composite blocks; the diameter bound is (ii) at step $1$ and the cube's diameter at the later steps; the coefficient $a_j$ is King's computation \cite[p.~653]{K1}, which uses only the form of the kernels.
\end{proof}

\begin{lemma}[flux through the sliver]\label{rot:lem:flux}
For two lattice curves with the same endpoints, within Hausdorff distance $w$ and of length $\le\ell$, $|\int_\gamma A-\int_{\gamma'}A|\le\|dA\|_\infty\,\mathrm{area}(\Sigma)$ for a surface $\Sigma$ they bound together with the matching segment, and in particular $\le C\|dA\|_\infty w\ell$. In the unit scale of level $r$ the paired systems of \Cref{rot:def:flows} bound a surface of area $\le1.5\eta'+33\eta'^2$ (\Cref{rot:lem:pair}), so the phases $U(A(\Gamma))$ of the two flows at matched endpoints differ by at most $C\eta'\|dA\|_\infty$, a relative $\eta'$ against the Wilson-line phase.
\end{lemma}
\begin{proof}
Stokes' theorem on the lattice (the sum of the plaquette fluxes over $\Sigma$). The pairing is a condition, not a convenience: if the enclosed loops did not shrink with the scale, the flux $\sim e_rpA\ell^2$ would be of the same order as the Wilson-line phase $e_rA\ell$ at the coarse momenta $p\ell=O(1)$ at which the block-spin comparison is made, and \Cref{rot:thm:A} would fail; \cite[(4.4)]{K2} asks only that the contour ``lie entirely inside $B(y)$'' and leaves its shape open.
\end{proof}

% r03_fill_K2a: proof of (K_0), part II: the Cameron-Martin transfer, the choice of the number of steps, effective actions, matching.

\subsubsection*{The Cameron--Martin transfer of the curvature source}

\begin{lemma}[Cameron--Martin transfer]\label{rot:lem:CM}
Let $\Lambda$ be a lattice of spacing $s=\eps_k$ with frame $(f_\alpha)$ descending to a torus $T$, $\gamma_\Lambda$ the normalised Gaussian photon measure of the action, with covariance $C_\Lambda=(-\Delta_\Lambda+\lmu^2)^{-1}\otimes\mathbf 1_3$ for the pairing $\ip uv_\Lambda=\sum_xs^3\sum_\alpha u_\alpha(x)v_\alpha(x)$, and write $\Zun_{\Lambda,T}(G,h)=\int d\gamma_\Lambda(A)\,e^{F_\Lambda(G)}W_\Lambda(A;h)$ with $W_\Lambda$ the matter integral, which depends on $A$ only through the bond phases. Then $F_\Lambda(G)=\ip A{\ell_\Lambda}_\Lambda$ with $\ell_{\Lambda,\beta}\defeq\sum_{\alpha\ne\beta}\partial^*_\alpha G_{\alpha\beta}$, $G_{\alpha\beta}(x)\defeq G_{c(x;\alpha\beta)}(f_\alpha,f_\beta)$, $\partial^*_\alpha u(x)\defeq s^{-1}(u(x-sf_\alpha)-u(x))$, and for real $z$
\[
\Zun_{\Lambda,T}(zG,h)=e^{z^2q_\Lambda(G)/2}\,\Zun_\Lambda[za_\Lambda](T;h),\qquad q_\Lambda(G)\defeq\ip{\ell_\Lambda}{C_\Lambda\ell_\Lambda}_\Lambda,\quad a_\Lambda\defeq C_\Lambda\ell_\Lambda,
\]
where $\Zun_\Lambda[b](T;h)\defeq\int d\gamma_\Lambda(A)\,W_\Lambda(A+b;h)$. Counterterms, Wick constant and wall are untouched. At $e=0$, $W_\Lambda$ does not depend on $A$, so $\langle e^{zF(G)}\rangle_{\fc,\Lambda}=e^{z^2q_\Lambda(G)/2}$ exactly.
\end{lemma}
\begin{proof}
$\sum_{\alpha<\beta}G_{\alpha\beta}(\partial_\alpha A_\beta-\partial_\beta A_\alpha)=\sum_{\alpha\ne\beta}G_{\alpha\beta}\partial_\alpha A_\beta$ by antisymmetry, and summation by parts on the torus gives $\ip A{\ell_\Lambda}$. Completing the square, $e^{z\ip A\ell}d\gamma(A)=e^{z^2\ip\ell{C\ell}/2}d\gamma(A-zC\ell)$; substitute $A\mapsto A+zC\ell$. Both sides are finite ($\lmu>0$, $\lambda>0$, finite lattice).
\end{proof}

\begin{lemma}[the Gaussian factor is frame-free]\label{rot:lem:q}
With $q(G)\defeq\int_T\ip{\dstar G}{(-\Delta+\lmu^2)^{-1}\dstar G}\,dx$: $|q_\Lambda(G)-q(G)|\le C(\lmu)s^2\|G\|_{H^3(T)}^2+C_N(G)s^N$ for every $N$ and every lattice $\Lambda=O(s\Z^3)$ descending to $T$.
\end{lemma}
\begin{proof}
In Fourier variables on $\Lambda/P$, with $k_\alpha\defeq k\cdot f_\alpha$, $\sigma_\alpha(k)\defeq\frac2s\sin\frac{sk_\alpha}2$ and $u_{\alpha\beta}(y)\defeq G_y(f_\alpha,f_\beta)$, Poisson summation on the coset of the plaquette centres gives $\hat\ell_{\Lambda,\beta}(k)=-\mathrm ie^{\mathrm isk_\beta/2}\sum_\alpha\sigma_\alpha(k)[\hat u_{\alpha\beta}(k)+r_{\alpha\beta}(k)]$ with alias terms $|r|\le C_Ns^N$ on the Brillouin zone, and $q_\Lambda=|T|^{-1}\sum_{k\in B_\Lambda\cap P^*}\sum_\beta|\hat\ell_{\Lambda,\beta}(k)|^2/(\sum_\gamma\sigma_\gamma(k)^2+\lmu^2)$. In the continuum, $q=|T|^{-1}\sum_{k\in P^*}(|k|^2+\lmu^2)^{-1}\sum_\beta|\sum_\alpha k_\alpha\hat u_{\alpha\beta}(k)|^2$, and $\sum_\beta|\sum_\alpha k_\alpha G(f_\alpha,f_\beta)|^2=\sum_\beta|G(k,f_\beta)|^2=|\iota_kG|^2$ is frame-free because $\sum_\alpha k_\alpha f_\alpha=k$. Finally $|\sigma_\alpha-k_\alpha|\le s^2|k|^3/24$, $|k|^2-s^2|k|^4/12\le\sum\sigma_\gamma^2\le|k|^2$, $\sum\sigma_\gamma^2\ge\frac4{\pi^2}|k|^2$ on $B_\Lambda$, and the continuum terms with $k\notin B_\Lambda$ sum to $O(s^N)$.
\end{proof}

\begin{lemma}[the shift field is regular]\label{rot:lem:shift}
Let $a^{\rm c}\defeq(-\Delta+\lmu^2)^{-1}\dstar G$, a smooth co-exact $1$-form on $T$. For every bond $(x,x+sf_\beta)$ of $\Lambda$, $|a_{\Lambda,\beta}(x)-(a^{\rm c}\cdot f_\beta)(x+\frac s2f_\beta)|\le C(G,\lmu)s^2$, and the same for lattice differences of any fixed order. In King's unit scaling at step $r$ the background $za_\Lambda$ is regular in the sense of \cite[Def.~3.2]{K1} for $\delta_r\le\delta_0(z,G,e)$, and so is $A^{(r)}+za_\Lambda$ whenever $A^{(r)}$ is.
\end{lemma}
\begin{proof}
$\hat a_{\Lambda,\beta}=(\sum_\gamma\sigma_\gamma^2+\lmu^2)^{-1}\hat\ell_{\Lambda,\beta}$; compare coefficient by coefficient with the sampled continuum field as in \Cref{rot:lem:q}. In unit scaling ($A_u(x_u)=\delta_r^{1/2}A(\delta_rx_u)$, $e_u=e\delta_r^{1/2}$, \cite[(2.20)]{K1}) a derivative of $za_\Lambda$ is $\le|z|\delta_r^{3/2}(\|\nabla a^{\rm c}\|_\infty+Cs^2)$, while \cite[Def.~3.2]{K1} allows $Ce^{\beta-1}\delta_r^{(\beta-1)/2}$ with a small $\beta>0$.
\end{proof}

\begin{remark}[what the transfer does for the comparison]\label{rot:rem:CMrg}
After \Cref{rot:lem:CM} the curvature source has disappeared from the photon: the model is King's with a scalar source and a fixed smooth external background $za_\Lambda$ in the matter coupling only. Running the matter averages in the total field, the background enters exactly where King's coarse background $A^{(r)}$ enters, through $A^{(r)}+za_\Lambda$, regular and small at every scale below $\delta_0(z,G)$ (\Cref{rot:lem:shift}); this is the content of the identities \cite[(3.82)--(3.87)]{K1}, which express the effective action at a background as the action at zero background with the background inserted in the lines and vertices. The two backgrounds $a_{\Xk}$ (axis bonds) and $a_{\Yk}$ (rotated bonds) are samples of one continuum $1$-form $a^{\rm c}$, so replacing one by the other at matched points is the same step as King's replacement of $g_{\mu\nu}(x')$ by $g_{\mu\nu}(x)$ \cite[p.~665]{K1}. The comparison of the effective actions below is therefore carried out at ``$h=0$ with background'', and no frame-paired source vertex \cite[(3.20)]{K1} appears.
\end{remark}

\begin{lemma}[lattice sums]\label{rot:lem:lsum}
For $h\in C^\infty(T)$, every lattice $\Lambda=O(s\Z^3)\supset P$ and every $N$: $|\sS_\Lambda(h)-\int_Th|\le C_N(h)s^N$, with $C_N(h)$ independent of $O$.
\end{lemma}
\begin{proof}
$\sS_\Lambda(h)=\sum_{k\in\Lambda^*}\hat h(k)$ by Poisson summation on $\Lambda/P$ ($P\subset\Lambda$ gives $\Lambda^*\subset P^*$); the term $k=0$ is $\int h$; every $k\in\Lambda^*\setminus0$ has $|k|\ge2\pi/s$; $|\hat h(k)|\le\|\Delta^Mh\|_{L^1}|k|^{-2M}$ and $\sum_{k\in\Lambda^*\setminus0}|k|^{-2M}\le C_Ms^{2M}$ for $2M>3$.
\end{proof}

\subsubsection*{The number of steps, the effective actions, the matching}

\begin{lemma}[the choice of $r$]\label{rot:lem:r}
Let $\gamma_*,\beta_*,\sigma>0$ be the exponents of \cite[Thm.~3.4]{K1}, $\beta_1$ the field-norm exponent of step (5) of the proof of \Cref{rot:thm:A}, $\theta_1$ the gain exponent of \Cref{rot:thm:H1S}, $\gamma\defeq\min(\gamma_*,\theta_1,\frac12)$ and $\beta\defeq\max(\beta_*,\beta_1+\frac12,5)$. Fix $0<\vr<\min\{\gamma/(\gamma+\beta),\frac19\}$ and put $k-r(k)\defeq\lfloor\vr k\rfloor$. Then, as $k\to\infty$: $\delta_r\to0$; $5^{-\gamma r}\delta_r^{-\beta}(1+k)^p\to0$ for every $p$; $5^{3k-5r}\to0$; $\delta_r^{-2}(\eps_k/\delta_r)^{1/4}\log(1/\eps_k)^2\to0$; and $\eps_k^{-q}e^{-c\delta_r^{-1}}\to0$ for all $q,c>0$. Moreover $p(\delta_r)^2\ge b_0^2(1+\vr k\ln5-\ln5)^{2\mathrm p}$ grows at least quadratically in $k$, where $p(\delta)=b_0(1+\log\delta^{-1})^{\mathrm p}$, $\mathrm p\ge1$, is King's small-field function \cite[(3.1)]{K1}.
\end{lemma}
\begin{proof}
$\delta_r=5^{-\lfloor\vr k\rfloor}$. The second quantity is $5^{-\gamma k+(\gamma+\beta)\lfloor\vr k\rfloor}(1+k)^p$; the third $\le5^{-2k+5\vr k+5}$; the fourth $\le5^{2\vr k+2-(1-\vr)k/4}(k\ln5)^2$, and $2\vr<(1-\vr)/4$ iff $\vr<\frac19$; the fifth because $\delta_r^{-1}\ge5^{\vr k-1}$ grows exponentially. The last claim is \cite[(3.1)]{K1} with $\mathrm p\ge1$. King's own choice \cite[(3.12)]{K1}, $\vr=\gamma/(2\beta+\gamma)$, satisfies every condition.
\end{proof}

\begin{definition}[effective actions; vertices]\label{rot:def:eff}
For the standard flow the $r$-step effective action is \cite[(3.14)]{K1} rewritten as \cite[(3.90)]{K1}:
\[
S^{(r)}=\tfrac12\ip{A_r}{\Delta^{(r)}A_r}+\ln N_r+\tfrac12\ip{\phi_r}{\Delta^{(r)}(A^{(r)})\phi_r}+\ln N_r(0)-\ln\big[Z_r(A^{(r)})Z_r(0)^{-1}\big]+P^{(r)}(A_r,\phi_r;s),
\]
where $P^{(r)}$ is the Taylor polynomial of order $\bar n=13$ in the running couplings of King's $E'_r$ \cite[(3.17)]{K1}, a finite sum of connected graphs with the vertices \cite[(3.18)--(3.26)]{K1}. The hybrid $\hs^{(r)}$ is the same expression for the hybrid flow. A curvature source enters only through the Cameron--Martin background (\Cref{rot:rem:CMrg}); a wall enters through the source vertex \cite[(3.21)]{K1}, the \emph{$h$-vertex}, $\theta(5^r\eps)^2\sum_x\eta^3h(x)\sum_i[(\phi'_i+\phi^{(r)}_i)(x)^2-(-\Delta^{\eta,\Omega}+m^2(5^r\eps)^2)^{-1}(x,x)]$, whose subtraction is in \cite{K1} the coincident torus propagator at the running scale and here, in the free-boundary convention which both flows inherit from the fine action by exact integration, the corresponding infinite-lattice coincident propagator, the same number for both lattices because $\Xk$ and $\Yk$ are isometric; the two readings differ by the image sum $\sum_{\gamma\in P\setminus0}$ of the running propagator, a $c$-number times $\sS(h)$, which is part of the bridge \Cref{rot:prop:CT}. Everything below is written in the free-boundary reading.
\end{definition}

\begin{definition}[matching; slices]\label{rot:def:match}
In the unit scale of level $r$ let $\Gamma_r=R^{-1}\Z^3$ (the coarse lattice), $\Gamma_1=R^{-1}(\eta\Z^3)$ ($=Y_{k-1}$ rescaled), $\Lambda_X=\eta'\Z^3$ and $\Lambda_Y=R^{-1}(\eta'\Z^3)$. For $x\in\Gamma_1$ the first-step cells are the tile $x+\eta'B_0\subset\Lambda_X$ and the centred cube $Q_Y(x)\subset\Lambda_Y$, $125$ sites each. Fix bijections $\tau_x:x+\eta'B_0\to Q_Y(x)$; they glue to a bijection $\tau:\Lambda_X\to\Lambda_Y$, and on every admissible torus to $\Xk/P\to\Yk/P$, with $|\tau x'-x'|\le c_\tau\eta'$, $c_\tau=\sqrt{13}+2\sqrt3<7.1$. Points $x''=\tau x'$ are \emph{matched}. For a flow $F\in\{X,Y\}$ let $S^F_0\defeq G^F_1$ be the first fluctuation covariance and $S^F_j\defeq H^F_jC^F_jH^{F*}_j$ the $(j+1)$-st, lifted to the fine lattice; then $G^F_r=\sum_{j<r}S^F_j$ exactly, and on the infinite lattice $C^F=G^F_n+\Gamma^F_n$ with $\Gamma^F_n\defeq H^F_n(A^F_n)^{-1}H^{F*}_n$ (Woodbury). Slice $j$ has length scale $\ell_j\defeq5^j\eta'$. The \emph{coarse slices} $j\ge r$ are the further steps of the same flow; for $F=X$ these are standard rotated steps on $Y_{k-r},Y_{k-r-1},\dots$
\end{definition}

\begin{lemma}[vertex matching; frame-free contractions]\label{rot:lem:vert}
\textup{(i)} $\sum_{y\in\Lambda_Y^n}\eta'^{3n}g(y)=\sum_{x\in\Lambda_X^n}\eta'^{3n}g(\tau x)$: every $Y$-graph is, exactly, a sum over $X$-vertices with $\tau$-transported kernels. \textup{(ii)} A smooth vertex function $w$ (couplings, the Cameron--Martin background, partitions of unity) sampled at $x'$ and at $\tau x'$ differs by $\le c_\tau\eta'\|\nabla w\|_\infty$. \textup{(iii)} Vector indices are contracted by the Euclidean inner product and derivative lines are compared as vectors; the current vertex $eA\cdot j$ (\cite[(3.22)]{K1}, $m=1$) and the seagull $e^2|A|^2|\phi|^2$ (the $m=2$ member of \cite[(3.23)]{K1}) are frame-free; \cite[(3.22)]{K1} with $m=2$ is the bond-diagonal derivative vertex $\frac{e^2\eta}2\sum_bA_b^2D_b\phi\cdot U\phi(b_+)$ and is not frame-free; the vertices with $m\ge3$ in (3.22) and $m\ge4$ in (3.23) carry explicit factors $\eta'^{m-2}$ or more and are bounded per flow.
\end{lemma}
\begin{proof}
(i) $\tau$ is a bijection of $\Lambda_X$ onto $\Lambda_Y$ and both lattices have the cell volume $\eta'^3$ per site, so $y=\tau x$ is an exact change of the summation variable in each factor; on an admissible torus the same holds for the bijection $\Xk/P\to\Yk/P$. (ii) The mean value theorem along the segment from $x'$ to $\tau x'$. (iii) In the unit scale of level $r$ the hopping term on the fine lattice of spacing $\eta'$ is $\sum_b\eta'\big[|\phi(b_-)|^2+|\phi(b_+)|^2-2\Re\big(\bar\phi(b_+)e^{\mathrm ie_r\eta'A_b}\phi(b_-)\big)\big]$, the bond integral of the unit-scaled field being $\eta'A_b$ \cite[(2.20)]{K1} and $e_r=e(5^r\eps)^{1/2}$ the running charge; its expansion in powers of $A_b$ carries $(e_r\eta')^m/m!$ at order $m$, one explicit $\eta'$ per power of $A$, which is how \cite[(3.22)--(3.23)]{K1} are organised. At $m=1$ the vertex is $\sum_x\eta'^3e_rA(x)\cdot j(x)$ with $j=\sum_\alpha j_\alpha f_\alpha$ the lattice current, $j_\alpha(x)=2\,\Im(\bar\phi(x+\eta'f_\alpha)\phi(x))/\eta'$; at $m=2$ the member of (3.23) is $\sum_x\eta'^3e_r^2|A(x)|^2|\phi(x)|^2$. The pairing $A\cdot j$ and the norm $|A|^2$ are the same functions of the vectors in every orthonormal frame; comparing the $X$-version with the $\tau$-transported $Y$-version therefore compares the vectors $A$, $j$, $\phi$ at matched points only. The member of (3.22) with $m=2$ contracts $A\otimes A\otimes D\phi$ with the rank-three tensor $\sum_\alpha f_\alpha\otimes f_\alpha\otimes f_\alpha$, which is not $O(3)$-invariant (there is no non-zero isotropic symmetric tensor of rank three), so the two versions differ at order one on the same vectors; its tadpole is handled by the summation by parts of \Cref{rot:lem:CK4}. The vertices with $m\ge3$, $m\ge4$ carry $\eta'^{m-2}$ relative to the two order-$e_r^2$ vertices, and each flow's own line bounds \cite[(3.63)--(3.65)]{K1} bound them without any comparison.
\end{proof}

\begin{lemma}[counterterms take the same value along both flows]\label{rot:lem:ctsame}
Let $g$ be a counterterm graph of $\dm$ or $E_1$ (free boundary conditions), an infinite-lattice graph with full lines $C$ localised at one vertex. Expanding each of its lines in the slices of flow $F$ \cite[(3.76)]{K1}, the sum over all slice assignments equals $g$ for both flows. Mixed assignments take other values; they occur only inside a telescoping sum whose endpoints are the pure assignments.
\end{lemma}
\begin{proof}
$\sum_jS^F_j=C^F$ pointwise; $C^Y$ is the image of $C^X$ under the isometry $R^{-1}$, and the localised vertex sum of an infinite-lattice graph is invariant under that isometry.
\end{proof}

% r03_fill_K2b: proof of (K_0), part III: the equal-spacing propagator package at zero gauge field.

\subsubsection*{The equal-spacing propagator package}

Unit scale of level $r\ge2$, infinite lattice, $A=0$; $V_\eta=R^{-1}[-\frac\pi\eta,\frac\pi\eta]^3$, $V_1=R^{-1}[-\pi,\pi]^3$; $u_{r-1}$ is the common symbol of the $(r-1)$-step rotated centred average; $u_X(q)=\frac1{125}\sum_{x\in B_0}\cos(\eta'q\cdot x)$ is the tile symbol and $u_Y(q)=\prod_\alpha W(\eta'q\cdot f_\alpha)$, $W(t)=\frac15\sum_{j=-2}^2\cos jt$, the cube symbol; $\Delta'_X(q)=4\eta'^{-2}\sum_\mu\sin^2\frac{\eta'q_\mu}2+m_r^2$ and $\Delta'_Y$ its rotated analogue, $m_r=m\delta_r$ the running mass; $L^F_2$ are the $125$ alias classes; $A^F$ is the coarse multiplier of flow $F$ (the symbol of $\Delta^{F(r)}$), $A^F\le a$; $\Delta_R$ is the rotated Laplacian at spacing $\eta$; $J_s\defeq\int_{V_\eta}|u_{r-1}(q)||q|^{-s}dq$, bounded uniformly in $r$ for $s=1$. Constants depend only on $a$ and are uniform in $r$, $k$, the torus and the mass $m_r^2\ge0$.

\begin{proposition}[free core at equal spacing]\label{rot:prop:fc}
$A^F(p)^{-1}=a_{r-1}^{-1}+\sum_{l_1\in\Gamma_r^*/\Gamma_1^*}u_{r-1}(p+l_1)^2[\eta^2a^{-1}+S_F(p+l_1)]$ with $S_F(q)\defeq\sum_{l\in L^F_2}u_F(q+l)^2/\Delta'_F(q+l)$, and $\sup_{V_\eta}|S_X-S_Y|\le0.198\,\eta^2$. Hence for all $r\ge2$, all $p$ and all $m_r\ge0$,
\[
|A^X(p)-A^Y(p)|\le0.20\,\eta^2A^X(p)A^Y(p)\le0.20\,a\,\eta^2A^Y(p),
\]
and the same for the photon (three copies at mass $\lmu\delta_r$).
\end{proposition}
\begin{proof}
The formula is the Woodbury identity for a composite average whose symbol is $u_{r-1}u_F$; the terms $\eta^2/a$ cancel in the difference and $\sum_{l_1}u_{r-1}^2=1$ gives the sup bound. Scale to $\eta=1$, $\eta'=\frac15$. \emph{Zones:} on $V_1$, $\Delta'_X\ge0.93718|q|^2+m^2$ and $\Delta'_Y\ge0.96753|q|^2+m^2$ (monotonicity of $\sin x/x$), while $\Delta'_F\le|q|^2+m^2$ and $(|q|^2+m^2)-\Delta'_F\le\eta'^2|q|^4/12$. \emph{Second moments:} $1-u_X\le\frac{\eta'^2}2q^{\mathrm T}M_0q$ with $M_0=\diag(\frac{52}{25},\frac{52}{25},2)$ (exact enumeration of the tile), so $1-u_X^2\le\frac{52}{25}\eta'^2|q|^2$; $1-W(t)\le t^2$ gives $1-u_Y^2\le2\eta'^2|q|^2$. The main terms ($l=0$) then differ by at most $0.0924$. \emph{Aliases:} every non-zero alias $q+l$ has a frame component of modulus in $[\pi,5\pi]$, where $100\sum\sin^2(r_\mu/10)\ge9.549$ (concavity of $t\mapsto\sin^2(\sqrt t/10)$ on $[0,25\pi^2]$); with Parseval $\sum_lu_F(q+l)^2=1$, $0\le\mathrm{Al}_F\le(1-u_F^2)/9.549\le0.1047$. Total $0.198$. Finally $A^F\le a_{r-1}\le a$. (The true constant is about $0.005$: both first moments vanish by inversion symmetry and the second moments differ by $4\%$.)
\end{proof}

\begin{lemma}[external lines]\label{rot:lem:H}
For $x'\in x+\eta'B_0$, $x''=\tau x'\in Q_Y(x)$ ($x\in\Gamma_1$) and $z\in\Gamma_r$,
\[
H^F(x^F,z)=(2\pi)^{-3}\int_{V_\eta}dq\,e^{\mathrm iq\cdot(x-z)}u_{r-1}(q)M^F(q)A^F(q),\qquad M^F(q)\defeq\sum_{l\in L^F_2}\frac{e^{\mathrm i(q+l)\cdot r_F}u_F(q+l)}{\Delta'_F(q+l)},
\]
with $r_X=x'-x$, $r_Y=x''-x$, $x^X=x'$, $x^Y=x''$, and
\[
|H^X(x',z)-H^Y(x'',z)|\le C_EJ_1\,\eta,\qquad |H^X(x',z)-H^Y(x'',z)|\le C(C_EJ_1\eta)^{1/2}e^{-\delta|x-z|/2},
\]
$C_E=345a+0.2a(C_*/20+6.9a)$, $C_*=\max\{43.6(1+\sqrt3\pi)^2,\frac{\pi^2}4a(1+\frac1\pi)^2\}$.
\end{lemma}
\begin{proof}
The representation: $H=D^{-1}Q^*\Delta^{(r)}$ (Woodbury), $Q^*$ has symbol $u_Fu_{r-1}$; decompose the fine Brillouin zone into the translates $V_\eta+l$, $l\in L^F_2$, and use the $\Gamma_1^*$-periodicity of $u_{r-1}$, $A^F$ and $e^{-\mathrm iq\cdot z}$, with $e^{\mathrm il\cdot x}=1$. \emph{Bounds on $V_\eta$}: for both flows $|M^F-\Delta_R^{-1}|\le172.1\,\eta|q|^{-1}$. For $Y$: $|e^{\mathrm iq\cdot s}u_Y(q)-1|\le|q||s|+(1-u_Y)\le0.910\,\eta|q|$ ($|s|\le2\sqrt3\eta'$, $1-u_Y\le\eta'^2|q|^2$), divided by $\Delta'_Y\ge\frac4{\pi^2}|q|^2$ gives $2.25\,\eta|q|^{-1}$; $|\Delta'^{-1}_Y-\Delta_R^{-1}|\le\min\{0.51\eta^2,\frac{\pi^2}4|q|^{-2}\}\le1.12\,\eta|q|^{-1}$; the $124$ aliases have $\Delta'_Y(q+l)\ge4/\eta^2$ and $|u_Y|\le1$, so their sum is $\le31\eta^2\le168.7\,\eta|q|^{-1}$ on $V_\eta$. For $X$ the same three bounds are $2.37+0.73+169$. Hence $|M^X-M^Y|\le345\,\eta|q|^{-1}$. \emph{Multipliers}: $E\defeq M^XA^X-M^YA^Y=(M^X-M^Y)A^X+M^Y(A^X-A^Y)$; the first term is $\le345a\eta|q|^{-1}$; for the second use \Cref{rot:prop:fc}, $|M^Y|\le\Delta_R^{-1}+172.1\eta|q|^{-1}$ and $\Delta_R^{-1}A^Y\le C_*(1+|q|)^{-2}$ (on $V_1$ from $A^Y(q)^{-1}\ge u_{r-1}^2u_Y^2\Delta'^{-1}_Y$ with $u_{r-1}^2\ge(2/\pi)^6$, $u_Y\ge0.95$, $\Delta'_Y/\Delta_R\le\pi^2/4$; off $V_1$ from $|q|\ge\pi$), so $|M^Y(A^X-A^Y)|\le0.2a(C_*/20+6.9a)\eta|q|^{-1}$. Integrate $|E|\le C_E\eta|q|^{-1}$ against $|u_{r-1}|$. The decaying form is the geometric mean with the separate decay bound $|H^F(x^F,z)|\le Ce^{-\delta|x^F-z|}$ of each chain (Combes--Thomas for the coercive coarse operators \cite{CT73}) and $|x^F-x|\le2\sqrt{13}\eta'$.
\end{proof}

\begin{lemma}[gradients and H\"older bounds at $A=0$]\label{rot:lem:grad}
$|\nabla H^F(x,z)|\le C_\nabla$ for the forward difference of each flow, hence $|H^F(x,z)-H^F(y,z)|\le C|x-y|^{1/2}(e^{-\delta|x-z|/2}+e^{-\delta|y-z|/2})$.
\end{lemma}
\begin{proof}
The difference multiplies the alias-$l$ summand by $d_F(q+l)$ with $|d_F(Q)|\le\min(|Q|,2\sqrt3/\eta')$; the main term is $\le C|q|(1+|q|)^{-2}$ and the aliases give $537a\,\eta$ integrated against $|u_{r-1}|$ with $J_0\le C\log^3(1/\eta)$ (\Cref{rot:lem:CK0}), $\eta\log^3(1/\eta)\le C$. Geometric mean with the decay.
\end{proof}

\begin{lemma}[coarse covariances]\label{rot:lem:Gam}
Let $B^F\defeq(A^F)^{-1}$ and $\Gamma^F\defeq H^FB^FH^{F*}$. \textup{(i)} $C^F=G^F_r+\Gamma^F$ exactly (Woodbury). \textup{(ii)} $\Gamma^F(x,y)=(2\pi)^{-3}\int_{V_1}B^F(p)\hat h^F_x(p)\overline{\hat h^F_y(p)}dp$ with $\hat h^F_x(p)\defeq\sum_zH^F(x,z)e^{-\mathrm ip\cdot z}$; $B^F(p)\le C_B|p|^{-2}$ on $V_1$ and $\sup|\Gamma^F|\le C_1$. \textup{(iii)} $|\Gamma^X(x',y')-\Gamma^Y(x'',y'')|\le C_\Gamma\eta^{1/2}$ for matched pairs. \textup{(iv)} $|\Gamma^F(x,y)-\Gamma^F(x,w)|\le C_2|y-w|^{1/2}$.
\end{lemma}
\begin{proof}
(ii) Parseval on $\Gamma_r$; $A^F(p)^{-1}=a_{r-1}^{-1}+\sum_{l_1}u_{r-1}(p+l_1)^2[\eta^2a^{-1}+S_F(p+l_1)]$ with $S_F(q)\le\frac{\pi^2}{4|q|^2}+31\eta^2$; the term $l_1=0$ gives $\frac{\pi^2}{4|p|^2}+C$ and the others are summable. (iii) $|B^X-B^Y|\le0.198\,\eta^2$ pointwise (\Cref{rot:prop:fc}), $|\hat h^X_{x'}-\hat h^Y_{x''}|\le C'_h\eta^{1/2}$ (\Cref{rot:lem:H}, decaying form); three differences. (iv) \Cref{rot:lem:grad}.
\end{proof}

\begin{lemma}[scale pieces $j\ge2$]\label{rot:lem:slice}
For $j\ge2$, in the unit scale of level $j$ ($\eta_j=5^{1-j}$) and matched points, $|S^X_j(x',y')-S^Y_j(x'',y'')|\le C\eta_j^{1/2}e^{-\delta|x'-y'|}$. In the unit scale of level $r$ ($j<r$): $\le C(5\eta'/\ell_j)^{1/2}\ell_j^{-1}e^{-\delta|x'-y'|/\ell_j}$. The same holds for the photon.
\end{lemma}
\begin{proof}
At level $j$: $|(A^X_j-A^Y_j)(z,w)|\le C\eta_je^{-\delta|z-w|/2}$ (geometric mean of the sup bound $0.20a^2\eta_j^2$ with the decay of each kernel); $C^F_j=(A^F_j+bQ^*Q)^{-1}$ decays for both chains; $C^X_j-C^Y_j=C^X_j(A^Y_j-A^X_j)C^Y_j$. Then $S^X_j-S^Y_j\circ\tau=(H^X_j-H^Y_j\circ\tau)C^X_jH^{X*}_j+H^Y_j\circ\tau(C^X_j-C^Y_j)H^{X*}_j+H^Y_j\circ\tau\,C^Y_j(H^X_j-H^Y_j\circ\tau)^*$; bound the three terms by \Cref{rot:lem:H} at level $j$ and sum the exponentials. Rescale with \cite[(2.20)]{K1}: scalar kernels scale as $\ell^{-1}$ in three dimensions.
\end{proof}

\begin{corollary}[uniform comparison form]\label{rot:cor:unif}
Let $D_j\defeq S^X_j-S^Y_j\circ\tau$ in the unit scale of level $r$ and $\theta\defeq\frac12$. \textup{(a)} Fine slices $0\le j\le r-1$: $|D_j(x',y')|\le C_D\eta'^{\theta}\ell_j^{-1-\theta}e^{-\delta|x'-y'|/\ell_j}$. \textup{(b)} Coarse slices and remainders $j\ge r$: $|D_j|\le C\eta^{1/2}\ell_j^{-1}e^{-\delta|x'-y'|/\ell_j}$ and $|\Gamma^X_n-\Gamma^Y_n\circ\tau|\le C\eta^{1/2}$ at level $n$. \textup{(c)} Given \Cref{rot:prop:BD}(a) for the hybrid flow and \cite[Thm.~3.3]{K1} for the standard flow, the derivative, H\"older and contour versions hold: the contour lines \cite[(3.27)]{K1} along the two paired contour systems differ by the kernel difference along the contour plus the flux of the slice kernel through the sliver between the paired partial contours, a gain $(\eta'/\ell_j)^{1/2}$ on slice $j$; the double contour lines, both ends on contours, have the same relative gain against their size $C\ell_j$.
\end{corollary}
\begin{proof}
(a) For $j\ge2$ this is \Cref{rot:lem:slice}. For $j\in\{0,1\}$ bound $|S^X_j|$ and $|S^Y_j\circ\tau|$ separately by $C\ell_j^{-1}e^{-\delta|\cdot|/\ell_j}$ and use $\ell_j\le5\eta'$, so $1\le5^\theta(\eta'/\ell_j)^\theta$: there is no gain on these lines, and the factor $\eta'^\theta$ is paid for by a degree loss $\theta$ (King's own device at the finest slice, \cite[p.~675]{K1}: ``When $j=0$, we bound each term separately''). (b) Beyond level $r$ both flows apply the same rotated steps; only the input $B^F$ differs, by the relative amount $0.2a\eta^2$ (\Cref{rot:prop:fc}), and every coarse multiplier is built from it by the same formulas with non-negative weights. (c) Interpolation of (a), (b) with the separate bounds of each flow; for the contours, \Cref{rot:lem:pair,rot:lem:flux}: the sliver has area $\le1.5\eta'\ell_{j+1}+33\eta'^2$, against the slice kernel's gradient $\le C\ell_j^{-2}$, plus the end segment $c_\tau\eta'$ and Riemann errors; for $j\in\{0,1\}$ there is no gain and the factor is paid by a degree loss, as in (a).
\end{proof}

\begin{corollary}[the full propagators of the two lattices]\label{rot:cor:C}
On the infinite lattices, for matched pairs and $u=x'-y'$ with $|u|\ge25\eps$ (physical units),
\[
|C^X(x'-y')-C^Y(x''-y'')|\le C_6\,\eps^{1/2}|u|^{-3/2},
\]
and the same for the photon and for $G^F_r$, $\Gamma^F$ separately. For all $u$ the difference is $\le2C_0(|u|+\eps)^{-1}$.
\end{corollary}
\begin{proof}
Pick $n\ge2$ with $\delta_n\le|u|<5\delta_n$ and write $C^F=\sum_{j<n}S^F_j+\Gamma^F_n$. The slices $j\in\{0,1\}$ are each $\le C\eps^{-1}e^{-\delta_*|u|/(5\eps)}\le C'\eps^{1/2}|u|^{-3/2}$; the slices $2\le j<n$ differ by $\le C\delta_j^{-1}(5\eps/\delta_j)^{1/2}e^{-\delta|u|/\delta_j}$ and sum to $\le C\eps^{1/2}|u|^{-3/2}$; the remainder differs by $\le C_\Gamma\delta_n^{-1}(5\eps/\delta_n)^{1/2}$. Pairs stay matched at every level because $\tau$ is defined by the first-step cells.
\end{proof}

Since $C^Y(x''-y'')=C^X(R(x''-y''))$, \Cref{rot:cor:C} is a quantitative rotation estimate for the massive lattice Green's function, with relative error $(\eps/|u|)^{1/2}$, obtained from the two block-spin decompositions alone.

% r03_fill_K3a: proof of (K_0), part IV: the comparison principle, the two-flow bound, the constant-background reduction.

\subsubsection*{The comparison principle and the two-flow bound}

A \emph{renormalised family} is a finite sum of King's renormalised graphs \cite[\S3.5]{K1}, Ba{\l}aban's classes of \cite[\S\S2--3]{Ba3} (\Cref{def:a07:class}), localised as in \cite[(1.33)]{Ba3}. It is in \emph{per-flow normal form} when, for each flow separately and before any element is mixed, every counterterm graph is expanded in the lines of the flow \cite[(3.76)]{K1}, every Taylor subtraction \cite[(3.83)]{K1} and Ward-identity rewriting has been carried out with that flow's exact identities, and the exact cancellations between members with identical element sets (\cite[(3.6)/(3.7), (2.21b)/(2.21g), (2.18)]{Ba3}) are co-switched, the identified elements of both members being switched together. Its \emph{elements} are lines, vertex functions and external fields.

\begin{proposition}[comparison principle]\label{rot:prop:cp}
Let $\cR$ be a renormalised family in per-flow normal form with $n_\cR$ elements and $E^F(\cR)$ its value along flow $F$. Then, with $\theta=\min(\frac12,\theta_1)$ and $\theta_1$ of \Cref{rot:thm:H1S},
\[
|E^X(\cR)-E^Y(\cR)|\le Cn_\cR\,p_r^2\delta_r^{-1/2}\,\eta'^{\theta}(1+k)^p\cdot[\text{right side of \cite[(1.33)]{Ba3} for }\cR],
\]
the prefactor $p_r^2\delta_r^{-1/2}$ being that of the contour elements at $e>0$ (absent at $e=0$).
\end{proposition}
\begin{proof}
By \Cref{rot:lem:vert}(i) both values are sums over $X$-vertices. Order the elements consistently across members and telescope: $E^X-E^Y=\sum_iE[e_1,\dots,e_{i-1}\text{ in $X$-version};\ e_i\to e_i^X-e_i^Y;\ e_{i+1},\dots\text{ in $Y$-version}]$. In the $i$-th term the element $e^X_i-e^Y_i$, divided by $C\eta'^\theta$, obeys the bounds of its type weakened by the degree $\theta$: fine lines by \Cref{rot:cor:unif}(a), coarse lines by (b) ($\eta^{1/2}\le\sqrt5\eta'^{1/2}$), vertex functions by \Cref{rot:lem:vert}(ii) with $\eta'=\eta'^\theta\eta'^{1-\theta}$, external fields $H^F\phi_r$ and their covariant derivatives by \Cref{rot:lem:H} and its derivative and H\"older versions; at $e>0$, contour lines \cite[(3.27)]{K1} by \Cref{rot:cor:unif}(c) (gain $(\eta'/\ell_j)^{1/2}$), double contour lines by \Cref{rot:prop:BD}(b) (the same relative gain, size $C\ell_j$), contour vertices \cite[(3.24)--(3.26)]{K1} by \Cref{rot:prop:BD}(b) (relative gain $\eta^{1/2}$, size $Ce_rp_r^2$) and the constant-background phases \cite[(3.46)]{K1} by \Cref{rot:prop:BD}(b) ($\le e\,p_r\lmu^{-1}c_\tau\delta_r^{-1/2}\eta'$), whence the prefactor. The weakened-element bound for each term is \Cref{rot:thm:H1S}, whose proof runs this telescoping with the subtraction rule described there.
\end{proof}

\begin{theorem}[the two-flow comparison bound]\label{rot:thm:H1S}
Let $\alpha\in(0,1)$ be the H\"older exponent of \cite[(3.8)]{K1} and put $\theta_1\defeq\alpha^2/(4(1+\alpha)^2)$. There are $p=p(\bar n)$ and $C$ such that for every renormalised family $\cR$ of order $\le\bar n$ in per-flow normal form (the $h$-families of \Cref{rot:lem:S1} other than $\cT_h$, and the families of step (5) of the proof of \Cref{rot:thm:A} at $e>0$ after the constant-background reduction of \Cref{rot:lem:CB}), for $0<\theta\le\theta_1$, uniformly in $r$, $k$, the torus and the masses, the two-flow difference obeys the bound of \Cref{rot:prop:cp}: in every term of its telescoping, the element replaced by the difference $e^X_i-e^Y_i$, which obeys the bounds of its type weakened by the degree $\theta$, \cite[(1.33)]{Ba3} holds with the extra factor $C(1+k)^p\eta'^{\theta}$ and the prefactor stated there. The exponent $\theta_1$ does not depend on $\bar n$, and any $\theta>0$ with any polynomial factor suffices for the consumers (\Cref{rot:lem:r}).
\end{theorem}
The proof is given at the end of this part, after the constant-background reduction (\Cref{rot:lem:CB}), which it uses at $e>0$. Its inputs are \Cref{rot:prop:BD}, \cite[Thm.~3.3]{K1} for the standard flow, \cite[Prop.~2.1, (2.8)]{Ba3} with the ordering device (2.7) there, the zero-momentum counterterm conditions \cite[(3.33)]{K1} through order four with the vanishing of the one-loop vacuum-polarisation moments \cite[(2.23)--(2.28)]{Ba3}, the local limit law for the simple random walk \cite[Thm.~2.5.6]{LL10}, and the $h$-census of \Cref{rot:lem:S1}. In outline: every line is written as one flow-independent lattice function sliced by the heat kernel; degrees are read per ordering of the scale-summed elements; each unit (a self-energy, vacuum-polarisation or triangle insertion) is bundled with one of its legs into a single element, and one subtraction rule is applied, in each flow's own frame, to every connected component of degree zero that forms in an initial segment, its local terms being common numbers of both flows; after the rule every component containing the weakened element has degree $\ge\min(\alpha,1)$, and the scale sums give the gain $\eta'^{\theta_1}$.

\begin{lemma}[the constant-background reduction]\label{rot:lem:CB}
In the localised form \cite[(3.55)]{K1} every family has its external vertices in a cube $D$ of side $\le Cp_r$. Let $D''$ be the union of the level-$r$ blocks within $p$ of $D$, $x_0\in D$, $B\defeq A_r(x_0)$ and $\tilde A\defeq A^{(r)}-B$ on all of $T$. For each flow $F$ the reduction runs on the infinite fine lattice $\Lambda_F$ of level $r$ with the $P_T$-periodic lifts $A_{\rm lift}$, $\tilde A_{\rm lift}$ (no periodic cell is formed); $F_T$ is the lifted fundamental cube centred at $D$, $F_1\defeq\{\dist(\cdot,D'')\le p/3\}$, and the vertex region $F_v$ is the union of the level-$r$ blocks within $p/3-3$ of $D''$. The sup-systole of $P_T$ in the unit scale of level $r$ is $s_T\ge40MJ\,5^{\lfloor\vr k\rfloor}\ge Cp_r$ for $k\ge k_0$.
\begin{enumerate}[label=\textup{(\arabic*)},leftmargin=*]
\item[(0)] \textup{(P)} For $x,y\in F_T$, $G_r(T,A^{(r)})(x,y)=\sum_{\gamma\in P_T}G_r(\Lambda_F,A_{\rm lift})(x,y+\gamma)$, and likewise for $D_AG_r$, $Q_r$ and $H^F$; the image part $\delta^{\rm im}$ (the terms $\gamma\ne0$) obeys $|(\delta^{\rm im}f)(x)|,|(D_A\delta^{\rm im}f)(x)|\le Ce^{-\delta_0s_T/4}\|f\|_\infty$ for $\supp f\subset B(x,s_T/4)$. \textup{(G)} On $\Lambda_F$ the constant $B$ is removed exactly by the gauge function $B\cdot x$, leaving King's phases \cite[(3.46)]{K1} at the vertices. \textup{(R)} $|\tilde A_r(w)|\le p_r(d_T(w,x_0)+1)$ and $|\partial\tilde A^{(r)}|\le Cp_r$ on all of $T$, so $\tilde A_{\rm lift}$, $\vartheta\tilde A_{\rm lift}$ and $B_{\rm bg}\defeq(1-\vartheta)\tilde A_{\rm lift}$ are regular. \textup{(N)} If $A_1$, $A_2$ are regular on $\Lambda_F$ and $A_1=A_2$ on the bonds inside $F_1$, then $\delta\defeq G_r(\Lambda_F,A_1)-G_r(\Lambda_F,A_2)$ obeys $|\delta(x,y)|\le Ce^{-\delta_1(|x-y|+\dist(x,F_1^c)+\dist(y,F_1^c))}$ with $C$ independent of $F_1$, and the same for $D_A\delta$ when the bond at $x$ lies inside $F_1$; in particular $G_r(\Lambda_F,B_{\rm bg})=G_r(\Lambda_F,0)+\delta^{\rm bg}$.
\item For each flow $F\in\{X,\tau Y\}$ and each renormalised family $\cR$ at background $A^{(r)}$ on $\supp\chi_r$: $\cR^F(A^{(r)})=\sum_{q\le q_0}\sum_{w_1,\dots,w_q}A_r(w_1)\cdots A_r(w_q)\,\cR^F_{q;w}(A=0)\prod_vU^F(B(\Gamma_v))+\cR^F_{\delta G}+\cR^F_{\rm rem}$, where $\cR^F_{q;w}$ is a renormalised family in per-flow normal form whose lines are the $A=0$ lines of the flow ($G^F_r(0)=C-\Gamma^F_r(0)$, $H^F$, $(A^F)^{-1}$) and whose vertices are \cite[(3.18)--(3.26)]{K1} at $B=0$ with the explicit factors $\eta$ of the pin and contour vertices; the coarse-field factors $A_r(w_i)$ are the same functions of the common coarse field for both flows; $U^F(B(\Gamma_v))=\exp[\mathrm ie(5^r\eps)^{1/2}qB(\Gamma^F_v)]$ are the constant-background phases \cite[(3.46)]{K1}; $\cR^F_{\delta G}$ collects the terms with a line $\delta^{\rm im}$ or $\delta^{\rm bg}$ or with a vertex outside $F_v$, and $\cR^F_{\rm rem}$ the Taylor remainders of \cite[(3.48), (3.52)]{K1}.
\item $|\cR^F_{\delta G}|+|\cR^F_{\rm rem}|\le C(\delta_r^\sigma+e^{-\delta'p(\delta_r)})|T|$ per flow.
\item The two-flow difference of the leading sum is controlled element by element: the $A=0$ lines by \Cref{rot:cor:unif} and the difference elements $\Delta\Gamma$, $\Delta H$, $\Delta(A^{-1})$ of size $O(\eta'^{1/2})$; for a constant field $B$ the line integral depends only on the endpoints, $B(\Gamma^F_v)=B\cdot(c(v)-v)$, so $U^X(B(\Gamma_v))$ and $U^{\tau Y}(B(\Gamma_{\tau v}))$ differ by the phase of $e(5^r\eps)^{1/2}B\cdot(\tau v-v)$, $|\tau v-v|\le c_\tau\eta'$, i.e.\ by $\le ep_r\lmu^{-1}c_\tau\delta_r^{-1/2}\eta'$ (\cite[(3.2)]{K1}); the contour vertices \cite[(3.24)--(3.26)]{K1} at $B=0$ differ entrywise by a relative $\eta^{1/2}$ and the contour lines, one-ended and double, obey the slice-$j$ bounds in both flows with mismatch $(\eta'/\ell_j)^{1/2}$ (\Cref{rot:prop:BD}(b)); every connected subgraph containing a contour line has Ba{\l}aban's degree $\ge1$; the explicit factors $A_r(w_i)$ are common.
\end{enumerate}
\end{lemma}
\begin{proof}
(1) is King's expansion \cite[(3.43)--(3.55)]{K1} applied to each flow separately; it uses only the gauge covariance of the flow's lines (the hybrid flow is gauge covariant: its block average is the frame-covariant photon average of \Cref{rot:rem:frameav} and \Cref{rot:lem:tile}, so a constant $B$ is mapped to the coarse constant $R^{-1}B$ for the $Y$ flow and $B$ for the $X$ flow, and the gauge transformation $A^{(r)}\mapsto A^{(r)}-B$ acts on both by the same unitary at each vertex), the universal-cover form (0) of the localisation, and the Taylor expansion around $\vartheta\tilde A_{\rm lift}=0$ with background $B_{\rm bg}$, \cite[(3.47)--(3.54)]{K1}, exact algebra in which the lines become $A=0$ lines and the powers of $\tilde A$ become explicit coarse-field factors sampled at block points. The paired contours (\Cref{rot:lem:pair}) make the contour phases of the two flows differ by the flux through an area $O(\eta')$ (\Cref{rot:lem:flux}), so that the contour vertices and lines are difference elements with the gains of (3) rather than $O(1)$ mismatches. (0): the block systems of both flows, tiles and fixed staircases included, are $P_T$-invariant (\Cref{rot:lem:adm}(iii) and the in-tile convention of \Cref{rot:def:flows}), so the operator $M_{A_{\rm lift}}+a_rQ_F^*Q_F$ on $\Lambda_F$ commutes with $P_T$; its inverse has a summable kernel (the bounds of (2)), its periodisation is a right inverse of the torus operator, which is positive definite, and so equals its inverse; for $|x-y|\le s_T/2$ and $\gamma\ne0$, $|x-y-\gamma|\ge|\gamma|/2$, whence the image bound. (G): on $\Lambda_F$ the gauge function is single-valued and both flows are gauge covariant for every gauge function. (R): \cite[(3.2)]{K1} on all of $T$, $\sum_z\nabla H^{\rm ph}(x,z)=0$ with the $r$-uniform decay of $\nabla H^{\rm ph}$, and the product rule with $|\nabla\vartheta|\le3/p$. (N): the random-walk representation of \Cref{rot:prop:BD}(a) with one domain: chains whose cells lie with a fixed margin inside $F_1$ are the same for $A_1$ and $A_2$, every other chain of length $n$ has $\sqrt3M_{\rm w}n\ge\dist(x,F_1^c)+\dist(\supp f,F_1^c)-c$ and $\sqrt3M_{\rm w}n\ge|x-\supp f|-c$, and the average of the two gives the rate $\delta_1=(2\sqrt3M_{\rm w})^{-1}$. Vertices in $F_v$ see $B_{\rm bg}=0$ on every bond and in-block contour they touch, so King's ``taking $A'=\theta\tilde A$ and $B=0$'' \cite[p.~662]{K1} holds there literally. (2) is King's bound \cite[p.~665]{K1} with the kernel bounds at the regular fields of (R): a $\delta^{\rm bg}$ line with an end within $p$ of $D$ gains $e^{-\delta_1p/3}$, otherwise the string from an external vertex gives $e^{-\delta_2p}$; a vertex outside $F_v$ gives $e^{-\delta_2(4p/3-3)}$ by tree decay; a $\delta^{\rm im}$ line gives $e^{-\delta_0s_T/4}$; half of every exponent pays the vertex sums, whence $\delta'=\delta_2/6$, and $e^{-\delta'p(\delta_r)}\le C\delta_r^\sigma$ since $\mathrm p>1$; the Taylor remainders are $\le(Ce_rp_r^2)^{\bar n+1}\le C\delta_r^\sigma$. The universal-cover form replaces the periodic-cell reading of \cite[p.~337]{K2}: the periodised restriction of $\tilde A$ to a cell of side $\asymp p_r$ jumps by $\asymp p_r^2$ across the seam and is not regular, so the kernel bounds do not apply to it. (3) is as stated; the inversion symmetry about the coarse site used in the oddness rule requires the centred block convention for both flows (\Cref{rot:conv:C}(i)), and the frame-covariant average keeps it for the tiles.
\end{proof}

\begin{corollary}[the one-flow instance]\label{rot:cor:OF}
For one flow $F$ (standard or hybrid), a renormalised family of order $\le\bar n$ with lines from the slices $j=0,\dots,k$ of that flow (\cite[(3.76)]{K1}, each slice obeying \cite[(3.63)--(3.65)]{K1}; the coarse lines of depths $r$ and $r+n$ of the same flow included) and one element weakened by degree $\theta<\alpha_0$ obeys \cite[(1.33)]{Ba3} with an extra factor $C(1+k)^p$. Consequently King's estimate \cite[Prop.~3.6]{K1}, asserted on p.~665 without a printed argument, holds for both flows.
\end{corollary}
\begin{proof}
The multiscale bound on the point-split graph (the two summation points of a block product as vertices, the block kernel a line of index $r$): the slice lines satisfy the slice bounds with $d_l=-1$, a contour line of \cite[(3.26)--(3.27)]{K1} with $d_l=-1+\#(\text{contour ends})$ (\cite[(3.63)--(3.65), Prop.~3.7, Thm.~3.3]{K1} for the standard flow, \Cref{rot:prop:BD}(a) for the hybrid flow); the per-flow renormalised structure gives degree $\ge\alpha_0>0$ for every connected all-scale-summed subgraph except the trees \cite[(2.4)]{Ba3} and the counterterm-attachment subgraphs of \cite[(3.78)]{K1} with degree $0$, each of which costs a factor $(1+r)$ (King's ``possible addition of some power of $\ln(L^k\eps)^{-1}$'', p.~666, is $(1+k-r)^p$; both are $\le(1+k)^p$, \Cref{rot:lem:r}). The weakened line is the replacement $c_{l_0}\mapsto c_{l_0}\ell_j^{-\theta}$: the degree of every subgraph containing $l_0$ drops by $\theta<\alpha_0$ and stays positive, and no new zero-degree subgraph appears; a weakened vertex changes its exponent by $-\theta$; a weakened external field changes a norm only; the H\"older remainders of \cite[(3.83)]{K1} raise $d_l$ by $\alpha$ on the line they act on. This is King's own weakened-element induction \cite[(3.67)--(3.70)]{K1}, with no two-flow input; the strict positivity of the degrees after Ba{\l}aban's resummed rewriting is \cite[(3.5)]{Ba3}, and no cancellation for a fixed ordering of the lines is used.
\end{proof}

% r03_h1s_a: proof of the two-flow comparison bound (Theorem rot:thm:H1S), part 1: the all-C form, the slices, the elements, degrees per ordering, common moments, the current vertex.

\subsubsection*{Proof of the two-flow comparison bound}

The proof of \Cref{rot:thm:H1S} occupies the rest of this part. Its steps: the lines of every family are written as flow-independent lattice functions sliced by the heat kernel, plus coarse elements (this paragraph); degrees are read per ordering of the scale-summed elements, so that only the connected components of initial segments matter (\Cref{rot:h1s:MS}); each flow's expression is put in a normal form and one subtraction rule is applied, in each flow's own frame, to every component of degree $\le0$ that forms (\Cref{rot:h1s:rule}); after the rule every component containing the weakened element has degree $\ge\min(\alpha,1)$ (\Cref{rot:h1s:Dp}); every element type obeys a two-flow difference law with a positive exponent (\Cref{rot:h1s:P,rot:h1s:E,rot:h1s:U}); and the telescoping over the elements, together with the scale sums, gives the bound with the exponent $\theta_1$ (\Cref{rot:h1s:T} and the proof at the end).

\paragraph{The all-$C$ form and the slices.}
After the constant-background reduction (\Cref{rot:lem:CB}; at $e=0$ nothing is to be reduced) every line of a family is an $A=0$ line of its flow on the infinite lattice: $G^F_r=C^F-\Gamma^F$ (\Cref{rot:lem:Gam}(i)), where $C^F$ is the free propagator of the flow's lattice, $C^X(x,y)=C(x-y)$ and $C^Y(x'',y'')=C(R^{-1}(x''-y''))$ with \emph{one} lattice function $C=(-\Delta_{\eta'}+m_r^2)^{-1}$ on $\eta'\Z^3$ (for the photon, $D=(-\Delta_{\eta'}+\lmu^2\delta_r^2)^{-1}\otimes\mathbf 1_3$). Expanding every line of every graph in $C$ and $\Gamma$ gives \emph{$C/\Gamma$ monomials}, an exact identity; in each monomial the $C$-lines are the same lattice function evaluated in each flow's own frame, the vertices are matched by \Cref{rot:lem:vert} and the counterterm numbers are the same (\Cref{rot:lem:ctsame}), while $\Gamma^F$, $H^F$, $B^F=(A^F)^{-1}$ and the external fields depend on the flow. Write $C$ as the resolvent of the continuous-time simple random walk on $\eta'\Z^3$ with rate $\eta'^{-2}$ per direction, $C(u)=\eta'^{-3}\int_0^\infty e^{-m_r^2s}P_s(u)\,ds$, and slice it in the time variable:
\[
C^{(0)}\defeq\eta'^{-3}\!\int_0^{\ell_1^2},\qquad C^{(j)}\defeq\eta'^{-3}\!\int_{\ell_j^2}^{\ell_{j+1}^2}\ (1\le j\le r-1),\qquad C^{\rm top}\defeq\eta'^{-3}\!\int_1^\infty,
\]
all of $e^{-m_r^2s}P_s(u)\,ds$ ($\ell_r=1$). Gaussian and Poisson tail bounds for the walk give, for every $n\ge0$,
\begin{equation}\label{eq:h1s:B}
|\partial^nC^{(j)}(u)|\le c_n\ell_j^{-1-n}e^{-\delta|u|/\ell_j}\quad(0\le j\le r-1),\qquad|\partial^nC^{\rm top}(u)|\le c_ne^{-c'm_r|u|}(1+|u|)^{-1},
\end{equation}
$\partial$ a lattice difference; a derivative is charged to its line. The decomposition is the same for both flows as lattice functions; the $Y$ flow evaluates it at $R^{-1}(\tau x-\tau y)$, and in $X$-coordinates the $Y$-version obeys \eqref{eq:h1s:B} with $\delta/2$, since $|R^{-1}(\tau x-\tau y)|\ge|x-y|-2c_\tau\eta'$. The heat-kernel slices replace the block-spin slices of \Cref{rot:lem:slice} in the comparison of the all-$C$ lines; both decompositions are exact, and the heat-kernel one is flow-independent as a lattice function, which is what the subtraction rule below needs. (The gain $\eta'^{1/2}$ of \Cref{rot:cor:unif}(a) is recovered, and improved, by \Cref{rot:h1s:P}.)

\paragraph{Elements.}
The \emph{elements} of a monomial, with their per-flow bounds in the unit scale of level $r$, are: (a) the fine slice lines $C^{(j)}$, $D^{(j)}$, $0\le j\le r-1$, with their lattice derivatives, bounded by \eqref{eq:h1s:B}; (a$'$) the top pieces; (b) the coarse elements $\Gamma^F$, $H^F$, $B^F$ (\Cref{rot:lem:Gam}, \Cref{rot:lem:H}), bounded and decaying at the coarse mass, with one derivative and an $\alpha$-H\"older gradient (\cite[Thm.~3.3]{K1} for the standard flow, \Cref{rot:prop:BD}(a) for the hybrid flow); (c) the external fields $E^F\in\{H^F\phi_r,H^F_AA_r$, their covariant derivatives, the Cameron--Martin background$\}$, bounded by the field norms \cite[(3.2)]{K1} on $\supp\chi_r$ and $C^{1,\alpha}$ in the same sense; (d) the vertex functions $w$ (couplings, partitions of unity, the background), smooth with $\|\nabla w\|\le C$ (\Cref{rot:lem:vert}(ii)); (e) at $e>0$ the contour elements of \Cref{rot:prop:BD}(b): contour lines \cite[(3.27)]{K1}, double contour lines, contour vertices \cite[(3.24)--(3.26)]{K1} and the constant-background phases \cite[(3.46)]{K1}; (f) the counterterm numbers $\dm$, $E_1$, the same for both flows. Their matched differences, $X$-version at $x$ against $Y$-version at $\tau x$: (b) $|\Gamma^X-\Gamma^Y\circ\tau|\le C_\Gamma\eta^{1/2}$ (\Cref{rot:lem:Gam}(iii)), $|B^X-B^Y|\le0.2\eta^2$ (\Cref{rot:prop:fc}), coarse slices $\le C\eta^{1/2}\ell_j^{-1}e^{-\delta|x-y|/\ell_j}$ (\Cref{rot:cor:unif}(b)); (c) $|E^X(x)-E^Y(\tau x)|\le C\eta^{1/2}\|E\|$ (\Cref{rot:lem:H}, decaying form, applied to the lift); (d) $|w(x)-w(\tau x)|\le c_\tau\eta'\|\nabla w\|$; (e) contour lines: relative gain $(\eta'/\ell_j)^{1/2}$ on slice $j$, also with one lattice derivative; contour vertices: relative gain $\eta^{1/2}$ against their size $Ce_rp_r^2$; phases: $\le e\,p_r\lmu^{-1}c_\tau\delta_r^{-1/2}\eta'$ (\Cref{rot:prop:BD}(b)); the prefactor $p_r^2\delta_r^{-1/2}$ of the theorem comes from here and from nowhere else. Gradients of (b) and (c) are compared in \Cref{rot:h1s:E}.

\paragraph{The structures.}
With the vertices \cite[(3.18)--(3.23)]{K1} at $B=0$ (\Cref{rot:lem:CB}(1)), the connected all-$C$ line sets of non-positive degree that occur are the following (the list is proved in \Cref{rot:h1s:CL}). SE: the order-$e^2$ scalar self-energy, lines $\{C(x,y),D(x,y)\}$, two current vertices, both derivatives internal. VP: the order-$e^2$ vacuum polarisation, lines $\{C(x,y),C(x,y)\}$, two derivatives on the two lines, together with its seagull tadpole and derivative tadpole (\cite[(3.22)]{K1} with $m=2$), which enter only through the Ward sum of \Cref{rot:h1s:VP}. TAD: the $\lambda$-tadpole $C^{(j)}(0)$ and the seagull scalar tadpole $D^{(j)}(0)$, numbers, identical for both flows and cancelled exactly, slice by slice, by their counterterm shares (\cite[(3.33)]{K1}, \Cref{rot:lem:ctsame}). T: the order-$e^3$ triangle $\{C(a,b),C(b,c),D(a,c)\}$, three current vertices. SC: the order-$e^3$ seagull--current loop $\{C(s,b),D(s,b)\}$. W: the one-particle-irreducible two-point structures of order four: scalar (the $\lambda^2$ sunset; the $e^4$ rainbow, crossed, insertion and seagull graphs) and photon (the $e^4$ self-energy insertion, vertex correction, seagull graphs and the double seagull). L24: Ba{\l}aban's graph \cite[(2.4)]{Ba3}, a single scalar line between two bare current vertices carrying both their derivatives, $\partial\partial C$, with two external photon legs. CT: the counterterm share $c_\sigma$ of a self-energy or tadpole at slice assignment $\sigma$ (a one-site insertion, $|c_\sigma|\le C\ell_{\sigma_{\max}}^{-1}$, a loop line of dimension $-1$ at scale $\sigma_{\max}$), and the order-four share of dimension $0$. The $e^2\lambda$ two-point members are numbers or vanish identically; structures with an odd number of photon legs vanish by charge conjugation; the $(2,1)$-structure at order $e\lambda$ vanishes; four-leg loops have positive degree; contour vertices carry the explicit weight $\ge\frac12$ of \cite[(3.30)]{K1} at their site.

\paragraph{Degrees.}
For a connected set $H$ of elements (lines with their endpoints; a compound, defined in \Cref{rot:h1s:rule}, counts with the vertices and line dimensions it contains) put
\[
D(H)\defeq\sum_{l\in H}d_l+\sum_{v\in V(H)}a_v+3(|V(H)|-1),\qquad d_l=-1-(\text{number of lattice derivatives on }l),
\]
$a_v$ the vertex weight: an explicit factor $\eta'^a$, $a\ge0$, at $v$ counts as $\ell_{j(v)}^a$ with $j(v)\defeq\min\{j_l:l\ni v\}$; a derivative acting on an external field costs nothing, since the externals have bounded gradients. This is Ba{\l}aban's degree \cite[(2.1)--(2.3)]{Ba3} with every derivative charged to its line. The \emph{bare degree} $D_0(H)$ is the same count with the derivative placement of the normal form (NF2) below and before any derivative is transferred. All-$C$ degrees are integers; the only non-integer degrees are vertex weights $\alpha\in(0,1)$ produced by second-order subtractions with an external leg (\Cref{rot:h1s:U}(i)), so every degree is $n+m\alpha$ with $m\ge0$.

\begin{lemma}[the multiscale bound per ordering]\label{rot:h1s:MS}
Let $G$ be a connected set of elements of a monomial, each scale-summed element $l$ carrying the bound $c_l\ell_{j_l}^{d_l}e^{-\delta|x_l-y_l|/\ell_{j_l}}$ at its scale $j_l$ (fixed elements carry a bounded decaying kernel and impose no condition), and fix one ordering class of \cite[(2.7)]{Ba3}: an ordering $l(1),\dots,l(m)$ of the scale-summed elements with the scale sums restricted to $j_{l(1)}\le\dots\le j_{l(m)}$, ties assigned to one class. Suppose that every connected component of every initial segment $G_i\defeq\{l(1),\dots,l(i)\}$ has $D\ge0$, and $D\ge\theta$ if it contains a distinguished element $e_0$ whose constant carries the factor $(\eta'/\ell_{j(e_0)})^\theta$; let $n_0$ be the number of such components with $D=0$. Then the sum over the scales in the class and over all vertices but one is $\le C(1+r)^{n_0}\prod_lc_l$, and with the distinguished element $\le C\eta'^\theta(1+r)^{n_0}\prod_lc_l$, $C=C(G)$. Only the connected components of initial segments are tested.
\end{lemma}
\begin{proof}
This is the proof of \cite[Prop.~2.1]{Ba3} read for one ordering class, which is how Ba{\l}aban runs it: the finest element $l(1)$ is summed first, shrinking its line to a point with the exponent $D(G_1)$; the sum over $j_1\le j_2$ is geometric when $D(G_1)>0$ and costs a factor $(1+r)$ when $D(G_1)=0$; the merged vertex receives the weight $D(G_1)$, and the induction continues with $G_2,G_3,\dots$, at each stage testing only ``the zero-degree subgraphs met along the way'', the connected components of the initial segments. The same per-ordering positivity is King's \cite[p.~665]{K1}. With the distinguished element, every component containing it has exponent $D-\theta>0$ for the geometric sums, and the factor $\eta'^\theta\ell_{j(e_0)}^{-\theta}$ is absorbed by the merged vertex weight $\ell_{j'}^{D-\theta}\le\ell_{j(v)}^{D-\theta}$ at every later stage, so $\eta'^\theta$ survives to the end. A compound enters as a line of its stated dimension at the scale of its leg; its internal scale sums are done in \Cref{rot:h1s:U}(i) uniformly in the positions of its outer vertices, so the first step applies to it verbatim.
\end{proof}

\begin{lemma}[common moments]\label{rot:h1s:Z0}
Let $H(x,t)$ be a kernel built in a flow's own frame from $C$, $D$, their lattice derivatives, lattice Taylor remainders, lattice sums, the counterterm numbers and the bond shifts of \Cref{rot:h1s:cv} (no externals, no coarse elements). Then for $a=0,1$,
\[
\sum_{t\in\Lambda_X}H^Y(\tau x,\tau t)\,(\tau t-\tau x)^{\otimes a}=R^{\otimes a}\sum_{t\in\Lambda_X}H^X(x,t)\,(t-x)^{\otimes a},
\]
the indices compared frame-free (\Cref{rot:lem:vert}(iii)): zero moments are the same numbers and first moments the same vectors in both flows. An $\Oh$-invariant moment is isotropic: a rank-one or rank-three tensor vanishes, a rank-two one is $c\,\delta_{\mu\nu}$. On the tori the kernels are periodisations of the infinite-lattice ones and the moments agree up to $O(e^{-cm_r\cdot50MJ})$ (\Cref{rot:thm:V}(b)).
\end{lemma}
\begin{proof}
$H^Y$ is the conjugate of $H^X$ by the isometry $R$ (one lattice function evaluated in two frames) and $\tau:\Lambda_X\to\Lambda_Y$ is a bijection (\Cref{rot:def:match}), so the left side is $\sum_{v\in R\eta'\Z^3}H(R^{-1}v)v^{\otimes a}=R^{\otimes a}\sum_{u\in\eta'\Z^3}H(u)u^{\otimes a}$. The lattice model is $\Oh$-invariant (King's action on $\eps\Z^3$ with bond fields, in midpoint coordinates for bond kernels), so an $\Oh$-invariant tensor of odd rank vanishes under $-\mathbf 1\in\Oh$ and one of rank two is a multiple of $\delta_{\mu\nu}$.
\end{proof}

\begin{lemma}[the current vertex and the placement of its derivative]\label{rot:h1s:cv}
The current vertex \cite[(3.22)]{K1} ($m=1$) on the bond $b=(x,x+\eta'e_\mu)$ is $e_rA_b\,[\phi(b_+)\cdot q\phi(b_-)]/\eta'$ with $q$ antisymmetric. For two fields $a$, $b$ (a leg and an internal line after Wick contraction),
\[
(D_\mu a)(b)\cdot q\,b(b_-)=-a(b_+)\cdot q\,(D_\mu b)(b)+\eta'^{-1}\big[(a\cdot qb)(b_+)-(a\cdot qb)(b_-)\big],
\]
and summed against $A_b$ over the bonds the last term is $-\sum_x(\partial^*_\mu A_\mu)(x)\,(a\cdot qb)(x)$, the first term of \cite[(2.8)]{Ba3}; if the vertex function $g$ at the vertex is position dependent, the second term of \cite[(2.8)]{Ba3}, $\sum_x(a\cdot qb)(x)\sum_\mu(\partial^{\eta'}_\mu g)(x)A_\mu(x-\eta'e_\mu)$, appears as well. Hence the one lattice derivative of a current vertex may be written on either of its two scalar fields, the other being read at the shifted end of the bond, at the cost of a term in which the derivative acts on the vertex's photon field (and one with an explicit factor $\eta'$ from $\partial^{\eta'}g=\eta'\nabla g+O(\eta'^2)$). The identity is exact in each flow's own frame; a shifted evaluation point is matched by $\tau$ with $c_\tau$ replaced by $c_\tau+1$ in every matched-point bound. A current vertex has exactly two scalar fields and one photon field; the seagull $e_r^2|A|^2|\phi|^2$ and the quartic $\lambda_r|\phi|^4$ carry no derivative.
\end{lemma}
\begin{proof}
$a(b_+)=a(b_-)+\eta'(D_\mu a)(b)$ and $b(b_-)=b(b_+)-\eta'(D_\mu b)(b)$ give $a(b_+)\cdot qb(b_-)/\eta'=(D_\mu a)\cdot qb(b_-)+(a\cdot qb)(b_-)/\eta'=-a(b_+)\cdot q(D_\mu b)+(a\cdot qb)(b_+)/\eta'$; subtract, and sum by parts over the bonds $(x,x+\eta'e_\mu)$. For $a=b=\phi$ the remainder is $\phi\cdot q\phi/\eta'=0$, which is why the one-field form of the vertex has no such term; for two distinct fields it does not vanish, since the photon covariance is $\delta_{\mu\nu}$ times a scalar and not transverse.
\end{proof}

% r03_h1s_b: proof of the two-flow comparison bound, part 2: the sets of non-positive bare degree, the normal form, the subtraction rule.

\begin{lemma}[the connected all-$C$ sets of non-positive bare degree]\label{rot:h1s:CL}
Let $H$ be a connected set of all-$C$ lines with every current-vertex derivative placed on an internal line of its unit when the vertex belongs to a unit (the placement (NF2) below), and let $n_e$, $n_s$, $n_\lambda$ be the numbers of its current, seagull and quartic vertices, $n_{\rm ct}$ the number of its counterterm insertions, $E$ the number of its slots (legs of $V(H)$ not on lines of $H$, external fields included) and $x$ the number of current-vertex derivatives of $V(H)$ placed outside $H$. Then
\[
D_0(H)=\tfrac12n_e+n_s+n_\lambda+n_{\rm ct}+\tfrac12E-3+x .
\]
Decompose $H$ into its two-edge-connected blocks $P$ joined by $\rho-1$ bridges $b$. Then $D_0(H)=\sum_P[\omega(P)+\mathrm{out}(P)]+\sum_b(2-\mathrm{in}(b))\ge\rho-2$, where $\omega(P)$ is the degree of $P$ with every derivative of its vertices placed inside $P$, $\mathrm{out}(P)$ the number of derivatives of vertices of $P$ placed outside $P$ and $\mathrm{in}(b)\le2$ the number of derivatives on the bridge, so that $\sum_P\mathrm{out}(P)\ge\sum_b\mathrm{in}(b)$; and $\omega(P)\ge-1$ for every block with legs, with equality exactly for the units SE and VP (the tadpoles being dropped), for a counterterm insertion and, by convention, for a bare current vertex ($\omega=-1$, $\mathrm{out}=1$). Consequently $D_0(H)\le0$ forces one of: $\rho=1$ and $H$ is a unit ($-1$) or one of T, SC, W, the double seagull, a $\lambda^2$ or $e^4$ or $e^2\lambda$ two-point structure ($0$); $\rho=2$ and $H$ is two blocks with $\omega=-1$ and every non-internal derivative on the bridge: $U$--$l$--$U'$ ($l$ undifferentiated), $U$--$\partial l$--$v$, $v$--$\partial\partial l$--$v'$ ($=$ L24), and with a counterterm insertion ct--$l$--$U$, ct--$\partial l$--$v$, ct--$l$--ct; and, in the local-term monomials of the rule below, the same with a vertex of weight $-1$ (the vacuum-polarisation leftover of \Cref{rot:h1s:VP}) joined by an undifferentiated photon line to a vacuum-polarisation compound whose leg is that line, or to another such vertex. Unions of blocks at a cut vertex (a quartic or seagull vertex) have both slots at one vertex and are common numbers. Every member is a two-point structure except T and SC (three slots) and the one-point numbers.
\end{lemma}
\begin{proof}
Count: $\sum_ld_l=-I-n_\partial$ with $I$ the number of lines and $n_\partial$ the number of internal derivatives, $n_\partial=n_e-x$; the scalar legs give $2I_\phi+E_\phi=2n_e+2n_s+4n_\lambda+2n_{\rm ct}$ and the photon legs $2I_A+E_A=n_e+2n_s$; $|V|=n_e+n_s+n_\lambda+n_{\rm ct}$, and the counterterm insertions carry the vertex weights $\sum_va_v=-n_{\rm ct}$. Substituting in $-I-n_\partial+\sum_va_v+3(|V|-1)$ gives the formula. Blocks and bridges: $3(|V(H)|-1)=\sum_P3(|V(P)|-1)+3(\rho-1)$ because the block--bridge graph is a tree; a scalar bridge has $d_b=-1-\mathrm{in}(b)$ and a photon bridge carries no derivative in the placement (NF2); a derivative of a vertex of $P$ placed on a bridge counts $+1$ in $\mathrm{out}(P)$ and $-1$ on the bridge, and every derivative on a bridge is the own derivative of one of its end vertices, each of which lies in a block, whence $\sum_P\mathrm{out}(P)\ge\sum_b\mathrm{in}(b)$ and $D_0(H)\ge-\rho+2(\rho-1)=\rho-2$. The formula applied to a block with all derivatives internal gives $\omega(P)=\frac12n_e+n_s+n_\lambda+\frac12E-3$ for a block without insertions, and a counterterm insertion alone is a block with $E=2$ and $\omega=-1$; hence $\omega(P)\ge-1$ with the stated equality cases ($E=2$: $(n_e,n_s,n_\lambda)=(2,0,0)$; $E=3$: a single current vertex, whose derivative is necessarily outside), and the enumeration of the pairs $(P,P')$ with $\omega=-1$ and $\mathrm{in}(b)\le2$ gives the $\rho=2$ list; $\rho\ge3$ is impossible. The connected sets of bare degree $0$ with all derivatives internal are those with $\frac12n_e+n_s+n_\lambda+\frac12E=3$: $E=2$ with $(4,0,0)$, $(0,0,2)$, $(2,0,1)$, $(0,2,0)$, $(2,1,0)$, $(0,1,1)$; $E=3$ with $(3,0,0)$, $(1,1,0)$, $(1,0,1)$; $E=4$ with $(2,0,0)$, which is L24 (two current vertices joined by one scalar line carrying both derivatives; joined by a photon line both derivatives are outside and the degree is $2$). Integrality: $E\equiv n_e\pmod2$.
\end{proof}

\begin{corollary}[placement at the terminals; the derivative count]\label{rot:h1s:nd}
\textup{(i)} In a connected set with $D_0(H)\le0$ no vertex of $H$ places its own derivative outside $H$. \textup{(ii)} After the rule of \Cref{rot:h1s:rule} every line carries at most one lattice derivative at each end (its end vertex's own derivative or one transferred derivative, never both), hence at most two in all. The leg of a compound carries at most one as formed; when the compound is the differenced factor of (R2b) at one of its outer vertices, its leg carries the difference (\Cref{rot:h1s:U}(iv)) and then at most two; the middle line of a double compound carries none; a photon line carries only transferred derivatives, at most one per end; a contour line receives at most one, at its bond end.
\end{corollary}
\begin{proof}
(i) With $x\ge1$ the formula gives $D_0(H)\ge-1+x\ge0$, with equality only if the set with internal placement is a unit; but a unit's vertices have their derivatives inside by (NF2). (ii) Own derivatives: a current vertex has one, placed on one of its two scalar lines. Transferred derivatives land on a line $l$ at a vertex $w$ of the component $H$ being processed, $l\notin H$ later than $H$, $D(H)=0$ at formation, hence $D_0(H)\le0$ and, by (i), $w$'s own derivative is inside $H$: $l$ carries no own derivative of $w$. A second transfer onto $l$ at $w$ would come from a component $H'\ni w$ formed before $l$ enters the segment; $H'\supseteq H$ (nested initial segments), so $D(H')\ge D_0(H')+1\ge1$ unless $H'$ is a bare unit, which is never a component (\Cref{rot:h1s:Dp}(b)); so $H'$ is not processed. The leg of a compound is attached to the unit at a vertex whose own derivative is internal and at which no rule acts (the unit's lines are not elements before the compound enters, and afterwards they are inside it), so the leg carries at most one derivative, at its far end; the middle line of a double compound has two such ends. Photon lines carry no own derivative (\Cref{rot:h1s:cv}). For contour lines see (iii) of the rule.
\end{proof}

\begin{definition}[normal form and the subtraction rule]\label{rot:h1s:rule}
In each flow, in its own frame, before any element is mixed:
\begin{enumerate}[label=\textup{(NF\arabic*)},leftmargin=*]
\item Tadpoles and their counterterm shares cancel exactly per slice (TAD): they are dropped.
\item At every vertex of a unit (SE, VP, T, SC) whose derivative sits on a leg, the derivative is written on the internal scalar line by \Cref{rot:h1s:cv}; the remainder term, in which the derivative acts on the vertex's photon field, is the unit's own photon line, so it stays inside the unit (for SE the kernel is the sum of the pieces $\partial\partial C\cdot D$, $\partial C\cdot\partial D$, $C\cdot\partial\partial D$, each with $\sum_ld_l=-4$). After this a unit's leg carries at most the far vertex's derivative.
\item The seagull and derivative tadpoles of the order-$e^2$ vacuum polarisation are eliminated through the exact Ward sum of \Cref{rot:h1s:VP}; what remains at a vacuum-polarisation unit is the bubble for every slice assignment plus one common vertex weight.
\item The telescoping is done per ordering class of \cite[(2.7)]{Ba3}. The two classes that differ only by swapping the two internal lines of a vacuum-polarisation unit are merged before telescoping (identical element sets; \Cref{rot:h1s:MS} is applied per class in absolute value, a factor $2$ per such unit), so a vacuum-polarisation slice assignment is the unordered pair $\{j_1,j_2\}$ and its kernel the line-symmetrised one. The two Wick contractions of a self-energy between the same two legs lie in one class and are one element (their sum is the symmetrised kernel of \Cref{rot:h1s:M}).
\end{enumerate}
Fix an ordering of the scale-summed elements and process the initial segments $G_1\subset G_2\subset\dots$. Let $\sigma_{\max}(U)$ be the coarsest scale of a unit $U$.
\begin{enumerate}[label=\textup{(R\arabic*)},leftmargin=*]
\item\emph{Compounds.} A unit $U$ (SE, VP; order two) or a degree-zero one-loop piece (T, SC; order one) one of whose vertices $y$ has all its non-$U$ elements later than $\sigma_{\max}(U)$ (an external field counts as later) is bundled into the \emph{compound} $(U,F_y)$: the own-frame Taylor subtraction, about $U$'s other vertex $x$ (in leg-site coordinates, the vertex sites shifted by the bonds of the currents), of the product $F_y$ of the elements at $y$, taken on the earliest line $f$ of $F_y$ (or on the external if $F_y$ has no line), to second order for SE, VP and to first order for T, SC; one leg per unit, preferring a derivative-free, unclaimed leg. In the reduced ordering the compound is one element at the position of $f$ with the vertex set $V(U)\cup V(f)$ and the line dimensions of $U$ and $f$: $D(\{(U,f)\})=2+d_f\in\{1,0\}$ (SE, VP), $=3+d_f\ge1$ (T, SC); with an external leg $D=\alpha$ resp. $1$ (\Cref{rot:h1s:U}). Two units forced onto the same leg $l$ form the \emph{double compound} $(U_1,l,U_2)$, one element at $l$'s position with $D=D_0(U_1\cup l\cup U_2)=0$. The local terms of the subtraction are exact and taken in the flow's own frame: for SE the zeroth-order term $c_\sigma f(x)g(x)$ cancels the co-switched counterterm share of the same $\sigma$ (\cite[(3.33)]{K1} through order four; \Cref{rot:lem:ctsame}) and the first-order term vanishes by \Cref{rot:h1s:M}; for VP the zeroth-order term enters the Ward sum of \Cref{rot:h1s:VP} and the first-order term vanishes by \Cref{rot:h1s:M}; for T the local term is $f(x)\sum_zT''(x,z)g(z)$ with the two-point kernel $T''(x,z)\defeq\sum_yT(x;y,z)$ of dimension $-3$, a degree-zero line processed by (R2); for SC the local term vanishes or carries an explicit factor $\eta'$ (the bond shift). A unit none of whose vertices qualifies (both legs earlier than $\sigma_{\max}$, the \emph{split} unit) is not an element: its lines are ordinary elements and its counterterm share stands alone as a loop line of dimension $-1$ at scale $\sigma_{\max}$ in a separate monomial.
\item\emph{Components of degree zero.} Whenever a newly formed connected component $H$ of the current initial segment has $D(H)=0$ (degrees after the earlier transfers), apply the first available clause.
\begin{enumerate}[label=\textup{(\alph*)},leftmargin=*]
\item If $H$ contains a bare current vertex $w$, or a vertex $w$ of a split unit, whose derivative sits on a line of $H$, whose other scalar field is not in $H$ and whose photon field is an external field or a photon line not in $H$: rewrite the vertex by \Cref{rot:h1s:cv}. This gives two monomials, (i) with the derivative on the other scalar field and (ii) with $\partial^*$ on the photon field; in both $D(H)$ rises to $1$, the recipient (later than $H$, since $H$ is a component) changing $d\mapsto d-1$, or at no cost if it is an external field; the term with $\partial^{\eta'}g$ carries an explicit $\eta'$, a vertex weight $1$. Each vertex is rewritten at most once. If $w$'s photon field is a line inside $H$, the clause is not available (the move would leave $D(H)=0$): this happens exactly for the outer vertices of the two-point structures W and for the leg-carrying vertices of a triangle, which go to (b) and (R1).
\item Else, if $H$ is a two-point structure (exactly two vertices $x$, $t$ of $H$ carry elements outside $H$ or externals): the first-order subtraction at the terminal $t$ about $x$ (either terminal may be chosen),
\[
\textstyle\sum_tH(x,t)G(t)=c\,G(x)+\sum_tH(x,t)[G(t)-G(x)],\qquad c\defeq\sum_tH(x,t),
\]
$G$ the product of the later elements at $t$, Leibniz-expanded so that exactly one factor is differenced. $c$ is a common number of both flows (\Cref{rot:h1s:Z0}), bounded by \Cref{rot:h1s:MS} applied per flow to $H$ with one vertex fixed ($|\sum_\sigma c_\sigma|\le C(1+r)^{n_0+1}$); for the sunset and the scalar $e^4$ structures it is $0$ exactly with the co-switched order-four counterterm share, for the $e^4$ vacuum polarisation it is $c\,\delta_{\mu\nu}$, for a triangle kernel $T''$ it is $0$ exactly (the summand is odd under the lattice inversion through a site: the current on a bond goes to minus the current on the inverted bond, the propagators are even and the three lines are distinguishable), and for a double compound standing alone it is $0$ exactly (translation invariance of the outer unit's kernel). $D(H)$ rises to $1$ (the factor $|t-x|$ absorbed on a line of $H$); the differenced factor changes $d\mapsto d-1$ if it is a line (for a compound, its leg, \Cref{rot:h1s:U}(iv)), at no cost if it is an external, a coarse element or a vertex function; the $Y$-displacement $|\tau t-\tau x-(t-x)|\le2c_\tau\eta'$ against the $Y$-gradient is an explicit-$\eta'$ vertex weight, a degree-zero component with gain $\eta'\ge\eta'^\theta$, counted among the logarithms.
\item Else $H$ has at most one vertex with outside elements or externals and is a common number (\Cref{rot:h1s:Z0} with $a=0$; both flows evaluate the same closed lattice sum), bounded per flow by \Cref{rot:h1s:MS}: a vertex weight $0$, never a switched element.
\end{enumerate}
\item\emph{Local-term monomials.} The monomials produced by the zeroth-order terms of (R1) and (R2b) (the unit or structure collapsed to a vertex of weight $0$, or $-1$ for the vacuum-polarisation leftover of \Cref{rot:h1s:VP}) are processed by the same rules; their degree-zero components are L24-type lines, the weight classes of \Cref{rot:h1s:CL}, and, when a triangle inside a W is collapsed, the kernels $\{\partial C,D\}$ and $\{C,\partial C\}$ between the collapsed point and the other terminal, all handled by (R2b) with common numbers.
\end{enumerate}
\textup{(iii)} At $e>0$ the recipient of (R2a)(ii) may be a contour line \cite[(3.27)]{K1} whose bond end sits at $w$; it receives this one derivative at its bond end and no other, since its contour end is the site of a contour vertex, which carries the weight $\ge\frac12$ and so never lies in a degree-zero component. Contour vertices are never recipients of any clause, for the same reason, and the constant-background phases \cite[(3.46)]{K1} are not photon fields. The lifted external background at $w$ is an external field with the gradient bound $\le Cp_r$ (\Cref{rot:lem:CB}(0)(R)).
\end{definition}

The rule is Ba{\l}aban's own device for the graph (2.4), ``the first few $G_1,G_2^{(2)},\dots$ are graphs (2.4)\dots replaced by graphs with degree $+1$'' by the integration by parts \cite[(2.8)]{Ba3}, applied to every degree-zero component that forms in an initial segment, with composite lines (compounds) allowed and with the first-order subtraction (R2b) for the components without a usable bare vertex. Two features are forced by the two flows: the units are bundled with a leg as one element, so that no line of a unit is ever processed alone inside an isolated unit (processing the lone $\partial\partial C$ of a self-energy with both legs later would lose the gain, since the Taylor order at the slot would be fixed before the second unit line enters), and the local terms are common numbers of both flows, so that no first moment of a flow-dependent object is ever mixed between the frames.

% r03_h1s_c: proof of the two-flow comparison bound, part 3: degree bookkeeping per ordering after the rule; the enumeration.

\begin{lemma}[degree bookkeeping]\label{rot:h1s:Dp}
Fix an all-$C$ component $G$ of a monomial, a derivative placement, an ordering of its scale-summed elements, and apply (NF1)--(NF4) and (R1)--(R3). Then in the reduced ordering every connected component $H$ of every initial segment satisfies
\[
D(H)\ge1,\quad\text{or}\quad D(H)=n+m\alpha\ \ (n\ge0,\ m\ge1;\ H\text{ contains a compound on an external leg}),\quad\text{or}\quad H\text{ is an explicit-$\eta'$ component}
\]
(a $Y$-displacement term of (R2b) or a $\partial^{\eta'}g$ term of (R2a), with $D(H)=1-\theta$ after the gain $\eta'\ge\eta'^\theta$ is taken out), with at most $n_0\le\bar n$ zero-degree components in the per-flow count (local vertices of weight $0$, explicit-$\eta'$ components). In particular $D(H)-\theta\ge\min(\alpha,1)-\theta>0$ for $\theta\le\theta_1$ whenever $H$ contains the weakened element.
\end{lemma}
\begin{proof}
(a) Degrees are a property of the multilinear form and flow-independent: both flows have the same lines, the same placements and the same transfers, since (NF2) and the rule are applied in each flow's own frame to the same combinatorial objects.

(b) \emph{No component has $D<0$.} By \Cref{rot:h1s:CL} a connected line set with $D_0<0$ is a unit $U$. In the reduced ordering $U$'s lines are either inside a compound, and then $U$ alone is not an element, or $U$ is split, and then both legs $f$, $g$ of $U$ are earlier than $\sigma_{\max}(U)$: the first initial segment containing both lines of $U$ contains $f$ and $g$, attached at $U$'s vertices, so the component is $\{f,g,U,\dots\}$ with $D\ge d_f+d_g-4+9\ge1$ ($d_f,d_g\ge-2$ by \Cref{rot:h1s:nd}(i)); the segment containing only $U$'s first line, the lone $\partial\partial C$ of a split self-energy with $D=0$ when it precedes both legs, is a degree-zero component covered by (c). Transfers never create $D<0$: a transfer lands on an element later than its source component $H_0$ and attached to a vertex of $H_0$, so every later component containing the recipient contains $H_0$, where the $+1$ and the $-1$ cancel (this is \cite[(2.8)--(2.9)]{Ba3}: ``the degrees of these graphs are $\ge D(G_0)$''); the components containing $H_0$ but not the recipient rose by $1$; no other component changed. Loop closure after a transfer is harmless for the same reason: if the recipient line later closes a loop, the component $H\cup l$ has $D(H\cup l)\ge D_0(H\cup l)\ge-1$, with $-1$ only for a bare unit; then $H$ is the first line of a split unit and $l$ its other line, the coarsest, so both legs are earlier and attached when $l$ enters, and the component is $\{f,g,U,\dots\}$ with $D\ge1$ as above.

(c) \emph{Components with $D(H)=0$ at formation.} Since every transfer inside $H$ with recipient inside $H$ is neutral and one with recipient outside is $+1$, $D(H)\ge D_0(H)$, so $D_0(H)\in\{-1,0\}$; $-1$ is excluded by (b), so $H$'s line set is in the list of \Cref{rot:h1s:CL} (compounds expanded; for a derived structure the degree is that of its line set). Case by case:
\begin{itemize}[nosep,leftmargin=*]
\item $H=(U,\partial f)$ with the far vertex $w$ of $f$ bare, or $H=$ L24, or $H=\{\partial f,{\rm ct}_\sigma\}$ (a split form), or $H=(U,\partial f)$ with the derivative transferred onto $f$ earlier: a bare vertex $w\in H$ has its derivative on a line of $H$, its other scalar field outside $H$ and its photon field outside $H$ (external, or a photon line later than $H$); (R2a) raises $D$ to $1$. The lone $\partial\partial C$ of a split self-energy: (R2a) at a vertex of the split unit (its photon is the unit's $D$ line, outside the one-line component), or, equivalently, (R2b) with $c=\sum_y\partial\partial C(x,y)=0$ exactly (a lattice total derivative), the recipient being the leg at a terminal, which precedes the second unit line in the split regime; then $\{\partial\partial C^{+1},\nabla f\}$ has $D=2$ and $\{\partial\partial C^{+1},\nabla f,D\}$ has $D=1$.
\item $H=T''$ (the triangle's two-point local kernel): the terminals' photon field is the triangle's own $D$ line, inside $T''$, so (R2a) is not available; (R2b) applies with $c=0$ exactly (inversion oddness); $D\mapsto1$.
\item $H=$ W, scalar (outer current vertices bare, derivatives internal): an outer vertex's photon is a line inside W; (R2b) with $c=0$ exactly against the co-switched order-four counterterm share; $D\mapsto1$. $H=$ W, photon, or the double seagull, the sunset, the $e^2\lambda$ members: (R2b) with the common number $c$; $D\mapsto1$. When two triangles inside a W are both bundled (the crossed and vertex-correction graphs: each triangle's outer vertex carries an external leg or a later line, so each qualifies under (R1)), their line sets overlap and the two compounds merge into one element whose line set is all of W: a two-point degree-zero component when it stands alone, handled by (R2b). The collapsed local terms inside W, $\{\partial C,D\}$ and $\{C,\partial C\}$ between the collapsed point and the other terminal, are degree-zero two-point kernels of the SC and VP type: (R2b) with the common number; $D\mapsto1$.
\item Chains $U_1$--$l$--$U_2$, ct--$l$--$U$, ct--$l$--ct with $D=0$: each $U_i$ is bundled on the chain line $l$ or on its outer leg (if neither vertex of $U_i$ qualifies, its outer leg is earlier and attached, and then $H$ contains it with $D\ge1$); if $U_1$ bundles on $l$ and $U_2$ on its outer leg $g$, $\{(U_1,l),(U_2,g)\}$ has $D=(2+d_l)+(2+d_g)\ge1$; if both bundle on $l$, $H$ is the double compound element, which with two units is never alone (its outer legs are earlier) and in chains of three or more units can stand alone (the middle unit forced onto a leg already claimed by a neighbour while both outer legs of the double compound are later): then it is a two-point degree-zero component with no bare vertex and (R2b) applies, with $c=0$ exactly; $D\mapsto1$. A counterterm share stands alone only when its unit is split, so $\{f,g,{\rm ct}\}\ge5+d_f+d_g\ge1$, and $\{\partial f,{\rm ct}\}$ alone never forms since $g$ is earlier too.
\item $H$ one-point (both slots at one vertex: the seagull with two external photons and a hanging self-energy; the all-$C$ components of vacuum monomials cut at their $\Gamma$-lines): a common number, (R2c).
\item The weight classes of the local-term monomials, a weight-$(-1)$ vertex joined by an undifferentiated photon line to a vacuum-polarisation compound whose leg is that line, or to another such vertex ($-1+(-4-1)+3\cdot2=0$, $-1-1-1+3=0$): two-point structures whose outer vertices carry external or later photon fields; (R2b) with the common number; $D\mapsto1$.
\end{itemize}
Nothing else occurs: the list of \Cref{rot:h1s:CL} is exhausted (T and SC are bundled when a vertex qualifies, and otherwise never alone: $\{f_{\rm earlier},T\}$ has $D=d_f+3\ge1$).

(d) \emph{Components with $D=\alpha$.} A compound $(U,E)$ on an external leg alone; any component containing it and a line has $D\ge\alpha+1$.

(e) \emph{Termination and the count.} Each application of (R2) raises one freshly formed component and rewrites one vertex or differences one factor; a vertex is rewritten at most once (after (R2a) at $w$ every later component containing the recipient contains $H_0$ with $D\ge1$); (R2b) is applied once per formed two-point component; the local-term monomials have fewer vertices. The per-flow zero-degree components are the weight-$0$ local vertices and the explicit-$\eta'$ components, at most one per rule application, so $n_0\le\bar n$.
\end{proof}

\begin{remark}[the enumeration]\label{rot:h1s:census}
The case analysis of \Cref{rot:h1s:Dp}(c) is a finite statement for each graph, and an exhaustive enumeration, in exact integer arithmetic, confirms it for every graph with at most two units among the two- and three-point structures through order four (the isolated self-energy and vacuum polarisation, the tree (2.4), the triangle and the seagull--current loop with their leg currents, a unit between two currents, the rainbow, crossed, insertion, vertex-correction and seagull graphs, the chains $U$--$l$--$U$ and $U$--$\partial l$--$w$--$l'$--$U$ with bare and unit vertices, the collapsed local-term graphs and the weight classes): for every derivative placement allowed by \Cref{rot:h1s:cv} and every ordering of the scale-summed lines, walking the initial segments and applying the rule as stated (both monomials of (R2a) followed separately), no component of negative degree forms, every degree-zero component falls under one of the clauses, the rule terminates, and no line carries more than two lattice derivatives; the same holds for the chains of three and four units, the rings and the counterterm monomials, the double compound standing alone in a fraction of the orderings of the three-unit chain and being handled by (R2b). The proof for every family of order $\le\bar n$ is \Cref{rot:h1s:CL} together with (b)--(e), which use only the block arithmetic and the nesting of the initial segments. The statement for chains of any number of units is degree arithmetic: a third unit cannot join a set of bare degree $\le0$ (\Cref{rot:h1s:CL}, $\rho\le2$), and each further unit enters as a compound at a later leg or inside a component of degree $\ge1$.
\end{remark}

% r03_h1s_d: proof of the two-flow comparison bound, part 4: first moments of the units, the vacuum polarisation, the slice law, gradients of externals, the compound as a weakened element.

\begin{lemma}[the first moments of the units vanish]\label{rot:h1s:M}
In each flow's own frame: \textup{(i)} for every slice assignment $\sigma$ the self-energy kernel $k_\sigma$, as it enters any family (the sum of the two Wick contractions between the same two legs, in leg-site coordinates), has $F_\sigma\defeq\sum_uk_\sigma(u)\,u=0$; \textup{(ii)} for every unordered slice assignment $\sigma=\{j_1,j_2\}$ of the vacuum polarisation, the line-symmetrised kernel $\Pi_{\sigma,\mu\nu}\defeq\Pi_{(j_1,j_2),\mu\nu}+\Pi_{(j_2,j_1),\mu\nu}$ has $F_{\sigma,\mu\nu\rho}\defeq\sum_y\Pi_{\sigma,\mu\nu}(x,y)(y-x)_\rho=0$, in site and in bond-midpoint coordinates alike. (For a labelled assignment $(j_1,j_2)$ with $j_1\ne j_2$ the first moment is not zero; this is why (NF4) merges the two classes.)
\end{lemma}
\begin{proof}
Why it is needed: the compound's expansion $\sum_yk_\sigma(x,y)f(y)=c_\sigma f(x)+F_\sigma\cdot\nabla f(x)+\sum_yk_\sigma(x,y)T_xf(y)$ removes the first-order term. For a slice-line leg a term $F_\sigma\cdot\nabla f(x)$ would be harmless even if $F_\sigma\ne0$ ($F^Y=RF^X$ by \Cref{rot:h1s:Z0}). For an external or coarse leg it would not: $F^X\cdot\nabla E^X(x)$ and $F^X\cdot R^{-1}\nabla E^Y(\tau x)$ differ by $F^X\cdot(1-R^{-1})\nabla E(x)$ up to small terms, a frame term of order $|F|\|E\|$ with no gain in $\eta'$. (i) After (NF2) the self-energy at slice $\sigma$ is a bilinear form $B_\sigma(\phi,\phi)=\sum_{x',y'}\phi(x')\cdot\tilde k_\sigma(x',y')\phi(y')$ in the two leg fields, $x'$, $y'$ the leg sites (vertex sites shifted by the bonds of the currents), with $\tilde k_\sigma$ proportional to the identity in the $O(2)$ indices ($q^2=-1$; no $q$-part by charge conjugation). In every family the self-energy sits between two leg elements $g$ at $x'$ and $f$ at $y'$ through both Wick contractions, which have identical element sets; their sum is $\sum g(x')[\tilde k_\sigma(x',y')+\tilde k_\sigma(y',x')]f(y')$. The kernel $k_\sigma\defeq\tilde k_\sigma+\tilde k_\sigma^T$ is translation invariant and symmetric, so $\kappa_\sigma(u)\defeq k_\sigma(x',x'+u)$ is even and $F_\sigma=0$, whatever the rewriting of \Cref{rot:h1s:cv} did inside the kernel; the expansion is about the leg site, so no shifted-site moment appears. (ii) The bubble at labelled slices is $\Pi_{(j_1,j_2),\mu\nu}(x,y)\propto C^{(j_1)}(x+e_\mu,y+e_\nu)C^{(j_2)}(x,y)-C^{(j_1)}(x+e_\mu,y)C^{(j_2)}(x,y+e_\nu)$; under $-\mathbf 1\in\Oh$ the two factor positions swap slices, so the labelled kernel is $\Oh$-covariant only up to $j_1\leftrightarrow j_2$ and the symmetrised one is the $\Oh$-covariant object, a kernel on bonds covariant in midpoint coordinates, $\Pi_\sigma(Rm,Rm')=R\Pi_\sigma(m,m')R^T$ for $R\in\Oh$. Its midpoint first moment is an $\Oh$-invariant rank-three tensor, hence $0$; its zero moment is $c_\sigma\delta_{\mu\nu}$. In site coordinates $m'-m=(y-x)+(e_\nu-e_\mu)\eta'/2$, so $F_{\sigma,\mu\nu\rho}=0-c_\sigma\delta_{\mu\nu}(e_\nu-e_\mu)_\rho\eta'/2=0$.
\end{proof}

\begin{lemma}[the vacuum polarisation]\label{rot:h1s:VP}
\textup{(a)} With every line sliced, the order-$e^2$ photon two-point function at zero momentum vanishes in each flow, \cite[(2.23)--(2.28)]{Ba3} (valid with free boundary conditions) and \cite[\S3.5]{K1}: $\sum_{\sigma}c_{\sigma,\mu\nu}+\sum_jc^{\rm sg}_j\delta_{\mu\nu}+c^{\rm dt}_{\mu\nu}+c^{\rm top}_{\mu\nu}=0$, where $c_\sigma\defeq\sum_y\Pi_\sigma(x,y)$ is the bubble's zero moment, $c^{\rm sg}_j$ the seagull tadpole, $c^{\rm dt}$ the derivative tadpole and $c^{\rm top}$ the top pieces, all common numbers, all sums finite ($\sum_\sigma|c_\sigma|\le C\eta'^{-1}$ at fixed $r$). \textup{(b)} The seagull and derivative tadpoles are single-vertex structures of degree $-1$ and may not appear as elements; (a) replaces them by $-\sum_\sigma c_\sigma-c^{\rm top}$. With the local terms $+c_\sigma f(x)g(x)$ of the bundled bubbles ((R1); $f$, $g$ the photon legs) the complete local content at a vacuum-polarisation vertex is $[-\sum_{\sigma\ \rm unbundled}c_\sigma-c^{\rm top}]f(x)g(x)$, a slice assignment being unbundled iff $\sigma_{\max}>J\defeq\max(j_f,j_g)$, plus the remainders of the bundled bubbles and the unbundled bubbles themselves, which are never alone ($\{f,\Pi_\sigma\}$ has $D=1$). Since $|c_\sigma|\le C\ell_{\sigma_{\min}}/\ell_{\sigma_{\max}}^2$ (placement $(1,1)$) and $\le C\ell_{\sigma_{\min}}^2/\ell_{\sigma_{\max}}^3$ (placement $(2,0)$, using $\sum_u\partial\partial C^{(j_1)}(u)=0$),
\[
\Big|\sum_{\sigma_{\max}>J}c_\sigma\Big|\le C\sum_{j_2>J}\ell_{j_2}^{-2}\sum_{j_1\le j_2}\ell_{j_1}\le C\ell_J^{-1},\qquad|c^{\rm top}|\le C:
\]
a vertex weight $-1$ at scale $J\ge j(x)$ plus a weight $0$. Photon legs carry no own derivative, so $\{f,\text{weight}\}$ has $D=1$ and $\{f,g,\text{weight}\}$ has $D=3$. The Ward identity is used once, per flow, for the sum over all $\sigma$; its fine end, $\sum_{\sigma_{\max}\le J}c_\sigma$, of order $\eta'^{-1}$, is exactly what the bundled remainders carry, and only the coarse end survives as a weight. With external photon legs the leftover is $0$ exactly. \textup{(c)} The order-$e^4$ vacuum polarisation, and every degree-zero two-point structure handled by (R2b), needs no identity: $c_{\mu\nu}\defeq\sum_t\Pi^{(4)}_W(x,t)$ is a common number (\Cref{rot:h1s:Z0}), isotropic, and $|\sum_\sigma c_\sigma|\le C(1+r)^{n_0+1}$ by \Cref{rot:h1s:MS} applied per flow to W with one vertex fixed (all proper components of W are already of degree $\ge0$ after the rule; $D(W)=0$ gives one logarithm). The local term $c\,A^F(x)\cdot A^F(x)$ is a weight-$0$ vertex term common to both flows; its difference is telescoped over the externals. \textup{(d)} In the scalar sector the bundled self-energy's zeroth-order term cancels against ${\rm ct}_\sigma$ exactly; the sunset and the scalar $e^4$ structures give $c=0$ exactly with the co-switched order-four share.
\end{lemma}

\begin{lemma}[the slice law]\label{rot:h1s:P}
Let $g_j\defeq C^{(j)}-C^{(j)}_c$ with $C^{(j)}_c(v)\defeq\int_{\ell_j^2}^{\ell_{j+1}^2}e^{-m_r^2s}(4\pi s)^{-3/2}e^{-|v|^2/4s}ds$ the isotropic continuum slice, $3\le j\le r-1$. Then for every $n\ge0$ lattice derivatives
\[
|\partial^ng_j(u)|\le C_n(\eta'/\ell_j)^{1/2^n}\ell_j^{-1-n}e^{-\delta_n|u|/\ell_j},
\]
and consequently, for matched points $|u'-R^{-1}u|\le4c_\tau\eta'$ (bond shifts included),
\[
|\partial^nC^{(j)}(u)-R^{\otimes n}\partial^nC^{(j)}(u')|\le C(\eta'/\ell_j)^{1/2^n}\ell_j^{-1-n}e^{-\delta_n|u|/(2\ell_j)},
\]
the same for the photon slices and, with $\ell_j$ replaced by $1$ and the mass decay, for the top piece. For $j\le2$ there is no gain and the factor $(\eta'/\ell_j)^\theta$ is paid by a degree loss $\theta$ (\Cref{rot:cor:unif}(a), King's device at the finest slices \cite[p.~675]{K1}). Constants are uniform in $m_r\in(0,1]$ and $r$.
\end{lemma}
\begin{proof}
Value clause $n=0$: the coordinates of the continuous-time walk are independent one-dimensional walks \cite[Prop.~1.2.2]{LL10}, and the local limit law \cite[Thm.~2.5.6]{LL10} gives, for $|u|\le s/(8C_0\eta')$ and $s\ge\ell_3^2$, $|\eta'^{-3}P_s(u)-(4\pi s)^{-3/2}e^{-|u|^2/4s}|\le c(\eta'/\sqrt s)(4\pi s)^{-3/2}e^{-|u|^2/8s}$; integrated over $s\in[\ell_j^2,\ell_{j+1}^2]$ this is $\le c(\eta'/\ell_j)\ell_j^{-1}e^{-\delta|u|/\ell_j}$, and beyond that range both kernels are bounded by their Gaussian and Poisson tails, $\le C_N(\eta'/\ell_j)^Ne^{-\delta|u|/\ell_j}$; the geometric mean with the decay of each kernel halves the rate and keeps the factor. Derivatives by interpolation: for a lattice function $h$ and $\rho=N\eta'$, $\partial_\mu h(u)=[h(u+\rho e_\mu)-h(u)]/\rho-\rho^{-1}\sum_{i=1}^{N-1}(\rho-i\eta')\partial_\mu\partial_\mu h(u+i\eta'e_\mu)\eta'$, so $|\partial_\mu h(u)|\le2\sup|h|/\rho+\rho\sup|\partial\partial h|$ over the segment. With $h=\partial^ng_j$, the crude bound $|\partial^{n+2}g_j|\le|\partial^{n+2}C^{(j)}|+|\partial^{n+2}C_c^{(j)}|\le C\ell_j^{-3-n}e^{-\delta|u|/\ell_j}$ from \eqref{eq:h1s:B} and the smoothness of $C_c$, and $\rho\defeq\ell_j\epsilon_n^{1/2}$ with $\epsilon_n\defeq C_n(\eta'/\ell_j)^{1/2^n}$ (so $\rho\ge\eta'$; round to a multiple of $\eta'$), one gets $|\partial^{n+1}g_j|\le C\epsilon_n^{1/2}\ell_j^{-2-n}$, the decay rate kept since the segment has length $\le\ell_j$; induction on $n$. Matched points: $\partial^nC^{(j)}(u)-R^{\otimes n}\partial^nC^{(j)}(u')=\partial^ng_j(u)+[\partial^nC_c^{(j)}(u)-R^{\otimes n}\partial^nC^{(j)}_c(u')]-R^{\otimes n}\partial^ng_j(u')$; the middle term is the rotation covariance of the isotropic $C_c$, $\nabla^nC_c(Rv)=R^{\otimes n}\nabla^nC_c(v)$, plus the displacement $|R^{-1}u-u'|\le4c_\tau\eta'$ against $|\nabla^{n+1}C_c|\le C\ell_j^{-2-n}$, and the lattice-difference-versus-derivative error of $C_c$, $\le C\eta'\ell_j^{-2-n}$; both are inside the claim. (The true law is $(\eta'/\ell_j)^2$ for the values and $\eta'/\ell_j$ for first differences; the exponents $1/2^n$ are far from sharp and suffice, since every gain is compared with $\theta_1\le\frac1{16}$.)
\end{proof}

\begin{lemma}[gradients of externals and coarse elements]\label{rot:h1s:E}
For $E\in\{H^F\phi_r,H^F_AA_r$, their covariant derivatives, the background$\}$ and for $\Gamma^F(\cdot,z)$, $H^F(\cdot,z)$: $|E^X(x)-E^Y(\tau x)|\le C\eta^{1/2}\|E\|$ (\Cref{rot:lem:H}, \Cref{rot:lem:Gam}(iii)), and both are restrictions of $C^{1,\alpha}$ functions $g^X$, $g^Y$ on $\R^3$ with $\|\nabla g^F\|_\infty\le C_\nabla\|E\|$ and $[\nabla g^F]_\alpha\le C\|E\|$ (\cite[Thm.~3.3, (3.8)]{K1} for the standard flow, \Cref{rot:prop:BD}(a) for the hybrid flow; for the lifts, the Fourier representation of \Cref{rot:lem:H} with smooth symbols). Then, with $\kappa\defeq\alpha/(2(1+\alpha))$, $|\nabla g^X(x)-\nabla g^Y(x)|\le C\eta^\kappa\|E\|$ for all $x\in\R^3$, in particular at $x$ and at $\tau x$. Consequently the second-order own-frame Taylor remainders $T^F_xE^F(y)\defeq E^F(y)-E^F(x)-(y-x)\cdot\nabla g^F(x)$ (in the $Y$ frame at $\tau x$, $\tau y$) differ by
\[
|T^X_xE^X(y)-T^Y_{\tau x}E^Y(\tau y)|\le C\|E\|\min\{|u|^{1+\alpha}+\eta'^{1+\alpha},\ \eta^{1/2}+|u|\eta^\kappa+\eta'\},\qquad u=y-x .
\]
For a fine slice line $f=\partial^aC^{(j_f)}(\cdot-z)$ in place of $E$ the same holds with $\eta^{1/2}$, $\eta^\kappa$, $\eta'$ replaced by $(\eta'/\ell_{j_f})^{1/2^a}\ell_{j_f}^{-1-a}$, $(\eta'/\ell_{j_f})^{1/2^{a+1}}\ell_{j_f}^{-2-a}$, $\eta'\ell_{j_f}^{-2-a}$, times $e^{-\delta\,{\rm dist}/(2\ell_{j_f})}$, and with the remainder bound $c|u|^{1+\alpha}\ell_{j_f}^{-2-a-\alpha}$ (\Cref{rot:h1s:P}; no interpolation needed).
\end{lemma}
\begin{proof}
$h\defeq g^X-g^Y$ has $\|h\|_\infty\le C\eta^{1/2}\|E\|$ (at lattice points by the value difference and $|g^Y(\tau x)-g^Y(x)|\le C_\nabla c_\tau\eta'\|E\|$, between them by the Lipschitz bound) and $[\nabla h]_\alpha\le2C\|E\|$; for a unit vector $e$ and $\rho>0$, $|e\cdot\nabla h(x)|\le|h(x+\rho e)-h(x)|/\rho+[\nabla h]_\alpha\rho^\alpha\le2C\eta^{1/2}\|E\|/\rho+2C\|E\|\rho^\alpha$; choose $\rho=\eta^{1/(2(1+\alpha))}$. The first remainder bound is the sum of the two remainders ($|\tau y-\tau x|\le|u|+2c_\tau\eta'$); the second is the value law at $y$ and $x$, the gradient law at $x$ and the displacement $|(\tau y-\tau x)-(y-x)|\le2c_\tau\eta'$ against $\|\nabla g^Y\|$.
\end{proof}

\begin{lemma}[the compound as a weakened element]\label{rot:h1s:U}
Let $B=(U,f)$ be a compound of (R1): $U$ a self-energy or vacuum polarisation with slice assignment $\sigma$ between the leg sites $x$ and $y$, placement-summed and symmetrised kernel $k_\sigma(x,y)=\kappa_\sigma(x-y)$ with $|\kappa_\sigma(u)|\le C\ell_{j_1}^{d_1}\ell_{j_2}^{d_2}e^{-\delta|u|/\ell_{\sigma_{\min}}}$, $d_1+d_2=-4$ (the product of two slices decays at the finest line); $f=\partial^aC^{(j_f)}(y-w)$, $a\in\{0,1,2\}$ ($a\le1$ for a compound as formed, \Cref{rot:h1s:nd}(ii); $a=2$ when the compound is the differenced factor of (R2b), see (iv)), $\sigma_{\max}\le j_f$; $K^F_\sigma$ the own-frame second-order remainder sums $K^X_\sigma(x,w)\defeq\sum_{y}\kappa_\sigma(x-y)T^X_xf^X(y)$, $K^Y_\sigma(\tau x,\tau w)\defeq\sum_{y'\in\Lambda_Y}\kappa_\sigma(R^{-1}(\tau x-y'))T^Y_{\tau x}f^Y(y')$, with lattice differences in the Taylor polynomials. Put $s\defeq\alpha/(1+\alpha)$, $\kappa\defeq\alpha/(2(1+\alpha))$.
\begin{enumerate}[label=\textup{(\roman*)},nosep,leftmargin=*]
\item $|\sum_{\sigma:\sigma_{\max}\le j_f}K_\sigma(x,w)|\le C\ell_{j_f}^{-2-a}e^{-\delta|x-w|/(4\ell_{j_f})}$ per flow: the compound is a line of dimension $-2-a$ at scale $j_f$. With an external or coarse leg $E$: $|\sum_{\sigma_{\max}\le j}K_\sigma(x)|\le C\ell_j^\alpha\|E\|$, a vertex weight $\alpha$ at scale $\sigma_{\max}$, summed geometrically.
\item $\sum_{\sigma_{\max}\le j_f}|K^X_\sigma(x,w)-K^Y_\sigma(\tau x,\tau w)|\le C(1+r)\eta'^{\theta_U}\ell_{j_f}^{-2-a-\theta_U}e^{-\delta|x-w|/(4\ell_{j_f})}$ with $\theta_U\defeq\min\{\theta,s/2^{a+1}\}$; with an external or coarse leg $\le C(1+r)\eta'^{\theta_U}\|E\|$ with $\theta_U\defeq\min\{\theta,\kappa s\}$. No frame vector appears.
\item For the double compound $(U_1,l,U_2)$ with $l=\partial^aC^{(j_l)}$, $a\in\{0,1\}$: per flow a line of dimension $-3-a$ at scale $j_l$; the difference gains $(\eta'/\ell_{j_l})^{s/8}$, with one logarithm per unit. The same holds when $l$ is a coarse element $\Gamma^F$: per flow $|\sum_{\sigma_1,\sigma_2}(U_1,\Gamma,U_2)(x,w)|\le Ce^{-m_r|x-w|/2}$ (dimension $0$ at the top scale), and the difference is $\le C(1+r)^2\eta^{\kappa s}e^{-m_r|x-w|/4}$.
\item First-order forms (the compounds $(T,f)$, $(SC,f)$ and the brackets of (R2b) whose differenced factor is a line carrying $a\le2$ derivatives): gain $(\eta'/\ell_{j_f})^{1/2^{a+2}}\ge(\eta'/\ell_{j_f})^{1/16}$ (slice leg) resp.\ $\eta^{1/4}$ (external leg), with one logarithm. When the differenced factor of (R2b) is a compound $(U,f)$, the bracket is bounded by (i)--(ii) with $a+1\le2$ in place of $a$: per flow $|t-x|$ times a line of dimension $-3-a$ at scale $j_f$, and the gain $\eta'^{\theta_U}$ with $\theta_U\ge\min\{\theta,s/8\}$.
\end{enumerate}
\end{lemma}
\begin{proof}
(i) $\sum_u\eta'^3|\kappa_\sigma(u)||u|^{1+\alpha}\le C\ell_{j_1}^{d_1}\ell_{j_2}^{d_2}\ell_{\sigma_{\min}}^{4+\alpha}\le C\ell_{\sigma_{\max}}^\alpha$ (the sum is dominated by $|u|\sim\ell_{\sigma_{\min}}$, where the coarser line is bounded by $\ell_{\sigma_{\max}}^{d_{\max}}$; more precisely it is $\le C\ell_{\sigma_{\min}}^{\alpha}(\ell_{\sigma_{\min}}/\ell_{\sigma_{\max}})^{|d_{\max}|}$, geometric in $\sigma_{\min}$ at fixed $\sigma_{\max}$); $|T_xf(y)|\le C\ell_{j_f}^{-2-a-\alpha}|u|^{1+\alpha}$ with the leg's decay $e^{-\delta|y-w|/(2\ell_{j_f})}\le e^{-\delta|x-w|/(4\ell_{j_f})}e^{\delta|u|/(4\ell_{j_f})}$, the last factor absorbed by the unit's decay at $\ell_{\sigma_{\min}}\le\ell_{j_f}$; $\sum_{\sigma_{\max}\le j_f}\sum_{\sigma_{\min}\le\sigma_{\max}}\ell_{\sigma_{\min}}^{\alpha}(\ell_{\sigma_{\min}}/\ell_{\sigma_{\max}})^{|d_{\max}|}\le C\ell_{j_f}^\alpha$. External leg: the same with $\|E\|$. (ii) Split $K^X_\sigma-K^Y_\sigma=\sum_y[\kappa_\sigma(x-y)-\kappa_\sigma(R^{-1}(\tau x-\tau y))]T^X_xf^X(y)+\sum_y\kappa_\sigma(R^{-1}(\tau x-\tau y))[T^X_xf^X(y)-T^Y_{\tau x}f^Y(\tau y)]=:A_\sigma+B_\sigma$. $A_\sigma$: telescoping the product of the unit's lines with \Cref{rot:h1s:P} ($n\le2$) for the differenced line and \eqref{eq:h1s:B} for the other gives $|\Delta\kappa_\sigma|\le C(\eta'/\ell_{\sigma_{\min}})^{1/4}\ell_{j_1}^{d_1}\ell_{j_2}^{d_2}e^{-\delta|u|/(2\ell_{\sigma_{\min}})}$ (with the $125^\theta$-form for $\sigma_{\min}\le2$), so as in (i) $\sum_{\sigma_{\max}\le j_f}A_\sigma\le C\eta'^\theta\ell_{j_f}^{\alpha-\theta}\ell_{j_f}^{-2-a-\alpha}$ for $\theta\le\frac14$, $\theta<\alpha$. $B_\sigma$: the bracket is $\le\min\{P,Q\}$ with $P\defeq C\ell_{j_f}^{-2-a-\alpha}(|u|^{1+\alpha}+\eta'^{1+\alpha})$ and, by \Cref{rot:h1s:E} for a slice leg, $Q\defeq b_0+b_1|u|+b_2$ with $b_0=C(\eta'/\ell_f)^{1/2^a}\ell_f^{-1-a}$, $b_1=C(\eta'/\ell_f)^{1/2^{a+1}}\ell_f^{-2-a}$, $b_2=C\eta'\ell_f^{-2-a}$ ($\ell_f\defeq\ell_{j_f}$). Since $(1+\alpha)(1-s)=1$, $\min\{P,Q\}\le P^{1-s}Q^s\le C\gamma_f\ell_f^{-2-a}(|u|+\eta')(1+|u|/\ell_f)^s$ with $\gamma_f\defeq(\eta'/\ell_f)^{s/2^{a+1}}$ (the $\ell_f$-exponent is $-(2+a+\alpha)(1-s)-(1+a)s=-2-a$); then $B_\sigma\le C\gamma_f\ell_f^{-2-a}e^{\dots}\sum_u\eta'^3|\kappa_\sigma(u)|(|u|+\eta')$, and the first-order sum of a degree-$(-1)$ kernel is $\le C$ per $\sigma$, hence $\le C(1+r)$ after the $\sigma$-sum: one logarithm. External leg: $b_0=C\eta^{1/2}\|E\|$, $b_1=C\eta^\kappa\|E\|$, $b_2=C\eta'\|E\|$, $\gamma_f=\eta^{\kappa s}\le5\eta'^{\kappa s}$. (iii) Slice $l$: $|T_xT_wl(y,z)|\le\ell_l^{-3-a-2\alpha}|u_1|^{1+\alpha}|u_2|^{1+\alpha}$ (mixed second-order lattice remainder, from \eqref{eq:h1s:B} at order $a+4$) gives, summed, $\ell_l^{-3-a}$ per flow. Difference: $P=\ell_l^{-3-a-2\alpha}|u_1|^{1+\alpha}|u_2|^{1+\alpha}$ and $Q\le C(\eta'/\ell_l)^{1/8}\ell_l^{-1-a}(1+|u_1|/\ell_l)(1+|u_2|/\ell_l)$, since the differenced factors $l,\nabla l,\nabla\nabla l$ carry $n\le a+2\le3$ derivatives (\Cref{rot:h1s:P}, exponent $\frac18$) plus the displacements; $P^{1-s}Q^s\le C(\eta'/\ell_l)^{s/8}\ell_l^{-3-a}|u_1||u_2|(\dots)^s$ as $(3+a+2\alpha)(1-s)+(1+a)s=3+a$; two first-order sums give $(1+r)^2$. Coarse $l=\Gamma^F=H^FB^FH^{F*}$ (\Cref{rot:lem:Gam}): $T_xT_w$ acts on the two $H$-factors separately, each $C^{1,\alpha}$ in its first variable with $|T_xH^X(\cdot,p)(y)|\le C|y-x|^{1+\alpha}e^{-\delta|x-p|}$, and $|B^F(p,q)|\le Ce^{-\delta|p-q|}$, so $|T_xT_w\Gamma^F(y,z)|\le C|u_1|^{1+\alpha}|u_2|^{1+\alpha}e^{-\delta|x-w|/2}$; the difference uses $|\Gamma^X-\Gamma^Y\circ\tau|\le C\eta^{1/2}$ and \Cref{rot:h1s:E} for each $H$-factor, $Q\le C\eta^\kappa(1+|u_1|)(1+|u_2|)$, whence $P^{1-s}Q^s\le C\eta^{\kappa s}|u_1||u_2|$. (iv) First-order forms: $P=\ell_f^{-2-a}|u|$ (slice; $C\|E\||u|$ external), $Q$ as above with $n\le a+1$ in the gradient; $P^{1/2}Q^{1/2}\le C(\eta'/\ell_f)^{1/2^{a+2}}\ell_f^{-1-a}|u|^{1/2}(\dots)$ resp. $C\eta^{1/4}\|E\||u|^{1/2}(\dots)$, and the first-order sum converges per $\sigma$ with the $(1+r)$ after the $\sigma$-sum. For the brackets of (R2b) the switched factor $\Delta[g(t)-g(x)]$ is $\le|t-x|$ times the slice law for $\nabla g$ plus $2c_\tau\eta'\sup|\nabla g^Y|$: the degree-one term with its gain and the explicit-$\eta'$ term (degree zero, one logarithm). When the differenced factor is a compound $(U,f)$ at an outer vertex $t\in\{x,w\}$, its lattice difference is the compound with the differenced leg: at $t=w$ directly, and at the unit vertex $t=x$ by the translation invariance of $\kappa_\sigma$, since shifting the summation variable gives $K_\sigma(x+\eta'e_\mu,w)-K_\sigma(x,w)=\eta'\sum_y\kappa_\sigma(x-y)\,T_x(\partial_\mu f)(y)$ exactly, the lattice differences commuting (and likewise on $\Lambda_Y$ with the step $\eta'Re_\mu$). So the difference of the bracket is a sum of $|t-x|/\eta'$ terms $\eta'\,[K^X_\sigma-K^Y_\sigma]$ with the leg $\partial_\mu f$, $a+1\le2$, each bounded by (ii) (\Cref{rot:h1s:E} for the slice leg with $a+1$, \Cref{rot:h1s:P} at $n\le3$; the unit lines keep $n\le2$): per flow $|t-x|\,\ell_{j_f}^{-3-a}$, gain $\eta'^{\theta_U}$ with $\theta_U=\min\{\theta,s/2^{a+2}\}\ge\min\{\theta,s/8\}$, one logarithm; the $Y$-displacement term is explicit-$\eta'$ as before.
\end{proof}

% r03_h1s_e: proof of the two-flow comparison bound, part 5: the telescoping, every term is small, the exponent theta_1, the sectors, the proof of the theorem.

\paragraph{The telescoping.}
After (NF1)--(NF4) and (R1)--(R3) in each flow's own frame, $E^F(\cR)$ for one ordering class is a finite sum of multilinear expressions in the elements: slice lines and top pieces (with at most two lattice derivatives, \Cref{rot:h1s:nd}), compounds $(U,f)$ and $(U_1,l,U_2)$ (own-frame remainder sums, \Cref{rot:h1s:U}), first-order compounds $(T,f)$, $(SC,f)$, the brackets $[G(t)-G(x)]$ of (R2b) with one differenced factor, the coarse elements $\Gamma^F$, $H^F$, $B^F$, the externals and vertex functions at $\tau$-matched points (shifts included), the contour elements at $e>0$, and the common numbers (the zero moments $c$, the order-four shares, the vacuum-polarisation leftover weight, the one-point structures), which are never switched. The two flows' expressions have the same combinatorics. Order the elements consistently (compounds first, then lines, coarse elements, externals, contour elements) and telescope:
\[
E^X(\cR)-E^Y(\cR)=\sum_iT_i,\qquad T_i\defeq\big[e_1,\dots,e_{i-1}\text{ in the $Y\circ\tau$ version};\ e_i\to e_i^X-e_i^Y\circ\tau;\ e_{i+1},\dots\text{ in the $X$ version}\big],
\]
an identity, every term being the same multilinear form with elements substituted, the co-switched pairs (a self-energy and its counterterm share, the two Wick contractions of \Cref{rot:h1s:M}, a scalar W and its order-four share) switched together; at most $n_\cR$ terms per class, and finitely many classes per family. Compounds are switched first so that a unit and its leg are compared as a whole and no mixed-frame first moment is ever written; externals are switched late so that the only frame-sensitive quantity left is the difference of the two flows' own-frame remainders, which \Cref{rot:h1s:E} bounds. Identities (the counterterm cancellations, the Ward sum, the moment subtractions) are used only at the two pure endpoints, in each flow's own frame; the mixed terms are bounded in absolute value.

\begin{lemma}[every telescoping term is small]\label{rot:h1s:T}
For every ordering class and every term $T_i$, with $\theta\le\theta_1=\alpha^2/(4(1+\alpha)^2)$,
\[
|T_i|\le C(1+k)^p\eta'^\theta\cdot[\text{the per-flow bound \cite[(1.33)]{Ba3} of the monomial}],\qquad p=p(\bar n),
\]
with the prefactor $p_r^2\delta_r^{-1/2}$ when $e_i$ is a contour element.
\end{lemma}
\begin{proof}
Apply \Cref{rot:h1s:MS} to $|T_i|$, the absolute value of a multilinear form whose degrees are those of \Cref{rot:h1s:Dp}(a).
(i) \emph{Per-flow bounds} of every element in $X$-coordinates, whichever flow it comes from: lines by \eqref{eq:h1s:B}, with $\delta/2$ for the $Y\circ\tau$ versions; compounds by \Cref{rot:h1s:U}(i),(iii), lines of dimension $-2-a$, resp.\ $-3-a$ ($0$ at the top scale for a $\Gamma$ middle), at their leg's scale, or vertex weights $\alpha$ for external legs; first-order compounds by \Cref{rot:h1s:U}(iv) (degree-one objects); the brackets of (R2b) by $|G(t)-G(x)|\le|t-x|\sup|\nabla g|$ plus the explicit-$\eta'$ displacement term; coarse elements, externals and vertex functions by their table; the common numbers by \Cref{rot:h1s:Z0} and \Cref{rot:h1s:MS} per flow ($\le C(1+r)^{n_0+1}$ each; the vacuum-polarisation leftover $\le C\ell_J^{-1}$ as a weight $-1$); at $e>0$ the contour elements by \Cref{rot:prop:BD}(b), a contour line that received a derivative being a line of dimension $-1$ at its slice (the contour sum carries the weight $\eta$ per bond against $\sim\ell_j/\eta$ bonds within the decay range, \cite[(3.63)--(3.64)]{K1} with the contour clause of \cite[Prop.~3.7]{K1}). Decay is re-anchored at merged vertices as in the first step of \Cref{rot:h1s:MS}.
(ii) \emph{The switched element}, divided by $C\eta'^\theta$, obeys the bound of its type weakened by $\ell_j^{-\theta}$ at its scale (or unweakened if it is fixed): a line with $n\le2$ derivatives by \Cref{rot:h1s:P}, gain $(\eta'/\ell_j)^{1/4}\ge(\eta'/\ell_j)^\theta$ (top piece: $\eta'^{1/4}$); a compound by \Cref{rot:h1s:U}(ii),(iii), gain $\eta'^{\theta_U}\ell_{j_f}^{-\theta_U}$ with $\theta_U\ge\theta$; a first-order compound or a bracket by \Cref{rot:h1s:U}(iv), gain $\ge(\eta'/\ell_j)^{1/16}$ resp.\ $\eta^{1/4}$ (a bracket whose differenced factor is a compound: $\eta'^{\theta_U}\ell_{j_f}^{-\theta_U}$ with $\theta_U\ge\min\{\theta,s/8\}$); a coarse element by $C\eta^{1/2}$; an external by $C\eta^{1/2}\|E\|$, its gradient as a recipient by $C\eta^\kappa\|E\|$ (\Cref{rot:h1s:E}); a vertex function by $c_\tau\eta'\|\nabla w\|$, and the $\partial^{\eta'}g$ term of (R2a), of size $O(\eta')$ in both flows, by $O(\eta')\le\eta'^\theta\ell_{j(v)}^{1-\theta}$; at $e>0$ a contour element by \Cref{rot:prop:BD}(b): $(\eta'/\ell_j)^{1/2}$ for contour lines, also with the one lattice derivative a recipient contour line carries, $\eta^{1/2}$ for contour vertices, $e\,p_r\lmu^{-1}c_\tau\delta_r^{-1/2}\eta'$ for the phases, whence the prefactor; contour vertices are never recipients (\Cref{rot:h1s:rule}(iii)). For $j\le2$ there is no gain and the factor is paid by the degree loss $\theta$, as in \Cref{rot:cor:unif}(a).
(iii) \emph{Degrees.} By \Cref{rot:h1s:Dp} every component of every initial segment of the reduced ordering has $D\ge1$, or $D\in\alpha+\N_0$, or is an explicit-$\eta'$ component with $D=1-\theta$ after $\eta'^\theta$ is taken out; fixed elements impose no condition. Hence $D(H)-\theta>0$ for every component containing the switched element, and the scale sums of \Cref{rot:h1s:MS} converge geometrically there.
(iv) \emph{Logarithms.} $n_0\le\bar n$ zero-degree components, $(1+r)$ per compound, $(1+r)^{n_0+1}$ per common number, $(1+r)^2$ per double compound, and the $(1+r)^{n_0}$ of \Cref{rot:h1s:MS} itself; with $\le\bar n$ elements, $p(\bar n)\le3\bar n+2$ suffices, and $(1+r)\le(1+k)$.
\Cref{rot:h1s:MS} then gives $|T_i|\le C(1+k)^p\prod_lc_l$ with $\prod_lc_l=\eta'^\theta\cdot[\text{per-flow constants}]$; the localisation \cite[(3.55)]{K1} costs at most $p_r^3$ per coarse line, absorbed in the field norms of step (5) of the proof of \Cref{rot:thm:A}.
\end{proof}

\paragraph{The exponent $\theta_1$.}
In every term the switched element carries $\eta'^{\theta'}\ell_{j}^{-\theta'}$ with its own exponent $\theta'$, and \Cref{rot:h1s:MS} at the stage where the component $H\ni e_i$ is shrunk sums $\sum_{j\le j'}\ell_j^{D(H)-\theta}\eta'^\theta\le C\ell_{j'}^{D(H)-\theta}\eta'^\theta$, geometric because $D(H)-\theta>0$; the remaining factor is the merged vertex's weight, and since all later components containing it again have $D-\theta>0$ the factor $\eta'^\theta$ survives. (When $D(H)=0$ the sum $\sum_j(\eta'/\ell_j)^\theta=\sum_j5^{-j\theta}$ is $O(1)$ and the gain is lost; \Cref{rot:h1s:Dp} excludes this.) With $s=\alpha/(1+\alpha)$, $\kappa=\alpha/(2(1+\alpha))$, the gains are: a switched line with $n\le2$ derivatives, $1/2^n\ge\frac14$; the top piece, $\frac14$; a compound with a slice leg carrying $a\le2$ derivatives ($a=2$ only for the differenced compound of (R2b)), $s/2^{a+1}\ge s/8$; with an external or coarse leg, $\kappa s$; the double compound, $s/8$ (slice middle) or $\kappa s$ (coarse middle); the first-order forms, $1/2^{a+2}\ge\frac1{16}$ (slice) or $\frac14$ (external); coarse elements, $\frac12$; vertex functions and explicit-$\eta'$ terms, $1$; contour lines, $\frac12$; contour vertices, $\frac12$; phases, $1$ with the prefactor. The arithmetic: $\kappa s=\alpha^2/(2(1+\alpha)^2)=2\theta_1$; $s/4-\theta_1=\alpha/(4(1+\alpha)^2)>0$; $s/8-\theta_1=\alpha(1-\alpha)/(8(1+\alpha)^2)>0$ for $\alpha<1$; $\frac1{16}-\theta_1=(1-\alpha)(1+3\alpha)/(16(1+\alpha)^2)>0$; and $\min(\alpha,\frac12)-\theta_1>0$ because $\theta_1<\alpha/4\le\alpha$ and $\theta_1<\frac1{16}$. So $\theta_1\defeq\alpha^2/(4(1+\alpha)^2)$ is at most every gain exponent (the weakest two being $s/8$ and $2\theta_1$), and every component containing the switched element keeps the margin $D-\theta_1\ge\min(\alpha,1)-\theta_1>0$. Nothing is nested: no compound is the leg of another (the double compound is one element, and the recipients of (R2) receive a single gradient), so $\theta_1$ does not depend on $\bar n$. The consumers (\Cref{rot:lem:r}, \Cref{rot:prop:cp}) use only $\theta_1>0$ and a polynomial $(1+k)^p$, so the exact value plays no role; what matters is that the set of element types and derivative counts is finite, that every type has a difference law with some positive exponent, and that the local terms are common numbers.

\begin{proof}[Proof of \Cref{rot:thm:H1S}]
Sum \Cref{rot:h1s:T} over the $\le n_\cR$ terms of each ordering class and over the finitely many classes (\cite[(2.7)]{Ba3}; their number depends on the family only). The sectors:
\emph{$e>0$.} The constant-background reduction (\Cref{rot:lem:CB}) is flow-independent algebra; after it the new elements are the contour elements of \Cref{rot:prop:BD}(b). Every connected subgraph containing a contour vertex has degree $\ge\frac12$ (\cite[(3.30)]{K1}; \Cref{rot:lem:CB}(3)), so no new structure of degree $\le0$ appears and \Cref{rot:h1s:Dp} is unchanged; a contour line may be the recipient of (R2a)(ii) at its bond end, and this is the one-derivative clause of \Cref{rot:prop:BD}(b); the phases are Lipschitz in the vertex position with the bound $e\,p_r\lmu^{-1}\delta_r^{-1/2}$, inside the prefactor. $e$ enters only through $e_r^2=e^2\delta_r$, so the bound holds at every $e>0$.
\emph{The sectors with legs at $A=0$.} Scalar two-point: orders $\lambda$, $\lambda^2$, $e^2$, $e^4$ with the counterterm cancellations of \Cref{rot:h1s:VP}(d), $e^2\lambda$ trivial; the $(2,1)$-sector at order $e^3$: T and SC as first-order compounds (no Ward identity of the vertex class is used: in the per-flow normal form the identities of this sector are identities on all-$C$ expressions, applied in each flow alone, and a reduction of this sector to a ``mixed zero-moment'' bound for current--current subgraphs with lines from different flows is not used, the mixed zero moment being a lattice-scale quantity of order $e_r^2/\eta'$); the vacuum polarisation at $e^2$ through \Cref{rot:h1s:VP}(a),(b) and at $e^4$ through (c); L24; everything else of positive degree (\cite[Cor.~2.3, (2.17)]{Ba3}).
\emph{The vacuum sector.} All-$C$ vacuum monomials cancel against $E_1$ exactly (\Cref{rot:lem:ctsame}, \cite[(3.35)]{K1}); the $\Gamma$-monomials, cut at their $\Gamma$-lines, leave all-$C$ components with at least two legs, renormalised by the rule with the $\Gamma$-lines as legs and recipients (one-point components are common numbers, (R2c)); \Cref{rot:h1s:T} applies per coarse site, times $|T|\delta_r^{-3}$ sites (\Cref{rot:thm:V}(e)).
\emph{The source vertex.} The $h$-vertex has degree $2$ (\Cref{rot:lem:S1}); the $h$-families other than $\cT_h$ contain no new structure of degree $\le0$.
Hence $|E^X(\cR)-E^Y(\cR)|\le Cn_\cR\,p_r^2\delta_r^{-1/2}\eta'^\theta(1+k)^p\cdot[\text{right side of \cite[(1.33)]{Ba3}}]$ for $0<\theta\le\theta_1$, uniformly in $r$, $k$, the torus (the tori enter through periodisation, \Cref{rot:h1s:Z0}) and the masses, and in every term of the telescoping the switched element obeys the bound of its type weakened by $\theta$, which is the statement.
\end{proof}


% r03_fill_K3b: proof of (K_0), part V: the comparison of the effective actions in total form, at h = 0 and with a wall.

\subsubsection*{The comparison of the effective actions at equal spacing}

\begin{theorem}[equal-spacing comparison, total form]\label{rot:thm:A}
Let $T\in\{T_1,T_2\}$, $G$ a real smooth $P$-periodic $2$-form, $z\in\R$, and $r=r(k)$ as in \Cref{rot:lem:r}. Then for $k\ge k_0(zG)$,
\[
\sup_{\supp\chi_r}\big|\hs^{(r)}_T(A_r,\phi_r;zG,0)-S^{(r)}_T(A_r,\phi_r;zG,0)\big|\le\vartheta_k|T|,\qquad
\vartheta_k\defeq C\big(5^{-\gamma r}\delta_r^{-\beta}(1+k)^p+\delta_r^\sigma\big)\to0,
\]
every $c$-number included: there is no vacuum constant and no source-linear constant at $h=0$.
\end{theorem}

\begin{theorem}[the vacuum constant]\label{rot:thm:V}
Let $v_k(T)\defeq S^{(r)}_T(0,0;0)-\hs^{(r)}_T(0,0;0)$. Then: \textup{(a)} $v_k=v^{\rm G}_k+v^{\rm I}_k$ with $v^{\rm G}=\ln N^Y_r-\ln\hat N_r+\ln N^Y_r(0)-\ln\hat N_r(0)$ and $|v^{\rm G}_k(T)|\le12.5\,a\,5^{3k-5r}|T|$; \textup{(b)} $E_1$ cancels exactly in $v^{\rm I}$, and $v^{\rm I}_k(T)=|T|\Delta p_k+\varpi_k(T)$ with $\Delta p_k$ a difference of infinite-lattice vacuum densities, the same for both tori, and $|\varpi_k(T)|\le C\eps^{-q}e^{-c\delta_r^{-1}}$; \textup{(c)} the order-$\lambda$ part obeys $|v^{\rm I}_\lambda|\le C\lambda\eta^{1/2}\delta_r^{-2}|T|+|\varpi_k|$; \textup{(d)} the order-$\lambda^2e^0$ part obeys $|v^{\rm I}_{\lambda^2}|\le C\lambda^2\delta_r^{-2}(\eps/\delta_r)^{1/4}\log(1/\eps)^2|T|+|\varpi_k|$; \textup{(e)} $|v^{\rm I}_k(T)|\le\vartheta_k|T|$. Hence $v_k(T)\to0$ on each torus.
\end{theorem}
\begin{proof}
(a) \cite[(3.89)]{K1} gives $E_0=\ln Z_r+\ln Z_r(0)+\ln N_r+\ln N_r(0)$; inserted in \cite[(3.14)]{K1} at $A_r=\phi_r=0$ this leaves $\ln N_r+\ln N_r(0)+P^{(r)}(0)$ for each flow, with $\ln N_r=-\frac32\sum_{p\in T^{(r)*}}\ln(\Delta^{(r)}_{\lmu}(p)/2\pi)$ and $\ln N_r(0)=-\frac N2\sum_p\ln(\Delta^{(r)}_m(p)/2\pi)$ over the dual of the common coarse torus. By \Cref{rot:prop:fc}, $|\ln(A^X(p)/A^Y(p))|\le0.20\,a\eta^2$ at every $p$ and every mass; sum over $|T|\delta_r^{-3}$ momenta and $3+N=5$ components with the factor $\frac12$; $\eta^2\delta_r^{-3}=25\cdot5^{-2r}5^{3(k-r)}$. (b) $E_1=|T|e_1(\eps)$ is the same $c$-number for both lattices and passes through every block-spin step unchanged. The remaining vertices at zero coarse fields are \cite[(3.18), (3.19), (3.22), (3.23)]{K1} and the products (3.25)$\times$(3.26), (3.26)$\times$(3.26), with fluctuation kernels as lines; each torus kernel is the $P$-periodisation of its infinite-lattice kernel, decaying at rate $\delta_*/\delta_r$ in unit scale, with prefactor $\le C\eta'^{-3}$; expanding every torus line as infinite-lattice line plus winding remainder, a graph with $n\le2\bar n$ lines differs from its infinite-lattice value by $\le C^n\eta'^{-3n}|T|\delta_r^{-3}\sum_{\gamma\ne0}e^{-\delta_*|\gamma|/2}\le C\eps^{-q}e^{-c\delta_r^{-1}}$, the shortest period being $50MJ\ge50\delta_r^{-1}$ in unit scale. The infinite-lattice values are $|T|p^F_k$ by $Y_{k-r}$-periodicity of both flows. (c) The order-$\lambda e^0$ vertices are (3.18) and the order-$\lambda$ part of (3.19), and $\dm_1=-4(N+2)\lambda C(0)$ (the vanishing of the one-particle-irreducible two-point function at zero momentum at order $\lambda$, \cite[(3.33)]{K1}). With Gaussian moments, $P^F(0)|_\lambda=\lambda N(N+2)\sum_x[(C(0)-G^F_T(x,x))^2-C(0)^2]$; $C(0)-G^F_r(x,x)=\Gamma^F(x,x)$ by \Cref{rot:lem:Gam}(i); the terms $-C(0)^2$ give equal totals; re-index the $Y$-sum by $\tau$ and use \Cref{rot:lem:Gam}(ii),(iii): $|\Gamma^X(x',x')^2-\Gamma^Y(\tau x',\tau x')^2|\le2C_1C_\Gamma\eta^{1/2}\delta_r^{-2}$. (d) With $V_1\defeq\lambda\sum{:}|\phi|^4{:}_C$ there are no tadpoles at order $\lambda^2$; $\frac12\dm_2=16(N+2)\lambda^2c_3$, $c_3=\sum_yC(y)^3$; writing ${:}F{:}_C=e^{-\frac12\Delta_\Gamma}{:}F{:}_G$ with $\Gamma=C-G$, the order-$\lambda^2e^0$ part of $v^{\rm I}$ is the $Y$-minus-$X$ difference of $N(N+2)\lambda^2[16T_1-24T_2+16T_3-4T_4-4(N+2)T_5]$ with $T_1=\sum_{x,y}C_{xy}^3(\Gamma_{xy}-\Gamma_{xx})$, $T_2=\sum C^2_{xy}\Gamma^2_{xy}$, $T_3=\sum C_{xy}\Gamma^3_{xy}$, $T_4=\sum\Gamma^4_{xy}$, $T_5=\sum\Gamma_{xx}\Gamma_{yy}(C_{xy}-\Gamma_{xy})^2$, up to $O(\eps^{-q}e^{-c\delta_r^{-1}})$ (\Cref{rot:prop:V2}); the pure-$C$ terms have equal totals. Each difference is bounded, in the three regions $|u|\le\rho\defeq(\eps\delta_r)^{1/2}$, $\rho<|u|\le R_\eps\defeq(40/m)\ln(1/\eps)$ and $|u|>R_\eps$, by the pointwise bounds $|C_{xy}|\le C_0(|u|+\eps)^{-1}$, $|\Gamma_{xy}|\le C'_1\min(\delta_r^{-1},(|u|+\eps)^{-1})$, $|\Gamma_{xy}-\Gamma_{xx}|\le C_2\delta_r^{-1}(|u|/\delta_r)^{1/2}$, the matched-pair bounds of \Cref{rot:cor:C} and \Cref{rot:lem:Gam}(iii), and Combes--Thomas decay beyond $R_\eps$; the result is $\le C\delta_r^{-2}(\eps/\delta_r)^{1/4}\ln(1/\eps)^2$ per unit volume. (e) Group the vacuum graphs, with $E_1$ expanded in each flow's slices (\Cref{rot:lem:ctsame}), into renormalised vacuum families in per-flow normal form; all lines are $A=0$ lines on the infinite lattice; the vacuum graphs reduce to families with at least two legs, to which \Cref{rot:prop:cp} applies per coarse site; multiply by $|T|\delta_r^{-3}$ and use \Cref{rot:lem:r}.
\end{proof}

\begin{proposition}[the order-$\lambda^2$ identity]\label{rot:prop:V2}
The decomposition of $v^{\rm I}_{\lambda^2}$ into $T_1,\dots,T_5$ in the proof of \Cref{rot:thm:V}(d) is an exact identity; in particular the counterterm $\dm_2$ enters only through $T_1$, in the form $\sum C^3[\Gamma(x,y)-\Gamma(x,x)]$.
\end{proposition}
\begin{proof}
Wick's theorem for the two Gaussians with covariances $C$ and $G=C-\Gamma$ at second order in $V_1$, with ${:}|\phi|^4{:}_C$ re-expanded in $G$-normal order: the connected contractions of two quartic vertices with $n$ lines of type $C$ and $4-n$ of type $\Gamma$, $n=0,\dots,4$, give $T_4,T_3,T_2,T_1$ and the pure-$C$ term, with the combinatorial weights $4$, $16$, $24$, $16$; the counterterm $\frac12\dm_2|\phi|^2$ contracted with the quartic vertex gives the $-c_3\sum\Gamma(x,x)$ half of $T_1$; the two one-loop pieces of the re-expansion contracted together give $T_5$. The form is the one in which each telescoping term is small (\Cref{rot:rem:tele}): with an unexpanded counterterm the intermediate terms would be of order $\lambda^2|T|\delta_r^{-1}$ and small only because $\sum_y\eps^3(C^{X3}-(C^Y\circ\tau)^3)=0$.
\end{proof}

\begin{remark}[why manifest forms]\label{rot:rem:tele}
The comparison must be applied to families in which every counterterm is expanded in the flow's own lines, and the same applies to Ward-identity cancellations, which hold for pure slice assignments only: King's ``replace one by one every factor'' \cite[p.~665]{K1} is correct on such manifest forms only (\Cref{rot:lem:ctsame}).
\end{remark}

\begin{proof}[Proof of \Cref{rot:thm:A}]
Fix $T$ and work on $\supp\chi_r$, with the bounds \cite[(3.2)]{K1}: $|A_r|\le p_r(\lmu\delta_r)^{-1}$, $|\partial A_r|\le p_r$, $|\phi_r|\le p_r(\lambda\delta_r)^{-1/4}$, $|D_{A_r}\phi_r|\le p_r$, $p_r\defeq p(\delta_r)$. Write both effective actions in the form of \Cref{rot:def:eff}, with the curvature source carried by the Cameron--Martin background (\Cref{rot:rem:CMrg}), and compare term by term; every bound is per coarse site, times $|T|\delta_r^{-3}$ sites.
\begin{enumerate}[leftmargin=*]
\item \emph{Photon quadratic form.} Three copies of the scalar coarse operator at mass $\lmu$, and $|\frac12\ip{A_r}{(\hat\Delta^{(r)}-\Delta^{(r)})A_r}|\le0.3a^2\eta^2p_r^2(\lmu\delta_r)^{-2}|T|\delta_r^{-3}$ by \Cref{rot:prop:fc} (King's (3.92) at equal spacing).
\item \emph{Normalisations.} \Cref{rot:thm:V}(a).
\item \emph{Scalar quadratic form at $A^{(r)}=0$.} As (1). Its $A^{(r)}$-dependent part has the form of a graph in $P^{(r)}$ with two external scalar fields \cite[p.~669]{K1} and is treated with (5).
\item \emph{The background determinant} $B3^F\defeq-\ln[Z^F_r(A^F)Z^F_r(0)^{-1}]$. By the scale factorisation of the determinant, valid for both flows by \Cref{rot:lem:tile}(iii),(iv), $B3^F=\frac12\sum_{j<r}\sum_x[g^F_{j,x}(A^F)-g^F_{j,x}(0)]$ with $g^F_{j,x}=\Tr_N[\ln K^F_j](x,x)$; the depths $j\le1$ are bounded per flow (per-site gauge and charge-conjugation parity), at cost $\le CNe^2p_r^2\delta_r^{-2}5^{-r}|T|$; the depths $j\ge2$ live on the common lattice $Y_{k-j}$ and are compared per site by an exact matched re-indexing, the telescoping of the resolvents along a cascade of exponential weights, a kernel sweep with the kernel bounds of \Cref{rot:prop:BD}(a) and \cite[Thm.~3.3]{K1}, and gradients compared as vectors, giving, for every $s_0>0$, $|B3^X-B3^Y|\le C_{s_0}Ne^2p_r^c(1+k)^p\delta_r^{-2}(5^{-\theta_Dr}+\delta_r^{s_0})|T|+C\eps^{-q}e^{-c/\delta_r}\le\vartheta_k|T|$, $\theta_D=\alpha/4$, by \Cref{rot:lem:r}. King's \cite[(3.94)--(3.98)]{K1} is not used: the per-$p$ middle line of (3.98) is false as printed, its seagull and bubble classes each diverging like $h^{-1}$ and cancelling only between classes.
\item \emph{The interaction $P^{(r)}$.} For each flow apply King's reduction \cite[(3.43)--(3.57)]{K1} in the form of \Cref{rot:lem:CB}: the constant-background gauge transformation (\Cref{rot:rem:frameav}); the localisation in universal-cover form; the expansions (3.47)--(3.54) around $\vartheta A^{(r)}=0$; the contour phases (\Cref{rot:conv:C}(ii), \Cref{rot:lem:shift}). After the reduction the background sits in the explicit factors of (3.55) and in the special lines $\delta^{\rm im}$, $\delta^{\rm bg}$, and every compared line is an $A=0$ line of \Cref{rot:cor:unif}; the contour vertices (3.24)--(3.26) and the phases (3.46) are compared by \Cref{rot:prop:BD}(b), with relative gains $\eta^{1/2}$ and $\eta'$ and the prefactor $p_r^2\delta_r^{-1/2}$; the explicit factors $A_r(w_1)\cdots A_r(w_q)$ are the same coarse field for both flows; the special lines, the terms with a vertex outside the vertex region and the Taylor remainders are bounded per flow by \Cref{rot:lem:CB}(2). Group the graphs into renormalised families in per-flow normal form; this is possible for both flows with the same combinatorics because each flow satisfies the identities used: counterterms (\Cref{rot:lem:ctsame}), inversion about coarse sites (\Cref{rot:conv:C}(i)), gauge covariance, and \cite[(3.87)]{K1}, an identity for each lattice's own full propagator. \Cref{rot:prop:cp} bounds each family difference by $Cp_r^2\delta_r^{-1/2}\eta'^{\theta}(1+k)^p$ times \cite[(1.33)]{Ba3}; the localisation sums give $C$ per coarse site, the field norms on $\supp\chi_r$ give at most $Cp_r^{\bar n}\delta_r^{-\beta_1}$ with $\beta_1$ fixed by \cite[(3.2)]{K1} and $\bar n$, and the number of families depends on $\bar n$ only.
\end{enumerate}
Summing, $|\hs^{(r)}-S^{(r)}|\le C(\eta'^{\theta}\delta_r^{-\beta}(1+k)^p+\delta_r^\sigma)|T|$ with $\eta^2\le25\eta'^\theta$ and the prefactor $p_r^2\delta_r^{-1/2}$ absorbed in $\delta_r^{-\beta}$ ($\beta\ge\beta_1+\frac12$ in \Cref{rot:lem:r}); with \Cref{rot:lem:r} this is $\vartheta_k|T|\to0$. The vacuum $c$-number is \Cref{rot:thm:V}(e). On $T_1$ the $Y$-flow lives on a torus skew in its own frame; this enters only through periodisation at $A=0$ and deck descent at $A\ne0$ (\Cref{rot:lem:F}).
\end{proof}

\subsubsection*{Walls: the source vertex}

\begin{lemma}[the $h$-census]\label{rot:lem:S1}
In King's degree (\cite[(3.30)--(3.31)]{K1}; $d=3$: a scalar end counts $-\frac12$, a differentiated end $-\frac32$, each factor $\eta$ counts $+1$, $D(G)=\sum_vD_G(v)-3$) the $h$-vertex \cite[(3.21)]{K1} has $D=3-\frac12\cdot2=2$ with both legs internal. Inserting one $h$-vertex into a scalar line of a graph $G_0$ gives $D(G_0+h)=D(G_0)+2$. Field-dependent $h$-free families have $D\ge-1$ with one exception, the current tadpole (one vertex \cite[(3.22)]{K1} with $m=1$ on a scalar loop, the photon leg external), which has $D=-2$: its value at $A=0$ vanishes ($\Tr q=0$), and in a background every insertion adds an external photon leg and raises the degree by $1$. So every non-vanishing field-dependent $h$-family has $D\ge1$. The only $h$-families with $D\le0$ are field-independent and linear in $h$: the free tadpole ($D=-1$; its residue at zero coarse fields, $-N\sum\eta^3h\,\Gamma^F(x,x)$, is Term~0 of \Cref{rot:prop:CK}), the order-$\lambda$ figure-eight with its $\dm_1$ partner, and the order-$e^2$ exchange and seagull tadpoles with their $\dm_2$ partner ($D=0$). Call this set $\cT_h$. The tree-level term $\delta_r^2\sum\eta^3h|\phi^{(r)}|^2$ has $D=0$ but no internal line; it is compared by \Cref{rot:lem:H} directly and gives $O(|T|\eps\delta_r^{-2}p_r^2)\to0$. The $X$/$\tau Y$ difference of the $h$-vertex is a smooth vertex function sampled at $x'$ and $\tau x'$ (\Cref{rot:lem:vert}(ii)), an explicit factor $\eta'$; $\sum_i\phi_i^2$ is frame-free.
\end{lemma}
\begin{proof}
Degree arithmetic from the vertex list \cite[(3.18)--(3.26)]{K1} (the seagull of order two included) and the minimal degree $-1$ of the field-dependent $h$-free families read from Ba{\l}aban's list \cite[(2.4)]{Ba3}; the exhaustive enumeration of the $h$-linear graphs of order $\le2$ in the couplings, with the merged counterterm vertices, gives $\cT_h$. For a one-particle-reducible family a bridge line adds $-1+3=+2$ to the degree, so its degree is at least the minimum over its one-particle-irreducible parts.
\end{proof}

\begin{lemma}[the source-linear constant is the same on both tori]\label{rot:lem:C}
Let $f^F(x)$ be the infinite-lattice one-point density of $\cT_h$ along flow $F$ (one $h$-vertex at $x$, all lines infinite-lattice kernels); it is $Y_{k-r}$-periodic. Let $\bar f^F$ be its cell average and $c_k\defeq-(\bar f^X-\bar f^Y)$, a number independent of $T$. For $T\in\{T_1,T_2\}$ the $\cT_h$-part of $\hs_T-S_T$ at zero coarse fields equals $-c_k\sS_{\Xk}(h)+\rho_T(h)$ with $|\rho_T(h)|\le C(\eps^{-q}e^{-c\delta_r^{-1}}+C_N\eps^{-q}\delta_r^N)\|h\|_{C^N}\to0$. Moreover $|c_k|\le C\eps^{-q}$.
\end{lemma}
\begin{proof}
Periodisation as in \Cref{rot:thm:V}(b), with one $h$-vertex summed over $T$. Write $x=z+t$ with $z\in Y_{k-r}/P$ and $t$ in a cell; the Riemann sum of the smooth periodic $h$ over the coarse lattice is exact to $O(C_N\delta_r^N)$, so $\sum_x\eps^3h(x)f^F(x)=\bar f^F\int h+O(\|f^F\|_\infty\delta_r^N)$, with $\|f^F\|_\infty\le C\eps^{-q}$ since each line is $\le C\eta'^{-3}$; and $\sS_{\Xk}(h)=\int h+O(C_N\eps^N)$ (\Cref{rot:lem:lsum}). Take $N$ large: with $\delta_r=5^{-\lfloor\vr k\rfloor}$, $N>(3+q-3\vr)/\vr$ suffices.
\end{proof}

\begin{theorem}[comparison with a wall]\label{rot:thm:Asrc}
For $T\in\{T_1,T_2\}$ and $s=(zG,h)$ real, on $\supp\chi_r$: $\big|\hs^{(r)}_T(\cdot;s)-S^{(r)}_T(\cdot;s)+c_k\sS_{\Xk}(h)\big|\le\vartheta_k(s)|T|\to0$, with $c_k$ of \Cref{rot:lem:C}.
\end{theorem}
\begin{proof}
Split $\hs(G,h)-S(G,h)=[\hs(G,0)-S(G,0)]+[\Delta_h\hs-\Delta_hS]$. The first bracket is \Cref{rot:thm:A}. Since \cite[(3.17)]{K1} depends on $h$ only through the $h$-vertex, the second bracket consists of the $h$-families of the two Taylor polynomials plus Taylor remainders bounded per flow by $C\delta_r^\sigma|T|$ \cite[p.~665]{K1}. The $\cT_h$ part is $-c_k\sS_{\Xk}(h)+o(1)$ by \Cref{rot:lem:C}; its $G$-dependence through the Cameron--Martin background is an $h$-family with background insertions, which by gauge invariance of each flow acts through the curvature of the background and so has positive degree. Every other $h$-family has $D\ge1$ (\Cref{rot:lem:S1}); group them in per-flow normal form with the same combinatorics in both flows, telescope element by element (lines: \Cref{rot:cor:unif}; external fields: \Cref{rot:lem:H}; the $h$-vertex: \Cref{rot:lem:vert}(ii)), apply \Cref{rot:prop:cp}, and sum as in the proof of \Cref{rot:thm:A}. The exponents of \Cref{rot:lem:r} are source-independent (\cite[Thm.~3.4]{K1}: ``$C$ depends on $g,h$'').
\end{proof}

% r03_fill_K4: proof of (K_0), part VI: the source-linear constant vanishes.

\subsubsection*{The source-linear constant vanishes}

Throughout, unit scale of level $r$ (\Cref{rot:def:eff,rot:def:match}): fine lattices $\Lambda_X=\eta'\Z^3$, $\Lambda_Y=R^{-1}(\eta'\Z^3)$, $\eta'=5^{-r}$, $\eta=5\eta'$; $\tau$ the cell-wise bijection with $|\tau x'-x'|\le c_\tau\eta'$; for $F\in\{X,Y\}$ the full propagator $C^F$ (mass $m_r=m\delta_r$; $C^Y=C^X\circ R$), the photon propagator $D^F_C$ (three scalar copies, frame-free), the $r$-step fluctuation kernels $G^F=\sum_{j<r}S^F_j$ and $D^F$, and the coarse covariances $\Gamma^F=C^F-G^F=H^FB^FH^{F*}$, $\Gamma^F_{\rm ph}=D^F_C-D^F$ (\Cref{rot:lem:Gam}(i)). Running couplings $\lambda_r=\lambda\delta_r$, $e_r^2=e^2\delta_r$; the $h$-vertex at zero coarse fields is $\delta_r^2\sum_x\eta'^3h(x)[|\phi'(x)|^2-C^F(x,x)]$, $N=2$ real components, with the free-boundary subtraction $C^F(x,x)=C(0)$ common to both flows. A unit-scale density $d$ corresponds to the physical density $\delta_r^{-1}d$. Constants depend on $a,m,\lmu,N,\tau$ only. $L_\eta\defeq C\log^3(1/\eta)$.

\emph{Kernel toolbox} ($u$ the displacement of the two arguments; ``both flows'' $=F\in\{X,Y\}$).
(Q1) $|\nabla^jC(u)|\le C_j(|u|+\eta')^{-1-j}e^{-cm_r|u|}$, $j\le3$, lattice differences along frame directions; the same for $D_C$; constants uniform in $\eta'\le1$ and in the mass ($m_r\in[0,m]$ for $C$, $\lmu\delta_r\in[0,\lmu]$ for $D_C$; $c$ halved for $j=3$). This is the standard bound for the massive free lattice Laplacian: scale to $\eta'=1$, $\mu=m_r\eta'$, symbol $E(q)+\mu^2$ with $E(q)=\sum_\nu4\sin^2(q_\nu/2)$; by the lattice symmetries let $|u_1|=|u|_\infty$, and shift the $q_1$-contour by $\mathrm i\kappa$, $\kappa=\min(\mu/2,1)$, where $\Re[E(q+\mathrm i\kappa e_1)+\mu^2]\ge\frac12(E(q)+\mu^2)$ ($\cosh\kappa-1\le\kappa^2$; the vertical sides cancel by periodicity): factor $e^{-\kappa|u_1|}$; on the shifted contour the integrand with $j$ differences obeys $|\partial_{q_1}^b(\cdot)|\le C_b(E+\mu^2)^{(j-2-b)/2}\le C'_b|q|^{j-2-b}$ uniformly in the mass ($j\le2$), and a cutoff at $|q|\sim|u_1|^{-1}$ with $j+2$ integrations by parts in $q_1$ on the high part gives $|u_1|^{-1-j}$ ($j=3$: an extra factor $1+\mu|u_1|$, absorbed by halving $c$). Both factors come from one integral, so the product of the two bounds holds.
(Q2) Matched-pair differences of $C$-type kernels: \Cref{rot:cor:C} gives $|C^X(y',z')-C^Y(\tau y',\tau z')|\le C_6\eta'^{1/2}|u|^{-3/2}$ for $|u|\ge25\eta'$ (the same for $D_C$); for a kernel carrying $j\le2$ differences, interpolation of this bound against the per-flow bounds $|\nabla^{j+1}C|\le C(|u|+\eta')^{-2-j}$ (difference quotients at an optimised step) gives the relative gains $(\eta'/|u|)^{1/4}$ ($j=1$) and $(\eta'/|u|)^{1/8}$ ($j=2$) on the per-flow size $(|u|+\eta')^{-1-j}$; inside $|u|<25\eta'$ the difference is bounded by the sum of the two per-flow sizes.
(Q3) Fluctuation kernels per flow: $|S^F_j(x,y)|\le C\ell_j^{-1}e^{-\delta|x-y|/\ell_j}$, $|\nabla S^F_j|\le C\ell_j^{-2}e^{-\delta|x-y|/\ell_j}$, $|\nabla S^F_j(x,y)-\nabla S^F_j(x',y)|\le C_{\alpha,r}\ell_j^{-2-\alpha}|x-x'|^\alpha e^{-\delta\min(|x-y|,|x'-y|)/\ell_j}$, $\ell_j=5^j\eta'$; hybrid flow: \Cref{rot:prop:BD}(a) at $A=0$, with $C_{\alpha,r}\le C_\alpha\log^{3\alpha}(2+\ell_j/\eta')$ on slice $j$; standard flow: \cite[Lemmas~2.1, 2.2, 2.4]{Ba4} and \cite[Thm.~3.3]{K1} at $A=0$. Hence $|\nabla^jG^F(u)|\le C(|u|+\eta')^{-1-j}e^{-\delta'|u|}$, $j=0,1$.
(Q4) Coarse covariances per flow: $\sup|\Gamma^F|\le C_1$ (\Cref{rot:lem:Gam}(ii)); $|\Gamma^F(x,y)|,|\nabla_y\Gamma^F(x,y)|\le C(1+|x-y|)^{-1}e^{-cm_r|x-y|}$, uniformly in $r$ and $m_r\in[0,m]$: for $|x-y|\ge1$ the triangle inequality between (Q1) ($j=0,1$) and the fixed-rate decay $|\nabla_y^iG^F(x,y)|\le C_Ge^{-\delta'|x-y|}$, $i\le1$, which (Q3) gives by summing the slices ($\ell_j\le\frac15$), with $e^{-\delta'\rho}\le c_{\delta'}\rho^{-2}e^{-(\delta'/2m)m_r\rho}$; for $|x-y|<1$ it is $C_1$ and the Lipschitz bound with $C'_\nabla\defeq(2\pi)^{-3}C_BC_hC_{\nabla h}\int_{V_1}|p|^{-2}dp$ from \Cref{rot:lem:Gam}(ii) and the bounds $|\hat h_x|\le C_h$, $|\nabla_y\hat h_y|\le C_{\nabla h}$ of the proof of \Cref{rot:lem:CK1}. H\"older $|\Gamma^F(x,y)-\Gamma^F(x,w)|\le C_2|y-w|^{1/2}$ (\Cref{rot:lem:Gam}(iv)) and Lipschitz with constant $\sqrt3C'_\nabla$. The same for $\Gamma_{\rm ph}$ (mass $\lmu\delta_r$).
(Q5) Second differences, both flows: $|\nabla\nabla H^F(x,z)|\le L_\eta$ and $|\nabla_y\nabla_y\Gamma^F(x,y)|\le CL_\eta$ (\Cref{rot:lem:CK1}).
(Q6) Matched-pair differences without derivatives: $|\Gamma^X(x',y')-\Gamma^Y(\tau x',\tau y')|\le C_\Gamma\eta^{1/2}$ (\Cref{rot:lem:Gam}(iii)), and $\le C_6\eta'^{1/2}|u|^{-3/2}$ for $|u|\ge25\eta'$ (\Cref{rot:cor:C}, for $G^F_r$ and $\Gamma^F$ separately); $|G^X(x',y')-G^Y(\tau x',\tau y')|\le C_6\eta'^{1/2}|u|^{-3/2}$ ($|u|\ge25\eta'$) and $\le2C_0(|u|+\eta')^{-1}$ always; the same for the photon kernels.
(Q7) Vertex matching: \Cref{rot:lem:vert}(i)--(iii). (Q8) Summation by parts on the infinite lattice and the lattice Leibniz rule $\nabla_\mu(FG)(y)=(\nabla_\mu F)(y)G(y+\eta'e_\mu)+F(y)(\nabla_\mu G)(y)$.

\begin{lemma}[the block-average symbol]\label{rot:lem:CK0}
The $(r-1)$-step centred cube average is the composition of $5$-point averages along each rotated axis, whose symbols telescope, $\prod_j\sin(5^{j+1}t/2)/(5\sin(5^jt/2))=\sin(5^{r-1}t/2)/(5^{r-1}\sin(t/2))$; hence $u_{r-1}(q)=\prod_\alpha\sin(q_\alpha/2)/(\eta^{-1}\sin(\eta q_\alpha/2))$, $q_\alpha=q\cdot f_\alpha$, on $V_\eta$. Consequently $J_0\defeq\int_{V_\eta}|u_{r-1}(q)|\,dq\le(C\log(1/\eta))^3$ (the $L^1$ norm of the Dirichlet kernel per axis), $J_1\le C$, and $\sum_{l_1}|u_{r-1}(p+l_1)|\le C\log^3(1/\eta)$, $\sum_{l_1\ne0}|u_{r-1}(p+l_1)|/|p+l_1|\le C$, uniformly in $p\in V_1$.
\end{lemma}
\begin{proof}
The telescoping of the Dirichlet kernels along each axis; the $L^1$ norm of the Dirichlet kernel of order $n$ is $\asymp\log n$.
\end{proof}

\begin{lemma}[second differences of the coarse elements, both flows]\label{rot:lem:CK1}
$|\nabla\nabla H^F(x,z)|\le CL_\eta$ for every pair of frame directions, $F\in\{X,Y\}$; and, uniformly in $m_r\ge0$, $|\nabla_y\nabla_y\Gamma^F(x,y)|\le CL_\eta$ and $|\nabla_y\nabla_w\Gamma^F(y,w)|\le C$. The same for $\Gamma_{\rm ph}$.
\end{lemma}
\begin{proof}
By \Cref{rot:lem:H}, $H^F(x,z)=(2\pi)^{-3}\int_{V_\eta}e^{\mathrm iq\cdot(x-z)}u_{r-1}(q)M^F(q)A^F(q)\,dq$ with $M^F(q)=\sum_le^{\mathrm i(q+l)\cdot r_F}u_F(q+l)/\Delta'_F(q+l)$. A second lattice difference in the fine variable multiplies the alias-$l$ summand by $d_F(q+l)^2$, $|d_F(Q)|\le|Q|$; since $\Delta'_F(Q)\ge c|Q|^2$ on every alias cell and $\sum_lu_F(q+l)^2=1$ (\Cref{rot:prop:fc}), $|\sum_ld_F^2u_F/\Delta'_F|\le C$, and the integral against $|u_{r-1}|$ costs $J_0\le CL_\eta$ (\Cref{rot:lem:CK0}). For $\Gamma$ write $\Gamma(x,y)=(2\pi)^{-3}\int_{V_1}B(p)\hat h_x(p)\overline{\hat h_y(p)}\,dp$, $\hat h_y(p)=\sum_zH(y,z)e^{-\mathrm ipz}$, which by Poisson summation equals $A^F(p)\sum_{Q=p+l_1+l}e^{\mathrm iQ\cdot y}u_{r-1}(p+l_1)u_F(Q)/\Delta'_F(Q)$; a fine difference in $y$ multiplies the $Q$-summand by $d_F(Q)$. Hence $|\hat h_y|\le C$, $|\nabla_y\hat h_y|\le C[|p|A^F/\Delta'_F+a\sum_{l_1\ne0}|u_{r-1}(p+l_1)|/|p+l_1|]\le C$ and $|\nabla_y\nabla_y\hat h_y|\le Ca\sum_{l_1}|u_{r-1}(p+l_1)|\sum_l|u_F|\le CL_\eta$, uniformly in $p\in V_1$ and $m_r\ge0$ (\Cref{rot:lem:CK0}). With $B^F(p)\le C_B|p|^{-2}$ (\Cref{rot:lem:Gam}(ii)), integrable on $V_1$ in three dimensions, both bounds follow.
\end{proof}

\begin{lemma}[matched-pair comparison of gradient vectors]\label{rot:lem:CK2}
Let $g$ on $\Lambda_X$ and $h$ on $\Lambda_Y$ satisfy, on a region, $|g(x')-h(\tau x')|\le A$ at matched pairs, $|\nabla g|,|\nabla h|\le C_\nabla$, and the H\"older modulus $|\nabla g(y)-\nabla g(y+v)|\le C_H|v|^\alpha$ (likewise $h$), $0<\alpha\le1$. With $\nabla^Xg(y')\defeq\sum_\mu(\nabla_\mu g)(y')e_\mu$ and $\nabla^Yh(\tau y')\defeq\sum_\mu(\nabla^Y_\mu h)(\tau y')f_\mu$ the lattice gradient vectors, for every $\rho\in\eta'\N$
\[
|\nabla^Xg(y')-\nabla^Yh(\tau y')|\le3\big[2A/\rho+2C_H\rho^\alpha+4c_\tau C_\nabla\eta'/\rho\big].
\]
\textup{(a)} For $\Gamma$-type kernels, $A=C_\Gamma\eta^{1/2}$ (Q6), $C_H=CL_\eta$ with $\alpha=1$ (Q5): $\rho=\eta^{1/4}$ gives $|\nabla^X\Gamma^X(x',\cdot)(y')-\nabla^Y\Gamma^Y(\tau x',\cdot)(\tau y')|\le C\eta^{1/4}L_\eta=:A_1$. \textup{(b)} Decaying form: $A=C_6\eta'^{1/2}|u|^{-3/2}$ for $|u|\ge25\eta'$ gives $\le C\eta'^{1/4}|u|^{-3/4}L_\eta$. \textup{(c)} Fluctuation kernels slice by slice: on slice $j$, $A_j=C\eta'^{1/2}\ell_j^{-3/2}e^{-\delta|u|/\ell_j}$, $C_{H,j}=C_\alpha\log^{3\alpha}(2+\ell_j/\eta')\,\ell_j^{-2-\alpha}$, $C_{\nabla,j}=C\ell_j^{-2}$; with $\alpha=\frac12$ and $\rho=\ell_j(\eta'/\ell_j)^{1/3}$, $|\nabla^XS^X_j(x',y')-\nabla^YS^Y_j(\tau x',\tau y')|\le C_\vartheta\ell_j^{-2}(\eta'/\ell_j)^{\vartheta}e^{-\delta|u|/(2\ell_j)}$ for every $\vartheta<\frac16$ and every $u$, and, summing over $j<r$, $|\nabla^XG^X(x',y')-\nabla^YG^Y(\tau x',\tau y')|\le C\eta'^{1/6}\log(1/\eta)\,(|u|+\eta')^{-13/6}e^{-\delta'|u|}$.
\end{lemma}
\begin{proof}
Lattice Taylor expansion: for $v\in\Lambda_Y-\Lambda_Y$, $h(\tau y'+v)-h(\tau y')=\nabla^Yh(\tau y')\cdot v+\varrho$, $|\varrho|\le3C_H|v|^{1+\alpha}$ (sum the increments along a lattice path of at most $3|v|/\eta'$ bonds). Take $v_\mu\defeq\tau(y'+\rho e_\mu)-\tau y'$, $|v_\mu-\rho e_\mu|\le2c_\tau\eta'$; then $e_\mu\cdot\nabla^Yh(\tau y')=\rho^{-1}[h(\tau(y'+\rho e_\mu))-h(\tau y')]+O(C_H\rho^\alpha)+O(C_\nabla c_\tau\eta'/\rho)$, and likewise $e_\mu\cdot\nabla^Xg(y')=\rho^{-1}[g(y'+\rho e_\mu)-g(y')]+O(C_H\rho^\alpha)$. Subtract: the bracket differences are matched-pair differences, each $\le A$; sum over the three components. In (c) the two main terms balance at the stated $\rho$, giving the exponent $\alpha/(2(1+\alpha))=\frac16$, the $\eta'$ term is smaller, and the sum over $j$ is dominated by $\ell_j\sim|u|$.
\end{proof}

% r03_fill_K4b: proof of (K_0), part VI (continued): the order-lambda and order-e^2 densities and the vanishing of the source-linear constant.

\begin{lemma}[order $\lambda$: the $h$-inserted figure-eight]\label{rot:lem:CK3}
$f^F_\lambda(x)=s_\lambda\,4N(N+2)\lambda_r\sum_{y\in\Lambda_F}\eta'^3G^F(x,y)^2\Gamma^F(y,y)$ with a sign $s_\lambda$ common to both flows, and $|f^X_\lambda(x')-f^Y_\lambda(\tau x')|\le C\lambda_r\eta^{1/2}$ (unit scale), i.e.\ $\le C\lambda\eta^{1/2}$ in physical units.
\end{lemma}
\begin{proof}
The order-$\lambda$ vertices are \cite[(3.18)]{K1}, $\lambda_r\sum_y\eta'^3|\phi(y)|^4$, and the order-$\lambda$ part of (3.19), $\frac12\dm_1\sum_y\eta'^3|\phi(y)|^2$ with $\dm_1=-4(N+2)\lambda_rC(0)$. The connected pairing of $|\phi(x)|^2$ with $|\phi(y)|^4$ gives $4N(N+2)G(x,y)^2G(y,y)$, with $|\phi(y)|^2$ it gives $2NG(x,y)^2$; the sum is $4N(N+2)\lambda_rG(x,y)^2[G(y,y)-C(0)]=-4N(N+2)\lambda_rG(x,y)^2\Gamma(y,y)$ (\Cref{rot:lem:Gam}(i)). Re-index the $Y$-sum by $\tau$ (Q7) and telescope: the term with $[\Gamma^X(y',y')-\Gamma^Y(\tau y',\tau y')]$ is bounded by $C_\Gamma\eta^{1/2}$ times the convergent sum $\sum_{y'}\eta'^3G^Y(\tau x',\tau y')^2\le C$; the term with $[G^X(x',y')^2-G^Y(\tau x',\tau y')^2]\Gamma^X(y',y')$ is bounded through (Q6) and (Q3): for $|u|<25\eta'$ the at most $25^3\cdot4\pi/3$ sites give $\le C\eta'$, and for $|u|\ge25\eta'$, $\sum\eta'^3C_6\eta'^{1/2}|u|^{-3/2}\cdot2C|u|^{-1}e^{-\delta'|u|}\le C\eta'^{1/2}$. Convert with $\lambda_r\delta_r^{-1}=\lambda$.
\end{proof}

\begin{lemma}[order $e^2$: kernel form]\label{rot:lem:CK4}
Let $\Pi^F[L_c,L_d](y,z)$ be the one-loop self-energy kernel with scalar loop line $L_c$ and photon line $L_d$ (exchange part from two current vertices, \cite[(3.22)]{K1} with $m=1$; seagull part from the $A^2$ vertex), $\Pi_C\defeq\Pi[C,D_C]$, $\dm_2\defeq-\hat\Pi_C(0)$ (the zero-momentum reading of \cite[(3.33)]{K1} at order $e^2$ in the free-boundary convention, the same number for both flows), and $\Pi_{C,R}(y,z)\defeq\Pi_C(y,z)+\dm_2\eta'^{-3}\delta_{yz}$, so that $\sum_z\eta'^3\Pi_{C,R}(y,z)=0$. Then, with a sign $s_e$ common to both flows,
\[
f^F_{e^2}(x)=s_e\Big\{\sum_{y,z}\eta'^6G^F(x,y)\,\Pi^F[G^F,D^F](y,z)\,G^F(z,x)+\dm_2\sum_y\eta'^3G^F(x,y)^2\Big\},
\]
$|\Pi_C(y,z)|\le Ce_r^2(|u|+\eta')^{-4}e^{-cm_r|u|}$ (Q1), and $|\Pi^X_C(y',z')-\Pi^Y_C(\tau y',\tau z')|\le Ce_r^2(\eta'/|u|)^{1/8}|u|^{-4}$ for $|u|\ge25\eta'$, $u=y-z$ (telescoping the two lines, each carrying at most two differences, with (Q2)). In each flow the exchange part equals a finite sum of terms $s_\iota\sum_{w,w'}\eta'^6L_a(x,w^+)L_b(x,w'^+)L_c(w^+,w'^+)L_d(w,w')$ with $w^+\in\{w,w+\eta'e_\mu\}$, $L_a,L_b,L_c\in\{G,\nabla_\mu G\}$, $L_d\in\{D,\nabla_\mu D\}$, exactly two differences per term and at most one per line, both along the same $\mu$ (a Euclidean inner product of two gradient vectors, \Cref{rot:lem:vert}(iii)); the $A^2$ part is frame-free.
\end{lemma}
\begin{proof}
Wick algebra at zero coarse fields; the graph with both $h$-legs at one current vertex vanishes because the current one-point function $\Im G(b_+,b_-)$ is zero. At each current vertex $\sum_s(-1)^sG(x,b_{\rm ext})G_c(b_{\rm loop},\cdot)=\eta'[(\nabla_\mu G)(x,w)G_c(w+\eta'e_\mu,\cdot)-G(x,w+\eta'e_\mu)(\nabla_\mu G_c)(w,\cdot)]$; in the (loop, loop) type the loop line carries two differences, which (Q8) moves onto $G(x,\cdot)$ or $D(\cdot,w')$ with arguments shifted by at most $\eta'$. The shifts are frame-dependent ($w+\eta'e_\mu$ against $\tau w+\eta'f_\mu$); the comparison charges one per-flow Taylor step per shifted argument, a relative gain $\eta'/\rho$ that integrates to $C\eta'(1+\log(1/\eta'))$ and is absorbed.
\end{proof}

\begin{lemma}[the $C/\Gamma$ split of the order-$e^2$ density]\label{rot:lem:CK5}
With $G=C-\Gamma$, $D=D_C-\Gamma_{\rm ph}$ and $\Pi_r\defeq\Pi[G,D]$, $f^F_{e^2}=s_e\{Z^F+A^F(x)+B^F(x)\}$ exactly, where $Z^F\defeq\sum\eta'^6C(x,y)\Pi_{C,R}(y,z)C(z,x)$, $A^F(x)\defeq\sum\eta'^6[G(x,y)G(x,z)-C(x,y)C(x,z)]\Pi_{C,R}(y,z)$ and $B^F(x)\defeq\sum\eta'^6G(x,y)G(x,z)[\Pi_r-\Pi_C](y,z)$. \textup{(i)} $Z^F$ is independent of $x$ and $Z^X=Z^Y$ (the all-$C$ part cancels exactly; it contains the logarithm of \Cref{rot:rem:log}). \textup{(ii)} $A^F(x)=\sum_{y\ne z}\eta'^6\Pi_C(y,z)[\Gamma(x,y)-\Gamma(x,z)]\{[C(x,y)-C(x,z)]-\frac12[\Gamma(x,y)-\Gamma(x,z)]\}$. \textup{(iii)} $\Pi_r-\Pi_C=-\Pi[\Gamma,D_C]-\Pi[C,\Gamma_{\rm ph}]+\Pi[\Gamma,\Gamma_{\rm ph}]-\Pi^{\rm sg}[\Gamma_{\rm ph}]$: every term has at least one internal $\Gamma$-line. The counterterm never appears alone: it is bound into $\Pi_{C,R}$, which occurs only in $Z$ and $A$. If King's scheme differs from the zero-momentum reading by a constant $\varrho=O(e_r^2)$, $f^F_{e^2}$ acquires $\varrho\sum_y\eta'^3G(x,y)^2$, whose two-flow difference is $\le C|\varrho|\eta'^{1/2}$ by the argument of \Cref{rot:lem:CK3}.
\end{lemma}
\begin{proof}
(i) is \Cref{rot:lem:ctsame}'s argument. (ii) $GG-CC=-\Gamma C-C\Gamma+\Gamma\Gamma$; symmetrise in $(y,z)$ using $\Pi(y,z)=\Pi(z,y)$ and the zero sum of $\Pi_{C,R}$ in $z$. (iii) is multilinearity.
\end{proof}

\begin{lemma}[the external-$\Gamma$ part]\label{rot:lem:CK6}
$|A^F(x)|\le Ce_r^2(1+m_r^{-1})$ and $|A^X(x')-A^Y(\tau x')|\le Ce_r^2\big[\eta^{1/4}L_\eta\log(1/\eta')+\eta'^{1/8}(1+\log(1/\eta'))^2(1+m_r^{-1})\big]$.
\end{lemma}
\begin{proof}
With $\rho=|x-y|\le\rho'=|x-z|$ and $u=y-z$: per flow $\mathsf a\defeq|\Gamma(x,y)-\Gamma(x,z)|\le C\min(|u|,1)(1+\rho)^{-1}e^{-c'm_r\rho}$ for $|u|\le\rho/2$, $\le2C(1+\rho)^{-1}e^{-c'm_r\rho}$ and $\le\sqrt3C'_\nabla|u|$ always ((Q4), $c'=c/2$); $\mathsf a_C\defeq|C(x,y)-C(x,z)|\le C|u|(\rho+\eta')^{-2}$ for $|u|\le\rho/2$, $\le2C(\rho+\eta')^{-1}$ always (Q1). Difference elements: $\delta\Pi\le Ce_r^2(\eta'/|u|)^{1/8}|u|^{-4}$ (\Cref{rot:lem:CK4}); $\delta\mathsf a\le C\min\{\eta^{1/2};\ \eta'^{1/2}(\rho^{-3/2}+\rho'^{-3/2})\ (\rho,\rho'\ge25\eta');\ A_1|u|+L_\eta|u|^2+\eta'\}$ ((Q6), \Cref{rot:cor:C}, \Cref{rot:lem:CK2}(a) with the lattice Taylor expansion of $\Gamma(x,\cdot)$ in each flow); $\delta\mathsf a_C\le C(\eta'/\rho)^{1/2}(\rho+\eta')^{-1}$ (Q2). Telescope (ii) of \Cref{rot:lem:CK5} and sum over the three regions $|u|\le\rho/2$; $|u|>\rho/2$, $\rho\le1$; $|u|>\rho/2$, $\rho>1$, bounding lattice sums by integrals and using $\sum_y\eta'^3(\rho+\eta')^{-2}e^{-cm_r\rho}\le C/m_r$, $\sum_y\eta'^3\rho^{-3}e^{-cm_r\rho}\le C\log(1/m_r)$; the regions $\rho,|u|<25\eta'$ are bounded by the per-flow sizes. The $\delta\Pi$ terms give $C\eta'^{1/8}(1+\log(1/\eta'))(1+m_r^{-1})$, the $\delta\mathsf a$ terms $C[\eta^{1/4}L_\eta\log(1/\eta')+\eta'^{1/2}\log(1/\eta')m_r^{-1/2}]$, the $\delta\mathsf a_C$ terms $C\eta'^{1/2}(1+\log(1/\eta'))^2$, and $\log(1/m_r)\le\log(1/\eta')$ absorbs the remaining logarithms.
\end{proof}

\begin{lemma}[the internal-$\Gamma$ part]\label{rot:lem:CK7}
$|B^F(x)|\le Ce_r^2$ and $|B^X(x')-B^Y(\tau x')|\le Ce_r^2\eta'^{1/6}(1+\log(1/\eta'))^2$.
\end{lemma}
\begin{proof}
By \Cref{rot:lem:CK4,rot:lem:CK5}(iii), $B^F$ is a finite sum of terms with externals $L_a,L_b\in\{G,\nabla G\}$, an internal pair with at least one $\Gamma$-type line and at most one difference per line, plus the $A^2$ terms with the photon tadpole $\Gamma_{\rm ph}(w,w)$. Per-flow sizes: $|L_a|\le C(\rho+\eta')^{-1-j_a}e^{-\delta'\rho}$ (Q3); a $C$-type internal line with $j$ differences $\le C(|u|+\eta')^{-1-j}$, a $\Gamma$-type one $\le C(1+|u|)^{-1}$ (Q1, Q4). The total scaling dimension is $6-(1+j_a)-(1+j_b)-(1+j_C)\ge1>0$ and every sub-region converges, whence the size bound; the region where $\rho$, $\rho'$ or $|u|$ is $\le M\eta'$ contributes $\le CM\eta'(1+\log(1/\eta'))$. Telescope over the lines. Outside that region the difference elements are $\delta G\le(\eta'/\rho)^{1/2}C\rho^{-1}$ (Q6), $\delta\nabla G\le(\eta'/\rho)^{1/6}\log(1/\eta)\,C\rho^{-2}$ (\Cref{rot:lem:CK2}(c)), $C$-type internal $\le(\eta'/|u|)^{1/2^{j+1}}C|u|^{-1-j}$, $j\le1$ (Q2), $\Gamma$-type internal $\delta\Gamma\le C\eta^{1/2}$, $\delta\nabla\Gamma\le A_1$ (\Cref{rot:lem:CK2}(a)). The gains $(\eta'/\rho)^\theta$ are integrated against the size integrand region by region: $\int\rho^2\rho^{-1}(\eta'/\rho)^{1/2}d\rho\le C\eta'^{1/2}$, $\int\rho^2\rho^{-2}(\eta'/\rho)^{1/6}\log(1/\eta)\,d\rho\le C\eta'^{1/6}\log(1/\eta)$, $\int|u|^2|u|^{-1-j}(\eta'/|u|)^{1/2^{j+1}}d|u|\le C\eta'^{1/4}$, and the $\Gamma$-type gain $C\eta^{1/4}L_\eta(1+|u|)$ is confined by the externals' decay; the joint sub-regions have the margin $1$ of the power counting and every gain exponent is $\le1$, so each is $\le C\eta'^{1/6}(1+\log(1/\eta'))^2$. Inside the near region the difference is at most the sum of the two per-flow sizes, $\le C\eta'(1+\log(1/\eta'))$. The photon-tadpole difference is $\le C_\Gamma\eta^{1/2}$ (Q6), and the $\eta'^2(\nabla G)^2$ term has size $\le Ce_r^2\eta'$.
\end{proof}

\begin{proposition}[the source-linear constant vanishes]\label{rot:prop:CK}
With $c_k$ as in \Cref{rot:lem:C} and $r=k-\lfloor\vr k\rfloor$, $\vr<\frac19$, for $k\ge k_0(a,m,\lmu,\vr)$
\[
|c_k|\le NC_\Gamma\,\eta^{1/2}\delta_r^{-1}+C_1\lambda\,\eta^{1/2}+C_2e^2\,\eta'^{1/8}(1+\delta_r^{-1})(1+\log(1/\eta'))^4\le C\,5^{-(1-9\vr)k/8}(1+k)^4,
\]
with constants depending on $a,m,\lmu,N,\tau$ only (not on $J$, $\psi$, $n$, the torus or $h$). Hence $c_k\to0$.
\end{proposition}
\begin{proof}
Both densities are $Y_{k-r}$-periodic (\Cref{rot:lem:C}) and the cell average of $f^Y$ equals the average of $f^Y\circ\tau$ over the matching $X$-cell (\Cref{rot:lem:vert}(i)), so $|\bar f^X-\bar f^Y|\le\sup_{x'}|f^X(x')-f^Y(\tau x')|$. By \Cref{rot:lem:S1}, $f^F=f^F_0+f^F_\lambda+f^F_{e^2}$. \emph{Term 0}: at zero coarse fields $\langle|\phi'(x)|^2\rangle=NG^F(x,x)=N(C(0)-\Gamma^F(x,x))$, so the free tadpole leaves the residue $f^F_0(x)=-N\Gamma^F(x,x)$ (the subtraction in \cite[(3.21)]{K1} removes the full-propagator expectation, not this), with matched-pair difference $\le NC_\Gamma\eta^{1/2}$ (Q6); in physical units $\Gamma$ carries $\delta_r^{-1}$, and $\eta^{1/2}\delta_r^{-1}=5^{1/2}5^{-(k-3\lfloor\vr k\rfloor)/2}\le C5^{-k/3}$. The order-$\lambda$ term is \Cref{rot:lem:CK3}, $\eta^{1/2}\le C5^{-4k/9}$; the order-$e^2$ term is \Cref{rot:lem:CK5,rot:lem:CK6,rot:lem:CK7} with $e_r^2\delta_r^{-1}=e^2$ and $m_r^{-1}=m^{-1}\delta_r^{-1}$, and $\eta'^{1/8}\delta_r^{-1}\le C5^{-(1-9\vr)k/8}$ because $\vr<\frac19$. The rate is not optimal; only $c_k\to0$ is used.
\end{proof}

% r03_fill_K5a: proof of (K_0), part VII: the inputs of the forward bound (kernel bounds of the hybrid flow, the two stability bounds, the polynomial lower bound at the wall).

\subsubsection*{The inputs of the forward bound}

Throughout $T\in\{T_1,T_2\}$ and $r=r(k)$ is the choice of \Cref{rot:lem:r}; $p(\delta)=b_0(1+\log\delta^{-1})^{\mathrm p}$ is King's small-field function.

\begin{proposition}[kernel bounds of the hybrid flow at a non-zero gauge field]\label{rot:prop:BD}
Let $r\ge2$ and let $A$ be regular in the sense of \cite[Def.~3.2]{K1} on $\Lambda_X$ or on $\Lambda_X/P$ for a torus $P$ fitted by both lattices, with $e_r^\beta\le c_0(a)$.
\begin{enumerate}[label=\textup{(\alph*)},leftmargin=*]
\item \textup{(kernel bounds)} The operators $G_r$, $\Delta^{(r)}$, $C^{(r)}$ of the hybrid flow with the paired contours of \Cref{rot:def:flows} obey the bounds of \cite[Thm.~3.3, (3.5)--(3.8)]{K1} in the periodic form (on the infinite lattice and on tori fitted by both lattices, each in its own quotient distance), with constants depending on $a$, $M$, $\alpha$ only, except the H\"older constant of (3.8), which is $C_\alpha\log^{3\alpha}(1/\eta')$ (\Cref{rot:hyb:rem:Hol}); the differences between two tori, or between a torus and the infinite lattice, with a common fundamental region $F$ that is a union of $(r+1)$-blocks on which the backgrounds agree up to a collar, obey the same bounds with the distance taken in the smaller $d_\wedge$ of the two quotient metrics and the additional factor $e^{-\delta(\dist(\{x,y\},\partial F)+\dist(\supp f,\partial F))}$ for $G_r$ and $e^{-\delta(\dist(z,\partial F)+\dist(w,\partial F))}$ for $\Delta^{(r)}$ and $C^{(r)}$ (\Cref{rot:hyb:lem:domdiff}), in place of the Neumann corrections; the next-step precision $\Delta^{(r)}+a5^{-2}\hat Q^*\hat Q$ obeys $(25\gamma_0^{-1}+a_r^{-1})^{-1}\le\Delta^{(r)}+a5^{-2}\hat Q^*\hat Q\le a_r+a/25$, with $\gamma_0=\gamma_0(a)$ the coercivity constant of \Cref{rot:hyb:lem:coer}, and has an exponentially decaying kernel, and so has its restriction to every subset $\Lambda$ of the coarse lattice, in particular to every union of big blocks, whose inverse $\hat C^{(r)}_\Lambda$ of \cite[(5.53)]{K2} obeys \cite[(5.54)]{K2} and the two-torus clause (\Cref{rot:hyb:prop:CLam}); the fine-slice lines obey \cite[(3.63)--(3.65)]{K1} in the covariant form at regular $A$ (covariant lattice derivatives and transported H\"older quotients, \cite[(2.10)--(2.11)]{Ba3}) and in the plain form at the reduced background of step (5) of the proof of \Cref{rot:thm:A}; and the level-$r$ elements obey $|H^X(x,z)|,|D_AH^X(x,z)|\le Ce^{-\delta_0|x-z|}$ with the H\"older clause, $|B^X(z,w)|=|(A^X)^{-1}(z,w)|\le C(1+|z-w|)^{-1}$ and $|\Gamma^X(x,y)|\le C(1+|x-y|)^{-1}$ with the fine-variable regularity of $H^X$, $B^X$ and $\Gamma^X$ decaying exponentially only at the coarse mass; all constants are uniform in $r$, $k$ and the torus. The same holds at every level $j\le r$ with constants uniform in $j$, and, in the form in which the hybrid lower bound uses it, at the levels $j\in\{0,1\}$. The same statements hold for the standard flow with King's uniform constants (\cite[Thm.~3.3]{K1}, proved in \cite{Ba4}).
\item \textup{(gains)} After the constant-background reduction of \Cref{rot:lem:CB} the difference elements $\Delta\Gamma\defeq\Gamma^X-\Gamma^Y\circ\tau$, $\Delta H$, $\Delta(A^{-1})$ at regular $A$ obey the bounds of \Cref{rot:cor:unif} with the factor $\eta'^{1/2}$; the contour vertices \cite[(3.24)--(3.26)]{K1} at the reduced background differ entrywise by a relative $\eta^{1/2}$ against their size $Ce_rp_r^2$; the constant-background phases \cite[(3.46)]{K1} at matched vertices differ by $\le e\,p_r\lmu^{-1}c_\tau\delta_r^{-1/2}\eta'$; and the per-slice pieces carried by the contour vertex \cite[(3.26)--(3.27)]{K1}, with one end or both ends on contours, obey the slice-$j$ bounds \cite[(3.63)--(3.65)]{K1} in both flows with a relative mismatch $(\eta'/\ell_j)^{1/2}$, and, for a one-ended line with one lattice derivative at its lattice end, with the relative mismatch $(\eta'/\ell_j)^{\theta_\partial}$, $\theta_\partial\defeq\alpha/(2(2+\alpha))\ge\theta_1$ with $\alpha$ the H\"older exponent of \cite[(3.8)]{K1} (\Cref{rot:hyb:lem:DS}).
\end{enumerate}
\end{proposition}
At $A=0$ the proposition is the propagator package (\Cref{rot:prop:fc,rot:lem:H,rot:lem:grad,rot:lem:Gam,rot:lem:slice,rot:cor:unif,rot:cor:C}). At $A\ne0$ the proof, which occupies the rest of this subsection, is the programme of \cite[\S5.3]{K2} written out: coercivity of $-\Delta_A+m_r^2+a_rQ(A)^*Q(A)$ at a regular background by a cover with axis cubes carrying a local gauge for the constant part of the field (the tile properties of \Cref{rot:lem:tile} and the paired contours of \Cref{rot:lem:pair} are what enters: blocks partition the lattice, have bounded diameter, carry contours of bounded length inside the block, and $QQ^*=\mathbf 1$ for every $A$, so that $Q^*Q$ is the projection onto the block-constant mode); Combes--Thomas decay in $L^2$ form \cite{CT73}; the domain-difference lemma for the two-torus clause; Ba{\l}aban's random walk \cite[\S2]{Ba4} with periodic cells fitted by both lattices (\cite[(5.55)--(5.56)]{K2}); the $L^p$ and H\"older bounds of \cite[Lemma~2.2]{Ba4} with the $A=0$ input, the slice clauses stated covariantly; and, for (b), King's constant-background reduction \cite[(3.43)--(3.55)]{K1} applied to both flows on their infinite lattices in the universal-cover form of \Cref{rot:lem:CB}, after which every difference element at matched points is of one of seven types (an $A=0$ line, with the gain $\eta'^{1/2}$ of \Cref{rot:cor:unif}; a contour line, with the kernel difference along the contour plus the flux of the slice kernel through the sliver of \Cref{rot:lem:flux}; a double contour line; a contour vertex; an external lifted field, \Cref{rot:lem:H}; a constant-background phase; a hopping vertex, \Cref{rot:lem:vert}, unchanged), the extension of the lattice kernels off the lattice being the cubic-Hermite cardinal interpolation in the first argument, which decays like the lattice kernel.

% r03_hyb_setup: proof of \Cref{rot:prop:BD}, part I: the hybrid block-spin operators at a gauge field; lengths of the contour pieces; the block-spin algebra for any contour system.
% Input from sec/r03_fill_K5a.tex after the statement of \Cref{rot:prop:BD}. Parts II-VII follow as sec/r03_hyb_coer.tex, r03_hyb_domdiff.tex, r03_hyb_A0.tex, r03_hyb_walk.tex, r03_hyb_ES.tex, r03_hyb_gain.tex.

\paragraph{The hybrid operators at a gauge field.}
Throughout the proof of \Cref{rot:prop:BD} we work in the unit scale of level $r\ge2$ (\cite[(2.20)]{K1}): the physical spacing $\delta_r=5^r\eps_k$ is the unit, the fine spacing is $\eta'=5^{-r}$, $\eta=5\eta'$, and $\ell_j\defeq5^j\eta'$ ($\ell_r=1$). We write $L\defeq5$ for the block factor where a formula is clearer with it, $d=3$, and $f_\alpha=R^{-1}e_\alpha$ for the frame of $\Yk$. The coarse lattice $\Gamma_r=R^{-1}\Z^3$ is common to both flows; $\Gamma_j\defeq R^{-1}(\ell_j\Z^3)$ for $0\le j\le r$, so that $\Gamma_0=\Lambda_Y=R^{-1}(\eta'\Z^3)$ is the fine lattice of the flow $Y$ and $\Gamma_1=R^{-1}(\eta\Z^3)$ is $Y_{k-1}$ rescaled; $\Lambda_X=\eta'\Z^3$ is the fine lattice of the flow $X$ (\Cref{rot:def:match}). Since $5R^{-1}$ is an integer matrix (\Cref{rot:lem:geom}(a)), $\Gamma_1\subset\Lambda_X\cap\Lambda_Y$.

For $1\le j\le r$ and $z\in\Gamma_j$ the \emph{level-$j$ block} of the flow $F$ is
\[
B^F_j(z)\defeq\Lambda_F\cap\big(z+R^{-1}(-\tfrac{\ell_j}2,\tfrac{\ell_j}2)^3\big).
\]
For $F=Y$ this is a cube of $5^{3j}$ points aligned with $(f_\alpha)$; for $F=X$ it is the digitised rotated cube, with no lattice point on its boundary (\Cref{rot:lem:tile}(i)), and $B^X_1(z)=z+\eta'B_0$ is the tile. Both are nested, $B^F_{j+1}(z)=\bigcup_{s\in\{-2,\dots,2\}^3}B^F_j(z+\ell_jR^{-1}s)$, and for $x\in B^F_r(z)$ we write $x_j=x_j(x)$ for the centre of the level-$j$ block containing $x$ ($x_r=z$). Blocks are centred at the sites of the next coarse lattice (\Cref{rot:conv:C}(i)): the root of every block is its centre, in both flows. The matching $\tau:\Lambda_X\to\Lambda_Y$ of \Cref{rot:def:match} maps the tile $B^X_1(x)$ onto the cube $B^Y_1(x)$ for every $x\in\Gamma_1$, with $|\tau x'-x'|\le c_\tau\eta'$, $c_\tau=\sqrt{13}+2\sqrt3<7.1$; points $x''=\tau x'$ are matched. The composite contours are those of \Cref{rot:def:flows}: for the flow $Y$ King's \cite[(2.12)]{K1}, $\Gamma^{Y,(r)}_{z,x}=\Gamma^Y_{z,x_{r-1}}\cup\dots\cup\Gamma^Y_{x_1,x}$, where $\Gamma^Y_{u,v}$ is the L-path that changes the $f_3$-coordinate first, then $f_2$, then $f_1$; for the flow $X$, $\Gamma^{X,(r)}_{z,x'}=\Gamma^X_{z,x_{r-1}}\cup\dots\cup\Gamma^X_{x_2,x_1}\cup\Gamma^X_{x_1,x'}$, where for $i\ge1$ the level-$i$ piece $\Gamma^X_{x_{i+1},x_i}$ is the L-path $\Gamma^Y_{x_{i+1},x_i}$ with every $\eta$-bond $(u,u+\eta f_\alpha)$, $u\in\Gamma_1$, replaced by the fixed staircase $u+\eta'\gamma_\alpha$ (the reversed bond by the reversed staircase), and $\Gamma^X_{x_1,x'}$ is the fixed in-tile path of at most $5$ bonds. The gauge field along a composite contour is read in King's convention \cite[(2.19)]{K1} (\Cref{rot:conv:C}(ii)).

\emph{Backgrounds.} A background is a real bond field $A$ on $\Lambda_F$, identified with a site vector field by $A_{(x,x+\eta'f_\alpha)}=A_\alpha(x)$ (\cite[p.~652]{K1}); $U(A(\Gamma))\defeq\exp[\mathrm ie_rA(\Gamma)]$ with $A(\Gamma)=\sum_{b\in\Gamma}\eta'A_b$ (oriented) and the running charge $e_r=e\,\delta_r^{1/2}$ (\cite[(3.22)]{K1}). \emph{Regular} means \cite[Def.~3.2]{K1}: $|\partial_\mu A_\nu(x)|\le Ce_r^{\beta-1}$ for all $x,\mu,\nu$ and some $\beta>0$. The small-field class at level $r$ is \cite[(3.2)]{K1}: a coarse field $A_r$ on $\Gamma_r$ with $|A_r|\le p_r(\lmu\delta_r)^{-1}$, $|\partial A_r|\le p_r$, $p_r=p(\delta_r)$, lifted to the flow $F$ by \cite[(2.18)]{K1}, $A^{(r),F}\defeq a_rG^{F,\mathrm{ph}}_rQ_F^*A_r$ (the photon fluctuation propagator and the frame-covariant photon average of \Cref{rot:rem:frameav}), plus the Cameron--Martin shift $za_F$ of \Cref{rot:lem:CM}. Both are regular for $\delta_r\le\delta_0$ (\cite[(3.45)]{K1}; \Cref{rot:lem:CM}).

\emph{The pinned operators.} For a flow $F$ and a background $A$ on $\Lambda_F$,
\[
Q_F(A)\phi(z)\defeq\sum_{x\in B^F_r(z)}\eta'^3\,U(A(\Gamma^{F,(r)}_{z,x}))^{-1}\phi(x)\quad(z\in\Gamma_r),\qquad
M_A\defeq-\Delta_A+m_r^2,\qquad G^F_r(A)\defeq\big(M_A+a_rQ_F(A)^*Q_F(A)\big)^{-1},
\]
with $-\Delta_A=D_A^*D_A$ the covariant Laplacian of the fine lattice, $(D_A\phi)(b)=\eta'^{-1}(U(-A_b)\phi(b_+)-\phi(b_-))$, $m_r=m\delta_r$, $a_r=a(1-5^{-2})(1-5^{-2r})^{-1}$, and the pairings $\ip uv=\sum\eta'^3\bar uv$ on $\Lambda_F$ and $\sum_z\bar uv$ on $\Gamma_r$ (unit block volume, $|B^F_r|\eta'^3=1$). This is \cite[(2.11), (2.13)]{K1} in unit scale; the convention $U^{-1}$ makes $Q_F(A)\phi$ transform with the coarse gauge function (\Cref{rot:hyb:lem:alg}). Put
\[
C_A\defeq M_A^{-1},\qquad u^F(A)\defeq C_AQ_F(A)^*,\qquad A^F(A)\defeq\big(a_r^{-1}+Q_F(A)C_AQ_F(A)^*\big)^{-1},\qquad B^F(A)\defeq A^F(A)^{-1},
\]
\[
H^F(A)\defeq a_rG^F_r(A)Q_F(A)^*=u^F(A)A^F(A),\qquad \Gamma^F(A)\defeq C_A-G^F_r(A)=u^F(A)A^F(A)u^F(A)^*=H^FB^FH^{F*}.
\]
By the Woodbury identity these relations are exact for every $A$. $A^F(A)$ is King's effective Laplacian $\Delta^{(r)}(A)=a_r-a_r^2Q_FG_rQ_F^*$ (\cite[(2.14)]{K1}; indeed $a_r(1-a_rQ_FG_rQ_F^*)=a_r(1+a_rQ_FC_AQ_F^*)^{-1}=A^F(A)$), $H^F$ is the lift operator of \cite[(2.18)]{K1} and \Cref{rot:lem:H}, and $C^{(r)}(A)=(A^F(A)+aL^{-2}Q^{(1)}(A)^*Q^{(1)}(A))^{-1}$ is the next fluctuation covariance (\cite[(2.16)]{K1}), $Q^{(1)}$ the one-step average of the coarse lattice. The photon kernels are three copies of the scalar kernels at $A=0$ and mass $\lmu\delta_r$ (\cite[p.~654]{K1}; \Cref{rot:rem:frameav}, \Cref{rot:prop:fc}): no photon kernel sees a background, so the photon sector of \Cref{rot:prop:BD} is the $A=0$ package of \Cref{rot:prop:fc,rot:lem:H,rot:lem:grad,rot:lem:Gam,rot:lem:slice,rot:cor:unif}.

\begin{lemma}[the pieces of the composite contours]\label{rot:hyb:lem:len}
Let $z\in\Gamma_r$, $x'\in B^X_r(z)$, $x''=\tau x'$, and let $x_i$ ($1\le i\le r$) be the common roots.
\begin{enumerate}[label=\textup{(\roman*)},leftmargin=*]
\item For $i\ge1$ the displacement is $x_i-x_{i+1}=\ell_iR^{-1}s$ with $s\in\{-2,\dots,2\}^3$, and every leg of $\Gamma^Y_{x_{i+1},x_i}$ is a concatenation of $|s_\alpha|5^{i-1}$ $\eta$-bonds between points of $\Gamma_1$. The level-$i$ piece has length $\le6\ell_i$ in the flow $Y$ and $\le\frac75\cdot6\ell_i=8.4\,\ell_i$ in the flow $X$; the level-$0$ pieces have length $\le6\eta'$ ($Y$) and $\le5\eta'$ ($X$). The partial contours $\Gamma^{F,(j+1)}_{x_{j+1},x}$ of \cite[(2.19)]{K1} (the pieces of levels $j,j-1,\dots,0$) have length $\le7.5\,\ell_j$ ($Y$) and $\le10.5\,\ell_j$ ($X$); the composite contours have length $\le1.5$ ($Y$) and $\le2.1$ ($X$).
\item Every level-$i$ piece lies inside the open rotated cube $x_{i+1}+R^{-1}(-\frac{\ell_{i+1}}2,\frac{\ell_{i+1}}2)^3$, hence inside $B^F_{i+1}(x_{i+1})$, and the level-$0$ piece inside its first-step cell; so every piece lies in its block and $\Gamma^{F,(r)}_{z,x}\subset B^F_r(z)$.
\item The two composite contours $\Gamma^{X,(r)}_{z,x'}$ and $\Gamma^{Y,(r)}_{z,x''}$ pass through the same points $x_{r-1},\dots,x_1$ and through the same endpoints of $\eta$-bonds; their level-$i$ pieces are within Hausdorff distance $\eta'$ for $i\ge1$, the level-$0$ pieces lie in the ball of radius $\sqrt{13}\eta'$ about $x_1$, and the closed polygon $\Gamma^{X,(r)}_{z,x'}\cup[x',x'']\cup(\Gamma^{Y,(r)}_{z,x''})^{-1}$ bounds a piecewise planar surface $\Sigma(z,x')$, a union of the $\eta$-bond slivers of \Cref{rot:lem:pair} and one polygon inside that ball, of area $\le1.5\eta'+33\eta'^2$. The sub-polygon formed by the partial contours up to level $j+1$ and the matching segment bounds a surface of area $\le1.5\eta'\ell_{j+1}+33\eta'^2$.
\end{enumerate}
\end{lemma}
\begin{proof}
(i) $x_i=\ell_iR^{-1}m_i$ with $m_i\in\Z^3$ the frame coordinates of the block centre and, by nesting, $m_i-5m_{i+1}=:s\in\{-2,\dots,2\}^3$ (the level-$i$ block of $x$ is one of the $125$ sub-blocks of its level-$(i+1)$ block; there are no ties since no lattice point lies on a block face). The L-path from $x_{i+1}$ to $x_i$ has legs $\ell_is_\alpha f_\alpha=5^{i-1}s_\alpha\cdot\eta f_\alpha$, each a run of $|s_\alpha|5^{i-1}$ $\eta$-bonds starting at a point of $\Gamma_1$ (because $x_{i+1}\in\Gamma_{i+1}\subset\Gamma_1$ and $\eta f_\alpha\in\Gamma_1$). Lengths: $\sum_\alpha|s_\alpha|\ell_i\le6\ell_i$; the staircases $\gamma_1,\gamma_2$ have $7$ fine steps for an $\eta$-bond of length $5\eta'$ and $\gamma_3$ has $5$; the in-tile piece has at most $5$ bonds (\Cref{rot:lem:tile}(ii)) and the level-$0$ $Y$-piece at most $6$. Summing, $\sum_{i=1}^{j}8.4\ell_i+5\eta'\le8.4\cdot\frac54\ell_j(1-5^{-j})+5\eta'<10.5\,\ell_j$ and $6\sum_{i=0}^{j}\ell_i\le7.5\,\ell_j$; for $j=r-1$ this gives $2.1$ and $1.5$ (\Cref{rot:lem:pair}). (ii) The endpoints of every $\eta$-bond of $\Gamma^Y_{x_{i+1},x_i}$ are points of $\Gamma_1$ inside the open cube $x_{i+1}+R^{-1}(-\frac{\ell_{i+1}}2,\frac{\ell_{i+1}}2)^3$; in the frame coordinates $(f_\alpha\cdot(\,\cdot\,-x_{i+1}))_\alpha$ they lie in $\eta\Z^3\cap(-\frac{\ell_{i+1}}2,\frac{\ell_{i+1}}2)^3$, and since $\frac{\ell_{i+1}}2=\frac{5^{i+1}}2\eta'$ is an odd multiple of $\frac{\eta'}2$ while the points are multiples of $5\eta'$, each coordinate is at distance $\ge\frac52\eta'$ from the faces. The segment between two such points keeps that distance, and the staircase deviates from the segment by $\le0.6\eta'<\frac52\eta'$ (\Cref{rot:lem:pair}), so it stays in the open cube; its points lie in $\Lambda_X$, hence in $B^X_{i+1}(x_{i+1})$. The in-tile piece lies in the tile (\Cref{rot:lem:tile}(ii)), which is a sub-block. (iii) The Hausdorff distances and the containment of the level-$0$ pieces are \Cref{rot:lem:pair} ($|x'-x_1|\le\sqrt{13}\eta'$, $|x''-x_1|\le2\sqrt3\eta'$, both level-$0$ pieces monotone inside their cells); the two composite contours pass through the same points $x_{r-1},\dots,x_1$ and through the same $\eta$-bond endpoints, so the closed polygon decomposes into the $\eta$-bond slivers (area $\le5\eta'^2$ each, at most $1.5/\eta=0.3/\eta'$ of them, total $\le1.5\eta'$) and the polygon $\Gamma^X_{x_1,x'}\cup[x',x'']\cup(\Gamma^Y_{x_1,x''})^{-1}$ inside the ball of radius $\sqrt{13}\eta'$, of perimeter $\le(5+7.1+6)\eta'$, which the cone from $x_1$ spans with area $\le\frac12\cdot18.1\cdot\sqrt{13}\,\eta'^2<33\eta'^2$. For the partial contours up to level $j+1$ the $Y$-length is $\le7.5\ell_j=1.5\ell_{j+1}$, so there are $\le1.5\ell_{j+1}/\eta$ $\eta$-bonds and the slivers have total area $\le5\eta'^2\cdot1.5\ell_{j+1}/\eta=1.5\eta'\ell_{j+1}$.
\end{proof}

\begin{lemma}[block-spin algebra for any contour system]\label{rot:hyb:lem:alg}
For every background $A$ on $\Lambda_X$ and every system of contours $\Gamma_{z,x}\subset B^X_r(z)$ from the root $z$ to the point $x$:
\begin{enumerate}[label=\textup{(\roman*)},leftmargin=*]
\item $Q_X(A)Q_X(A)^*=\mathbf 1$, and $P_X(A)\defeq Q_X(A)^*Q_X(A)=\sum_zw_zw_z^\dagger$ with $w_z(x)\defeq U(A(\Gamma_{z,x}))\mathbf 1_{B^X_r(z)}(x)$ is an orthogonal projection ($\|w_z\|=1$, orthogonal supports);
\item for a gauge function $\lambda$ on $\Lambda_X$, $Q_X(A+\partial\lambda)(e^{\mathrm ie_r\lambda}\phi)=e^{\mathrm ie_r\lambda|_{\Gamma_r}}Q_X(A)\phi$, hence $G^X_r(A+\partial\lambda)=e^{\mathrm ie_r\lambda}G^X_r(A)e^{-\mathrm ie_r\lambda}$, and likewise for $H^X$, $A^X$, $B^X$, $\Gamma^X$, $C^{(r)}$ with the coarse phases;
\item for a constant field $A\equiv B_0$, $U(B_0(\Gamma_{z,x}))=e^{\mathrm ie_rB_0\cdot(x-z)}$ for every contour, and $G^X_r(B_0)(x,y)=e^{\mathrm ie_rB_0\cdot x}G^X_r(0)(x,y)e^{-\mathrm ie_rB_0\cdot y}$ on $\Lambda_X$ (on a torus: the twisted periodisation of \Cref{rot:hyb:lem:A0}(c));
\item with the contours of \Cref{rot:def:flows}, $Q_X(A)=Q^{\mathrm{rot}}_{r-1}(A)\circ\hat Q(A)$ under the convention \cite[(2.19)]{K1}, and the composite blocks are unions of tiles (\Cref{rot:lem:tile}(iv)).
\end{enumerate}
\end{lemma}
\begin{proof}
(i) is \Cref{rot:lem:tile}(iii),(iv): $|U|=1$ and $|B^X_r(z)|\eta'^3=1$. (ii) $A(\Gamma_{z,x})\mapsto A(\Gamma_{z,x})+\lambda(x)-\lambda(z)$ for any path from $z$ to $x$, so $U(A(\Gamma_{z,x}))^{-1}e^{\mathrm ie_r\lambda(x)}=e^{\mathrm ie_r\lambda(z)}U(A(\Gamma_{z,x}))^{-1}$, and $-\Delta_{A+\partial\lambda}=e^{\mathrm ie_r\lambda}(-\Delta_A)e^{-\mathrm ie_r\lambda}$. (iii) A constant field is $\partial(B_0\cdot x)$ on $\Lambda_X$. (iv) The level-$j$ pieces are the contours of the steps of the flow $Y$ with the bonds re-drawn; iterating the one-step transformations produces exactly the phases of \cite[(2.19)]{K1} along the composite contour (\Cref{rot:conv:C}(ii); \cite[(3.33)--(3.34)]{Ba1}).
\end{proof}

The pairing of the contours is a property of \Cref{rot:def:flows}, not a further hypothesis; a hybrid transformation whose first-step contour were a fine-axis shortest path from the level-$r$ root across the whole block would be a different transformation, for which the contour phases of the two flows differ at order one (\Cref{rot:lem:flux}).

% r03_hyb_coer: proof of \Cref{rot:prop:BD}, part II: coercivity at a regular background, L^2 bounds, Combes--Thomas decay, the decay of C^{(r)}.

\paragraph{Coercivity and decay in $L^2$ form.}
The proof of \Cref{rot:prop:BD}(a) follows \cite[\S2]{Ba4} and \cite[\S5.3]{K2} with two changes: the lower bound on the quadratic form is proved directly at a regular background by a cover with axis cubes (so that no small-domain reduction is needed for the $L^2$ theory), and the local cells of the random walk are periodic cells fitted by both lattices (\cite[(5.55)]{K2}), on which the $A=0$ kernels are the twisted periodisations of the explicit kernels of \Cref{rot:lem:H,rot:lem:slice}. On the tori of the chain the periods lie in $10\Z^3$ in the unit scale of every level $r\le k$ ($P_J=50MJ\,5^{k-r}\Z^3$ and $R^{-1}P_J=10MJ\,5^{k-r}(5R^{-1})\Z^3$, \Cref{rot:lem:adm}(iii)), with sup-systole $s_\infty(P)\ge50$ for $r\le k-1$.

\begin{lemma}[coercivity at a regular background]\label{rot:hyb:lem:coer}
Let $\Omega$ be $\Lambda_X$ or $\Lambda_X/P$ with $P$ fitted by both lattices and $s_\infty(P)\ge10$ in unit scale. There are $c_0(a)>0$ and $\gamma_0(a)>0$ such that for every regular $A$ on $\Omega$ with $e_r^\beta\le c_0$ and every $\phi$,
\[
\ip{\phi}{(-\Delta_A+m_r^2+a_rP_X(A))\phi}\ \ge\ \gamma_0\|\phi\|^2 .
\]
Explicitly $\gamma_0=\frac1{2m_c}\min\{\frac18,\frac{a'}{64(8a'+1)}\}$ with $a'=a(1-5^{-2})\le a_r\le a$ and $m_c$ the multiplicity of the cover below; $m_c=8$ when $P\subset2\Z^3$, as on the tori of the chain.
\end{lemma}
\begin{proof}
\emph{Cover.} Let $\{\Delta_i\}$ be the axis cubes $y_i+[0,4)^3$ with $y_i\in\Lambda_c\defeq2\Z^3+P$ ($P=0$ on $\Lambda_X$); $\Lambda_c$ is $P$-invariant, so the family descends to a finite family on the torus, and a cube does not wrap since every non-zero period has $|\gamma|_\infty\ge10$. Every point, every bond and every block lies in at most $m_c$ cubes, $m_c=8$ if $P\subset2\Z^3$ and $m_c\le64$ in general ($\Lambda_c\subset\Z^3$). Every $\Delta_i$ contains a full block: the covering radius of $\Gamma_r=R^{-1}\Z^3$ is $\frac{\sqrt3}2$, so some $z_i\in\Gamma_r$ has $|z_i-c_i|\le\frac{\sqrt3}2$ ($c_i$ the centre of $\Delta_i$), and $B^X_r(z_i)\subset\{|x-z_i|<\frac{\sqrt3}2\}\subset\{|x-c_i|<\sqrt3\}\subset\Delta_i$ because $\sqrt3<2$. Dropping the non-negative terms of the quadratic form that belong to bonds joining different cubes, to the mass and to the pins of the other blocks, and dividing by the multiplicity,
\[
\ip{\phi}{(-\Delta_A+m_r^2+a_rP_X(A))\phi}\ \ge\ \tfrac1{m_c}\sum_i\Big[\sum_{b\subset\Delta_i}\eta'^3|D_A\phi(b)|^2+a_r|\ip{w_{z_i}}{\phi}|^2\Big]=:\tfrac1{m_c}\sum_iq_i(\phi).
\]
\emph{Local gauge.} On the cube $\Delta_i$ (contractible) write $A=A_0^{(i)}+A'^{(i)}$ with $A_0^{(i)}\defeq A(c_i)$ constant and, by regularity along lattice paths of $\ell^1$-length $\le12$, $|A'^{(i)}|\le12\,Ce_r^{\beta-1}$ on $\Delta_i$. The gauge function $\lambda_i(x)=A_0^{(i)}\cdot x$ removes $A_0^{(i)}$ from both terms of $q_i$ (\Cref{rot:hyb:lem:alg}(ii),(iii), applied to the bonds and the block inside $\Delta_i$), so $q_i(\phi)=q_i^{A'}(\tilde\phi)$ with $\tilde\phi=e^{-\mathrm ie_r\lambda_i}\phi$, $|\tilde\phi|=|\phi|$.
\emph{Perturbation.} $D_{A'}\phi(b)=U(-A'_b)D_0\phi(b)+\eta'^{-1}(U(-A'_b)-1)\phi(b_-)$, hence $\|D_{A'}\phi\|_{\Delta_i}\ge\|D_0\phi\|_{\Delta_i}-e_r|A'|_\infty\|\phi\|_{\Delta_i}$ and $\|D_{A'}\phi\|^2\ge\frac34\|D_0\phi\|^2-3e_r^2|A'|_\infty^2\|\phi\|^2$. For the pin, $w_z(A')-w_z(0)$ has sup norm $\le e_r|A'(\Gamma)|\le2.1\,e_r|A'|_\infty$ (\Cref{rot:hyb:lem:len}(i)), so $\big||\ip{w_z(A')}{\phi}|^2-|\ip{w_z(0)}{\phi}|^2\big|\le(4.2e_r|A'|_\infty+4.41e_r^2|A'|_\infty^2)\|\phi\|_B^2$.
\emph{The free cube.} With $\psi\defeq\tilde\phi-c$ on $\Delta_i$, $c$ the mean, $\sum_{b\subset\Delta_i}\eta'^3|D_0\tilde\phi(b)|^2\ge\lambda_1\|\psi\|^2_{\Delta_i}$ with $\lambda_1=4\eta'^{-2}\sin^2(\pi\eta'/8)\ge\frac14$ (the first Neumann eigenvalue of the cube of $4/\eta'$ points per side; $\sin t\ge\frac2\pi t$), and, since $|B_i|\eta'^3=1$, $|\ip{w_{z_i}}{\tilde\phi}|^2=|c+\ip{\mathbf 1_{B_i}}{\psi}|^2\ge(1-t)|c|^2-(\frac1t-1)\|\psi\|^2_{B_i}$ for $0<t<1$. Choosing $\frac1t-1=\frac{\lambda_1}{2a_r}$,
\[
\lambda_1\|\psi\|^2+a_r|\ip{w_{z_i}}{\tilde\phi}|^2\ \ge\ \tfrac{\lambda_1}2\|\psi\|^2+\tfrac{a_r\lambda_1}{2a_r+\lambda_1}|c|^2\ \ge\ \gamma_1\|\tilde\phi\|^2_{\Delta_i},\qquad
\gamma_1\defeq\min\Big\{\tfrac{\lambda_1}2,\tfrac{a_r\lambda_1}{64(2a_r+\lambda_1)}\Big\},
\]
using $\|\tilde\phi\|^2_{\Delta_i}=\|\psi\|^2+64|c|^2$. \emph{Assembly.} $q_i(\phi)\ge\frac34\gamma_1\|\phi\|^2_{\Delta_i}-(3e_r^2|A'|^2+a_r(4.2e_r|A'|+4.41e_r^2|A'|^2))\|\phi\|^2_{\Delta_i}\ge\frac12\gamma_1\|\phi\|^2_{\Delta_i}$ once $e_r^\beta\le c_0(a)$, and $\sum_i\|\phi\|^2_{\Delta_i}\ge\|\phi\|^2$. With $\lambda_1\ge\frac14$, $\gamma_1\ge\min\{\frac18,\frac{a'}{64(8a'+1)}\}$.
\end{proof}

This is \cite[(5.59)--(5.60)]{K2} (``cover $\square$ by cubes $\{\Delta_i\}$ on the $\eta$-lattice with sides parallel to the coordinate planes \dots\ each cube contains fully at least one block'') written out, with the local gauge step that \cite[p.~345]{K2} leaves to ``$A=B+A'$ with $B$ constant and $A'$ small'' made explicit; it is also \cite[(1.8)]{Ba4} for the hybrid flow. No Poincar\'e inequality on the digitised rotated cube is needed: only axis cubes of the fine lattice enter.

\begin{corollary}[$L^2$ bounds]\label{rot:hyb:cor:L2}
Under the hypotheses of \Cref{rot:hyb:lem:coer}, $T\defeq-\Delta_A+m_r^2+a_rP_X(A)\ge\gamma_0$, so $\|G^X_r(A)\|\le\gamma_0^{-1}$, $\|D_AG^X_r(A)^{1/2}\|\le1$, $\|G^X_r(A)^{1/2}D_A^*\|\le1$, $\|D_AG^X_r(A)D_A^*\|\le1$ and $0\le G^X_r(A)\le\gamma_0^{-1}$ (\cite[Lemma~2.1]{Ba4}, \cite[Lemma~5.7]{K2}). The same holds, on $\Lambda_X$ and on a torus with $s_\infty(P)\ge10L$ in unit scale (true on the tori of the chain for $r\le k-1$), for the operator with the level-$(r+1)$ blocks, $-\Delta_A+m_r^2+a_{r+1}L^{-2}P_{r+1}(A)$ (the pin at scale $L$), with $\gamma_0$ replaced by $\gamma_0L^{-2}$.
\end{corollary}
\begin{proof}
$D_A^*D_A\le T$. For the $(r+1)$-blocks repeat the proof with cubes of side $4L$, corners in $2L\Z^3+P$ (multiplicity $8$ if $P\subset10\Z^3$, as for $P_J$ and $R^{-1}P_J$; periods with $|\gamma|_\infty\ge10L$) and blocks of $L^3$ unit cells (\cite[p.~588]{Ba4}: ``replace $\min\{\pi^2,a_k\}$ by $\min\{\pi^2L^{-2},a_{k+1}L^{-2}\}$'').
\end{proof}

\begin{proposition}[Combes--Thomas decay]\label{rot:hyb:prop:CT}
Under the hypotheses of \Cref{rot:hyb:lem:coer} there are $\theta_0,c_1>0$ depending on $a$ only such that for all $f,f'$ supported in unit cubes,
\[
|\ip f{G^X_r(A)f'}|,\ |\ip f{D_AG^X_r(A)f'}|,\ |\ip f{G^X_r(A)D_A^*f'}|\ \le\ c_1e^{-\theta_0\dist(\supp f,\supp f')}\|f\|_2\|f'\|_2 ,
\]
$\dist$ the quotient distance of $\Omega$. Hence (\cite[(3.6)]{K1}) $|A^X(A)(z,w)|=|a_r\delta_{zw}-a_r^2\ip{w_z}{G^X_r(A)w_w}|\le Ce^{-\theta_0|z-w|}$. The same holds for the operator with the $(r+1)$-pin of \Cref{rot:hyb:cor:L2}, with $\theta_0,c_1$ depending on $a$ and $L$.
\end{proposition}
\begin{proof}
For a weight $\psi$ with $|\psi(x)-\psi(x')|\le|x-x'|$ and $0\le t\le\theta_0$ put $T_\psi\defeq e^{t\psi}Te^{-t\psi}$. On $\Lambda_X$ one may take the linear weight $\psi=\hat n\cdot x$; on a torus $\vartheta\cdot x$ is not a function, and $\psi=\min\{\dist(\cdot,\supp f'),R\}$ is used, with the quotient distance; the estimates below use only the Lipschitz bound on bonds and inside blocks. The hopping part changes by a form-bounded perturbation: $e^{t\psi}D_Ae^{-t\psi}=D_A+B_t$, with $B_t$ the bond multiplication by $\eta'^{-1}(e^{-t(\psi(b_+)-\psi(b_-))}-1)$ composed with the covariant shift (bond functions weighted at $b_-$), of norm $\|B_t\|\le\sqrt3\,te^{t\eta'}$ (every site is the end $b_+$ of three bonds); with $e^{t\psi}D_A^*e^{-t\psi}=(D_A+B_{-t})^*$ this gives $|\ip\phi{(T_\psi-T)\phi'}|\le Ct(\|D_A\phi\|\|\phi'\|+\|\phi\|\|D_A\phi'\|)+Ct^2\|\phi\|\|\phi'\|$ for $t\eta'\le1$. The pin changes by $\sum_zw_z(x)\overline{w_z(y)}(e^{t(\psi(x)-\psi(y))}-1)$ with $|\psi(x)-\psi(y)|\le|x-y|\le\sqrt3$ inside a block, an operator of norm $\le2\sqrt3\,te^{\sqrt3t}$ (one rank-one piece per block, orthogonal supports). Hence $\|T^{-1/2}(T_\psi-T)T^{-1/2}\|\le Ct\gamma_0^{-1/2}+Ct^2\gamma_0^{-1}\le\frac12$ for $t\le\theta_0(a)$, using $\|D_AT^{-1/2}\|\le1$ (\Cref{rot:hyb:cor:L2}); so $T_\psi=T^{1/2}(1+T^{-1/2}(T_\psi-T)T^{-1/2})T^{1/2}$ is invertible with $\|T_\psi^{-1}\|\le2\gamma_0^{-1}$. Then $\ip f{G^X_rf'}=\ip{e^{-t\psi}f}{T_\psi^{-1}e^{t\psi}f'}$, since $G^X_r=e^{-t\psi}T_\psi^{-1}e^{t\psi}$; choose $t=\theta_0$ and $\psi=\min\{\dist(\cdot,\supp f'),R\}$ with $R\ge\dist(\supp f,\supp f')$ (or, on $\Lambda_X$, $\psi=\hat n\cdot x$ with $\hat n$ the unit vector from $\supp f'$ to $\supp f$), and use $\sup_{\supp f}e^{-t\psi}\sup_{\supp f'}e^{t\psi}=e^{-\theta_0\dist(\supp f,\supp f')}$ for the distance weight ($\psi=0$ on $\supp f'$), $\le e^{-\theta_0(\dist-2\sqrt3)}$ for the linear one. The derivative versions use $\|D_AT_\psi^{-1}\|,\|T_\psi^{-1}D_A^*\|\le C$ (the same algebra). For $A^X$: $w_z,w_w$ have unit norm and supports in blocks of diameter $\le\sqrt3$. The same proof with the $(r+1)$-pin (blocks of diameter $\le\sqrt3L$, \Cref{rot:hyb:cor:L2}) gives the bounds for that operator.
\end{proof}

\begin{corollary}[{\cite[(3.6)]{K1} for $C^{(r)}$}]\label{rot:hyb:cor:Ck}
$|C^{(r),X}(A)(z,w)|\le Ce^{-\delta_0d(z,w)}$ in the quotient distance of the domain. For two domains $\Omega_1$, $\Omega_2$ as in \Cref{rot:hyb:lem:domdiff} (a common fundamental region $F$, a union of $(r+1)$-blocks, on which the backgrounds agree up to a collar of width $W_0$) and $y,y'\in\Gamma_r\cap F$,
\[
|C^{(r)}(\Omega_1,A_1;y,y')-C^{(r)}(\Omega_2,A_2;y,y')|\le Ce^{-\delta(d_\wedge(y,y')+\dist(y,\partial F)+\dist(y',\partial F))},
\]
$d_\wedge$ the smaller of the two quotient distances; the same for $\Delta^{(r)}=A^X$.
\end{corollary}
\begin{proof}
\cite[\S3]{Ba4} writes $C^{(r)}(\Omega,A;y,y')$ through the Gaussian integral (3.1), factorises it over the next-level blocks after the gauge transformation ``$\psi(y)\to U(A(\Gamma_{y,z}))\psi(y)$'' (p.~587), and arrives at the identity (3.8): $C^{(r)}$ is a bounded combination of $\ip{q_r(y)}{G_r(\Omega,\Lambda,A)q_r(y')}$ and $\ip{q_{r+1}(y)}{\cdot\,q_{r+1}(y')}$, with $q_r(y)$, $q_{r+1}(y)$ the block modes (3.7) of unit and $L^{-d/2}$ norm, where $G_r(\Omega,\Lambda,A)$ is the fluctuation propagator pinned on the $(r+1)$-blocks (3.5). The two inputs of this algebra are that the blocks partition the lattice and that every contour runs from the block root to the point, both true for the hybrid flow (\Cref{rot:hyb:lem:alg}(i),(iv), \Cref{rot:lem:tile}); the decay of $G_r(\Omega,\Lambda,A)$ in the $L^2$-weighted form is \Cref{rot:hyb:prop:CT} for the $(r+1)$-pin. This is what \cite[pp.~345--346]{K2} says: ``To prove Proposition 5.6, we use the method presented in Chap.~3 of [Ba4]. This method requires only $L^2$-bounds on the operators $G_k(\Omega,A)$.'' The two-domain bound is not \cite[(3.9)--(3.10)]{Ba4}, which change the pin region inside one domain: on each domain (3.6) gives $C^{(r)}=V^*G_r(\Omega,\Omega^{(r)},A)V-(\beta/a_r)P+a_r^{-1}$ with $G_r(\Omega,\Omega^{(r)},A)$ the $(r+1)$-pinned operator (3.5), $V=Q_r^*(1-\beta P)$, $\beta=a_{r+1}L^{-2}/a_r$ ($Q_r=Q_X(A)$, $Q=Q^{(1)}(A)$, $P=Q^*Q$), whose local terms and block functions agree on the common region, so that $\delta C^{(r)}(y,y')=\ip{v(y)}{\delta G\,v(y')}$ there; $\delta G$ for the $(r+1)$-pin obeys the weighted bound \eqref{rot:hyb:eq:B1} of the geometric resolvent identity $\ip{f'}{(G_2-G_1)f}=\ip{f'}{G_1[H_2,\chi]G_2f}$ with \Cref{rot:hyb:prop:CT} on each side (\Cref{rot:hyb:lem:domdiff}).
\end{proof}

% r03_hyb_domdiff: proof of \Cref{rot:prop:BD}, part III: the domain-difference lemma (two tori, or a torus against the infinite lattice) and the restricted precision inverses on unions of big blocks.

\begin{lemma}[domain difference]\label{rot:hyb:lem:domdiff}
Let $\Omega_1=\Lambda_X/P_1$ be a torus fitted by both lattices with $s_\infty(P_1)\ge10L$ in unit scale, and $F$ a fundamental region of $P_1$ that is a union of $(r+1)$-blocks. Let $\Omega_2$ be $\Lambda_X$, or a torus $\Lambda_X/P_2$ fitted by both lattices into which $F$ embeds (the projection is injective on $F$); $F$ is identified with its images. Let $A_i$ be regular on $\Omega_i$ with $e_r^\beta\le c_0(a)$, and let $A_1=A_2$ on every bond with both ends in $F_0\defeq\{x\in F:\dist(x,\partial F)\ge W_0\}$, for a fixed $W_0\ge0$. Put $H_i\defeq-\Delta_{A_i}+m_r^2+a_{r+1}L^{-2}P_{r+1}(A_i)$ on $\Omega_i$ (the $(r+1)$-pin, $P_{r+1}=Q_{r+1}^*Q_{r+1}$ with the level-$(r+1)$ blocks) and $G_i\defeq H_i^{-1}$. Let $d_i$ be the quotient metric of $\Omega_i$ and $d_\wedge\defeq\min\{d_1,d_2\}$ on $F$; distances to $\partial F$ and to $F_0^c$ are taken in $\Lambda_X$, for a lift of $F$. There are $C,\delta>0$ depending on $a,L,W_0$ only such that for $f,f'$ supported in unit cubes in $F$
\begin{equation}\label{rot:hyb:eq:B1}
|\ip{f'}{(G_1-G_2)f}|\le C\exp\!\big[-\delta\big(d_\wedge(\supp f,\supp f')+\dist(\supp f,\partial F)+\dist(\supp f',\partial F)\big)\big]\|f\|\|f'\|,
\end{equation}
and, for $y,y'\in\Gamma_r\cap F$,
\begin{equation}\label{rot:hyb:eq:B2}
|C^{(r)}(\Omega_1,A_1;y,y')-C^{(r)}(\Omega_2,A_2;y,y')|\le C\exp\!\big[-\delta\big(d_\wedge(y,y')+\dist(y,\partial F)+\dist(y',\partial F)\big)\big].
\end{equation}
The same bound \eqref{rot:hyb:eq:B2} holds for $\Delta^{(r)}=A^X$, from \eqref{rot:hyb:eq:B1} at the $r$-pin. In particular $\delta C^{(r)}$ and $\delta\Delta^{(r)}$ carry the factor $e^{-\delta\dist(\{y,y'\},\partial F)}$ of \cite[Thm.~3.3]{K1}, in the metric $d_\wedge$.
\end{lemma}

The metric is forced: for $y$ near one face of $F$ and $y'$ near the opposite face, $d_1(y,y')=O(1)$, so $C^{(r)}(\Omega_1,A_1;y,y')$ is of order one, while $C^{(r)}(\Omega_2,A_2;y,y')$ is small when $\Omega_2=\Lambda_X$; no bound of the form \eqref{rot:hyb:eq:B2} with the distance $|y-y'|$ of the lift in place of $d_\wedge$ can hold. The kernel bounds on each domain separately are in the quotient distance of that domain (\Cref{rot:hyb:prop:CT}).

\begin{proof}
\emph{Step 0 (decay on each domain).} By \Cref{rot:hyb:prop:CT} at the $(r+1)$-pin, with the distance weight on a torus and the linear weight on $\Lambda_X$, for $f,f'$ supported in unit cubes and $i=1,2$,
$|\ip f{G_if'}|,\ |\ip f{D_{A_i}G_if'}|,\ |\ip f{G_iD_{A_i}^*f'}|\le c_1e^{-\theta_0d_i(\supp f,\supp f')}\|f\|\|f'\|$, $\theta_0=\theta_0(a,L)$.

\emph{Step 1 (locality).} Put $r_1\defeq1+\sqrt3L$ and $S\defeq\{x\in F_0:\dist(x,F_0^c)\ge r_1\}$. For $u$ supported in $S$, $H_1u=H_2u$ (both supported in $F_0$, identified). Indeed $(H_iu)(x)$ involves $u$ at the bond neighbours of $x$ and on the $(r+1)$-block of $x$, the bond variables of the bonds meeting $\supp u$ (both ends in $F_0$), and the block functions $w_z$ of the blocks meeting $S$; such a block lies in $F$ ($F$ is a union of $(r+1)$-blocks) within $\sqrt3L$ of $S$, hence in $F_0$, where $A_1=A_2$; its contours lie inside it (\Cref{rot:hyb:lem:alg}), so the contour phases of $w_z$ agree.

\emph{Step 2 (geometric resolvent identity).} Let $\chi(x)\defeq\min\{1,\max\{0,\dist(x,F_0^c)-r_1\}\}$ (Lipschitz constant $1$), so $\supp\chi\subset S$ and $\chi=1$ on $S''\defeq\{x:\dist(x,F_0^c)\ge r_1+1\}$. Let $f,f'$ be supported in $S''$. Then $\chi G_2f$ is supported in $S$, so by Step~1 $H_1(\chi G_2f)=H_2(\chi G_2f)=\chi f+[H_2,\chi]G_2f=f+[H_2,\chi]G_2f$. The function $[H_2,\chi]G_2f$ vanishes at points farther than $r_1$ from the set where $\chi$ is not locally constant, so it is supported in $F_0$ and is a function on $\Omega_1$. Applying $G_1$ and pairing with $f'$ ($\chi f'=f'$, $\chi$ real),
\begin{equation}\label{rot:hyb:eq:B3}
\ip{f'}{G_2f}-\ip{f'}{G_1f}=\ip{f'}{G_1[H_2,\chi]G_2f}.
\end{equation}

\emph{Step 3 (the commutator).} With $(D_A\phi)(b)=\eta'^{-1}(U(-A_b)\phi(b_+)-\phi(b_-))$, $D_A(\chi u)(b)=\chi(b_-)(D_Au)(b)+(Mu)(b)$, $(Mu)(b)\defeq\eta'^{-1}(\chi(b_+)-\chi(b_-))U(-A_b)u(b_+)$; taking adjoints, $\chi D_A^*-D_A^*\chi_-=M^*$, hence $[-\Delta_A,\chi]=D_A^*M-M^*D_A$. Since $|\chi(b_+)-\chi(b_-)|\le\eta'$ and every site is the end $b_+$ of three bonds, $\|M\|\le\sqrt3$. For the pin, $P_{r+1}=\sum_zw_zw_z^*$ with $\|w_z\|=1$ and $\supp w_z\subset B_{r+1}(z)$ (\Cref{rot:hyb:lem:alg}), so $[P_{r+1},\chi]=\sum_z\big(w_zw_z^*(\chi-\chi(z))-(\chi-\chi(z))w_zw_z^*\big)$, each term acting in $\ell^2(B_{r+1}(z))$ with norm $\le2\sup_{B_{r+1}(z)}|\chi-\chi(z)|\le\sqrt3L$, and zero unless $\chi$ is non-constant on the block; $m_r^2$ commutes with $\chi$. All these terms live on bonds and blocks meeting the slab $\{r_1\le\dist(\cdot,F_0^c)\le r_1+1\}$, hence on the set $Y\defeq\{x:\tfrac12\le\dist(x,F_0^c)\le r_1+2+\sqrt3L\}\subset F_0$ (the pieces lie within $\sqrt3L+\eta'$ of the slab and $r_1-\sqrt3L-\eta'>\tfrac12$), on whose bonds $D_{A_1}=D_{A_2}$.

\emph{Step 4 (bound of \eqref{rot:hyb:eq:B3}).} Write $\ip{f'}{G_1D^*MG_2f}=\ip{DG_1f'}{MG_2f}$, $\ip{f'}{G_1M^*DG_2f}=\ip{MG_1f'}{DG_2f}$ and the pin terms as $\ip{G_1f'}{[w_zw_z^*,\chi]G_2f}$. Split $Y$ into unit cubes $Q$ (every bond and block meets at most $(\sqrt3L+2)^3$ of them) and apply the Cauchy--Schwarz inequality cube by cube, with $\|\mathbf 1_QDG_2f\|=\sup\{|\ip h{DG_2f}|:h\in\ell^2(Q),\|h\|=1\}$ and the same without $D$, and Step~0 on $\Omega_1$ for the left factor and on $\Omega_2$ for the right factor. Since $\dist(x,\partial F)\le\dist(x,F_0^c)+W_0$, $Y\subset K\defeq\{x\in F:\dist(x,\partial F)\le W_0+r_2\}$ with $r_2\defeq2r_1+2$, and
\begin{equation}\label{rot:hyb:eq:B4}
|\ip{f'}{(G_2-G_1)f}|\le C_L\sum_{Q\subset K}e^{-\theta_0(d_1(\supp f',Q)+d_2(Q,\supp f))}\|f\|\|f'\|.
\end{equation}
\emph{Step 5 (geometry).} For $x\in F$ and $q\in F$, a shortest lattice path in $\Omega_i$ from $x$ to $q$, lifted from $x$, either stays in the lift of $F$ (the projection is injective on $F$), and then has length $\ge|x-q|\ge\dist(x,\partial F)-\dist(q,\partial F)$, or it leaves it, and then has length $\ge\dist(x,\partial F)$. Hence $d_i(x,q)\ge\dist(x,\partial F)-\dist(q,\partial F)$, and for $q\in K$, $d_i(x,q)\ge\dist(x,\partial F)-W_0-r_2-1$ (the $1$ for the unit cubes). Splitting $\theta_0(d_1+d_2)=\frac{\theta_0}2d_1+\frac{\theta_0}2(d_1+d_2)$ in \eqref{rot:hyb:eq:B4},
\[
\theta_0\big(d_1(\supp f',Q)+d_2(Q,\supp f)\big)\ge\tfrac{\theta_0}2d_1(\supp f',Q)+\tfrac{\theta_0}2\big(\dist(\supp f,\partial F)+\dist(\supp f',\partial F)\big)-\theta_0(W_0+r_2+3),
\]
and $\sum_Qe^{-(\theta_0/2)d_1(\supp f',Q)}\le C(\theta_0)$, since the unit cubes of a fundamental region at $d_1$-distance in $[n,n+1)$ from a unit cube number at most $C(n+2)^3$ (choose for each the lift nearest to a fixed lift of $\supp f'$). So, for $f,f'$ supported in $S''$,
\[
\text{(ii)}\qquad|\ip{f'}{(G_1-G_2)f}|\le C'e^{-(\theta_0/2)(\dist(\supp f,\partial F)+\dist(\supp f',\partial F))}\|f\|\|f'\|.
\]
Since $\dist(x,F_0^c)\ge\dist(x,\partial F)-W_0$, a unit cube not contained in $S''$ lies within $r_3\defeq W_0+r_1+3$ of $\partial F$. If $\supp f$ is such a cube, then by the first inequality of this step $d_\wedge(\supp f,\supp f')\ge\dist(\supp f',\partial F)-r_3-2$, so $\frac{\theta_0}2(\dist(\supp f,\partial F)+\dist(\supp f',\partial F))\le\frac{\theta_0}2d_\wedge(\supp f,\supp f')+\theta_0(r_3+1)$, and (ii) follows from (i) below with $C'$ replaced by $\max\{C',2c_1e^{\theta_0(r_3+1)}\}$; the same if $\supp f'$ is such a cube. So (ii) holds for all $f,f'$ supported in unit cubes in $F$.

\emph{Step 6 (proof of \eqref{rot:hyb:eq:B1}).} By Step~0, (i) $|\ip{f'}{(G_1-G_2)f}|\le|\ip{f'}{G_1f}|+|\ip{f'}{G_2f}|\le2c_1e^{-\theta_0d_\wedge(\supp f,\supp f')}\|f\|\|f'\|$. Then $\min\{\text{(i)},\text{(ii)}\}\le\text{(i)}^{1/2}\text{(ii)}^{1/2}$, whose exponent is $\frac{\theta_0}2d_\wedge+\frac{\theta_0}4(\dist(\supp f,\partial F)+\dist(\supp f',\partial F))$; this is \eqref{rot:hyb:eq:B1} with $\delta=\theta_0/4$ and $C=(2c_1C')^{1/2}$.

\emph{Step 7 (proof of \eqref{rot:hyb:eq:B2}).} Take $\Lambda=\Omega^{(r)}$ in \cite[(3.5)]{Ba4}; then $B^r(\Lambda)=\Omega$ and $G_r(\Omega_i,\Omega_i^{(r)},A_i)=G_i$. In the conventions above ($Q_rQ_r^*=\mathbf 1$, $QQ^*=\mathbf 1$, $P=Q^*Q$, $P_{r+1}=Q_r^*PQ_r$, with $Q_r=Q_X(A)$ and $Q=Q^{(1)}(A)$), \cite[(3.6)]{Ba4} reads
\[
C^{(r)}(\Omega_i,A_i)=V_i^*G_iV_i-\tfrac{\beta}{a_r}P(A_i)+a_r^{-1},\qquad V_i\defeq Q_r^*(A_i)\big(1-\beta P(A_i)\big),\qquad\beta\defeq\frac{a_{r+1}L^{-2}}{a_r},
\]
with $a_{r+1}=aa_r/(aL^{-2}+a_r)$; it is obtained by integrating the coarse variable out of the Gaussian of \cite[(3.1)]{Ba4}, is an algebraic identity between bounded operators, and holds on $\Lambda_X$ as on a torus. The vector $v_i(y)\defeq V_i\delta_y$ is supported in the $(r+1)$-block of $y$, is built from contours inside that block (\Cref{rot:hyb:lem:alg}), and has $\|v_i(y)\|\le1+\beta\le2$.
\emph{Case 1:} the $(r+1)$-blocks of $y$ and $y'$ lie in $F_0$. Then $v_1(y)=v_2(y)=:v(y)$, the same for $y'$, and $P(A_1;y,y')=P(A_2;y,y')$, since $A_1=A_2$ there; so $\delta C^{(r)}(y,y')=\ip{v(y)}{(G_1-G_2)v(y')}$. Splitting $v(y)$, $v(y')$ into at most $(\sqrt3L+2)^3$ unit-cube pieces and applying \eqref{rot:hyb:eq:B1} to each pair, $|\delta C^{(r)}(y,y')|\le4(\sqrt3L+2)^6e^{6\delta\sqrt3L}C\,e^{-\delta(d_\wedge(y,y')+\dist(y,\partial F)+\dist(y',\partial F))}$.
\emph{Case 2:} a $(r+1)$-block of $y$ or $y'$ meets $F\setminus F_0$, say that of $y$, so $\dist(y,\partial F)\le r_4\defeq W_0+\sqrt3L$. Then $|\delta C^{(r)}|\le|C^{(r)}(\Omega_1)|+|C^{(r)}(\Omega_2)|\le2C_0e^{-\delta_0d_\wedge(y,y')}$ by the single-domain decay of \Cref{rot:hyb:cor:Ck} on each domain (in its quotient distance), and $d_\wedge(y,y')\ge\dist(y',\partial F)-r_4$ by Step~5, so $e^{-\delta_0d_\wedge}\le e^{\delta_0r_4}e^{-(\delta_0/2)(d_\wedge(y,y')+\dist(y,\partial F)+\dist(y',\partial F))}$. Both cases give \eqref{rot:hyb:eq:B2} with $\min\{\delta,\delta_0/2\}$ in place of $\delta$.
For $\Delta^{(r)}=a_r-a_r^2Q_rG^X_rQ_r^*$ (\cite[(2.14)]{K1}, the $r$-pin) the same argument applies with $w_z$ in place of $v(y)$, with \eqref{rot:hyb:eq:B1} at the $r$-pin (Steps 0--6 with $r_1=1+\sqrt3$, \Cref{rot:hyb:lem:coer}).
\end{proof}

\begin{proposition}[{\cite[Prop.~5.6]{K2} for the hybrid flow}]\label{rot:hyb:prop:CLam}
Let $X\defeq A^X(A)+aL^{-2}P(A)$ on the coarse lattice $\Gamma_r$ of the domain, $P=Q^{(1)}(A)^*Q^{(1)}(A)$. Then $(L^2\gamma_0^{-1}+a_r^{-1})^{-1}\le X\le a_r+aL^{-2}$ and $|X(y,y')|\le Ce^{-\theta d(y,y')}$; for every subset $\Lambda$ of the coarse lattice, in particular every union of big blocks, the same bounds hold for the compression $X|_\Lambda$, and $\hat C^{(r)}_\Lambda\defeq(X|_\Lambda)^{-1}$ (\cite[(5.53)]{K2}) obeys $|\hat C^{(r)}_\Lambda(y,y')|\le Ce^{-t_*d(y,y')}$ (\cite[(5.54)]{K2}), $|\hat C^{(r)}_\Lambda(y,y')-C^{(r)}(y,y')|\le Ce^{-(t_*/6)(d(y,y')+d(y,\Lambda^c)+d(y',\Lambda^c))}$, and, under the hypotheses of \Cref{rot:hyb:lem:domdiff} with $\Lambda\subset F$, the two-domain bound \eqref{rot:hyb:eq:B2} in the metric $d_\wedge$. The constants depend on $a$, $L$ (and $W_0$) only.
\end{proposition}
\begin{proof}
Upper bound: $A^X=(a_r^{-1}+Q_rC_AQ_r^*)^{-1}\le a_r$ and $P\le1$. Lower bound: \cite[(3.6)]{Ba4} in the form of Step~7 above gives $X^{-1}=C^{(r)}=V^*G_{r+1}V-(\beta/a_r)P+a_r^{-1}\le\|V\|^2\|G_{r+1}\|+a_r^{-1}$ with $\|V\|\le1$ and $\|G_{r+1}\|\le L^2\gamma_0^{-1}$ (\Cref{rot:hyb:cor:L2}, the $(r+1)$-pin). Decay of $X$: \Cref{rot:hyb:prop:CT} for $A^X$ and the finite range of $P$ (its kernel vanishes unless $y,y'$ lie in one $L$-block, of diameter $<\sqrt3L$). Compression to $\Lambda$ keeps both bounds and the entries. Then, for a $1$-Lipschitz $\psi$ on $\Lambda$, $e^{t\psi}X|_\Lambda e^{-t\psi}-X|_\Lambda$ has entries bounded by $Ct\,d(y,y')e^{-(\theta/2)d(y,y')}$ for $t\le\theta/2$, hence norm $\le Ct$ (Schur test), so $e^{t\psi}X|_\Lambda e^{-t\psi}$ is invertible with inverse of norm $\le2(L^2\gamma_0^{-1}+a_r^{-1})$ for $t\le t_*$, $t_*$ depending on $a,L$ only; with $\psi=\min\{d(\cdot,y'),d(y,y')\}$ this gives $|\hat C^{(r)}_\Lambda(y,y')|\le Ce^{-t_*d(y,y')}$ (the entry is bounded by the operator norm on the counting measure of $\Gamma_r$). The difference with $C^{(r)}$ is $\hat C^{(r)}_\Lambda X_{\Lambda\Lambda^c}C^{(r)}_{\Lambda^c\Lambda}$ (the block identity for the inverse of a compression), whose decay follows from the three factors' decay and the geometric mean as in Step~6; the two-domain bound is the resolvent identity on $\ell^2(\Lambda)$ with \eqref{rot:hyb:eq:B2} for $\Delta^{(r)}$ ($\delta P=0$ off the collar), Step~5 and Step~6's interpolation.
\end{proof}

% r03_hyb_A0: proof of \Cref{rot:prop:BD}, part IV: the A=0 input (the hybrid kernels on Lambda_X and on periodic cells) and the pointwise and Hoelder bounds at a regular background on a cell.

\paragraph{The $A=0$ input and the bounds on a cell.}

\begin{lemma}[$A=0$ bounds for the hybrid flow]\label{rot:hyb:lem:A0}
On $\Lambda_X$, with $S^X_j$ the slice kernels of \Cref{rot:def:match} ($G^X_r(0)=\sum_{j<r}S^X_j$ exactly):
\begin{enumerate}[label=\textup{(\alph*)},leftmargin=*]
\item $|S^X_j(x,y)|\le C\ell_j^{-1}e^{-\delta|x-y|/\ell_j}$, $|\nabla_xS^X_j(x,y)|\le C\ell_j^{-2}e^{-\delta|x-y|/\ell_j}$, and for $0<\alpha<1$
\[
|\nabla_xS^X_j(x,y)-\nabla_xS^X_j(x',y)|\le C_\alpha\log^{3\alpha}(2+\ell_j/\eta')\,\ell_j^{-2-\alpha}|x-x'|^\alpha e^{-\delta\min(|x-y|,|x'-y|)/\ell_j}
\]
(the logarithm is the slice's own, $\le\log^{3\alpha}(1/\eta')$);
\item consequently $\|G^X_r(0)f\|_\infty+\|\nabla G^X_r(0)f\|_\infty\le C\|f\|_\infty$ with $C$ uniform in $r$, $[\nabla G^X_r(0)f]_\alpha\le C_\alpha\log^{3\alpha}(1/\eta')\|f\|_\infty$, and $\|G^X_r(0)f\|_q\le C\|f\|_p$ for $0\le\frac1p-\frac1q<\frac23$, with the exponential localisation $e^{-\delta_0\dist(x,\supp f)}$ (\cite[(2.16)--(2.17)]{Ba4});
\item on a periodic cell $\square=\Lambda_X/P$ fitted by both lattices, of side $\ge2M$, and for a constant field $B_0$,
\[
G^X_r(\square,B_0)(x,y)=e^{\mathrm ie_rB_0\cdot(x-y)}\sum_{\gamma\in P}e^{\mathrm ie_rB_0\cdot\gamma}G^X_r(0)(x,y+\gamma)
\]
(twisted periodisation), and (b) holds on $\square$ for $G^X_r(\square,B_0)$ with constants independent of $B_0$ and $M$, the H\"older constant carrying the same $\log^{3\alpha}(1/\eta')$.
\end{enumerate}
\end{lemma}
\begin{proof}
(a) In the unit scale of level $j$, $S^X_j=H^X_jC^X_jH^{X*}_j$ (\Cref{rot:def:match}). The decay $|H^X_j(x,z)|\le Ce^{-\delta|x-z|}$ and the gradient bound $|\nabla H^X_j|\le C_\nabla$ are \Cref{rot:lem:H,rot:lem:grad} (the decay is the Combes--Thomas decay of the coercive coarse operators, \cite{CT73}); the decay of the coarse kernel $C^X_j=(A^X_j+bQ^*Q)^{-1}$ is the one used in the proof of \Cref{rot:lem:slice} (its symbol is analytic and bounded on a strip because $A^X_j+bQ^*Q\ge c$). \emph{The H\"older derivative.} The second difference of $H^X$ in $x$ multiplies the alias-$l$ summand of $M^X$ in the representation of \Cref{rot:lem:H} by $d(q+l)^2$ with $|d(Q)|\le|Q|$; since $\Delta'_X(Q)\ge c|Q|^2$ on each alias cell and $\sum_lu_X(q+l)^2=1$, the alias sum is $\le C\sum_l|u_X(q+l)|\le C$, and integrating against $|u_{r-1}|$ over $V_\eta$ costs $J_0\le C\log^3(1/\eta)$ (\Cref{rot:lem:CK0}). So $|\nabla\nabla H^X|\le\Lambda_\eta\defeq C\log^3(1/\eta)$ and, by interpolation with $|\nabla H^X|\le C_\nabla$,
\[
|\nabla H^X(x,z)-\nabla H^X(x',z)|\le\min\{2C_\nabla,\Lambda_\eta|x-x'|\}\le(2C_\nabla)^{1-\alpha}\Lambda_\eta^{\alpha}|x-x'|^\alpha\le C_\alpha\log^{3\alpha}(1/\eta)|x-x'|^\alpha
\]
for every $\alpha<1$; the geometric mean with the decay gives the localised form. Composing with $C^X_j$ and $H^{X*}_j$ (exponentially decaying kernels) and rescaling by \cite[(2.20)]{K1} (scalar kernels scale as $\ell^{-1}$ in $d=3$, derivatives as $\ell^{-2}$) gives (a); at level $j$ the logarithm is $\log(1/\eta_j)$ with $\eta_j=5^{1-j}$, i.e.\ $\le\log(2+\ell_j/\eta')$. \emph{Summation.} $\|S^X_jf\|_\infty\le C\ell_j^{-1}\ell_j^3\|f\|_\infty$, $\|\nabla S^X_jf\|_\infty\le C\ell_j\|f\|_\infty$, $[\nabla S^X_jf]_\alpha\le C_\alpha\log^{3\alpha}(2+\ell_j/\eta')\,\ell_j^{1-\alpha}\|f\|_\infty$; summing over $j<r$ converges for $\alpha<1$, the H\"older term to $\le C_\alpha\sum_{j<r}5^{-(r-j)(1-\alpha)}(1+j\ln5)^{3\alpha}\le C_\alpha\log^{3\alpha}(1/\eta')$ (this is \cite[(2.38)--(2.39)]{Ba4} with the slice kernels in place of $G_jQ_j^*C^{(j)}Q_jG_j$). The $L^p\to L^q$ bounds follow from the kernel bound $|G^X_r(0)(x,y)|\le C\min(\eta'^{-1},|x-y|^{-1})e^{-\delta|x-y|}$ by Young's inequality and the Riesz--Thorin theorem (\cite[p.~583]{Ba4}). (c) Both sides are the inverse of $-\Delta_{B_0}+m_r^2+a_rP_X(B_0)$ on $P$-periodic functions: $P\subset\Gamma_r$ makes the hybrid block system $P$-periodic, and $e^{\mathrm ie_rB_0\cdot x}$ conjugates the operator to the $A=0$ operator on functions with the twisted periodicity $\psi(x+\gamma)=e^{-\mathrm ie_rB_0\cdot\gamma}\psi(x)$, whose kernel is the twisted image sum (\cite[(B.3)]{K2} with phases). The images with $\gamma\ne0$ contribute $\le Ce^{-\delta_0(2M-\sqrt3)}$ to every bound in (b), uniformly in $B_0$.
\end{proof}

\begin{lemma}[{\cite[Lemma~2.2]{Ba4}, \cite[Lemma~5.8]{K2}, for the hybrid flow}]\label{rot:hyb:lem:HOL}
Let $\square$ be a periodic cell fitted by both lattices of side $2M$, and $\tilde A=A_0+A'$ regular on $\square$, $A_0$ constant, with $|A'|_\infty,|\partial A'|_\infty\le c'e_r^{\beta-1}$ and $A'$ supported away from the boundary of the fundamental region. Then for $e_r^\beta\le c_0(a,M)$, with $c_0,c_1,c_2$ uniform in $r$,
\[
\|G^X_r(\square,\tilde A)f\|_{1,0}\le c_1\|f\|_\infty,\qquad\|G^X_r(\square,\tilde A)f\|_{1,\alpha}\le c_1C_\alpha\log^{3\alpha}(1/\eta')\|f\|_\infty,
\]
\[
\|G^X_r(\square,\tilde A)f\|_q,\ \|D_{\tilde A}G^X_rf\|_q,\ \|G^X_rD_{\tilde A}^*f\|_q\le c_2\|f\|_p
\]
for $1\le p,q\le\infty$ with $\frac1p-\frac1q<\frac23$, where $\|\cdot\|_{1,\alpha}$ is the covariant H\"older norm \cite[(2.14)]{Ba4} with the field $\tilde A$ and $\|\cdot\|_{1,0}$ the sup norm plus the sup of the covariant derivative.
\end{lemma}
\begin{proof}
\cite[(2.24)--(2.26), (2.31)--(2.33)]{Ba4}: $G_r(\square,\tilde A)=\sum_{n\ge0}G_r(\square,A_0)[V_rG_r(\square,A_0)]^n$ with $V_r=V_r(A',A_0)$ the vertex operator of \cite[(3.16)]{Ba1} ($=$ \cite[(3.53)]{K1}); $V_r$ is a first-order differential operator with coefficients $O(e_r|A'|,e_r|\partial A'|)$ (\cite[(2.32)]{Ba4}, using that $A'$ vanishes near the boundary of the fundamental region) plus the pin part $a_r[Q(A_0)^*F_{2,r}(A',A_0)+F_{2,r}^*Q(A_0)+F_{2,r}^*F_{2,r}]$ with $(F_{2,r}(A',A_0)\phi)(y)=\sum_{x\in B(y)}\eta'^3(U(A'(\Gamma_{y,x}))-1)U(A_0(\Gamma_{y,x}))^{-1}\phi(x)$, $\|F_{2,r}\|_{\infty\to\infty}\le\sup|U(A'(\Gamma))-1|\le2.1\,e_r|A'|_\infty$ by \Cref{rot:hyb:lem:len}(i). Hence $\|V_rg\|_\infty\le C(M)e_r^\beta\|g\|_{1,0}$. By \Cref{rot:hyb:lem:A0}(b),(c), $\|G_r(\square,A_0)g\|_{1,0}\le c\|g\|_\infty$ with $c$ uniform in $r$ and $A_0$, and $\|G_r(\square,A_0)g\|_{1,\alpha}\le cC_\alpha\log^{3\alpha}(1/\eta')\|g\|_\infty$ (the constant background is a twisted periodisation of the $A=0$ kernel, and the H\"older norm (2.14) with field $A_0$ is the plain H\"older norm after the local gauge $e^{\mathrm ie_rA_0\cdot x}$ on each unit cube). \emph{Order of summation.} Put $g_0\defeq f$, $g_n\defeq V_rG_r(\square,A_0)g_{n-1}$. Then $\|g_n\|_\infty\le C(M)e_r^\beta\|G_r(\square,A_0)g_{n-1}\|_{1,0}\le(C(M)c\,e_r^\beta)^n\|f\|_\infty$, so for $C(M)c\,e_r^\beta\le\frac12$, a condition uniform in $r$, the sum $g\defeq\sum_ng_n$ converges in the sup norm with $\|g\|_\infty\le2\|f\|_\infty$, and $G_r(\square,\tilde A)f=\sum_nG_r(\square,A_0)g_n=G_r(\square,A_0)g$ (the cell is finite, every operator is bounded, and the sums may be exchanged). Hence $\|G_r(\square,\tilde A)f\|_{1,0}\le2c\|f\|_\infty$ and $\|G_r(\square,\tilde A)f\|_{1,\alpha}\le2cC_\alpha\log^{3\alpha}(1/\eta')\|f\|_\infty$: the H\"older constant enters once, at the end, and never the smallness condition. (Summing in $\|\cdot\|_{1,\alpha}$, as \cite[(2.33)]{Ba4} may with its uniform constant, would turn the condition into $e_r^\beta\log^{3\alpha}(1/\eta')\le c_0$, which depends on $r$.) The norms (2.14) with the fields $A_0$ and $\tilde A$ differ on the cell by at most $Ce_r^\beta\|\cdot\|_{1,0}$ ($D_{\tilde A}=U(A')D_{A_0}+F_{1,r}(A')$ and $|U(A'(\Gamma_{y,x}))-1|\le Ce_r|A'|_\infty|x-y|$ for the bounded-length paths of (2.14)), which changes $c_1$ by a factor $1+O(e_r^\beta)$. The $L^p\to L^q$ bounds are summed in the same order, with $D_{\tilde A}=U(A')D_{A_0}+F_{1,r}(A')$.
\end{proof}

\begin{remark}[the H\"older constant]\label{rot:hyb:rem:Hol}
In \Cref{rot:prop:BD}(a) the H\"older clause \cite[(3.8)]{K1} of $G^X_r(A)$, and the H\"older clauses of $H^X(A)$ and $\Gamma^X(A)$, hold with the constant $C_\alpha\log^{3\alpha}(1/\eta')$, $0<\alpha<1$ (\Cref{rot:hyb:lem:A0}, \Cref{rot:hyb:lem:HOL}, and the random walk below); on slice $j$ the constant is $C_\alpha\log^{3\alpha}(2+\ell_j/\eta')$ (\Cref{rot:hyb:lem:A0}(a)). The sup, decay and one-derivative bounds (3.6)--(3.7) and the $\|\cdot\|_{1,0}$ bounds have constants uniform in $r$, and so does every smallness condition. \cite[(3.8)]{K1} asks for some $0<\alpha<1$ with some constant (\cite[p.~656]{K1}). Every use of the clause in this section is either as a gain, where $\eta'^\alpha\log^{3\alpha}(1/\eta')\le C_{\alpha,\alpha'}\eta'^{\alpha'}$ and, slice by slice, $(\eta'/\ell_j)^\alpha\log^{3\alpha}(2+\ell_j/\eta')\le C_{\alpha,\alpha'}(\eta'/\ell_j)^{\alpha'}$ for $\alpha'<\alpha$ (\Cref{rot:thm:H1S}, whose exponent $\theta_1$ is fixed by $\alpha$ alone), or inside a per-flow bound that already allows a factor $(1+k)^p$ (\Cref{rot:cor:OF}), where $\log^{3\alpha}(1/\eta')=(r\ln5)^{3\alpha}\le C(1+k)^{3\alpha}$.
\end{remark}

% r03_hyb_walk: proof of \Cref{rot:prop:BD}, part V: the random walk with periodic cells, the boundary clause, the remaining kernels, the slice clauses, the levels 0 and 1; this completes part (a).

\begin{proof}[Proof of \Cref{rot:prop:BD}(a)]
\emph{Random walk.} Let $\Omega$ be $\Lambda_X$ or a torus fitted by both lattices with $P\subset M\Z^3$, $A$ regular on $\Omega$. For $j\in\Z^3$ let $\square_j$ be the periodic cell of side $2M$ ``centred'' at $Mj$ and fitted by both lattices (which is possible iff $25\mid2M$, as the tori of the chain allow; \cite[p.~345]{K2}: ``blocks $\{\square_j\}$ which are small scale versions of $D_J$, that is they can be fitted by both the usual lattice and the rotated lattice''), $h_j$ the partition of unity of \cite[\S2]{Ba4} ($h_j=1$ on the $M$-cube around $Mj$, supported in the $\frac43M$-cube, $\sum_jh_j^2=1$), and $\tilde A_j=A_0^{(j)}+\theta_jA'^{(j)}$ with $A_0^{(j)}=A(Mj)$ and the cutoff $\theta_j$ of \cite[p.~575]{Ba4}, so that $\tilde A_j=A$ on $\supp h_j$, is regular with constant $c'=O(Mc)$, and is constant near the boundary of the fundamental region of $\square_j$. With $G_0\defeq\sum_jh_jG^X_r(\square_j,\tilde A_j)h_j$ one computes, exactly as \cite[(2.3)--(2.9)]{Ba4} and \cite[(5.8)--(5.12)]{K2},
\[
(-\Delta_A+m_r^2+a_rP_X(A))G_0=\mathbf 1-\cR,\qquad \cR=\sum_jK_jG^X_r(\square_j,\tilde A_j)h_j,
\]
\[
(K_j\phi)(x)=\mathbf 1_{\square_j}(x)\Big[\sum_{b\in\mathrm{st}(x)}(\partial h_j)(b)(D_A\phi)(b)+(\Delta h_j)(x)\phi(x)-(R_r(A,\partial h_j)\phi)(x)\Big],
\]
where $R_r(A,\partial h)\phi(x)=\sum_{x'\in B(z(x))}\eta'^3U(A(\Gamma_{z,x}))U(A(\Gamma_{z,x'}))^{-1}(\partial h)(\Gamma_{x,x'})\phi(x')$ collects the commutator of the pin with the multiplication by $h_j$ (\cite[(2.8)]{Ba4}); $(\partial h)(\Gamma_{x,x'})=h(x')-h(x)$ for any path, so $|R_r(A,\partial h_j)\phi(x)|\le\sup|\partial h_j|\cdot2.1\cdot\sup_B|\phi|$ by \Cref{rot:hyb:lem:len}(i). The three identities used --- the product rule for $-\Delta_A$ (\cite[(2.5)]{Ba4}), the commutator of the pin with a smooth function (\cite[(2.8)]{Ba4}), and the agreement of the cell operator with the global one on $\supp h_j$ (\cite[(2.6)]{Ba4}; for a periodic cell this holds because $\supp h_j$ stays at distance $\ge\sqrt3$ from the boundary of the fundamental region, so no block and no bond of $\supp h_j$ wraps) --- hold for any block system with contours from the root to the point. Hence $G^X_r(\Omega,A)=\sum_{n\ge0}G_0\cR^n$ in $L^2$ (\Cref{rot:hyb:cor:L2} gives $\|\cR\|_{2\to2}\le O(M^{-1})$ as in \cite{Ba4}, since $\sup|\partial h_j|=O(M^{-1})$ and $\sup|\Delta h_j|=O(M^{-2})$), i.e.\ the path representation \cite[(2.13)]{Ba4}, \cite[(5.56)]{K2}.

\emph{Bounds.} \cite[(2.18)--(2.22)]{Ba4} verbatim: for $x,x'$ in $\supp h_{\omega_0}$ and $f$ supported in $\square_{\omega_n}$, each factor $\|K_{\omega_i}G^X_r(\square_{\omega_i},\tilde A_{\omega_i})h_{\omega_i}\|_{\infty\to\infty}$ is $\le O(M^{-1})c_1$ by \Cref{rot:hyb:lem:HOL} and the bound on $K_j$ ($c_1$ the uniform $\|\cdot\|_{1,0}$ constant of \Cref{rot:hyb:lem:HOL}: $K_j$ takes one covariant derivative of its argument, so the choice of $M$ below is uniform in $r$), the first factor is bounded in $\|\cdot\|_{1,\alpha}$ by \Cref{rot:hyb:lem:HOL} (the only place where $C_\alpha\log^{3\alpha}(1/\eta')$ enters, and only in the clause (3.8)), the number of paths of length $n$ is $\le2^d(3^d)^{n-1}2^d$, and $n\ge M^{-1}\dist(\{x,x'\},\supp f)-2$; for $M$ fixed with $3^dc_1O(1)M^{-1}\le e^{-1}$ the sum converges and gives \cite[(3.7)--(3.8)]{K1} with $\delta_0=(2M)^{-1}$. The tail with $n>n_0$ is handled with the $L^p$ bounds (2.17) as in \cite[(2.21)]{Ba4}. The bounds (3.6) for $\Delta^{(r)}$ and $C^{(r)}$ are \Cref{rot:hyb:prop:CT} and \Cref{rot:hyb:cor:Ck}; the bounds for the next-step precision $\Delta^{(r)}+aL^{-2}\hat Q^*\hat Q$ and its restricted inverses $\hat C^{(r)}_\Lambda$ are \Cref{rot:hyb:prop:CLam}.

\emph{Two-torus clause.} For two domains $\Omega_1$, $\Omega_2$ sharing a fundamental region $F$ as in \Cref{rot:hyb:lem:domdiff} (or a torus against $\Lambda_X$), the representations coincide term by term for the paths whose cells do not meet the boundary of $F$, and the difference is bounded, as on \cite[p.~579]{Ba4}, by the same sum restricted to paths meeting the boundary; such a path of $n$ cells also joins $x$ to $\supp f$ in its own domain, so $\sqrt3Mn\ge\max\{d_i(\{x,x'\},\supp f),\dist(\{x,x'\},\partial F)+\dist(\supp f,\partial F)\}-c$ for $i=1,2$, and the average of the two exponents gives the factor $e^{-(2\sqrt3M)^{-1}(d_\wedge(\{x,x'\},\supp f)+\dist(\{x,x'\},\partial F)+\dist(\supp f,\partial F))}$ for the difference of the $G^X_r$ clauses (3.7)--(3.8), with $d_\wedge\le d_i$ the smaller of the two quotient distances. For $\Delta^{(r)}$, $C^{(r)}$ and $\hat C^{(r)}_\Lambda$ the two-torus clause is \Cref{rot:hyb:lem:domdiff} and \Cref{rot:hyb:prop:CLam}. The metric $d_\wedge$ cannot be replaced by the distance of the lift (the remark after \Cref{rot:hyb:lem:domdiff}); the uses of the clause in this section (\Cref{rot:lem:CB}(N), the kernel sweep of \Cref{rot:thm:Asrc}, and \Cref{rot:lem:TC}(IV)) need only the decay on each torus in its own quotient distance together with the boundary factor, which is what is proved.

\emph{The remaining kernels.} $H^X(A)=a_rG^X_r(A)Q_X(A)^*$: apply (3.7)--(3.8) to $f=w_z$ ($\|w_z\|_\infty=1$, support in a block of diameter $\le\sqrt3$). $B^X(A)=a_r^{-1}+Q_X(A)C_AQ_X(A)^*$: by the diamagnetic inequality $|C_A(x,y)|\le C_0(x,y)=(-\Delta+m_r^2)^{-1}(x,y)$ entrywise ($-\Delta_A+m_r^2=c-T_A$ with $|T_A|=T_0$ entrywise and $(c-T_0)^{-1}=\sum c^{-n-1}T_0^n\ge0$), so $|B^X(z,w)|\le a_r^{-1}\delta_{zw}+\sum_{x\in B(z),y\in B(w)}\eta'^6C_0(x-y)\le C(1+|z-w|)^{-1}$ ($d=3$). $\Gamma^X(A)=H^XB^XH^{X*}$: the fine-variable regularity is that of $H^X$, and $|\Gamma^X(x,y)|\le\sum_{z,w}Ce^{-\delta_0|x-z|}(1+|z-w|)^{-1}e^{-\delta_0|y-w|}\le C(1+|x-y|)^{-1}$. $B^X(A)$ and $\Gamma^X(A)$ are not unit-decaying: $\Gamma^X(x,y)\approx C_A(x,y)$ for $|x-y|\gg1$, at the coarse mass $m_r=m\delta_r\to0$; \cite[Thm.~3.3]{K1} never claims otherwise (it concerns $G_r$, $\Delta^{(r)}$, $C^{(r)}$). Where a renormalised family carries a line $\Gamma^F$ or $(A^F)^{-1}$ at the top scale, either the coarse-slice decomposition of \Cref{rot:cor:unif}(b) is used, or, in the localised form \cite[(3.55)]{K1}, all vertex positions lie in the region $F_v$ of \Cref{rot:lem:CB}, of diameter $\le Cp_r$, where a line bounded by $C(1+|x-y|)^{-1}$ costs at most a factor $Cp_r^3$, absorbed in the field-norm polynomial of \Cref{rot:prop:cp}.

\emph{The slice clauses.} Nothing in the argument above uses the level: the same proof at level $j\le r$, in the unit scale of level $j$, gives (3.6)--(3.8) for the step-$j$ operators $G_j$, $\Delta^{(j)}$, $C^{(j)}$, $H_j$ of the hybrid flow with constants uniform in $j$ (the first step, with the tiles as blocks, is covered by \Cref{rot:hyb:lem:coer} and \Cref{rot:hyb:lem:A0} with $B^X_1$ in place of $B^X_r$; the steps $j\ge2$ are the standard rotated steps acting on composite blocks, \Cref{rot:hyb:lem:alg}(iv)). King's derivation of the slice bounds \cite[(3.63)--(3.65)]{K1} from Theorem~3.3 (\cite[Prop.~3.7]{K1}: the slice $S_j=H_jC_jH_j^*$ at a regular background, the one- and two-ended contour clauses summed along contours of bounded length) therefore applies verbatim to the hybrid flow, with the covariant lattice derivatives and transported H\"older quotients of \cite[(2.10)--(2.11)]{Ba3} at regular $A$; after the constant-background reduction of \Cref{rot:lem:CB} the compared lines are the $A=0$ slices of \Cref{rot:hyb:lem:A0}, so there the plain form holds. For the contour clauses the contour lengths are those of \Cref{rot:hyb:lem:len}(i).

\emph{Levels $0$ and $1$.} In the form in which the hybrid lower bound (\Cref{rot:prop:H4}) consumes (a) at its first steps (precision bounds, precision decay, sup-decay of $C^{(0)}$ and $C^{(1)}$), (a) holds also at $j\in\{0,1\}$ by the same proofs, which use $r\ge2$ nowhere: at $j=1$ the level-$1$ block is the tile, inside the ball of radius $\frac{\sqrt3}2$ about its centre, with an in-tile contour of length $\le1$ in the unit scale of level $1$, so \Cref{rot:hyb:lem:coer}, \Cref{rot:hyb:cor:L2}, \Cref{rot:hyb:prop:CT} and \Cref{rot:hyb:lem:domdiff} read verbatim (the first Neumann eigenvalue of the cube is then $100\sin^2(\pi/40)>\frac14$), and the fine spacing $\eta'=\frac15$ is a constant; at $j=0$ the precision $-\Delta_A+m_0^2+aL^{-2}\hat Q^*\hat Q$ has finite range and its lower bound is the $(r+1)$-pin clause of \Cref{rot:hyb:cor:L2} at $r=0$.

\emph{The Neumann variant.} \cite[Thm.~3.3]{K1} is printed for Neumann boundary conditions on a rectangular parallelepiped of large blocks. For the hybrid flow such a domain is a digitised rotated parallelepiped, on which neither the reflection representation \cite[(2.42)]{Ba4} nor the axis-cube cover near the boundary is available; the Neumann variant is not claimed and not used: \cite[\S\S4--5]{K2} decouple with periodic cells in the random walk, the localisation of \Cref{rot:lem:CB} runs on the infinite lattice in the universal-cover form, and the only Neumann object of the chain, the boundary factor of $\delta G_r$, is replaced by the two-torus clause above and the same-domain clause \Cref{rot:lem:CB}(N).

\emph{The standard flow.} For the flow $Y$ every statement of (a) is \cite[Thm.~3.3]{K1}, proved in \cite{Ba4} on $\eta'\Z^3$ in the frame of $\Yk$, with King's uniform constants; its two-torus clause and restricted inverses are \cite[Prop.~5.6]{K2} and \Cref{rot:hyb:lem:domdiff}, whose proof uses nothing specific to the hybrid blocks. This completes the proof of (a).
\end{proof}

% r03_hyb_ES: proof of \Cref{rot:prop:BD}, part VI: the continuum extension of the Y kernels (cubic-Hermite cardinal interpolation), its decay, its multiplier form at the X sites, and the decaying forms at the X sites.

\paragraph{The continuum extension of the $Y$ kernels.}
The comparisons of \Cref{rot:hyb:lem:CV,rot:hyb:lem:CLN} below compare a sum over $X$ bonds with a sum over $Y$ bonds through \Cref{rot:hyb:lem:CL}, which needs a $C^1$ field on the block that equals the $Y$ kernel on the $Y$ lattice. The Fourier formula of \Cref{rot:lem:H} evaluated off the lattice (the sharp-zone extension) is not admissible for this: the symbol does not vanish on the faces of the zone, so the extension and its gradient have algebraic tails of order $\eta'^3/|r-z|$ and $\eta'^2/|r-z|$ off the lattice, not summable over a torus in $d=3$. The field used is the cubic-Hermite cardinal interpolation of the lattice kernels: it is $C^1$, equals the lattice kernel on $\Lambda_Y$, has exponentially decaying value and gradient by the lattice bounds alone, and obeys the multiplier comparison of \Cref{rot:lem:H} at the $X$ sites with an explicit constant.

\begin{definition}[cubic-Hermite cardinal extension]\label{rot:hyb:def:ES}
$S(u)\defeq3u^2-2u^3$, $\psi(s)\defeq S(1-|s|)\mathbf 1_{|s|<1}=(1-|s|)^2(1+2|s|)\mathbf 1_{|s|<1}$, $\Psi(v)\defeq\prod_\alpha\psi(v\cdot f_\alpha/\eta')$ (the frame of $\Yk$). For $F$ on $\Lambda_Y=R^{-1}(\eta'\Z^3)$ put $(I_SF)(r)\defeq\sum_{y\in\Lambda_Y}\Psi(r-y)F(y)$, $r\in\R^3$. The extensions are $E^S(r,z)\defeq I_S[H^Y(\cdot,z)](r)$ for the level-$r$ lift (scalar and photon), $E^S_j(r,z)\defeq I_S[H^Y_j(\cdot,z)](r)$ for the level-$j$ lift, $\tilde S^S_j(r,y)\defeq\sum_{z,z'}E^S_j(r,z)C^Y_j(z,z')\overline{H^Y_j(y,z')}=I_S[S^Y_j(\cdot,y)](r)$ for $y\in\Lambda_Y$ ($S^Y_j=H^Y_jC^Y_jH^{Y*}_j$, \Cref{rot:def:match}), likewise for the photon slices, and $E^S\varphi\defeq I_S[H^Y\varphi]$ for a coarse field $\varphi$ (componentwise in any frame: $I_S$ is scalar and linear). The extension acts in the first argument only; second arguments stay on $\Lambda_Y$.
\end{definition}

\begin{lemma}[cell identity]\label{rot:hyb:lem:ESH}
Write the frame coordinates of $r$ as $f_\alpha\cdot r=\eta'(n_\alpha+u_\alpha)$, $n\in\Z^3$, $u\in[0,1)^3$, with cell corners $y_\sigma\defeq R^{-1}\eta'(n+\sigma)$, $\sigma\in\{0,1\}^3$, and weights $\omega_\sigma(u)\defeq\prod_\alpha\omega_{\sigma_\alpha}(u_\alpha)$, $\omega_0\defeq1-S$, $\omega_1\defeq S$ (so $\omega_\sigma\ge0$, $\sum_\sigma\omega_\sigma=1$). Then
\textup{(i)} $I_SF(r)=\sum_\sigma\omega_\sigma(u)F(y_\sigma)$;
\textup{(ii)} $\partial_{f_1}I_SF(r)=S'(u_1)\sum_{\sigma'\in\{0,1\}^2}\omega_{\sigma'}(u_2,u_3)\,(\partial_1F)(y_{(0,\sigma')})$ with the lattice forward difference $\partial_1F(y)\defeq[F(y+\eta'f_1)-F(y)]/\eta'$, and likewise for $f_2,f_3$;
\textup{(iii)} $|I_SF(r)|\le\max_\sigma|F(y_\sigma)|$ and $|\partial_{f_\alpha}I_SF(r)|\le\frac32\max|\partial_\alpha F|$ over the four cell edges parallel to $f_\alpha$, so $|\nabla I_SF(r)|\le\frac{3\sqrt3}2\max_{\mathrm{edges}}|\partial F|$;
\textup{(iv)} $I_SF\in C^1(\R^3)$ and $I_SF=F$ on $\Lambda_Y$.
\end{lemma}
\begin{proof}
Only the corners with $|(r-y)\cdot f_\alpha|<\eta'$ for all $\alpha$ contribute, i.e.\ the $y_\sigma$; for them $\psi((r-y_\sigma)\cdot f_\alpha/\eta')=\psi(u_\alpha-\sigma_\alpha)$, which is $S(1-u_\alpha)=1-S(u_\alpha)$ for $\sigma_\alpha=0$ and $S(u_\alpha)$ for $\sigma_\alpha=1$ ($S(u)+S(1-u)=1$): this is (i). Differentiating (i) in $r\cdot f_1=\eta'(n_1+u_1)$ gives $\eta'^{-1}S'(u_1)\sum_{\sigma'}\omega_{\sigma'}[F(y_{(1,\sigma')})-F(y_{(0,\sigma')})]$, which is (ii). (iii) follows from $\omega\ge0$, $\sum\omega=1$, $0\le S'(u)=6u(1-u)\le\frac32$. (iv) $\psi'(s)=-6s(1-|s|)$ on $|s|<1$ vanishes at $s=0,\pm1$, so $\psi\in C^1(\R)$, and the sum is locally finite; $\psi(m)=\delta_{m0}$ for $m\in\Z$.
\end{proof}

\begin{lemma}[decay of the extensions]\label{rot:hyb:lem:ESB}
The lattice inputs are \cite[(3.6), (3.7), (3.63)]{K1} for the standard flow: $|H^Y(y,z)|,|\partial H^Y(y,z)|\le Ce^{-\delta_0|y-z|}$ ((3.7) with $f=Q^*\delta_z$, supported in the block of $z$, $\|f\|_\infty\le1$; for the photon, ``the same bounds for the vector field operators'', \cite[p.~656]{K1}); at level $j$, in the unit scale of level $j$, $|H^Y_j|,|\partial H^Y_j|\le Ce^{-\delta_0|\cdot|}$ and $|C^Y_j(z,z')|\le Ce^{-\delta_0|z-z'|}$. Then for $0<\delta<\delta_0$, uniformly in $j,r,\eta'$, the torus and $m_j\ge0$, for all $r\in\R^3$, $z$, and $y\in\Lambda_Y$, in the unit scale of level $r$:
\begin{enumerate}[label=\textup{(\alph*)},leftmargin=*]
\item $|E^S(r,z)|,\ \frac{2}{3\sqrt3}|\nabla_rE^S(r,z)|\le Ce^{\sqrt3\delta_0\eta'}e^{-\delta_0|r-z|}$, and the same for $E^S_j$ in the unit scale of level $j$;
\item $|\tilde S^S_j(r,y)|\le C\ell_j^{-1}e^{-\delta|r-y|/\ell_j}$, $|\nabla_r\tilde S^S_j(r,y)|,\ |\partial_y\tilde S^S_j(r,y)|\le C\ell_j^{-2}e^{-\delta|r-y|/\ell_j}$ and $|\nabla_r\partial_y\tilde S^S_j(r,y)|\le C\ell_j^{-3}e^{-\delta|r-y|/\ell_j}$ ($\partial_y$ a lattice difference in the second argument);
\item $E^S(\cdot,z)$, $\tilde S^S_j(\cdot,y)$, $E^S\varphi$ are $C^1$ and equal $H^Y(\cdot,z)$, $S^Y_j(\cdot,y)$, $H^Y\varphi$ on $\Lambda_Y$.
\end{enumerate}
\end{lemma}
\begin{proof}
(a) \Cref{rot:hyb:lem:ESH}(iii) with $F=H^Y(\cdot,z)$ and $|y_\sigma-r|\le\sqrt3\eta'$. (b) In the unit scale of level $j$ put $F(y')\defeq S^Y_j(y',y)$, resp.\ $\partial_yS^Y_j(y',y)$ ($\partial_y$ acts on the second argument and commutes with $I_S$). With $S^Y_j=H^Y_jC^Y_jH^{Y*}_j$, $F$ and its lattice gradient in $y'$ are sums $\sum_{z,z'}$ of products of three kernels with at most one lattice difference each, so with $e^{-\delta_0(|y'-z|+|z-z'|+|z'-y|)}\le e^{-\delta|y'-y|}e^{-(\delta_0-\delta)(|y'-z|+|z-z'|)}$ they are $\le C^3K_\delta e^{-\delta|y'-y|}$, $K_\delta\defeq(\sum_{v\in\Z^3}e^{-(\delta_0-\delta)(|v|-\sqrt3)})^2$. \Cref{rot:hyb:lem:ESH}(iii) at the spacing $\eta'/\ell_j\le\frac15$ adds $e^{\sqrt3\delta/5}$ and the factor $\frac{3\sqrt3}2$ for the gradient. Rescale by \cite[(2.20)]{K1}: kernel $\ell_j^{-1}$, each derivative $\ell_j^{-1}$. (c) \Cref{rot:hyb:lem:ESH}(iv). Only lattice inputs with at most one derivative occur; no same-argument second derivative of a block-spin kernel is used.
\end{proof}

\begin{lemma}[multiplier form at the $X$ sites]\label{rot:hyb:lem:ADMS}
Unit scale of level $r$; $x\in\Gamma_1$, $s\in\R^3$, $|s|\le\sqrt{13}\eta'$. In the notation of \Cref{rot:lem:H},
\[
E^S(x+s,z)=(2\pi)^{-3}\!\int_{V_\eta}\!dq\,e^{\mathrm iq\cdot(x-z)}u_{r-1}(q)M^Y_S(q;s)A^Y(q),\qquad M^Y_S(q;s)\defeq\sum_{l\in L^Y_2}\frac{w_l(q;s)\,u_Y(q+l)}{\Delta'_Y(q+l)},\qquad
w_l(q;s)\defeq\sum_{L\in\Lambda_Y^*}\chi(q+l+L)e^{\mathrm i(q+l+L)\cdot s},
\]
$\chi(Q)\defeq\prod_\alpha\hat\psi(\eta'Q\cdot f_\alpha)$, $\hat\psi(t)=\int\psi(\sigma)e^{-\mathrm it\sigma}d\sigma=12(2-2\cos t-t\sin t)/t^4$, $\Lambda_Y^*=\frac{2\pi}{\eta'}R^{-1}\Z^3$; and on $V_\eta$: $|w_l|\le W^3\le1.20$ for $l\ne0$, $|w_0(q;s)-e^{\mathrm iq\cdot s}|\le0.106\,\eta|q|$, $|M^Y_S(q;s)-\Delta_R^{-1}|\le207\,\eta|q|^{-1}$. Consequently, for $x'=x+s\in\Lambda_X$ in the tile of $x$, $|H^X(x',z)-E^S(x',z)|\le C^S_EJ_1\eta$ with $C^S_E\defeq380a+0.2a(C_*/20+8.3a)$ ($C_*$ of \Cref{rot:lem:H}), and the same at level $j$ in the unit scale of level $j$.
\end{lemma}
\begin{proof}
Poisson summation over $\Lambda_Y$ ($\hat\Psi(Q)=\eta'^3\chi(Q)$): $\sum_y\Psi(r-y)e^{\mathrm iQ\cdot y}=\sum_L\chi(Q+L)e^{\mathrm i(Q+L)\cdot r}$. Inserted in $H^Y(y,z)=(2\pi)^{-3}\int_{V_{\eta'}}e^{\mathrm iQ\cdot(y-z)}h(Q)dQ$ ($h$ periodic under $\Lambda_Y^*$, $e^{\mathrm iL\cdot z}=1$) this gives $E^S(r,z)=(2\pi)^{-3}\int_{\R^3}e^{\mathrm iQ\cdot(r-z)}\chi(Q)h(Q)dQ$; splitting $\R^3$ into the translates $V_\eta+l+L$ and using $e^{\mathrm i(l+L)\cdot x}=1$ gives the formula. Bounds: $\psi\ge0$ gives $|\hat\psi|\le\hat\psi(0)=1$; $\hat\psi(2\pi n)=0$ for $n\ne0$; for $|\tau|\le\pi/5$ and $n\ne0$, $|2-2\cos\tau-(2\pi n+\tau)\sin\tau|\le\tau^4/12+2\pi|n||\tau|\le1.004\cdot2\pi|n||\tau|$ and $|2\pi n+\tau|\ge0.9\cdot2\pi|n|$, so $|\hat\psi(2\pi n+\tau)|\le0.074|\tau|/|n|^3$ and $T(\tau)\defeq\sum_{n\ne0}|\hat\psi(\tau+2\pi n)|\le2\zeta(3)\cdot0.074|\tau|\le0.18|\tau|$ (below the cruder $0.213|\tau|$ is used); $1-\hat\psi(t)\le t^2/15$ on $[0,\pi]$ (the series $1-t^2/15+t^4/560-\dots$); $W\defeq\sup_t\sum_n|\hat\psi(t+2\pi n)|\le1.06$ (a one-dimensional bound: the sum is $1$ at the integers and its oscillation on a period is controlled by the bound on $T$). For $l\ne0$: $|w_l|\le\prod_\alpha\sum_n|\hat\psi(t_\alpha+2\pi n)|\le W^3$, $t_\alpha\defeq\eta'(q+l)\cdot f_\alpha$. For $l=0$ ($|t_\alpha|\le\pi/5$ on $V_\eta$, $t\defeq\eta'q$): $|w_0-e^{\mathrm iq\cdot s}|\le(1-\chi(q))+[\prod_\alpha(1+T(t_\alpha))-1]\le|t|^2/15+(e^{0.213\sum_\alpha|t_\alpha|}-1)\le(0.073+0.454)|t|\le0.106\,\eta|q|$ (convexity on $\sum|t_\alpha|\le3\pi/5$). Main term: $|w_0u_Y-1|\le0.106\eta|q|+|q||s|+(1-u_Y)\le(0.106+0.722+0.218)\eta|q|$, over $\Delta'_Y\ge\frac4{\pi^2}|q|^2$: $\le2.59\eta|q|^{-1}$; $|\Delta'^{-1}_Y-\Delta_R^{-1}|\le1.12\eta|q|^{-1}$ (proof of \Cref{rot:lem:H}); the $124$ aliases: $\le124\,W^3\eta^2/4\le37.2\eta^2\le203\eta|q|^{-1}$ ($|q|\le\sqrt3\pi/\eta$). So $|M^Y_S-\Delta_R^{-1}|\le207\eta|q|^{-1}$, and with $|M^X-\Delta_R^{-1}|\le172.1\eta|q|^{-1}$ at the same offset $s$ (proof of \Cref{rot:lem:H}), $|M^X-M^Y_S|\le380\eta|q|^{-1}$. Then the multiplier step of \Cref{rot:lem:H} verbatim: $|M^XA^X-M^Y_SA^Y|\le380a\eta|q|^{-1}+0.2a\eta^2A^Y|M^Y_S|\le C^S_E\eta|q|^{-1}$; integrate against $|u_{r-1}|$.
\end{proof}

\begin{corollary}[decaying forms at the $X$ sites]\label{rot:hyb:cor:DF}
Unit scale of level $r$; $x'\in\Lambda_X$ in the tile of $x\in\Gamma_1$, $x''=\tau x'$; $y'\in\Lambda_X$, $y''=\tau y'$.
\begin{enumerate}[label=\textup{(\alph*)},leftmargin=*]
\item (lifts) $|H^X(x',z)-E^S(x',z)|\le C(C^S_EJ_1\eta)^{1/2}e^{-\delta|x-z|/2}$ and $|H^Y(x'',z)-E^S(x',z)|\le C(375a\,J_1\eta)^{1/2}e^{-\delta|x-z|/2}$; the same for the photon lifts and for $E^S_j$ at level $j$ (unit scale of level $j$, $\eta\to\eta_j=5^{1-j}$).
\item (slices) for $2\le j<r$: $|S^X_j(x',y')-\tilde S^S_j(x',y'')|\le C(\eta'/\ell_j)^{1/2}\ell_j^{-1}e^{-\delta|x'-y'|/\ell_j}$.
\item (slices, the derivative at the lattice end) for $2\le j<r$: $|\partial_{y''}S^Y_j(x'',y'')-\partial_{y''}\tilde S^S_j(x',y'')|\le C(\eta'/\ell_j)^{1/2}\ell_j^{-2}e^{-\delta|x'-y'|/\ell_j}$.
\end{enumerate}
For $j\in\{0,1\}$ each term is bounded separately by \Cref{rot:hyb:lem:ESB} and the per-flow bounds of \Cref{rot:cor:unif}(a); there is no gain ($\ell_j\le5\eta'$).
\end{corollary}
\begin{proof}
(a) \Cref{rot:hyb:lem:ADMS} gives the sup bound $C^S_EJ_1\eta$; the second form is the same computation with $M^Y(q)$ at the cube offset $r_Y=x''-x$ in place of $M^X$: $|w_0u_Y-e^{\mathrm iq\cdot r_Y}u_Y|\le(0.106+c_\tau/5)\eta|q|$ and $\sum_{l\ne0}|w_l-e^{\mathrm i(q+l)\cdot r_Y}|\le124(W^3+1)$, giving $\le375a\,J_1\eta$. The decay is \Cref{rot:hyb:lem:ESB}(a) for $E^S$ and the separate bounds $|H^F(x^F,z)|\le Ce^{-\delta|x^F-z|}$ of \Cref{rot:lem:H}; take the geometric mean ($|x'-x|\le\sqrt{13}\eta'$). (b) In the unit scale of level $j$,
\[
S^X_j(x',y')-\tilde S^S_j(x',y'')=\sum_{z,z'}\big[(H^X_j-E^S_j)(x',z)\,C^X_j(z,z')\overline{H^X_j(y',z')}+E^S_j(x',z)(C^X_j-C^Y_j)(z,z')\overline{H^X_j(y',z')}+E^S_j(x',z)C^Y_j(z,z')\overline{(H^X_j(y',z')-H^Y_j(y'',z'))}\big]:
\]
(a), the bound $|C^X_j-C^Y_j|\le C\eta_je^{-\delta|z-z'|/2}$ of the proof of \Cref{rot:lem:slice}, \Cref{rot:hyb:lem:ESB}(a) and the decaying form of \Cref{rot:lem:H} at matched lattice points; sum the exponentials and rescale ($\eta_j=5\eta'/\ell_j$). This is \Cref{rot:lem:slice} with $E^S_j$ in place of $H^Y_j\circ\tau$ in the first argument. (c) The difference is $\sum_{z,z'}(H^Y_j(x'',z)-E^S_j(x',z))C^Y_j(z,z')\overline{\partial H^Y_j(y'',z')}$: the second form of (a), \cite[(3.6)]{K1}, and \cite[(3.7)]{K1} for the lattice gradient at $y''$; one more factor $\ell_j^{-1}$ on rescaling. (Alternatively, by \Cref{rot:hyb:lem:ESB}(c) the difference is $\cA(x'')-\cA(x')$ for $\cA\defeq\partial_{y''}\tilde S^S_j(\cdot,y'')$, so it is $\le c_\tau\eta'\sup|\nabla\cA|\le C(\eta'/\ell_j)\ell_j^{-2}e^{-\delta|x'-y'|/\ell_j}$ by \Cref{rot:hyb:lem:ESB}(b).)
\end{proof}

% r03_hyb_dslice: the slice derivative at matched points (the input of the one-derivative clause of the contour lines).

\begin{lemma}[the slice derivative at matched points]\label{rot:hyb:lem:DS}
Let $\alpha\in(0,1)$ be the H\"older exponent of \cite[(3.8)]{K1} (as in \Cref{rot:thm:H1S}) and put
\[
\theta_\partial\defeq\frac{\alpha}{2(2+\alpha)},\qquad\text{so that}\qquad\theta_1=\frac{\alpha^2}{4(1+\alpha)^2}\le\theta_\partial\le\tfrac14 .
\]
For $2\le j<r$, in the unit scale of level $r$, let $\nabla^X_yS^X_j(x',y')\defeq\sum_\mu\partial^X_\mu S^X_j(x',y')\,e_\mu$ be the vector of forward lattice differences of the $X$ slice in its lattice end $y'$, and $\nabla^Y_yS^Y_j(x'',y'')\defeq\sum_\nu\partial^Y_\nu S^Y_j(x'',y'')\,f_\nu$ the same for the $Y$ slice along the $Y$ frame $f_\nu=R^{-1}e_\nu$, both read as vectors in the common frame (\Cref{rot:lem:vert}(iii)). Then for matched $(x',y')$, $(x'',y'')=(\tau x',\tau y')$,
\[
\big|\nabla^X_yS^X_j(x',y')-\nabla^Y_yS^Y_j(x'',y'')\big|\ \le\ C\Big(\frac{\eta'}{\ell_j}\Big)^{\theta_\partial}\ell_j^{-2}e^{-\delta|x'-y'|/\ell_j},
\]
with $C=C(\alpha)$ uniform in $j$, $r$ and the torus; the same for the photon slices. For $j\in\{0,1\}$ each gradient is bounded separately by $C\ell_j^{-2}e^{-\delta|x'-y'|/\ell_j}$ (\cite[(3.63)]{K1} for $Y$, \Cref{rot:hyb:lem:A0}(a) for $X$) and $(\eta'/\ell_j)^{\theta_\partial}\ge5^{-1/4}$: there is no gain, and the factor is paid by a degree loss, as in \Cref{rot:cor:unif}(a).
\end{lemma}
\begin{proof}
Write $D_j\defeq S^X_j-S^Y_j\circ\tau$ as in \Cref{rot:cor:unif}, so that $|D_j(p',q')|\le C_D(\eta'/\ell_j)^{1/2}\ell_j^{-1}e^{-\delta|p'-q'|/\ell_j}$ for every pair of matched points (\Cref{rot:lem:slice}; $j\ge2$). The per-flow inputs are the gradient bound $|\nabla_yS^F_j(x,y)|\le C\ell_j^{-2}e^{-\delta|x-y|/\ell_j}$ and the H\"older clause of the gradient in the lattice end,
\[
\big|\nabla_yS^F_j(x,y)-\nabla_yS^F_j(x,q)\big|\le C'\ell_j^{-2}\Big(\frac{|y-q|}{\ell_j}\Big)^{\alpha'}e^{-\delta\min(|x-y|,|x-q|)/\ell_j}\qquad(|y-q|\le2\ell_j),\qquad\alpha'\defeq\alpha/2,
\]
for both flows. For $F=Y$ this is \cite[(3.63)]{K1} (\cite[Prop.~3.7]{K1}, from \cite[Thm.~3.3, (3.8)]{K1} by scaling) with King's constant, and $\alpha'<\alpha$ costs nothing on $|y-q|\le2\ell_j$. For $F=X$ it is \Cref{rot:hyb:lem:A0}(a) read in the lattice end ($S^X_j=H^X_jC^X_jH^{X*}_j$ is symmetric), whose factor $C_\alpha\log^{3\alpha}(2+\ell_j/\eta')\,(|y-q|/\ell_j)^{\alpha}$ is $\le C'(|y-q|/\ell_j)^{\alpha'}$: with $t\defeq|y-q|/\ell_j\in[\eta'/\ell_j,2]$ one has $\log(2+\ell_j/\eta')\le\log(2+1/t)$, and $t^{\alpha/2}\log^{3\alpha}(2+1/t)$ is bounded on $(0,2]$ (the device of \Cref{rot:hyb:rem:Hol}).

Fix a unit vector $e_\mu$ of the $X$ frame and a multiple $\rho=N\eta'$ of $\eta'$ with $\eta'\le\rho\le\ell_j$, to be chosen, and put $p'\defeq y'+\rho e_\mu\in\Lambda_X$, $p''\defeq\tau p'$. All points below lie within $2\ell_j$ of $y'$, so the decay factors of the inputs differ from $e^{-\delta|x'-y'|/\ell_j}$ by constants, which are not written.

\emph{The $X$ side.} Telescoping the forward differences along the segment, $S^X_j(x',p')-S^X_j(x',y')=\sum_{i=0}^{N-1}\eta'\,\partial^X_\mu S^X_j(x',y'+i\eta'e_\mu)$, and replacing each difference by the one at $y'$ with the H\"older clause,
\[
\Big|\rho^{-1}\big[S^X_j(x',p')-S^X_j(x',y')\big]-e_\mu\cdot\nabla^X_yS^X_j(x',y')\Big|\le C\ell_j^{-2}\Big(\frac{\rho}{\ell_j}\Big)^{\alpha'}e^{-\delta|x'-y'|/\ell_j}.
\]

\emph{The $Y$ side.} $p''-y''$ is a vector of the $Y$ lattice, $p''-y''=\sum_\nu n_\nu\eta'f_\nu$ with $n_\nu\eta'=(p''-y'')\cdot f_\nu$, and $p''-y''=\rho e_\mu+w$ with $w=(\tau p'-p')-(\tau y'-y')$, $|w|\le2c_\tau\eta'$ (\Cref{rot:def:match}). Join $y''$ to $p''$ by the lattice path that takes the $|n_1|$ steps along $\pm f_1$ first, then those along $\pm f_2$, then those along $\pm f_3$; it has at most $\sqrt3(\rho+2c_\tau\eta')/\eta'+3$ steps, and every point of it is within $\rho+2c_\tau\eta'+\eta'$ of $y''$. Telescoping along the path (a step along $-f_\nu$ from $q$ is $-\eta'\,\partial^Y_\nu S^Y_j(x'',q-\eta'f_\nu)$) and replacing each difference by the one at $y''$ with the H\"older clause,
\[
S^Y_j(x'',p'')-S^Y_j(x'',y'')=(p''-y'')\cdot\nabla^Y_yS^Y_j(x'',y'')+\mathrm{Rem},\qquad|\mathrm{Rem}|\le C\rho\,\ell_j^{-2}\Big(\frac{\rho+2c_\tau\eta'}{\ell_j}\Big)^{\alpha'}e^{-\delta|x'-y'|/\ell_j},
\]
and $(p''-y'')\cdot\nabla^Y_yS^Y_j=\rho\,e_\mu\cdot\nabla^Y_yS^Y_j+w\cdot\nabla^Y_yS^Y_j$ with $|w\cdot\nabla^Y_yS^Y_j(x'',y'')|\le2c_\tau\eta'\,C\ell_j^{-2}e^{-\delta|x'-y'|/\ell_j}$. Hence, with $(a+b)^{\alpha'}\le a^{\alpha'}+b^{\alpha'}$,
\[
\Big|\rho^{-1}\big[S^Y_j(x'',p'')-S^Y_j(x'',y'')\big]-e_\mu\cdot\nabla^Y_yS^Y_j(x'',y'')\Big|\le C\ell_j^{-2}\Big[\Big(\frac{\rho}{\ell_j}\Big)^{\alpha'}+\Big(\frac{\eta'}{\ell_j}\Big)^{\alpha'}+\frac{\eta'}{\rho}\Big]e^{-\delta|x'-y'|/\ell_j}.
\]

\emph{The comparison.} The pairs $(x',y'),(x'',y'')$ and $(x',p'),(x'',p'')$ are matched, so
\[
\rho^{-1}\Big|\big[S^X_j(x',p')-S^X_j(x',y')\big]-\big[S^Y_j(x'',p'')-S^Y_j(x'',y'')\big]\Big|\le\rho^{-1}\big(|D_j(x',p')|+|D_j(x',y')|\big)\le C\Big(\frac{\eta'}{\ell_j}\Big)^{1/2}\frac{\ell_j}{\rho}\,\ell_j^{-2}e^{-\delta|x'-y'|/\ell_j}.
\]
Adding the three displays, with $s\defeq\rho/\ell_j\in[\eta'/\ell_j,1]$,
\[
\big|e_\mu\cdot\big(\nabla^X_yS^X_j(x',y')-\nabla^Y_yS^Y_j(x'',y'')\big)\big|\le C\ell_j^{-2}e^{-\delta|x'-y'|/\ell_j}\Big[\Big(\frac{\eta'}{\ell_j}\Big)^{1/2}\frac1s+s^{\alpha'}+\Big(\frac{\eta'}{\ell_j}\Big)^{\alpha'}+\frac{\eta'}{\ell_j}\,\frac1s\Big].
\]
Choose $s$ as the smallest multiple of $\eta'/\ell_j$ that is $\ge(\eta'/\ell_j)^{1/(2(1+\alpha'))}$; since $1/(2(1+\alpha'))<1$ it lies in $[\eta'/\ell_j,1]$ and is $\le2(\eta'/\ell_j)^{1/(2(1+\alpha'))}$. The first two terms are then $\le C(\eta'/\ell_j)^{\alpha'/(2(1+\alpha'))}$, and so are the last two, because $\alpha'\ge\alpha'/(2(1+\alpha'))$ and $1-\frac1{2(1+\alpha')}=\frac{1+2\alpha'}{2(1+\alpha')}\ge\frac{\alpha'}{2(1+\alpha')}$. With $\alpha'=\alpha/2$, $\frac{\alpha'}{2(1+\alpha')}=\frac{\alpha}{2(2+\alpha)}=\theta_\partial$. Summing over the three $e_\mu$ gives the vector bound; the photon slices are the scalar ones times the identity. The inequalities $\theta_1\le\theta_\partial\le\frac14$ are $\alpha(2+\alpha)\le2(1+\alpha)^2$ and $2\alpha\le2+\alpha$.
\end{proof}

% r03_hyb_gain: proof of \Cref{rot:prop:BD}, part VII: the gain of every difference element at a non-zero gauge field (the lift of a constant field and the phases, the contour-line comparison, contour vertices, contour lines, double contour lines, the list); this completes part (b).

\paragraph{The difference elements at a non-zero gauge field.}
After the constant-background reduction of \Cref{rot:lem:CB}, applied to both flows on their infinite lattices with the torus lines as exact periodisations, every propagator that is compared between the flows is an $A=0$ line $G_r(\Lambda_F,0)=\sum_{j<r}S^F_j$, exactly the kernels of \Cref{rot:prop:fc,rot:lem:H,rot:lem:slice,rot:cor:unif}, with all vertices in the region $F_v$, and the background enters only through the explicit coarse fields $A_r(w_1)\cdots A_r(w_q)$ of \cite[(3.55)]{K1} (the same for both flows), the constant-background phases \cite[(3.46)]{K1}, the contour vertices \cite[(3.24)--(3.26)]{K1} with their contour lines \cite[(3.27)]{K1}, the external lifted fields $\tilde A^{(r),F}$, and the per-flow remainders, which are bounded by \Cref{rot:prop:BD}(a) at the regular fields of \Cref{rot:lem:CB}(R). The elements that are not $A=0$ lines are compared here.

\begin{lemma}[the lift of a constant field; the phases at matched vertices]\label{rot:hyb:lem:bph}
The lift \cite[(2.18)]{K1} of a constant coarse field $B$ to either flow is the constant fine field $\vr_rB$, $\vr_r\defeq a_r/(a_r+\lmu^2\delta_r^2)\le1$, the same number in both flows (\cite[(3.43)]{K1}). For matched $v'\in\Lambda_X$, $v''=\tau v'$ and the common root $y\in\Gamma_r$,
\[
\big|U(\vr_rB(\Gamma^X_{y,v'}))-U(\vr_rB(\Gamma^Y_{y,v''}))\big|\le e_r|B|\,c_\tau\eta'\le e\,p_r\lmu^{-1}c_\tau\,\delta_r^{-1/2}\eta'
\]
for $B$ in the small-field class.
\end{lemma}
\begin{proof}
$Q_F^*\mathbf 1_{\Gamma_r}=\mathbf 1_{\Lambda_F}$ for the frame-covariant photon average (\Cref{rot:lem:tile}(iii): $(\hat Q^*\psi)(x)=\sum_\alpha\psi_\alpha(y(x))f_\alpha$, which for the frame components $\psi_\alpha=f_\alpha\cdot B$ of a constant $B$ is $B$), and $(-\Delta+\lmu^2\delta_r^2)\mathbf 1=\lmu^2\delta_r^2\mathbf 1$, so $(M+a_rQ_F^*Q_F)\mathbf 1=(\lmu^2\delta_r^2+a_r)\mathbf 1$ and $a_rG^{F,\mathrm{ph}}_rQ_F^*B=\vr_rB$ for both flows. Both phases are $e^{\mathrm ie_r\vr_rB\cdot(v-y)}$ with $v=v'$, $v''$ (\Cref{rot:hyb:lem:alg}(iii)); $|v'-v''|\le c_\tau\eta'$, $\vr_r\le1$, $|B|\le p_r(\lmu\delta_r)^{-1}$ (\cite[(3.2)]{K1}), $e_r=e\delta_r^{1/2}$.
\end{proof}

\begin{lemma}[contour-line comparison]\label{rot:hyb:lem:CL}
Let $\cA$ be a $C^1$ vector field on a neighbourhood of the level-$r$ block $z+R^{-1}[-\frac12,\frac12]^3$ with $|\cA|\le P_0$, $|\nabla\cA|\le P_1$ there. Let $a^X_b\defeq\cA(x_b)\cdot e_\mu$ for the $\Lambda_X$-bond $b=(x_b,x_b+\eta'e_\mu)$ and $a^Y_{b''}\defeq\cA(x_{b''})\cdot f_\alpha$ for the $\Lambda_Y$-bond $b''=(x_{b''},x_{b''}+\eta'f_\alpha)$ (site sampling, King's identification of bond and site fields). Then for $x'\in B^X_r(z)$, $x''=\tau x'$,
\[
\big|a^X(\Gamma^{X,(r)}_{z,x'})-a^Y(\Gamma^{Y,(r)}_{z,x''})\big|\ \le\ 10\,(P_0+P_1)\,\eta' .
\]
The same holds for the partial contours $\Gamma^{(j+1)}_{x_{j+1},x}$ of \cite[(2.19)]{K1}, with $P_0,P_1$ taken on the level-$(j+1)$ block, the sliver area $\le1.5\eta'\ell_{j+1}+33\eta'^2$ in the flux term and the lengths $\le1.5\ell_{j+1}$ ($Y$), $\le2.1\ell_{j+1}$ ($X$) in the Riemann term: the bound is $\le(7.1P_0+(6.6\ell_{j+1}+66\eta')P_1)\eta'\le C(P_0+P_1\ell_{j+1})\eta'$.
\end{lemma}
\begin{proof}
\emph{Riemann step.} For a bond $b$ of length $\eta'$ sampled at its initial point, $|\eta'\,\cA(x_b)\cdot\hat b-\int_b\cA\cdot dl|\le\eta'\cdot\eta'P_1$. Summing over the $\le2.1/\eta'$ bonds of $\Gamma^X$ and the $\le1.5/\eta'$ bonds of $\Gamma^Y$: $|a^X(\Gamma^X)-\int_{\Gamma^X}\cA|+|a^Y(\Gamma^Y)-\int_{\Gamma^Y}\cA|\le3.6\,P_1\eta'$. \emph{Stokes step.} $\int_{\Gamma^X}\cA-\int_{\Gamma^Y}\cA=\oint_{\partial\Sigma}\cA-\int_{[x',x'']}\cA=\int_\Sigma\curl\cA\cdot dS-\int_{[x',x'']}\cA$ with the surface $\Sigma=\Sigma(z,x')$ of \Cref{rot:hyb:lem:len}(iii) (piecewise planar, so Stokes' theorem applies piece by piece and the interior edges cancel); $|\curl\cA|\le2P_1$, $\mathrm{area}(\Sigma)\le1.5\eta'+33\eta'^2\le3\eta'$ for $r\ge2$, and $|[x',x'']|\le c_\tau\eta'$. Total $\le(3.6+6)P_1\eta'+7.1P_0\eta'\le10(P_0+P_1)\eta'$. For partial contours the same decomposition applies to the sub-polygon up to level $j+1$ (\Cref{rot:hyb:lem:len}(i),(iii)).
\end{proof}

\begin{lemma}[contour vertices at the lifted background]\label{rot:hyb:lem:CV}
Let $A'^F=\vartheta\tilde A^{(r),F}+za_F$ be the reduced background of the flow $F$ (\Cref{rot:lem:CB}). Then for matched $(y,x')$, $(y,x'')$,
\[
\big|U(A'^X(\Gamma^{X,(r)}_{y,x'}))-U(A'^Y(\Gamma^{Y,(r)}_{y,x''}))\big|\ \le\ e_r\big|A'^X(\Gamma^X_{y,x'})-A'^Y(\Gamma^Y_{y,x''})\big|\ \le\ C\,e_r\,p_r^2\,\eta^{1/2},
\]
and hence the kernels of $F^X_{2,r}(A'^X,0)$ and of $F^Y_{2,r}(A'^Y,0)\circ\tau$ (\cite[(3.51)]{K1}, the contour vertices (3.24)--(3.26)) differ entrywise by $\le Ce_rp_r^2\eta^{1/2}$ against their own size $\le Ce_rp_r^2$: a relative gain $\eta^{1/2}$, and an absolute gain $\eta^{1/2}e\,p_r^2\delta_r^{1/2}$. The products (3.24)--(3.26)$\times$(3.26) at the same $y$ inherit the bound.
\end{lemma}
\begin{proof}
Let $\cA\defeq\vartheta\,E^S_{\mathrm{ph}}\tilde A_r+za^{\mathrm c}$, where $E^S_{\mathrm{ph}}\tilde A_r\defeq I_S[H^{Y,\mathrm{ph}}\tilde A_r]$ is the cubic-Hermite extension of the photon lift of $\tilde A_r$ along the flow $Y$ (\Cref{rot:hyb:def:ES}; componentwise in any frame) and $a^{\mathrm c}$ the smooth Cameron--Martin form (\Cref{rot:lem:CM}). $\cA$ is $C^1$ (\Cref{rot:hyb:lem:ESH}(iv)). By \Cref{rot:hyb:lem:ESH}(iii) and \Cref{rot:lem:CB}(R), $|E^S_{\mathrm{ph}}\tilde A_r(r)|\le\max_\sigma|\tilde A^{(r),Y}(y_\sigma)|\le Cp_r(|r-x_0|+\sqrt3\eta')$ and $|\nabla E^S_{\mathrm{ph}}\tilde A_r(r)|\le\frac{3\sqrt3}2\max_{\mathrm{edges}}|\partial\tilde A^{(r),Y}|\le Cp_r$; with \Cref{rot:lem:CM} for $a^{\mathrm c}$, $P_0,P_1\le Cp_r^2$ on the vertex region (of diameter $\le Cp_r$). Then $a^Y_{b''}=\cA(x_{b''})\cdot f_\alpha$ exactly (\Cref{rot:hyb:lem:ESB}(c); \Cref{rot:lem:CM} for $a_Y$), while
\[
|a^X_b-\cA(x_b)\cdot e_\mu|\le\vartheta\sum_z|H^{X,\mathrm{ph}}(x_b,z)-E^S_{\mathrm{ph}}(x_b,z)||\tilde A_r(z)|+|a_{X,\mu}(x_b)-a^{\mathrm c}\cdot e_\mu|\le C(C^S_EJ_1\eta)^{1/2}p_r^2+C\eta'^2
\]
by \Cref{rot:hyb:cor:DF}(a) at the $X$ site $x_b$ (the decaying form; $|\tilde A_r(z)|\le p_r(|z-x_0|+1)$, the contributions from far away as in \Cref{rot:lem:CB}(R)) and \Cref{rot:lem:CM}. Summing the lift mismatch with weight $\eta'$ over the $\le2.1/\eta'$ bonds of $\Gamma^X$ gives $\le C\eta^{1/2}p_r^2$, and \Cref{rot:hyb:lem:CL} gives $\le10(P_0+P_1)\eta'\le Cp_r^2\eta'$ for the rest. Multiply by $e_r$ for the phases ($|U(a)-U(a')|\le e_r|a-a'|$), and by $\eta'^3$ per summand of (3.51).
\end{proof}

\begin{lemma}[{the contour lines \cite[(3.27)]{K1}}]\label{rot:hyb:lem:CLN}
For $0\le j<r$ let $S^{F,\mathrm{ph}}_j$ be the stage-$j$ photon slice kernel (an $A=0$ kernel: the scalar slice $S^F_j$ times the identity, frame-free), and let $\cL^F_j(x;y)\defeq\sum_{b_1\in\Gamma^{F,(j+1)}_{x_{j+1},x}}\eta'\,S^F_j(x_{b_1},y)\,\hat b_1\in\R^3$ ($\hat b_1$ the direction of $b_1$, $y\in\Lambda_F$) be King's contour line (3.27) as a vector: its component along the direction of the photon bond $b=(y,y+\eta'\hat e)$ is $G_{(j)}(\Gamma^{(j+1)}_{x_{j+1},x},b)$ of \cite[(3.27)]{K1}. Lines are compared as vectors (\Cref{rot:lem:vert}(iii)). Then for matched $(x',y')$, $(x'',y'')=(\tau x',\tau y')$,
\[
\big|\cL^X_j(x';y')-\cL^Y_j(x'';y'')\big|\ \le\ C\Big(\frac{\eta'}{\ell_j}\Big)^{\theta}e^{-\delta\dist(\Gamma^{(j+1)},y')/\ell_j},\qquad\theta=\tfrac12,
\]
against the unweakened bound $Ce^{-\delta\dist/\ell_j}$ of each line: a slice-$j$ line weakened by the degree $\theta$ in the sense of \Cref{rot:thm:H1S}, with the prefactor $\eta'^\theta$. With one lattice derivative at the lattice end $y$ (the gradients of the two lines read as vectors in the common frame, \Cref{rot:hyb:lem:DS}) the same holds with one more factor $\ell_j^{-1}$ and, in place of $(\eta'/\ell_j)^{1/2}$, the gain $(\eta'/\ell_j)^{\theta_\partial}$ of \Cref{rot:hyb:lem:DS}, $\theta_\partial=\alpha/(2(2+\alpha))\ge\theta_1$: a slice-$j$ line with one derivative, weakened by the degree $\theta_\partial$.
\end{lemma}
\begin{proof}
Keep the second argument of the $Y$ kernel at the lattice point $y''$ throughout (no extension in the second argument is used) and write
\[
\cL^X_j(x';y')-\cL^Y_j(x'';y'')=\sum_{b_1\in\Gamma^X}\eta'\big[S^X_j(x_{b_1},y')-\tilde S^S_j(x_{b_1},y'')\big]\hat b_1
+\Big[\sum_{b_1\in\Gamma^X}\eta'\tilde S^S_j(x_{b_1},y'')\hat b_1-\sum_{b''_1\in\Gamma^Y}\eta'S^Y_j(x_{b''_1},y'')\hat b''_1\Big].
\]
\emph{First sum} ($2\le j<r$): \Cref{rot:hyb:cor:DF}(b) at the $X$ sites $x_{b_1}$, per bond $\le C(\eta'/\ell_j)^{1/2}\ell_j^{-1}e^{-\delta|x_{b_1}-y'|/\ell_j}$; with weight $\eta'$ over the $\le10.5\ell_j/\eta'$ bonds of the partial $X$ contour (\Cref{rot:hyb:lem:len}(i)): $\le C(\eta'/\ell_j)^{1/2}e^{-\delta\dist(\Gamma^X,y')/\ell_j}$. \emph{Bracket}: for a unit vector $u$, $u\cdot[\,\cdot\,]=a^X(\Gamma^X)-a^Y(\Gamma^Y)$ for the field $\cA_u(r)\defeq\tilde S^S_j(r,y'')\,u$, which is $C^1$ (\Cref{rot:hyb:lem:ESB}(c)), whose $Y$ samples are exactly $S^Y_j(x_{b''},y'')\,u\cdot f_\alpha$, and which has $P_0\le C\ell_j^{-1}e^{-\delta\dist(\Gamma,y'')/\ell_j}$, $P_1\le C\ell_j^{-2}e^{-\delta\dist(\Gamma,y'')/\ell_j}$ on the level-$(j+1)$ block (\Cref{rot:hyb:lem:ESB}(b); the block's diameter $\sqrt3\ell_{j+1}$ and $|y'-y''|\le c_\tau\eta'$ change only the constant). \Cref{rot:hyb:lem:CL} for partial contours gives $\le C(P_0+P_1\ell_{j+1})\eta'\le C(\eta'/\ell_j)e^{-\delta\dist/\ell_j}$ (flux, end segment and Riemann terms together: the kernel difference along the contour plus the flux of the slice kernel through the sliver, as \Cref{rot:cor:unif}(c) states); $u$ over an orthonormal basis gives the vector bound. The matching $y'\leftrightarrow y''$ of the second argument sits inside \Cref{rot:hyb:cor:DF}(b). \emph{The derivative version} ($2\le j<r$). With the lattice derivative on the lattice end, split the first sum at the matched $Y$ point: $\sum_{b_1\in\Gamma^X}\eta'\big[\nabla^X_yS^X_j(x_{b_1},y')-\nabla^Y_yS^Y_j(\tau x_{b_1},y'')\big]\hat b_1+\sum_{b_1\in\Gamma^X}\eta'\big[\nabla^Y_yS^Y_j(\tau x_{b_1},y'')-\nabla^Y_y\tilde S^S_j(x_{b_1},y'')\big]\hat b_1$ (gradients as vectors, the matrix $[\,\cdot\,]\hat b_1$ contracted on the line's index). The first sum is \Cref{rot:hyb:lem:DS} at the matched pairs $(x_{b_1},y')$, $(\tau x_{b_1},y'')$, per bond $\le C(\eta'/\ell_j)^{\theta_\partial}\ell_j^{-2}e^{-\delta|x_{b_1}-y'|/\ell_j}$, hence $\le C(\eta'/\ell_j)^{\theta_\partial}\ell_j^{-1}e^{-\delta\dist(\Gamma^X,y')/\ell_j}$ after the sum over the $\le10.5\ell_j/\eta'$ bonds with weight $\eta'$; the second sum is \Cref{rot:hyb:cor:DF}(c) per bond, $\le C(\eta'/\ell_j)^{1/2}\ell_j^{-1}e^{-\delta\dist/\ell_j}$ after the same sum; and the bracket is \Cref{rot:hyb:lem:CL} for the $C^1$ field $\partial_{y''}\tilde S^S_j(\cdot,y'')\,u$ (\Cref{rot:hyb:lem:ESB}(b),(c): $P_0\le C\ell_j^{-2}e^{-\delta\dist/\ell_j}$, $P_1\le C\ell_j^{-3}e^{-\delta\dist/\ell_j}$), $\le C(\eta'/\ell_j)\ell_j^{-1}e^{-\delta\dist/\ell_j}$. The weakest gain is the first sum's, $(\eta'/\ell_j)^{\theta_\partial}$. \emph{$j\in\{0,1\}$}: each line, also with the derivative, is bounded separately by $Ce^{-\delta\dist/\ell_j}$ (times $\ell_j^{-1}$ with the derivative) (\cite[(3.64)]{K1} for $Y$; \Cref{rot:hyb:lem:A0}(a) summed along $\le10.5\ell_j$ for $X$), and $(\eta'/\ell_j)^{1/2}\ge5^{-1/2}$: there is no gain and the factor $\eta'^\theta$ is paid by a degree loss, exactly as in \Cref{rot:cor:unif}(a).
\end{proof}

\begin{lemma}[double contour lines]\label{rot:hyb:lem:DCL}
For $0\le j<r$ and $x_1,x_2\in\Lambda_F$ let
\[
\cD^F_j(x_1;x_2)\defeq\sum_{b_1\in\Gamma^{F,(j+1)}_{x_{1,j+1},x_1}}\sum_{b_2\in\Gamma^{F,(j+1)}_{x_{2,j+1},x_2}}\eta'^2S^F_j(x_{b_1},x_{b_2})\,\hat b_1\cdot\hat b_2,
\]
the line joining two contour factors of \cite[(3.26)]{K1} (at one vertex with $m\ge2$, at the two factors of a product (3.26)$\times$(3.26), or at two vertices). Then $|\cD^F_j(x_1;x_2)|\le C\ell_je^{-\delta|x_{1,j+1}-x_{2,j+1}|/\ell_j}$ for both flows, and at matched points $|\cD^X_j(x_1';x_2')-\cD^Y_j(x_1'';x_2'')|\le C(\eta'/\ell_j)^{1/2}\ell_je^{-\delta|x_{1,j+1}-x_{2,j+1}|/\ell_j}$. In Ba{\l}aban's degree $D(H)=\sum_ld_l+\sum_va_v+3(|V(H)|-1)$ (\cite[(2.1)--(2.2)]{Ba3}, with the two summation points of a block product as vertices, as in \Cref{rot:cor:OF}) a double line has $d_{l_0}=+1$, and every connected all-scale-summed subgraph $H$ containing it has $D(H)\ge1$: the double line is never the binding constraint of the weakened-element bound.
\end{lemma}
\begin{proof}
$\cD^F_j(x_1;x_2)=\sum_{b_2}\eta'\hat b_2\cdot\cL^F_j(x_1;x_{b_2})$. The size bound: lengths $\le10.5\ell_j$ ($X$), $\le7.5\ell_j$ ($Y$) (\Cref{rot:hyb:lem:len}(i)), \Cref{rot:hyb:lem:A0}(a), \cite[(3.63)]{K1}; since both contours lie in open rotated cubes of side $5\ell_j$ about their roots, $e^{-\delta|x_{b_1}-x_{b_2}|/\ell_j}\le e^{5\sqrt3\delta}e^{-\delta|x_{1,j+1}-x_{2,j+1}|/\ell_j}$. For $j\ge2$ put $\cV(r)\defeq\sum_{b_1''}\eta'\,I_S[S^Y_j(x_{b_1''},\cdot)](r)\,\hat b_1''$ (extension in the second argument); then
\[
\cD^X-\cD^Y=\sum_{b_2}\eta'\hat b_2\cdot[\cL^X_j(x_1';x_{b_2})-\cL^Y_j(x_1'';\tau x_{b_2})]+\sum_{b_2}\eta'\hat b_2\cdot[\cV(\tau x_{b_2})-\cV(x_{b_2})]+[a^X(\Gamma^X_2)-a^Y(\Gamma^Y_2)]_{\cA=\cV}
\]
exactly (\Cref{rot:lem:vert}(i); $\tau x_{b_2}\in\Lambda_Y$); the three terms are bounded by \Cref{rot:hyb:lem:CLN} summed along $\Gamma^X_2$, by $c_\tau\eta'\sup|\nabla\cV|$ with $|\nabla\cV|\le C\ell_j^{-1}$ (\Cref{rot:hyb:lem:ESH}(iii), \Cref{rot:hyb:lem:ESB}(b)), and by \Cref{rot:hyb:lem:CL} for partial contours with $P_0\le C$, $P_1\le C\ell_j^{-1}$. For $j\le1$ use the size bounds and $(\eta'/\ell_j)^{1/2}\ge5^{-1/2}$. Degree: the size $C\ell_j$ against $\ell_j^{-1}$ for a slice line ($d_l=-1$) and $\ell_j^0$ for a one-ended contour line ($d_l=0$) gives $d_{l_0}=+1$ (\Cref{rot:cor:OF}: $d_l=-1+\#(\text{contour ends})$). For a connected all-scale-summed $H$ containing $l_0$: if $H\setminus l_0$ is connected (or $l_0$ is a self-loop), $D(H)=D(H\setminus l_0)+1\ge1$, because $D\ge0$ for every all-scale-summed subgraph that does not contain $l_0$ (the per-flow renormalised structure, \Cref{rot:cor:OF}); if $l_0$ is a bridge between $H_1$ and $H_2$, $D(H)=D(H_1)+D(H_2)+1+3\ge4$. So $D(H)-\theta\ge\frac12$ for $\theta=\frac12$.
\end{proof}

\begin{proposition}[every difference element at $A\ne0$]\label{rot:hyb:prop:list}
In the setting of \Cref{rot:prop:BD}(b), after \Cref{rot:lem:CB}, order the elements of a renormalised family consistently across the two flows and telescope as in \Cref{rot:prop:cp}. Every difference element $e^X_i-e^Y_i$ is of one of the following types, with the stated bound:
\begin{enumerate}[label=\textup{(\roman*)},leftmargin=*]
\item an $A=0$ line: a fine slice $S^F_j$, $j<r$ (\Cref{rot:cor:unif}(a): $C_D\eta'^{\theta}\ell_j^{-1-\theta}e^{-\delta|x'-y'|/\ell_j}$, $\theta=\frac12$), a coarse slice or $\Gamma^F_n$ (\Cref{rot:cor:unif}(b): $C\eta^{1/2}\ell_j^{-1}e^{-\delta|x'-y'|/\ell_j}$), the coarse multiplier (\Cref{rot:prop:fc}: $0.2a\eta^2$ relative), on $\Lambda_F$;
\item a contour line (3.27): \Cref{rot:hyb:lem:CLN}, gain $(\eta'/\ell_j)^{1/2}$; with one lattice derivative at its lattice end, gain $(\eta'/\ell_j)^{\theta_\partial}$, $\theta_\partial=\alpha/(2(2+\alpha))$ (\Cref{rot:hyb:lem:DS});
\item a contour vertex (3.24)--(3.26): \Cref{rot:hyb:lem:CV}, entrywise $\le Ce_rp_r^2\eta^{1/2}$, relative gain $\eta^{1/2}$;
\item an external lifted field $\tilde A^{(r),F}$, $H^F\phi_r$ or its covariant derivative: \Cref{rot:lem:H,rot:lem:grad}, $\le C\eta^{1/2}$ times the field norm;
\item a constant-background phase (3.46): \Cref{rot:hyb:lem:bph}, $\le e\,p_r\lmu^{-1}c_\tau\delta_r^{-1/2}\eta'$;
\item a hopping vertex of $V_r(A',0)$ or a vertex function (couplings, partition of unity, the Cameron--Martin form): \Cref{rot:lem:vert}(ii),(iii), unchanged;
\item a double contour line, joining two contour factors of (3.26): \Cref{rot:hyb:lem:DCL}, size $\le C\ell_j$, relative gain $(\eta'/\ell_j)^{1/2}$ on slice $j$.
\end{enumerate}
Each weakened element therefore obeys the bound of its type weakened by the degree $\theta=\frac12$ (by $\theta_\partial\ge\theta_1$ for a contour line carrying a lattice derivative, which is all \Cref{rot:thm:H1S} asks) with a constant $\le C\,p_r^{2}\,\delta_r^{-1/2}$ (the worst of (iii) and (v)); the per-flow remainders (\Cref{rot:lem:CB}(2)) contribute $(Ce_rp_r^2)^{n+1}+e^{-\delta'p_r}$ per graph.
\end{proposition}
\begin{proof}
\Cref{rot:lem:CB} lists what carries the background after the reduction; the compared lines are $A=0$ lines, the contour lines, one-ended (3.27) and double (both ends on contours), are the only lines whose shape depends on the flow, and the only background-dependent vertices are the contour vertices (3.24)--(3.26) (\cite[p.~658]{K1}: ``The remaining vertices occur when the renormalization transformation of the scalar field is expanded about $A'=0$'') and the phases (3.46). The external coarse fields $A_r(w_i)$ of (3.55) are common. The bounds are the cited lemmas.
\end{proof}

\begin{proof}[Proof of \Cref{rot:prop:BD}(b)]
The difference elements $\Delta\Gamma$, $\Delta H$, $\Delta(A^{-1})$ at regular $A$ are, after the reduction, the $A=0$ differences of \Cref{rot:cor:unif} (types (i) and (iv) of \Cref{rot:hyb:prop:list}), with the factor $\eta'^{1/2}$. The contour vertices at the reduced background differ entrywise by a relative $\eta^{1/2}$ against their size $Ce_rp_r^2$ (\Cref{rot:hyb:lem:CV}); the constant-background phases at matched vertices differ by $\le e\,p_r\lmu^{-1}c_\tau\delta_r^{-1/2}\eta'$ (\Cref{rot:hyb:lem:bph}); the per-slice pieces carried by the contour vertex, with one end on a contour (\cite[(3.27)]{K1}) or both ends on contours, obey the slice-$j$ bounds \cite[(3.63)--(3.65)]{K1} in both flows ((a) for the hybrid flow, \cite[Prop.~3.7]{K1} for the standard flow) with the relative mismatch $(\eta'/\ell_j)^{1/2}$ (\Cref{rot:hyb:lem:CLN,rot:hyb:lem:DCL}), the one-ended lines also with one lattice derivative at the lattice end, there with the relative mismatch $(\eta'/\ell_j)^{\theta_\partial}$ of \Cref{rot:hyb:lem:DS} (\Cref{rot:hyb:lem:CLN}; a double line has no lattice end). A constant background is pure gauge on the infinite lattice and gives the same phase, up to the matching segment, for every contour (\Cref{rot:hyb:lem:bph}); what would break (b) are mis-paired contours or mis-rooted blocks, both excluded by \Cref{rot:def:flows} and \Cref{rot:conv:C}(i). The gains are not uniform in the size of the constant part of the background: they carry the small-field bound \cite[(3.2)]{K1}, as King's own reduction does.
\end{proof}


\begin{proposition}[King's upper bound for the rotated lattice]\label{rot:prop:H3}
For $\Yk$ on $T$ and every real smooth periodic source $s=(zG,h)$, with the free-boundary convention of \Cref{rot:def:tori} and the centred blocks of \Cref{rot:conv:C},
\[
\Zun_{\Yk,T}(s)\le e^{C\delta_r^\sigma|T|}\int\chi_r\,e^{-S^{(r)}_T}+e^{-p(\delta_r)^2+C|T|}\,\eps^{-q},
\]
$C=C(s,|T|)$, $q=q(s)<\infty$, with $q=0$ when $h=0$.
\end{proposition}
\begin{proof}
On $T_2=R^{-1}\DJ$, which is aligned in $Y$'s frame, the passive rotation $R$ (\Cref{rot:lem:walledalg}(b)) maps $(\Yk,T_2)$ to $(\Xk,\DJ)$ with its standard frame, so the statement is King's Theorem~3.1 \cite[(3.4)]{K1}, with the additive large-field term in the form in which King uses it (\cite[(3.11)]{K1}, \cite[(4.7)]{K2}), on King's torus at block factor $5$: this is \Cref{thm:a06:BALA}(iii) with \Cref{hyp:k01:BALint}, read with the torus-convention counterterms replaced by the free-boundary ones at a constant-shifted source (\Cref{rot:prop:CT}(b)) and with corner-rooted blocks replaced by centred ones (\Cref{rot:conv:C}(i)); King's constant depends on $g$, $h$ and $|T|$ and is uniform under a shift of the scalar source by a bounded constant, which is the uniformity of \cite[Thm.~3.1]{K1} in the mass over a compact mass set (\cite[p.~655]{K1}). On $T_1=\DJ$ the frame torus of $\Yk$ is $R\DJ$, skew in $Y$'s frame at every scale; this instance is \Cref{rot:prop:SK}. For a walled source ($\int h\ne0$) the constant of the large-field term inherits \cite[(2.23)]{K1}, which fails as printed when $\int h\ne0$ (\Cref{rot:rem:log}); what the inductive proof supplies is a bounded $h$-linear $c$-number per step, so the large-field term is read with the polynomial envelope factor $\eps^{-q}$, which \Cref{rot:prop:Kforward} absorbs.
\end{proof}

% r03_r1_setting: the cumulant bound for the step Gaussian, part I: the setting, the lattice facts, the constants, and the kernel and decoupling lemmas. Input from sec/r03_fill_K5a.tex before Theorem rot:thm:R1.

\subsubsection*{The cumulant bound for the step Gaussian}

The fluctuation integral of every block-spin step of both flows has the form $\int\chi_b\,e^{V}d\mu_P$: $\mu_P$ is a centred Gaussian whose precision is coercive and has exponentially decaying off-diagonal entries, $V$ is a real polynomial with exponentially decaying coefficients, and $\chi_b$ is a product over sites of indicators of a fixed convex body scaled by the threshold $b$. \Cref{rot:thm:R1} below is the order-$\bar n$ cumulant expansion of such an integral, with an error proportional to the volume of the support of $V$ and, in the lower bound, a normalisation term proportional to $|\Lambda|$. The printed cluster-expansion lemma of Benfatto and coauthors, which King cites for this step, is stated for the nearest-neighbour Markov field with a cutoff growing with the distance to the support of the interaction; its proof uses the Markov property and the nearest-neighbour formulas for conditional Gaussians in four places. The proof below keeps its architecture (cells with corridors, four shifted pavements, the cutoff moved from the corridors into the cells) and replaces those four uses by the decoupling and conditioning lemmas of this part, which hold for any coercive precision with exponentially decaying entries.

\paragraph{The setting.} Let $\mathbb{L}$ be one of: the lattice $\Z^3$; a cubic torus $(\Z/N\Z)^3$; a skew torus $\Z^3/mG_\pm$, $m\in\N$, where $G_+\defeq\langle(40,-30,0),(30,40,0),(0,0,50)\rangle$ and $G_-$ is its mirror image under $x_2\mapsto-x_2$. These are the frame period lattices of the skew tori of \Cref{rot:prop:SK} in the units of a step lattice: $\Gamma'=R^{\mathrm T}(50MJ\,\Z^3)=MJ\,G_+$ for $R\DJ$ in $Y$'s frame and $RP_J=MJ\,G_-$ for $R\DJ$ in $X$'s frame, so $m=MJ$ divided by the spacing of the step, an integer because every block lattice contains $P_J$ (\Cref{rot:lem:adm}(iii)). Distances are quotient Euclidean distances, $|x-y|\defeq\min_g|\tilde x-\tilde y-g|$ over the period lattice (a metric), and $\dist(x,S)\defeq\min_{y\in S}|x-y|$. $\Lambda\subset\mathbb{L}$ is finite, $z=(z_x)_{x\in\Lambda}$ with $z_x\in\R^n$, and a complex Hermitian operator on $\C^\Lambda$ is read in its real form (blocks $\bigl(\begin{smallmatrix}\Re&-\Im\\\Im&\Re\end{smallmatrix}\bigr)$), which preserves the spectrum and the block operator norms. The hypotheses are:
\begin{enumerate}[label=\textup{(R\arabic*)},leftmargin=*]
\item $\mu_P\defeq N(0,P^{-1})$ with $P$ real symmetric on $(\R^n)^\Lambda$ and $P\ge\gamma_0>0$ (no upper bound is needed);
\item $\|P(x,y)\|\le c_Pe^{-\delta_P|x-y|}$ for $x\ne y$ (block operator norms); the diagonal blocks are arbitrary;
\item $V=\sum_{Y\subset J}V_Y$ with $J\subset\Lambda$, where $V_Y$ is a real polynomial in $(z_y)_{y\in Y}$ of degree at most $D$, without constant term, each of whose monomials involves every site of $Y$ (its \emph{support}); $\ell(Y)$ is the length of a shortest connected graph on a vertex set containing $Y$ (edges between points of $\mathbb{L}$, of length their distance); $|V_Y|_1$ is the sum of the moduli of the coefficients; for $\theta\ge0$ and for any polynomial $F=\sum_YF_Y$ sorted by supports in this way,
\[
\|F\|_\theta\defeq\sup_{x\in\mathbb{L}}\sum_{Y\ni x}e^{\theta\ell(Y)}|F_Y|_1,\qquad\mathsf A\defeq\|V\|_\varsigma\ \text{ for a fixed rate }\varsigma>0,\qquad V_R\defeq\sum_{Y\subset R}V_Y;
\]
\item $\chi_\beta(z)\defeq\prod_{x\in\Lambda}\mathbf 1\{z_x\in\beta K\}$ with $K\subset\R^n$ a closed symmetric convex body, $\bar B(1)\subset K\subset\bar B(c_K)$: the Euclidean ball ($c_K=1$) or Ba{\l}aban's two-component cutoff $K=\{|z^{A}|\le1,\ |z^{\phi}|\le1\}$, $c_K=\sqrt2$, which restricts the photon and the scalar fluctuation at each site separately.
\end{enumerate}
$\kappa_\mu(F_1,\dots,F_k)$ denotes the joint cumulant under a law $\mu$ and $\kappa_\mu(F;k)\defeq\kappa_\mu(F,\dots,F)$ the $k$-th cumulant, the truncated expectation $\langle F^k\rangle^T$. For $X\subset\Lambda$, $P_{XX}$ is the principal block, $G_X\defeq(P_{XX})^{-1}$ and $C\defeq P^{-1}=G_\Lambda$; for $Y\defeq\Lambda\setminus X$ the conditional law of $z_X$ given $z_Y$ under $\mu_P$ is $N(-G_XP_{XY}z_Y,\,G_X)$, by completing the square in the exponent. $\chi^X_\beta\defeq\prod_{x\in X}\mathbf 1\{z_x\in\beta K\}$.

\begin{lemma}[lattice facts]\label{rot:lem:R1lat}
\textup{(i)} $x\mapsto|x-y_0|$ and $x\mapsto\dist(x,S)$ are $1$-Lipschitz. \textup{(ii)} For every $f\ge0$ and $x\in\mathbb{L}$, $\sum_{y\in\mathbb{L}}f(|x-y|)\le\sum_{z\in\Z^3}f(|z|)$. \textup{(iii)} For $x,y\in Y$, $|x-y|\le\ell(Y)$; a shortest connected graph for $Y$ joins any two of its vertices by a path of length at most $\ell(Y)$; and $\ell(Y')\le\ell(Y)$ for $Y'\subset Y$. \textup{(iv)} For $\theta\le\varsigma$ and $r\ge0$ the part of $V$ made of the monomials of length $\ell(Y)\ge r$ has $\|\cdot\|_\theta\le\mathsf Ae^{-(\varsigma-\theta)r}$; in particular $\sum_{Y\ni x,\ \ell(Y)\ge r}|V_Y|_1\le\mathsf Ae^{-\varsigma r}$.
\end{lemma}
\begin{proof}
(i) is the triangle inequality of the quotient metric. (ii) Fix a lift $\tilde x$ and for each $y$ choose the lift $\tilde y$ nearest to $\tilde x$, so that $f(|x-y|)=f(|\tilde x-\tilde y|)$; distinct $y$ have distinct lifts, so the sum is a sub-sum of $\sum_{\tilde z\in\Z^3}f(|\tilde x-\tilde z|)$. (iii) A path in the graph is no longer than the graph, and by the triangle inequality the distance of its endpoints is at most its length; a graph for $Y$ is a graph for $Y'$. (iv) is the definition of the norm, since $e^{\theta\ell}=e^{\varsigma\ell}e^{-(\varsigma-\theta)\ell}$.
\end{proof}

\paragraph{Constants.} Put $\tilde c(\theta)\defeq\sum_{z\in\Z^3}e^{-\theta|z|}$, $\mathsf K\defeq\sum_{z\in\Z^3}|z|e^{-\delta_P|z|/2}$, $c_d\defeq\tilde c(\delta_P/2)$, and
\begin{equation}\label{rot:eq:R1const}
\mu_1\defeq\min\Big\{\frac{\delta_P}2,\frac{\gamma_0}{2c_P\mathsf K}\Big\},\quad g_0\defeq\frac2{\gamma_0},\quad\mathsf M_1\defeq\frac{2c_P}{\gamma_0}\tilde c(\mu_1/2)^2,\quad c_{\mathrm D}\defeq\frac{4c_P}{\gamma_0^2}\tilde c(\mu_1/2)^2,\quad\gamma_K\defeq\min\Big\{\frac12,\frac1{2c_K\mathsf M_1}\Big\},\quad\bar\mu\defeq\min\{\mu_1,\varsigma\}.
\end{equation}
By \Cref{rot:lem:R1lat}(ii) every sum over $\mathbb{L}$ of a function of the distance is bounded by the corresponding sum over $\Z^3$, so these constants serve on every $\mathbb{L}$. They depend on $(\gamma_0,c_P,\delta_P,c_K,\varsigma)$ only, and no lower bound on $c_P$ is used anywhere ($c_P\to0$ gives $\mu_1=\delta_P/2$, $\mathsf M_1\to0$, $\gamma_K=\frac12$).

\begin{lemma}[Combes--Thomas bound]\label{rot:lem:R1ct}
For $X\subset\Lambda$ and $r\ge0$, $\|(P_{XX}+r)^{-1}(x,y)\|\le2(\gamma_0+r)^{-1}e^{-\mu_1|x-y|}$. In particular $\|G_X(x,y)\|\le g_0e^{-\mu_1|x-y|}$ for every $X$, and $\|C(x,y)\|\le g_0e^{-\mu_1|x-y|}$.
\end{lemma}
\begin{proof}
Fix $y_0\in X$ and $f\defeq|\cdot-y_0|$, and conjugate: $P^f\defeq e^{\mu_1f}P_{XX}e^{-\mu_1f}$. Then $(P^f-P_{XX})(x,y)=P(x,y)\big(e^{\mu_1(f(x)-f(y))}-1\big)$ vanishes on the diagonal and, by (R2), \Cref{rot:lem:R1lat}(i) and $e^t-1\le te^t$ for $t\ge0$, has block norm at most $c_Pe^{-\delta_P|x-y|}\mu_1|x-y|e^{\mu_1|x-y|}\le c_P\mu_1|x-y|e^{-\delta_P|x-y|/2}$, since $\mu_1\le\delta_P/2$; by Schur's test and \Cref{rot:lem:R1lat}(ii) its row and column sums are at most $c_P\mu_1\mathsf K\le\gamma_0/2$, so $\|P^f-P_{XX}\|\le\gamma_0/2$. Hence $\Re\langle u,(P^f+r)u\rangle\ge(\gamma_0/2+r)\|u\|^2$ for every $u$ ($P_{XX}\ge\gamma_0$ as a principal block), so $\|(P^f+r)^{-1}\|\le(\gamma_0/2+r)^{-1}\le2(\gamma_0+r)^{-1}$, and $(P_{XX}+r)^{-1}(x,y_0)=e^{-\mu_1f(x)}(P^f+r)^{-1}(x,y_0)e^{\mu_1f(y_0)}$ has norm at most $2(\gamma_0+r)^{-1}e^{-\mu_1|x-y_0|}$.
\end{proof}

\begin{lemma}[conditioning]\label{rot:lem:R1sch}
Let $X\subset\Lambda$ and $Y\defeq\Lambda\setminus X$. Then $0<G_X\le C_{XX}\le\gamma_0^{-1}$, so every component variance of $N(\cdot,G_X)$ and of $\mu_P$ is at most $\gamma_0^{-1}$. The conditional mean $h\defeq-G_XP_{XY}z_Y$ obeys
\[
|h(x)|\le\mathsf M_1e^{-(\mu_1/2)\dist(x,Y)}\|z_Y\|_\infty\qquad(x\in X),
\]
in particular $\|h\|_\infty\le\beta/2$ whenever $\|z_Y\|_\infty\le c_K\gamma_K\beta$. For $x,y\in X$,
\[
\|(C-G_X)(x,y)\|\le c_{\mathrm D}\exp\Big(-\frac{\mu_1}8\big(\dist(x,Y)+\dist(y,Y)\big)-\frac{\mu_1}4|x-y|\Big).
\]
\end{lemma}
\begin{proof}
$C_{XX}=(P_{XX}-P_{XY}P_{YY}^{-1}P_{YX})^{-1}\ge P_{XX}^{-1}=G_X$, because the subtracted operator is non-negative and inversion reverses the order of positive operators; and $C_{XX}\le\|C\|\le\gamma_0^{-1}$. For the mean, \Cref{rot:lem:R1ct} and (R2) (the pairs $(u,v)\in X\times Y$ are off-diagonal) give $|h(x)|\le\sum_{u\in X,v\in Y}g_0e^{-\mu_1|x-u|}c_Pe^{-\delta_P|u-v|}|z_v|$; as $\delta_P\ge2\mu_1$ and $|x-u|+|u-v|\ge|x-v|\ge\dist(x,Y)$, the exponent is at most $-\frac{\mu_1}2\dist(x,Y)-\frac{\mu_1}2(|x-u|+|u-v|)$, and the double sum of the last factor is at most $\tilde c(\mu_1/2)^2$ by \Cref{rot:lem:R1lat}(ii). The particular case is $\mathsf M_1c_K\gamma_K\le\frac12$. For the covariance difference, row $X$ of $PC=\mathbf 1$ reads $P_{XX}C_{XX}+P_{XY}C_{YX}=\mathbf 1$, so $C_{XX}-G_X=-G_XP_{XY}C_{YX}$, with kernel $-\sum_{u\in X,v\in Y}G_X(x,u)P(u,v)C(v,y)$, bounded by $g_0^2c_P\sum_{u,v}e^{-\mu_1L}$ with $L\defeq|x-u|+|u-v|+|v-y|$ (using $\delta_P\ge\mu_1$). Since $L\ge|x-y|$, $L\ge|x-v|\ge\dist(x,Y)$ and $L\ge|v-y|\ge\dist(y,Y)$, write $e^{-\mu_1L}=e^{-\mu_1L/4}e^{-\mu_1L/4}e^{-\mu_1L/2}\le e^{-\frac{\mu_1}8(\dist(x,Y)+\dist(y,Y))}e^{-\frac{\mu_1}4|x-y|}e^{-\frac{\mu_1}2(|x-u|+|u-v|)}$ and sum the last factor over $u,v$.
\end{proof}

\begin{lemma}[decoupling across corridors]\label{rot:lem:R1dec}
Let $\Lambda=\Gamma\sqcup D$ with $D=\bigsqcup_iD_i$ and $\dist(D_i,D_j)\ge w_s$ for $i\ne j$. Let $\mu_P(\cdot\,|\,z_\Gamma)$ be the conditional law of $z_D$ given $z_\Gamma$ and $\nu_{\mathrm{loc}}(\cdot\,|\,z_\Gamma)\defeq\bigotimes_iN(-G_{D_i}P_{D_i\Gamma}z_\Gamma,\,G_{D_i})$ its decoupled version; put $\omega\defeq c_Pc_de^{-\delta_Pw_s/2}$ and suppose $\omega\le\gamma_0/2$. On $\{\|z_D\|_\infty\le\beta,\ \|z_\Gamma\|_\infty\le\beta\}$,
\[
\Big|\log\frac{d\mu_P(\cdot\,|\,z_\Gamma)}{d\nu_{\mathrm{loc}}(\cdot\,|\,z_\Gamma)}\Big|\le\omega_{\mathrm{tot}}\defeq\omega|D|\Big[\frac{\beta^2}2+\frac n{\gamma_0}+\frac{(c_P\tilde c(\delta_P)\beta)^2}{2\gamma_0^2}\Big],
\]
hence $e^{-\omega_{\mathrm{tot}}}\int F\,d\nu_{\mathrm{loc}}\le\int F\,d\mu_P(\cdot\,|\,z_\Gamma)\le e^{\omega_{\mathrm{tot}}}\int F\,d\nu_{\mathrm{loc}}$ for every $F\ge0$ supported in that set. Each factor of $\nu_{\mathrm{loc}}$ is a law of the type of \Cref{rot:lem:R1sch}, with $X=D_i$, $Y=\Lambda\setminus D_i$ and $z_{D_j}=0$ for $j\ne i$.
\end{lemma}
\begin{proof}
Put $u\defeq P_{D\Gamma}z_\Gamma$, $B\defeq\bigoplus_iP_{D_iD_i}$ and $O\defeq P_{DD}-B$ (the blocks $P_{D_iD_j}$, $i\ne j$). The two laws are $N(-P_{DD}^{-1}u,P_{DD}^{-1})$ and $N(-B^{-1}u,B^{-1})$, since $(B^{-1}u)_{D_i}=G_{D_i}P_{D_i\Gamma}z_\Gamma$. The logarithm of the density of $N(-A^{-1}u,A^{-1})$ is $\frac12\log\det(A/2\pi)-\frac12\langle z,Az\rangle-\langle z,u\rangle-\frac12\langle u,A^{-1}u\rangle$; the linear terms agree for $A=P_{DD}$ and $A=B$, so the log-ratio is $\frac12\log\det(P_{DD}B^{-1})-\frac12\langle z_D,Oz_D\rangle-\frac12\langle u,(P_{DD}^{-1}-B^{-1})u\rangle$. Every entry of $O$ is an off-diagonal entry of $P$ between points at distance at least $w_s$, so by (R2) and \Cref{rot:lem:R1lat}(ii) its row and column sums are at most $c_Pe^{-\delta_Pw_s/2}\tilde c(\delta_P/2)=\omega$, and $\|O\|\le\omega$. First term: $|\frac12\langle z_D,Oz_D\rangle|\le\frac12\beta^2\sum_{x\in D}\sum_y\|O(x,y)\|\le\frac12\beta^2|D|\omega$. Second term: $\det(P_{DD}B^{-1})=\det(\mathbf 1+B^{-1/2}OB^{-1/2})$; the symmetric operator $B^{-1/2}OB^{-1/2}$ has $n|D|$ eigenvalues, each of modulus at most $\omega/\gamma_0\le\frac12$ ($B\ge\gamma_0$ as a direct sum of principal blocks), and $|\log(1+t)|\le2|t|$ for $|t|\le\frac12$; so the term is at most $n|D|\omega/\gamma_0$. Third term: $P_{DD}^{-1}-B^{-1}=-P_{DD}^{-1}OB^{-1}$ has norm at most $\omega/\gamma_0^2$, and $|u_x|\le\sum_{v\in\Gamma}c_Pe^{-\delta_P|x-v|}\beta\le c_P\tilde c(\delta_P)\beta$, so $\|u\|^2\le|D|(c_P\tilde c(\delta_P)\beta)^2$. The three bounds add up to $\omega_{\mathrm{tot}}$.
\end{proof}

% r03_r1_cumul: the cumulant bound for the step Gaussian, part II: tails, the normalisation of the cutoff, cumulants of bounded variables, the cutoff in the exponent and in cumulants, the connected Wick expansion and the comparison of two Gaussian laws.

\begin{lemma}[tails]\label{rot:lem:R1tail}
Let $\nu=N(h,G)$ on $(\R^n)^X$ with $\|h\|_\infty\le\beta/2$ and $G\le\gamma_0^{-1}$. Then $\nu(\exists x\in X:\ |z_x|>\beta)\le\epsilon(X,\beta)\defeq2n|X|e^{-\gamma_0\beta^2/(8n)}$; since $\bar B(\beta)\subset\beta K$, also $\nu(1-\chi^X_\beta)\le\epsilon(X,\beta)$.
\end{lemma}
\begin{proof}
$|z_x|>\beta$ forces $|\zeta_x|>\beta/2$ for $\zeta\defeq z-h$, hence $|\zeta^a_x|>\beta/(2\sqrt n)$ for some component $a$; $\zeta^a_x\sim N(0,\sigma^2)$ with $\sigma^2\le\gamma_0^{-1}$, $\Pr(|N(0,\sigma^2)|>t)\le2e^{-t^2/(2\sigma^2)}$, and a union bound over the $n|X|$ components gives the claim.
\end{proof}

\begin{lemma}[the normalisation of the cutoff]\label{rot:lem:R1gci}
Let $\mu$ be a centred Gaussian on $(\R^n)^\Lambda$ all of whose component variances are at most $\gamma_0^{-1}$, and let $2e^{-\gamma_0\beta^2/(2n)}\le\frac12$. Then $\mu(\chi_\beta)\ge\exp\big(-4n|\Lambda|e^{-\gamma_0\beta^2/(2n)}\big)$.
\end{lemma}
\begin{proof}
$\{z_x\in\beta K\}\supset\{|z_x|\le\beta\}\supset\bigcap_{a\le n}\{|z^a_x|\le\beta/\sqrt n\}$. By \v{S}id\'ak's inequality (Z.~\v{S}id\'ak, J.\ Amer.\ Statist.\ Assoc.\ 62 (1967), 626--633; C.~G.~Khatri, Ann.\ Math.\ Statist.\ 38 (1967), 1853--1867): for a centred Gaussian vector $(\xi_1,\dots,\xi_m)$ and $t_1,\dots,t_m>0$, $\Pr(\bigcap_i\{|\xi_i|\le t_i\})\ge\prod_i\Pr(|\xi_i|\le t_i)$; this is the case of one slab against a product of slabs of the Gaussian correlation inequality, and it is applied here to the $n|\Lambda|$ components of the real centred vector $z$. Each factor is at least $1-2e^{-\gamma_0\beta^2/(2n)}$, and $\log(1-u)\ge-2u$ for $u\le\frac12$.
\end{proof}

\begin{lemma}[cumulants of bounded variables]\label{rot:lem:R1cum}
For any probability measure $\nu$ and real variables with $|X_j|\le M_j$, $|\kappa_\nu(X_1,\dots,X_k)|\le k!\,2^{k-1}\prod_jM_j$. If $|X|\le M$, then $\big|\log\int e^{X}d\nu-\sum_{k=1}^{\bar n}\kappa_\nu(X;k)/k!\big|\le2^{\bar n}M^{\bar n+1}\le(2M)^{\bar n+1}$.
\end{lemma}
\begin{proof}
In the M\"obius formula $\kappa_\nu(X_1,\dots,X_k)=\sum_\pi(-1)^{|\pi|-1}(|\pi|-1)!\prod_{B\in\pi}\E_\nu\prod_{j\in B}X_j$, over the set partitions $\pi$ of $\{1,\dots,k\}$, each product of moments is at most $\prod_jM_j$, and $\sum_\pi(|\pi|-1)!$ is at most the number of ordered set partitions, which is at most $k!\,2^{k-1}$ (an ordered partition is a permutation cut into consecutive blocks). For the second claim, $f(s)\defeq\log\int e^{sX}d\nu$ has $f^{(m)}(s)=\kappa_{\nu_s}(X;m)$ with $\nu_s\propto e^{sX}\nu$ a probability measure under which still $|X|\le M$; Taylor's formula at $s=0$ with the Lagrange remainder $f^{(\bar n+1)}(s_*)/(\bar n+1)!$ and the first claim give the bound.
\end{proof}

\begin{lemma}[the cutoff in the exponent]\label{rot:lem:R1ext}
Let $\nu$ be a probability measure, $\chi$ an indicator with $\nu(1-\chi)\le\epsilon$, $W$ measurable with $|W|\le M$, and $e^{2M}\epsilon\le\frac12$. Then \textup{(i)} $\int\chi e^Wd\nu\ge e^{-2e^{2M}\epsilon}\int e^Wd\nu$, and \textup{(ii)} $\int\chi e^Wd\nu\le\int e^{\chi W}d\nu\le e^{2e^M\epsilon}\int\chi e^Wd\nu$.
\end{lemma}
\begin{proof}
(i) $\int(1-\chi)e^W\le e^M\epsilon\le e^{2M}\epsilon\int e^W$, because $\int e^W\ge e^{-M}$; and $1-u\ge e^{-2u}$ for $u\le\frac12$. (ii) On $\{\chi=0\}$, $e^{\chi W}=1\ge0=\chi e^W$, so $\int e^{\chi W}=\int\chi e^W+\nu(1-\chi)$; and $\int\chi e^W\ge e^{-M}(1-\epsilon)\ge e^{-M}/2$, so $\nu(1-\chi)\le2\epsilon e^M\int\chi e^W$.
\end{proof}
The lemma is applied with $W=\Psi\chi$, for which $\chi e^W=\chi e^\Psi$.

\begin{lemma}[connected Wick expansion]\label{rot:lem:R1wick}
Let $\zeta$ be a centred Gaussian vector with covariance $G$ and let $F_j=\zeta^{\alpha_j}$, $j=1,\dots,k$, be monomials, each a multiset of legs. Then $\kappa(F_1,\dots,F_k)=\sum_\pi\prod_{\{u,v\}\in\pi}G(u,v)$, the sum over the pairings $\pi$ of all legs that connect the $k$ factors (the graph on $\{1,\dots,k\}$ with an edge for each pair joining legs of two different factors is connected; pairs inside a factor are allowed). For $N(h,G)$ the same holds after the shift $z=h+\zeta$, the constants dropping out of joint cumulants with $k\ge2$.
\end{lemma}
\begin{proof}
Isserlis' formula gives $\E\prod_{j\in B}F_j=\sum\prod G$ over all pairings of the legs of the factors in $B$. Grouping each pairing by the partition of $B$ into the connected components of its factor graph gives $\E\prod_{j\in B}F_j=\sum_\sigma\prod_{b\in\sigma}\mathsf C(b)$, with $\mathsf C(b)$ the sum over the connected pairings of the legs of $b$; the moment--cumulant relation $\E\prod_{j\in B}F_j=\sum_\sigma\prod_{b\in\sigma}\kappa(b)$ determines the $\kappa(b)$ from the moments recursively, so $\kappa=\mathsf C$.
\end{proof}

\begin{lemma}[anchored connected bound]\label{rot:lem:R1conn}
Let $\nu=N(h,G)$ on $(\R^n)^X$ with $\|G(x,y)\|\le g_0e^{-\theta|x-y|}$ and $\|h\|_\infty\le\beta$. Let $F_1,\dots,F_k$ ($k\ge2$) be polynomials of degree at most $D$ in the variables of $X$, sorted by supports as in (R3), with $\|F_j\|_{\theta'}\le a_j$ for some $\theta'\le\theta$, and let $S_1,S_2\subset\mathbb{L}$ be such that every monomial of $F_1$ has support meeting $S_1$ and every monomial of $F_2$ has support inside $S_2$. Then
\[
|\kappa_\nu(F_1,\dots,F_k)|\le|S_1\cap X|\,K_k(\theta')\,(1+\beta)^{kD}\prod_ja_j\,e^{-(\theta'/2)\dist(S_1,S_2)},\qquad K_k(\theta')\defeq(kD)^{kD/2}\max\{1,g_0\}^{kD/2}\,\tilde c(\theta'/2)^{k-1},
\]
and the same bound without the distance factor when no $S_2$ is given (the anchored bound). For $k=1$, $|\E_\nu F_1|\le|S_1\cap X|\,D^{D/2}\max\{1,g_0\}^{D/2}(1+\beta)^Da_1$. For centred $\nu$ the factor $(1+\beta)^{kD}$ is absent.
\end{lemma}
\begin{proof}
(1) Shift $z=h+\zeta$: a monomial $z^\alpha$ with support $Y$ becomes $\sum_{\alpha'\le\alpha}\binom{\alpha}{\alpha'}h^{\alpha-\alpha'}\zeta^{\alpha'}$, whose coefficients have modulus sum at most $\prod(1+|h|)^{\alpha}\le(1+\beta)^D$; the pieces are kept grouped under the original support $Y$ (their sites lie in $Y$), so the shifted $F'_j$ has $\|F'_j\|_{\theta'}\le(1+\beta)^Da_j$, and the constant pieces drop out for $k\ge2$. For $h=0$ there is no shift and no factor. (2) By \Cref{rot:lem:R1wick}, $|\kappa_\nu(F_1,\dots,F_k)|\le\sum_{Y_1,\dots,Y_k}\sum_{\text{monomials}}\prod|\text{coefficient}|\sum_{\text{connected }\pi}\prod\|G\|$; the factors have at most $kD$ legs in all, so for fixed monomials there are at most $(kD)^{kD/2}$ pairings. (3) For a connected $\pi$ fix a spanning tree of its factor graph, rooted at factor $1$, with $k-1$ tree pairs. Sum the root over the sites $a\in S_1\cap X$ and the supports $Y_1\ni a$: at most $|S_1\cap X|\,\|F'_1\|_{\theta'}$ times the rest. Each child factor $j$ is entered through a tree pair $(y,y')$ with $y$ in the parent's support and $y'\in Y_j$, which gives $\sum_{y'}g_0e^{-\theta|y-y'|}\sum_{Y_j\ni y'}|F'_{j,Y_j}|_1$; the pairs not in the tree are bounded by $g_0$ each. (4) Write $|F'_{j,Y}|_1=\big(|F'_{j,Y}|_1e^{\theta'\ell(Y)}\big)e^{-\theta'\ell(Y)}$ and $e^{-\theta|y-y'|}\le e^{-(\theta'/2)|y-y'|}e^{-(\theta'/2)|y-y'|}$. The tree path from factor $1$ to factor $2$ lifts to a path in $\mathbb{L}$ from $a\in S_1$ to a point of $Y_2\subset S_2$, made of tree pairs and of paths inside shortest connecting graphs of the traversed supports $Y_j$, each of length at most $\ell(Y_j)$ (\Cref{rot:lem:R1lat}(iii)); hence $e^{-(\theta'/2)(\sum_j\ell(Y_j)+\sum_{\text{tree}}|y-y'|)}\le e^{-(\theta'/2)\dist(S_1,S_2)}$. The remaining halves give $\tilde c(\theta'/2)$ per tree pair (\Cref{rot:lem:R1lat}(ii)) and $\|F'_j\|_{\theta'}$ per factor. For $k=1$ the same count applies without a tree.
\end{proof}

\begin{lemma}[cumulants with and without the cutoff]\label{rot:lem:R1chi}
Let $\nu=N(h,G)$ on $(\R^n)^X$ with $\|h\|_\infty\le\beta/2$, $G\le\gamma_0^{-1}$ and $\beta\ge2$; put $\chi\defeq\chi^X_\beta$ and $\epsilon\defeq\nu(1-\chi)$. Let $F_1,\dots,F_k$ ($k\le\bar n$) be polynomials in the variables of $X$ of degree at most $D$ with $\|F_j\|_0\le a_j$, and assume the moment condition
\begin{equation}\label{rot:eq:R1MC}
2^q\,|X|\,(2n)^{1/2}\big(n^q(2q-1)!!\big)^{1/2}\gamma_0^{-q/2}\,e^{-\gamma_0\beta^2/(4n)}\le(2\beta)^q\qquad(1\le q\le2\bar nD).
\end{equation}
Then $|\kappa_\nu(\chi F_1,\dots,\chi F_k)-\kappa_\nu(F_1,\dots,F_k)|\le k\cdot k!\cdot4^k\prod_j(a_j|X|)\,(2\beta)^{kD}\,\epsilon^{1/2}$.
\end{lemma}
\begin{proof}
Let $Z\defeq\max\{1,\max_{x\in X}|z_x|\}$ and $\zeta\defeq z-h$. If $\max_x|\zeta_x|\le\beta$ then $Z\le2\beta$, otherwise $Z\le2\max_x|\zeta_x|$; so $\E Z^q\le(2\beta)^q+2^q\sum_x\E\big[|\zeta_x|^q\mathbf 1\{|\zeta_x|>\beta\}\big]\le2(2\beta)^q$ by the Cauchy--Schwarz inequality, $\E|\zeta_x|^{2q}\le n^q(2q-1)!!\gamma_0^{-q}$ (H\"older's inequality over the $n$ components and the Gaussian moments), $\Pr(|\zeta_x|>\beta)\le2ne^{-\gamma_0\beta^2/(2n)}$ and \eqref{rot:eq:R1MC}. Since $|F_j|\le a_j|X|Z^D$, every product of moments obeys $|\E\prod_{j\in B}F_j|\le2\prod_{j\in B}(a_j|X|)(2\beta)^{|B|D}$, and $|\E[(1-\chi)\prod_{j\in B}F_j]|\le\epsilon^{1/2}(\E\prod_{j\in B}F_j^2)^{1/2}\le\sqrt2\,\epsilon^{1/2}\prod_{j\in B}(a_j|X|)(2\beta)^{|B|D}$. In the M\"obius formula ($\chi^{|B|}=\chi$ on each block) replace the blocks one at a time: each partition contributes at most $k\cdot\sqrt2\cdot2^{k-1}\epsilon^{1/2}\prod_j(a_j|X|)(2\beta)^{kD}$, and $\sum_\pi(|\pi|-1)!\le k!\,2^{k-1}$.
\end{proof}
The moment condition \eqref{rot:eq:R1MC} is monotone: once it holds for a pair $(|X|,\beta)$ it holds for smaller $|X|$ and larger $\beta$.

\begin{lemma}[comparing two Gaussian laws]\label{rot:lem:R1repl}
Let $\nu,\tilde\nu$ be Gaussian laws of $(z_x)_{x\in S}$, $S$ finite, with means $h,\tilde h$ of sup norm at most $\beta/2$ and covariances $G,\tilde G$ all of whose blocks have norm at most $\gamma_0^{-1}$, and let $\sup_x|h(x)-\tilde h(x)|\le\eta_1$, $\sup_{x,y}\|(G-\tilde G)(x,y)\|\le\eta_2$. For polynomials $F_1,\dots,F_k$ in the variables of $S$ of degree at most $D$ with $\|F_j\|_0\le a_j$,
\[
|\kappa_\nu(F_1,\dots,F_k)-\kappa_{\tilde\nu}(F_1,\dots,F_k)|\le k\cdot k!\cdot4^k\,(kD)^{kD+1}\,(1+\beta+\gamma_0^{-1})^{kD}\prod_j(a_j|S|)\,\max\{\eta_1,\eta_2\}.
\]
\end{lemma}
\begin{proof}
A moment of $L\le kD$ coordinate variables is, by Isserlis' formula for a Gaussian with mean, a sum over at most $L^L$ partial pairings of products of at most $L$ factors (covariance entries or means), each of modulus at most $1+\beta+\gamma_0^{-1}$ under both laws; replacing the factors one at a time, two corresponding products differ by at most $L\max\{\eta_1,\eta_2\}(1+\beta+\gamma_0^{-1})^L$. The coefficient sums contribute $\prod_j(a_j|S|)$, and the M\"obius formula with the telescoping of \Cref{rot:lem:R1chi} gives the bound.
\end{proof}

% r03_r1_geo: the cumulant bound for the step Gaussian, part III: the four shifted pavements on Z^3, on the cubic tori and on the skew tori, and the layers of a cell.

\begin{lemma}[pavements]\label{rot:lem:R1geo}
Let $b\ge b_{\mathrm{geo}}\defeq576$, $q_b\defeq\lceil b^2\rceil$, $w_b\defeq\lceil b^{3/2}\rceil$ and $s_b\defeq\lfloor q_b/4\rfloor$, so that $s_b\ge5w_b+2$, $q_b\ge6w_b$ and $3q_b\le4b^2$; put $v_b\defeq(4b^2)^3=64\,b^6$. On each $\mathbb{L}$ of the setting there are four partitions $\mathcal P^i$, $i=0,1,2,3$, of $\mathbb{L}$ into \emph{cells} $\Box$, with \emph{levels} $\lambda_i:\mathbb{L}\to[0,\infty]$, such that
\begin{enumerate}[label=\textup{(G\arabic*)},leftmargin=*]
\item $\lambda_i$ is $1$-Lipschitz;
\item if $x,y$ lie in different cells of $\mathcal P^i$, then $|x-y|\ge\lambda_i(x)+\lambda_i(y)$;
\item $\bigcap_{i=0}^3\{x:\ \lambda_i(x)<2.5w_b+1\}=\emptyset$;
\item every cell has at most $v_b$ sites; in particular, if some $\mathcal P^i$ consists of a single cell, then $|\mathbb{L}|\le v_b$.
\end{enumerate}
\end{lemma}
\begin{proof}
\emph{Numbers.} For $b\ge576$: $s_b\ge b^2/4-1\ge5b^{3/2}+7\ge5w_b+2$ (the middle inequality is $\sqrt b\ge20+32b^{-3/2}$, true for $\sqrt b\ge24$), $q_b\ge b^2\ge6b^{3/2}+6\ge6w_b$, and $3q_b\le3b^2+3\le4b^2$.

\emph{Coordinates.} Each $\mathbb{L}$ has three coordinate cycles $\mathbb I_a$, $a=1,2,3$, with a well-defined map $x\mapsto x_a\in\mathbb I_a$ and $N_a\defeq|\mathbb I_a|\in\N\cup\{\infty\}$: for $\Z^3$, $\mathbb I_a=\Z$; for $(\Z/N\Z)^3$, $\mathbb I_a=\Z/N\Z$; for $\Z^3/mG_\pm$, $\mathbb I_a\defeq\Z/N_a\Z$ with $N_a\Z\defeq\pi_a(mG_\pm)$ the projection of the period lattice on the coordinate $a$, that is $N_1=N_2=10m$ and $N_3=50m$ (the entries $40,\mp30,0$ and $\pm30,40,0$ have greatest common divisor $10$); the map is well defined because the period lattice $G\defeq mG_\pm$ is contained in $G'\defeq\bigoplus_aN_a\Z$. In every case $\dist_{\mathbb I_a}(x_a,y_a)\le|\tilde x_a-\tilde y_a|\le|\tilde x-\tilde y|=|x-y|$ for lifts realising the quotient distance. For the skew tori we record: $\det G=(40\cdot40+30\cdot30)\cdot50\cdot m^3=125000\,m^3=(50m)^3$, $[G':G]=125000/5000=25$, and $G\cap\{g_3=0\}=m\langle(40,\mp30,0),(\pm30,40,0)\rangle$ has index $2500m^2$ in $\Z^2\times\{0\}$.

\emph{Cuts.} A coordinate $a$ is \emph{cut} if $N_a\ge q_b$ (always for $\Z^3$). For a cut coordinate choose a set $B^0_a\subset\mathbb I_a+\frac12$ of cuts all of whose (cyclic) gaps lie in $[q_b,2q_b)$: for $\mathbb I_a=\Z$ take $q_b\Z+\frac12$; on a cycle take $p\defeq\lfloor N_a/q_b\rfloor\ge1$ cuts with gaps $\lfloor N_a/p\rfloor$ or $\lceil N_a/p\rceil$, which lie in $[q_b,2q_b-1]$. Put $B^i_a\defeq B^0_a+is_b$, $i=0,\dots,3$. Cuts of $B^i_a$ and of $B^{i'}_a$, $i\ne i'$, are at cyclic distance at least $s_b$: two shifts of the same cut differ by $|i-i'|s_b\in[s_b,3s_b]$, and $N_a-3s_b\ge q_b-3s_b\ge s_b$; shifts of different cuts differ by at least $q_b-3s_b\ge s_b$. For $x\in\mathbb{L}$ put $\lambda_i(x)\defeq\min_{a\ \mathrm{cut}}\dist_{\mathbb I_a}(x_a,B^i_a)$, and $\lambda_i\defeq+\infty$ if no coordinate is cut. The \emph{cut set} $C^i$ is the set of points $\xi$ of the continuum torus ($\R^3$, $\R^3/N\Z^3$ or $\R^3/G$) with $\xi_a\in B^i_a$ for some cut coordinate $a$, and the cells of $\mathcal P^i$ are the intersections of $\mathbb{L}$ with the connected components of the complement of $C^i$; no site lies on a cut, so the cells partition $\mathbb{L}$.

\emph{(G1)} is the Lipschitz property of $\dist_{\mathbb I_a}(\cdot,B^i_a)$ together with $\dist_{\mathbb I_a}(x_a,y_a)\le|x-y|$. \emph{(G3)}: if $\lambda_i(x)<2.5w_b+1$ for all four $i$, each minimum is attained at a cut coordinate, and since there are at most three, two indices $i\ne i'$ use the same coordinate $a$; then a cut of $B^i_a$ and a cut of $B^{i'}_a$ are at cyclic distance less than $2(2.5w_b+1)=5w_b+2\le s_b$, a contradiction. If no coordinate is cut the four sets are empty. \emph{(G2)}: lift $x,y$ to $\tilde x,\tilde y\in\R^3$ with $|\tilde x-\tilde y|=|x-y|$. The lift $\tilde C^i$ of the cut set is the union of the planes $\{\xi_a=c\}$ over the cut coordinates $a$ and the lifts $c$ of the cuts of $B^i_a$; it is $G'$-periodic, hence $G$-periodic, the components of $\R^3\setminus\tilde C^i$ are the products of open gap intervals in the cut coordinates and of $\R$ in the others, and the projection maps them onto the components of the complement of $C^i$. If the segment $[\tilde x,\tilde y]$ met none of the planes it would lie in one component, and $x,y$ would lie in one cell; so it crosses a plane $\{\xi_a=c\}$ with $c$ between $\tilde x_a$ and $\tilde y_a$, and $|x-y|\ge|\tilde x_a-\tilde y_a|=|\tilde x_a-c|+|c-\tilde y_a|\ge\dist_{\mathbb I_a}(x_a,B^i_a)+\dist_{\mathbb I_a}(y_a,B^i_a)\ge\lambda_i(x)+\lambda_i(y)$. \emph{(G4)}: a component $Q$ of $\R^3\setminus\tilde C^i$ projects onto its cell injectively modulo $G_Q\defeq\{g\in G:\ Q+g=Q\}=\{g\in G:\ g_a=0\text{ for every cut }a\}$: if $Q$ and $Q+g$ meet they coincide, $\tilde C^i$ being $G$-invariant. On $\Z^3$, and whenever all three coordinates are cut, $G_Q=0$ and the cell is a box with sides at most $2q_b-1$, hence has at most $(2q_b)^3$ sites. On a cubic torus the only other case is that no coordinate is cut, when $|\mathbb{L}|=N^3<q_b^3$. On $\Z^3/mG_\pm$, where $N_3=5N_1$, two further cases occur: only the third coordinate is cut ($10m<q_b\le50m$), the cell is the product of $\Z^2/(G\cap\{g_3=0\})$ with an interval of fewer than $2q_b$ integers and has fewer than $2500m^2\cdot2q_b=25\,(10m)^2\cdot2q_b<50\,q_b^3$ sites; or no coordinate is cut ($50m<q_b$) and $|\mathbb{L}|=\det G=(50m)^3<q_b^3$. With $q_b\le b^2+1\le\frac{257}{256}b^2$ the three counts are below $8.1\,b^6$, $50.6\,b^6$ and $1.02\,b^6$, hence below $v_b=64\,b^6$.
\end{proof}

\paragraph{Layers.} For a cell $\Box$ of $\mathcal P^i$ write $\lambda\defeq\lambda_i$ and define, all intersected with $\Box\cap\Lambda$:
\[
\Gamma_1(\Box)\defeq\{\lambda<w_b\},\quad\Gamma_2(\Box)\defeq\{w_b\le\lambda<2w_b\},\quad\Gamma_3(\Box)\defeq\{1.5w_b\le\lambda<2w_b\},\quad\Box'\defeq\{\lambda\ge2w_b\},\quad\Gamma_4(\Box)\defeq\{2w_b\le\lambda<2.5w_b\},
\]
$\mathrm{Deep}(\Box)\defeq\{\lambda\ge2.5w_b\}$ and $D_\Box\defeq\Gamma_2(\Box)\cup\Box'=\{\lambda\ge w_b\}$. By (G1) and (G2), two points of one cell at levels $\lambda(x)\ge\lambda(y)$ are at distance at least $\lambda(x)-\lambda(y)$, and two points of different cells at distance at least $\lambda(x)+\lambda(y)$. Hence: the sets $D_\Box$ of distinct cells are at distance at least $2w_b$; $\dist(D_\Box,\Gamma_1(\tilde\Box))\ge w_b$ and $\dist(\mathrm{Deep}(\Box),\tilde\Box\cap\{\lambda\ge1.5w_b\})\ge4w_b$ for $\tilde\Box\ne\Box$; and $\dist(\mathrm{Deep}(\Box),\Lambda\setminus\Box')\ge\frac12w_b$, $\dist(\Box\cap\{\lambda\ge1.5w_b\},\Lambda\setminus D_\Box)\ge\frac12w_b$, $\dist(\mathrm{Deep}(\Box),\Lambda\setminus(\Box\cap\{\lambda\ge1.5w_b\}))\ge w_b$, $\dist(\Gamma_4(\Box),\Box\cap\{\lambda<1.5w_b\})\ge\frac12w_b$, $\dist(\Box',\Gamma_2(\Box)\setminus\Gamma_3(\Box))\ge\frac12w_b$, $\dist(\Gamma_1(\Box),\Box')\ge w_b$.

% r03_r1_stage: the cumulant bound for the step Gaussian, part IV: one pavement stage (the sorting of the monomials, the stage lemma and the bookkeeping lemma).

\paragraph{One stage.} Fix $i$, the pavement $\mathcal P^i$ with its layers, a set $J'\subset J$ (write $V'\defeq V_{J'}$) and a threshold $\beta\in[b_{\min},b_{\max}]$, where $b_{\min}\defeq\gamma_K^4b$ and $b_{\max}\defeq\gamma_K^{-4}b$. Put $\Phi(\beta;R)\defeq\int\chi_\beta e^{V_R}d\mu_P$, $\chi^\Box_\beta\defeq\chi^{D_\Box}_\beta$, $\mathcal B'\defeq\{\Box\in\mathcal P^i:\ \Box\cap J'\ne\emptyset\}$ (the \emph{$J$-cells}), $N'\defeq|\mathcal B'|\le|J'|$, $D\defeq\bigcup_{\Box\in\mathcal B'}D_\Box$, $\Gamma\defeq\Lambda\setminus D$ and $J''\defeq J'\cap\{\lambda_i<2.5w_b\}$. The monomials of $V'$, all of which have support in $J'$, are sorted as follows.
\begin{itemize}[leftmargin=*]
\item $V_\Gamma$: support in $\Gamma$. For $\Box\in\mathcal B'$, $\Psi_\Box$: support in $D_\Box$, or support in $\Gamma_1(\Box)\cup\Gamma_2(\Box)$ meeting both. $R_0\defeq V'-V_\Gamma-\sum_\Box\Psi_\Box$.
\item $\Psi_1$: support in $\Box\cap\{\lambda\ge1.5w_b\}$ and meeting $\Box'$. $\Psi_2$: support in $\Gamma_1(\Box)\cup\Gamma_2(\Box)$ and meeting $\Gamma_2(\Box)$. $\Psi_3\defeq\Psi_\Box-\Psi_1-\Psi_2$: the monomials of $\Psi_\Box$ meeting $\Box'$ and $\Gamma_2(\Box)\setminus\Gamma_3(\Box)$.
\item $\Psi_1'$: the monomials of $\Psi_1$ with support in $\Gamma_3(\Box)\cup\Gamma_4(\Box)$ (they meet $\Gamma_4(\Box)$). $\Psi_1''\defeq\Psi_1-\Psi_1'$: the monomials of $\Psi_1$ meeting $\mathrm{Deep}(\Box)$. $W_\Box\defeq\Psi_1'+\Psi_2$ and $R_1\defeq V_{J''}-V_\Gamma-\sum_\Box W_\Box$.
\end{itemize}
$\Psi_2$ depends on $z_{\Gamma_1(\Box)}$ as well as on $z_{D_\Box}$. Under the conditional laws below it is read as a polynomial in $z_{D_\Box}$ whose coefficients depend on $z_{\Gamma_1(\Box)}$: every such monomial keeps at least one variable in $D_\Box$, its support shrinks to $Y\cap D_\Box$ (so its length does not increase, \Cref{rot:lem:R1lat}(iii)), and on $\{\|z_\Gamma\|_\infty\le c_K\beta\}$ its coefficient sums grow by at most the factor $(c_K\beta)^D$. The other pieces involve $z_{D_\Box}$ only.

\emph{Length facts.} $V_\Gamma$, $\Psi_1'$ and $\Psi_2$ consist of monomials of $V_{J''}$ (their supports lie in $J'\cap\{\lambda<w_b\}$, in $\Gamma_3\cup\Gamma_4$ and in $\{\lambda<2w_b\}$ respectively; a site of $J'$ in a cell $\Box$ with $\lambda\ge w_b$ lies in $D_\Box$, since then $\Box\in\mathcal B'$), and the three classes are disjoint, so $R_1$ is a sub-sum of the monomials of $V_{J''}$. Every monomial of $R_0$ meets some $D_\Box$, $\Box\in\mathcal B'$, and has $\ell\ge w_b$: if it lies in $\Box$ it meets $\Gamma_1(\Box)$ (its support is not inside $D_\Box$) and $\Box'$ (its support is not inside $\Gamma_1\cup\Gamma_2$); otherwise it meets a second cell, in $D_{\tilde\Box}$ or in $\Gamma_1(\tilde\Box)$. Every monomial of $\Psi_3$ has $\ell\ge\frac12w_b$. Every monomial of $R_1$ meets $\Gamma_2(\Box)\cup\Gamma_4(\Box)$ for some $\Box\in\mathcal B'$ and has $\ell\ge\frac12w_b$: it meets some $D_\Box$ inside $\{\lambda<2.5w_b\}$; if it lies in $\Box$ and meets $\Gamma_4(\Box)$, it is not inside $\Gamma_3\cup\Gamma_4$ (it is not in $\Psi_1'$), so it meets $\{\lambda<1.5w_b\}$; if it lies in $\Box$ and misses $\Gamma_4(\Box)$, its support is inside $\Gamma_1\cup\Gamma_2$ and meets $\Gamma_2$, so it belongs to $\Psi_2$, which is excluded; if it meets a second cell, $\ell\ge w_b$. On $\{\|z\|_\infty\le c_K\beta\}$, with $\beta\le b_{\max}$ and $|D_\Box|\le v_b$, \Cref{rot:lem:R1lat}(iv) gives
\[
|R_0|,\ |R_1|\le N'v_b\mathsf A(c_Kb_{\max})^De^{-\varsigma w_b/2},\qquad|\Psi_\Box|,\ |W_\Box|\le M\defeq\mathsf Av_b(c_Kb_{\max})^D,\qquad|\Psi_3|\le Me^{-\varsigma w_b/2},
\]
and $M=\mathsf m_*\,\mathsf Ab^{D+6}$ with $\mathsf m_*\defeq64\,c_K^D\gamma_K^{-4D}$, so $M\le\mathsf m_*$ under $\mathsf Ab^{D+6}\le1$.

\paragraph{Thresholds.} Put $\epsilon_\beta\defeq2nv_be^{-\gamma_0\beta^2/(8n)}$, with $b_{\min}$, $b_{\max}$ and $\mathsf m_*$ as above. The threshold of \Cref{rot:thm:R1} is $b_*\defeq\max\{576,\,b_{\mathrm{dec}},\,b_{\mathrm{tail}},\,b_{\mathrm{mc}},\,b_{\mathrm{gci}}\}$, where $b_{\mathrm{dec}}$, $b_{\mathrm{tail}}$, $b_{\mathrm{mc}}$ and $b_{\mathrm{gci}}$ are the least numbers beyond which, respectively, $c_Pc_de^{-\delta_Pb^{3/2}}\le\gamma_0/2$; $2nv_be^{2\mathsf m_*}e^{-\gamma_0\gamma_K^8b^2/(8n)}\le\frac12$ and $\gamma_K^4b\ge2$; the moment condition \eqref{rot:eq:R1MC} holds with $|X|=v_b$ and $\beta=b_{\min}$; and $2e^{-\gamma_0\gamma_K^8b^2/(2n)}\le\frac12$. Each of these inequalities, once true, stays true for larger $b$, and all of them depend only on $(n,\bar n,D,\gamma_0,c_P,\delta_P,c_K)$.

\begin{lemma}[one stage]\label{rot:lem:R1stg}
Let $b\ge b_*$ and $\beta\in[b_{\min},b_{\max}]$, with $\beta/\gamma_K\le b_{\max}$ in \textup{(U)}. Then
\begin{align*}
\textup{(L)}\qquad&\Phi(\beta;J')\ge\Phi(\gamma_K\beta;J'')\,\exp\Big\{\sum_{k=1}^{\bar n}\frac1{k!}\big[\kappa_{\mu_P}(V_{J'};k)-\kappa_{\mu_P}(V_{J''};k)\big]-N'E_*\Big\},\\
\textup{(U)}\qquad&\Phi(\beta;J')\le\Phi(\beta/\gamma_K;J'')\,\exp\Big\{\sum_{k=1}^{\bar n}\frac1{k!}\big[\kappa_{\mu_P}(V_{J'};k)-\kappa_{\mu_P}(V_{J''};k)\big]+N'E_*\Big\},
\end{align*}
with the per-cell error $E_*=E_*(b,\mathsf A)\defeq e_R+e_{\mathrm{dec}}+e_{\mathrm{ext}}+e_{\mathrm{cum}}+e_3+e_{\mathrm{mix}}+e_{\mathrm{rep}}+e_{\mathrm{bk}}$ defined in the proof and in \Cref{rot:lem:R1bk}. All its terms except $e_{\mathrm{cum}}=2(2M)^{\bar n+1}$ are increasing in $\mathsf A$ and carry one of the factors $e^{-\varsigma w_b/2}$, $e^{-\delta_Pw_b}$, $e^{-\bar\mu w_b/4}$, $e^{-\mu_1w_b/8}$, $\epsilon(v_b,b_{\min})^{1/2}$.
\end{lemma}
\begin{proof}
Recall $\epsilon_\beta=\epsilon(v_b,\beta)=2nv_be^{-\gamma_0\beta^2/(8n)}$, which decreases in $\beta$.

\emph{Lower bound.} (L1) $V'=V_\Gamma+\sum_\Box\Psi_\Box+R_0$, so on $\supp\chi_\beta$, $e^{V'}\ge e^{-|R_0|_{\max}}e^{V_\Gamma+\sum_\Box\Psi_\Box}$. (L2) Shrink the cutoff on $\Gamma$: $\chi_\beta\ge\chi^\Gamma_{\gamma_K\beta}\chi^D_\beta$, as $\gamma_K\beta K\subset\beta K$. Disintegrate $\mu_P$ into its marginal on $z_\Gamma$ and the conditional law of $z_D$. The $z_D$-integrand $\prod_\Box\chi^\Box_\beta e^{\Psi_\Box}$ is non-negative and supported in $\{\|z_D\|_\infty\le c_K\beta\}$, and $\|z_\Gamma\|_\infty\le c_K\gamma_K\beta\le c_K\beta$ on $\supp\chi^\Gamma_{\gamma_K\beta}$; \Cref{rot:lem:R1dec} with the pieces $D_\Box$, $\Box\in\mathcal B'$ (separation $2w_b$, so $\omega=c_Pc_de^{-\delta_Pw_b}\le\gamma_0/2$ because $b\ge b_*$), gives
\[
\int\mu_P(dz_D\,|\,z_\Gamma)\prod_\Box\chi^\Box_\beta e^{\Psi_\Box}\ge e^{-\omega_{\mathrm{tot}}}\prod_{\Box\in\mathcal B'}\int\nu_\Box(dz_{D_\Box})\,\chi^\Box_\beta e^{\Psi_\Box},\qquad\nu_\Box\defeq N(h_\Box,G_{D_\Box}),\quad h_\Box\defeq-G_{D_\Box}P_{D_\Box\Gamma}z_\Gamma,
\]
with $\omega_{\mathrm{tot}}\le N'v_bc_Pc_de^{-\delta_Pw_b}\big[(c_K\beta)^2/2+n/\gamma_0+(c_P\tilde c(\delta_P)c_K\beta)^2/(2\gamma_0^2)\big]$; only the $J$-cells are decoupled, which is why the error is proportional to $N'$. (L3) By \Cref{rot:lem:R1sch}, $\|h_\Box\|_\infty\le\beta/2$ and $G_{D_\Box}\le\gamma_0^{-1}$; by \Cref{rot:lem:R1tail}, $\nu_\Box(1-\chi^\Box_\beta)\le\epsilon_\beta$; by \Cref{rot:lem:R1ext}(i) with $W\defeq\Psi_\Box\chi^\Box_\beta$, $|W|\le M$ (and $e^{2M}\epsilon_\beta\le\frac12$ because $b\ge b_*$), $\int\nu_\Box\chi^\Box_\beta e^{\Psi_\Box}\ge e^{-2e^{2M}\epsilon_\beta}\int\nu_\Box e^{W}$. (L4) By \Cref{rot:lem:R1cum}, $\log\int\nu_\Box e^{W}\ge\sum_{k\le\bar n}\kappa_{\nu_\Box}(W;k)/k!-(2M)^{\bar n+1}$. (L5) Expand $W=\Psi_1'\chi+\Psi_1''\chi+\Psi_2\chi+\Psi_3\chi$ ($\chi=\chi^\Box_\beta$) multilinearly: $\kappa_{\nu_\Box}(W;k)$ is the sum over the at most $4^k$ words of length $k$ in these four letters of their joint cumulants. Words containing $\Psi_3\chi$: by \Cref{rot:lem:R1cum} each is at most $k!\,2^{k-1}M^ke^{-\varsigma w_b/2}$, so they contribute at most $e_3\defeq\sum_{k\le\bar n}(8M)^ke^{-\varsigma w_b/2}$. (L6) Words with at least one $\Psi_1''\chi$ and one $\Psi_2\chi$ (so $k\ge2$): \Cref{rot:lem:R1chi} removes the cutoffs (with $X=D_\Box$, $|X|\le v_b$, $a_j\le\mathsf A(c_K\beta)^D$, and the moment condition holding because $b\ge b_*$), then \Cref{rot:lem:R1conn} with $S_1\defeq\mathrm{Deep}(\Box)$ (met by every monomial of $\Psi_1''$), $S_2\defeq\Box\cap\{\lambda<2w_b\}$ (containing the supports of $\Psi_2$ read in $z_{D_\Box}$), $\dist(S_1,S_2)\ge\frac12w_b$, $\theta=\mu_1$ (\Cref{rot:lem:R1ct}) and $\theta'=\bar\mu$ ($\|\cdot\|_{\bar\mu}\le\|\cdot\|_\varsigma$) bounds each by $v_bK_k(\bar\mu)(1+\beta)^{kD}(\mathsf A(c_K\beta)^D)^ke^{-\bar\mu w_b/4}$; in total these words contribute at most
\[
e_{\mathrm{mix}}\defeq\sum_{k=2}^{\bar n}\frac{4^k}{k!}\Big[k\,k!\,4^k\big(\mathsf A(c_Kb_{\max})^Dv_b\big)^k(2b_{\max})^{kD}\epsilon_{b_{\min}}^{1/2}+v_bK_k(\bar\mu)(1+b_{\max})^{kD}\big(\mathsf A(c_Kb_{\max})^D\big)^ke^{-\bar\mu w_b/4}\Big].
\]
(L7) Words with at least one $\Psi_1''\chi$ and neither $\Psi_2\chi$ nor $\Psi_3\chi$ (all factors $\Psi_1'\chi$ or $\Psi_1''\chi$, with supports in $S\defeq\Box\cap\{\lambda\ge1.5w_b\}$): \Cref{rot:lem:R1chi} removes the cutoffs, and \Cref{rot:lem:R1repl} compares the marginal of $\nu_\Box$ on $S$ (mean $h_\Box|_S$, covariance $(G_{D_\Box})_{SS}$) with that of $\mu_P$ (centred, covariance $C_{SS}$): since $\dist(S,\Lambda\setminus D_\Box)\ge\frac12w_b$, \Cref{rot:lem:R1sch} gives $\eta_1\le\mathsf M_1c_K\gamma_K\beta e^{-\mu_1w_b/4}$ and $\eta_2\le c_{\mathrm D}e^{-\mu_1w_b/8}$. This replaces these words by the $z_\Gamma$-independent quantity
\[
c_\Box\defeq\sum_{k=1}^{\bar n}\frac1{k!}\big[\kappa_{\mu_P}(\Psi_1'+\Psi_1'';k)-\kappa_{\mu_P}(\Psi_1';k)\big]
\]
with an error at most
\[
e_{\mathrm{rep}}\defeq\sum_{k=1}^{\bar n}\frac{2^k}{k!}\,k\,k!\,4^k\big(\mathsf A(c_Kb_{\max})^Dv_b\big)^k\Big[(2b_{\max})^{kD}\epsilon_{b_{\min}}^{1/2}+(kD)^{kD+1}(1+b_{\max}+\gamma_0^{-1})^{kD}\big(\mathsf M_1c_Kb_{\max}+c_{\mathrm D}\big)e^{-\mu_1w_b/8}\Big].
\]
What is left of $\sum_k\kappa_{\nu_\Box}(W;k)/k!$ are the words in $\Psi_1'\chi$ and $\Psi_2\chi$ alone, that is $\sum_k\kappa_{\nu_\Box}(W_\Box\chi;k)/k!$. (L8) By \Cref{rot:lem:R1cum} backwards and $e^{\chi W_\Box}\ge\chi e^{W_\Box}$, $\sum_k\kappa_{\nu_\Box}(W_\Box\chi;k)/k!\ge\log\int\nu_\Box\chi^\Box_\beta e^{W_\Box}-(2M)^{\bar n+1}$. (L9) $\prod_\Box\int\nu_\Box\chi^\Box_\beta e^{W_\Box}=\int\nu_{\mathrm{loc}}\prod_\Box\chi^\Box_\beta e^{W_\Box}\ge e^{-\omega_{\mathrm{tot}}}\int\mu_P(dz_D\,|\,z_\Gamma)\prod_\Box\chi^\Box_\beta e^{W_\Box}$ by \Cref{rot:lem:R1dec} backwards (the integrand is non-negative and supported in the required set). The $c_\Box$ do not depend on $z_\Gamma$, so they come out of the $z_\Gamma$ integral; with $\chi^D_\beta\ge\chi^D_{\gamma_K\beta}$ and $V_{J''}=V_\Gamma+\sum_\Box W_\Box+R_1$,
\[
\int d\mu_P\,\chi^\Gamma_{\gamma_K\beta}\chi^D_\beta\,e^{V_\Gamma+\sum_\Box W_\Box}\ge e^{-|R_1|_{\max}}\int d\mu_P\,\chi_{\gamma_K\beta}\,e^{V_{J''}}=e^{-|R_1|_{\max}}\,\Phi(\gamma_K\beta;J'').
\]
(L10) By \Cref{rot:lem:R1bk}, $\sum_{\Box\in\mathcal B'}c_\Box\ge\sum_{k\le\bar n}\frac1{k!}[\kappa_{\mu_P}(V_{J'};k)-\kappa_{\mu_P}(V_{J''};k)]-N'e_{\mathrm{bk}}$. Collecting the errors, with $e_R\defeq2v_b\mathsf A(c_Kb_{\max})^De^{-\varsigma w_b/2}$ (for $R_0$ and $R_1$), $e_{\mathrm{dec}}\defeq2v_bc_Pc_de^{-\delta_Pw_b}\big[(c_Kb_{\max})^2/2+n/\gamma_0+(c_P\tilde c(\delta_P)c_Kb_{\max})^2/(2\gamma_0^2)\big]$ (two uses of \Cref{rot:lem:R1dec}), $e_{\mathrm{ext}}\defeq4e^{2M}\epsilon_{b_{\min}}$ and $e_{\mathrm{cum}}\defeq2(2M)^{\bar n+1}$, proves (L).

\emph{Upper bound.} The same steps with the following changes. (U2) Enlarge the cutoff on $D$: $\chi_\beta\le\chi^\Gamma_\beta\chi^D_{\beta'}$ with $\beta'\defeq\beta/\gamma_K$; the $z_D$-integrand is non-negative and supported in $\{\|z_D\|_\infty\le c_K\beta'\}$, and $\|z_\Gamma\|_\infty\le c_K\beta=c_K\gamma_K\beta'$, so \Cref{rot:lem:R1sch} gives $\|h_\Box\|_\infty\le\beta'/2$ and the upper half of \Cref{rot:lem:R1dec} applies with the radius $c_K\beta'$. (U3) $\int\nu_\Box\chi^\Box_{\beta'}e^{\Psi_\Box}\le\int\nu_\Box e^{\Psi_\Box\chi^\Box_{\beta'}}$. (U4)--(U7) as (L4)--(L7) with $\beta'$ in place of $\beta$, using the upper half of \Cref{rot:lem:R1cum}. (U8) $\sum_k\kappa_{\nu_\Box}(W_\Box\chi;k)/k!\le\log\int\nu_\Box e^{W_\Box\chi}+(2M)^{\bar n+1}\le\log\int\nu_\Box\chi^\Box_{\beta'}e^{W_\Box}+2e^{M}\epsilon_{\beta'}+(2M)^{\bar n+1}$ by \Cref{rot:lem:R1ext}(ii). (U9) The upper half of \Cref{rot:lem:R1dec} backwards, then $\chi^\Gamma_\beta\chi^D_{\beta'}\le\chi_{\beta'}$ and $e^{V_\Gamma+\sum_\Box W_\Box}\le e^{|R_1|_{\max}}e^{V_{J''}}$. (U10) The upper half of \Cref{rot:lem:R1bk}. Since all thresholds used lie in $[b_{\min},b_{\max}]$, every error is again bounded by $E_*$.
\end{proof}

\begin{lemma}[bookkeeping]\label{rot:lem:R1bk}
$\big|\sum_{\Box\in\mathcal B'}c_\Box-\sum_{k=1}^{\bar n}\frac1{k!}\big[\kappa_{\mu_P}(V_{J'};k)-\kappa_{\mu_P}(V_{J''};k)\big]\big|\le N'e_{\mathrm{bk}}$ with $e_{\mathrm{bk}}\defeq\sum_{k=1}^{\bar n}\frac{6^k}{k!}\,v_bK_k(\bar\mu/2)\,\mathsf A^ke^{-\bar\mu w_b/4}$.
\end{lemma}
\begin{proof}
Here the law is the centred $\mu_P$, so \Cref{rot:lem:R1conn} applies with the covariance decay $g_0e^{-\mu_1|x-y|}$ of \Cref{rot:lem:R1ct} and without the factor $(1+\beta)^{kD}$. (a) $V_{J'}=V_{J''}+\sum_{\Box\in\mathcal B'}X_\Box+R_2$ with $X_\Box\defeq\Psi_1''$ of the cell $\Box$: a monomial of $V_{J'}$ not inside $J''$ has a site of $J'$ at level at least $2.5w_b$, so it meets $\mathrm{Deep}(\Box)$ for a $J$-cell $\Box$; if its support lies in $\Box\cap\{\lambda\ge1.5w_b\}$ it belongs to $\Psi_1$ and not to $\Psi_1'$, that is to $X_\Box$; otherwise it has $\ell\ge\dist(\mathrm{Deep}(\Box),\Lambda\setminus(\Box\cap\{\lambda\ge1.5w_b\}))\ge w_b$ and is put into $R_2$. (b) Expand $\kappa_{\mu_P}(V_{J''}+X+R_2;k)$, $X\defeq\sum_\Box X_\Box$, into words. Words with at least one $R_2$: \Cref{rot:lem:R1conn} anchored at $S_1\defeq\bigcup_{\Box\in\mathcal B'}\mathrm{Deep}(\Box)$ (met by every monomial of $R_2$; $|S_1\cap\Lambda|\le N'v_b$) with $\theta'=\bar\mu/2$ and $\|R_2\|_{\bar\mu/2}\le\mathsf Ae^{-\bar\mu w_b/2}$ (\Cref{rot:lem:R1lat}(iv)), the other factors having $\|\cdot\|_{\bar\mu/2}\le\mathsf A$: at most $N'v_bK_k(\bar\mu/2)\mathsf A^ke^{-\bar\mu w_b/2}$ each. Words with some $X$ and no $R_2$: sort them by the cell $\Box$ of the first $X$ factor and write every other $X$ factor as $X_\Box+X_{\ne\Box}$. Words containing an $X_{\ne\Box}$: \Cref{rot:lem:R1conn} with $S_1=\mathrm{Deep}(\Box)$, $S_2\defeq\bigcup_{\tilde\Box\ne\Box}(\tilde\Box\cap\{\lambda\ge1.5w_b\})$, $\dist(S_1,S_2)\ge4w_b$ and $\theta'=\bar\mu$: at most $v_bK_k(\bar\mu)\mathsf A^ke^{-2\bar\mu w_b}$ each. Words with only $X_\Box$ and $V_{J''}$ factors: for this cell write $V_{J''}=\Psi_1'+Y^{\mathrm{long}}+Y^{\mathrm{far}}$, where $Y^{\mathrm{long}}$ consists of the monomials of $V_{J''}$ meeting $\Gamma_4(\Box)$ that are not in $\Psi_1'$ (they are not inside $\Gamma_3\cup\Gamma_4$, so they meet $\{\lambda<1.5w_b\}$ or another cell, and have $\ell\ge\frac12w_b$) and $Y^{\mathrm{far}}$ of the rest (supports in $J''\setminus\Gamma_4(\Box)$, at distance at least $\frac12w_b$ from $\mathrm{Deep}(\Box)$). The words in $X_\Box$ and $\Psi_1'$ alone with at least one $X_\Box$ add up to exactly $\kappa_{\mu_P}(\Psi_1'+\Psi_1'';k)-\kappa_{\mu_P}(\Psi_1';k)$, which is $k!$ times the $k$-th term of $c_\Box$. Words with a $Y^{\mathrm{far}}$: \Cref{rot:lem:R1conn} with $S_1=\mathrm{Deep}(\Box)$, $S_2=J''\setminus\Gamma_4(\Box)$, $\theta'=\bar\mu$: at most $v_bK_k(\bar\mu)\mathsf A^ke^{-\bar\mu w_b/4}$ each. Words with a $Y^{\mathrm{long}}$ and no $Y^{\mathrm{far}}$: anchored at $\mathrm{Deep}(\Box)$ with $\theta'=\bar\mu/2$ and $\|Y^{\mathrm{long}}\|_{\bar\mu/2}\le\mathsf Ae^{-\bar\mu w_b/4}$: at most $v_bK_k(\bar\mu/2)\mathsf A^ke^{-\bar\mu w_b/4}$ each. Every slot of a word takes one of six roles ($\Psi_1'$, $Y^{\mathrm{far}}$, $Y^{\mathrm{long}}$, $X_\Box$, $X_{\ne\Box}$, $R_2$), so there are at most $6^k$ words for each $k$; each word outside $c_\Box$ is bounded, per cell, by $v_bK_k(\bar\mu/2)\mathsf A^ke^{-\bar\mu w_b/4}$ (as $K_k(\bar\mu)\le K_k(\bar\mu/2)$ and all rates are at least $\bar\mu w_b/4$), and the words made of $V_{J''}$ alone add up to $\kappa_{\mu_P}(V_{J''};k)$.
\end{proof}


\begin{theorem}[the cumulant bound for the step Gaussian]\label{rot:thm:R1}
In the setting \textup{(R1)}--\textup{(R4)} above (any $\mathbb{L}$ among $\Z^3$, the cubic tori and the skew tori, any finite $\Lambda\subset\mathbb{L}$, the Euclidean or the two-component cutoff) with $b_*$ the threshold fixed before \Cref{rot:lem:R1stg}, there are $S<\infty$ and $\rho_3>0$, all three depending only on $(n,\bar n,D,\varsigma,\gamma_0,c_P,\delta_P,c_K)$, such that for $b\ge b_*$ and $\mathsf A\,b^{D+6}\le1$
\begin{align*}
\log\int\chi_b\,e^{V}d\mu_P&\le\sum_{k=1}^{\bar n}\frac{\kappa_{\mu_P}(V;k)}{k!}+S|J|\big[(\mathsf A\,b^{D+6})^{\bar n+1}+e^{-\rho_3b^{3/2}}\big],\\
\log\int\chi_b\,e^{V}d\mu_P&\ge\sum_{k=1}^{\bar n}\frac{\kappa_{\mu_P}(V;k)}{k!}-S|J|\big[(\mathsf A\,b^{D+6})^{\bar n+1}+e^{-\rho_3b^{3/2}}\big]-4n|\Lambda|\,e^{-\gamma_0\gamma_K^{8}b^2/(2n)},
\end{align*}
where $\kappa_{\mu_P}(V;k)$ are the truncated expectations of $V$ under $\mu_P$ without cutoff and $\gamma_K=\min\{\frac12,(2c_K\mathsf M_1)^{-1}\}$ is the constant of \eqref{rot:eq:R1const}. The constants do not depend on $\Lambda$, on $J$, on the torus, or on $P$ and $V$ otherwise, and no translation invariance is used. For the step Gaussians of both flows, at every step $j$ below the stopping scale, on $T^{(j)}$, on a principal block of it or on a Neumann sub-domain that is a union of $j$-blocks, in the $X$-frame and in the $Y$-frame steps alike, the hypotheses hold with $\gamma_0$ Ba{\l}aban's coercivity constant \cite[(2.33)]{Ba1}, $\delta_P$ the decay rate of the step precision divided by $\sqrt3$, and $c_P$ a constant of $(a,M)$.
\end{theorem}
\begin{proof}
\emph{The hypotheses for the step Gaussians.} The lower bound $\Delta^{(j)}+a5^{-2}Q^*Q\ge\gamma_0$ is \cite[(2.33)]{Ba1} (for the hybrid flow by \Cref{rot:prop:BD}(a)), with constants independent of $A$, $j$ and the domain; this is (R1), which needs no upper bound. The kernel bound (R2) is the exponential decay of $(M_j)^{-1}$ by Combes--Thomas conjugation (for the hybrid flow, the kernel bound of \Cref{rot:prop:BD}(a)) (whose first-order hopping part is anti-Hermitian and whose remainder has norm $\le3d\delta^2$ uniformly in the spacing), with $\dist(B(u),B(v))\ge|u-v|-\sqrt3$, and the block term $a5^{-2}Q^*Q$ has range $\sqrt3\cdot5$; only the off-diagonal entries are constrained, so a diagonal term (such as the shell mass of the shell model) leaves the constants unchanged; Ba{\l}aban's bounds are stated in the sup metric, and $|x-y|\le\sqrt3\,d^*(x,y)$ converts them to the quotient Euclidean distance at the rate divided by $\sqrt3$; principal blocks and Neumann sub-domains inherit (R1) and (R2) as sub-matrices. The interaction at step $j$ is a real polynomial of fixed degree in the fluctuation variables (the analytic remainder vertices of order $\bar n+1$ are taken out of the exponent on the support of the cutoff as $e^{\pm C\ell_j^{7-\beta}|\Lambda|}$, as \cite[(2.31)]{Ba2} does), with $\mathsf A=\|V\|_\varsigma\le s_0\,O(1)\ell_j^{\beta_0}$ by \Cref{rot:cor:OF} at $\varsigma$ half of Ba{\l}aban's coefficient decay rate: a coefficient bound $O(1)\ell_j^{\beta_0}e^{-\delta_0d(x_1,\dots,x_p)}$, with $d$ the length of the shortest graph connecting the points \cite[(3.58)]{Ba1}, gives $\|V\|_{\delta_0/2}\le s_0O(1)\ell_j^{\beta_0}$ with $s_0\defeq\sum_{p\le D}p\,n^p\,\tilde c(\delta_0/(2p))^{p-1}$, because $d(x_1,\dots,x_p)\ge(p-1)^{-1}\sum_{i\ge2}|x_i-x_1|$; this is (R3). The cutoff is Ba{\l}aban's, a product over sites of $\mathbf 1\{|A'(x)|\le b\}\mathbf 1\{|\phi'(x)|\le b\}$ \cite[(3.43), (3.50)]{Ba1}, which is (R4) with $c_K=\sqrt2$. Finally $\mathsf Ab^{D+6}\le1$ and $b\ge b_*$ hold at $b=p_1(\ell_j)=b_1(1+\log\ell_j^{-1})^{p_1}$ for $\ell_j$ below Ba{\l}aban's final spacing, since $\ell_j^{\beta_0}$ beats every power of $1+\log\ell_j^{-1}$; the exponent $p_1\ge1$ is a choice within Ba{\l}aban's freedom, under which $e^{-\rho_3b^{3/2}}$ and $e^{-\gamma_0\gamma_K^8b^2/(2n)}$ are $O(\ell_j^{\kappa})$ for every $\kappa$.

% r03_r1_proof: the cumulant bound for the step Gaussian, part V: the thresholds, the small-volume case and the assembly over the four pavements. Input inside the proof of Theorem rot:thm:R1 in sec/r03_fill_K5a.tex.
\emph{The thresholds.} The threshold $b_*$ is the one fixed before \Cref{rot:lem:R1stg}; it depends only on $(n,\bar n,D,\gamma_0,c_P,\delta_P,c_K)$.

\emph{Small volume.} If some pavement consists of a single cell, then $|\Lambda|\le v_b$ by (G4) and no pavement is needed: with $\nu=\mu_P$ (centred), $\epsilon\defeq\mu_P(1-\chi_b)\le\epsilon_b$ by \Cref{rot:lem:R1tail} and $M\defeq\mathsf A|\Lambda|(c_Kb)^D\le\mathsf m_*$, \Cref{rot:lem:R1ext} moves the cutoff into the exponent at a cost of at most $2e^{2M}\epsilon$, \Cref{rot:lem:R1cum} expands to order $\bar n$ at a cost of $(2M)^{\bar n+1}$, and \Cref{rot:lem:R1chi} replaces $\kappa_{\mu_P}(V\chi_b;k)$ by $\kappa_{\mu_P}(V;k)$ at a cost of $k\,k!\,4^k(\mathsf A|\Lambda|)^k(2b)^{kD}\epsilon^{1/2}/k!$ for each $k$; the total is at most $E_*\le|J|E_*$ in both directions. (If $J=\emptyset$ then $V=0$ and the statement reduces to $0\ge\log\mu_P(\chi_b)\ge-4n|\Lambda|e^{-\gamma_0b^2/(2n)}$, which is \Cref{rot:lem:R1gci}.)

\emph{Assembly.} Otherwise iterate \Cref{rot:lem:R1stg} over $i=0,1,2,3$ with $J_0\defeq J$, $J_{i+1}\defeq J_i\cap\{\lambda_i<2.5w_b\}$ (the set $J''$ of stage $i$), and the thresholds $\beta_i\defeq\gamma_K^ib$ for the lower bound and $\beta_i\defeq\gamma_K^{-i}b$ for the upper bound, all of which lie in $[b_{\min},b_{\max}]$ together with $\beta_3/\gamma_K$. By (G3), $J_4\subset J\cap\bigcap_i\{\lambda_i<2.5w_b\}=\emptyset$, so $V_{J_4}=0$ and the cumulant increments of the four stages telescope to $\sum_{k\le\bar n}\kappa_{\mu_P}(V_J;k)/k!$; the numbers of $J$-cells add up to $\sum_iN_i\le4|J|$. The terminal factor is $\Phi(\gamma_K^4b;\emptyset)=\mu_P(\chi_{b_{\min}})\ge\exp\big(-4n|\Lambda|e^{-\gamma_0\gamma_K^8b^2/(2n)}\big)$ in the lower bound, by \Cref{rot:lem:R1gci} ($\mu_P$ is centred with component variances at most $\|C\|\le\gamma_0^{-1}$), and $\Phi(\gamma_K^{-4}b;\emptyset)=\mu_P(\chi_{b_{\max}})\le1$ in the upper bound. Hence
\[
\log\int\chi_be^Vd\mu_P-\sum_{k=1}^{\bar n}\frac{\kappa_{\mu_P}(V;k)}{k!}\ \le\ 4|J|E_*(b,\mathsf A),\qquad\log\int\chi_be^Vd\mu_P-\sum_{k=1}^{\bar n}\frac{\kappa_{\mu_P}(V;k)}{k!}\ \ge\ -4|J|E_*(b,\mathsf A)-4n|\Lambda|e^{-\gamma_0\gamma_K^8b^2/(2n)}.
\]

\emph{The constants.} $e_{\mathrm{cum}}=2(2M)^{\bar n+1}=2(2\mathsf m_*)^{\bar n+1}(\mathsf Ab^{D+6})^{\bar n+1}$ exactly. Every other term of $E_*$ is increasing in $\mathsf A$, so under $\mathsf Ab^{D+6}\le1$ it is at most its value at $\mathsf A=b^{-D-6}$; call the sum of these values $Q(b)$. It is a finite sum of polynomials in $b$ times one of $e^{-\varsigma w_b/2}$, $e^{-\delta_Pw_b}$, $e^{-\bar\mu w_b/4}$, $e^{-\mu_1w_b/8}$, $e^{-2\bar\mu w_b}$ and $\epsilon_{b_{\min}}^{1/2}$, each of which is at most a polynomial in $b$ times $e^{-\bar\mu b^{3/2}/8}$ (since $w_b\ge b^{3/2}$, $\bar\mu\le\mu_1\le\delta_P/2$ and $\bar\mu\le\varsigma$). With $\rho_3\defeq\bar\mu/16$ and
\[
S\defeq4\max\Big\{2(2\mathsf m_*)^{\bar n+1},\ \sup_{b\ge b_*}e^{\rho_3b^{3/2}}Q(b)\Big\}<\infty
\]
one has $4|J|E_*(b,\mathsf A)\le S|J|\big[(\mathsf Ab^{D+6})^{\bar n+1}+e^{-\rho_3b^{3/2}}\big]$, which is the statement. $S$, $b_*$, $\rho_3$ and $\gamma_K$ depend only on $(n,\bar n,D,\varsigma,\gamma_0,c_P,\delta_P,c_K)$; nothing depends on $\Lambda$, on $J$, on $\mathbb{L}$ or its period lattice, on $P$ and $V$ beyond the stated bounds, or on any translation invariance, of which none is used; no lower bound on $c_P$ and no upper bound on $P$ enters.

\end{proof}

% r03_fill_K5c: proof of (K_0), part VII (continued): the hybrid lower bound, the polynomial lower bound and the floor shift.

\begin{proposition}[the hybrid lower bound]\label{rot:prop:H4}
For the hybrid flow on $T$ and every real smooth periodic source $s$,
\[
\Zun_{\Xk,T}(s)\ge e^{-C(1+k)^p\delta_r^\sigma|T|}\int\chi_r\,e^{-\hs^{(r)}_T}.
\]
\end{proposition}
\begin{proof}
This is the analogue of \cite[Prop.~4.1, (4.5)]{K2} for the hybrid flow (the print has ``$+C$'' in the exponent; in a lower bound it must be ``$-C$''; the factor $(1+k)^p$ is that of \Cref{rot:cor:OF} for the hybrid flow and is harmless, $(1+k)^p\delta_r^\sigma\to0$ by \Cref{rot:lem:r}). Its inputs are: the block-spin algebra of the tile average (\Cref{rot:lem:tile}, \Cref{rot:lem:pair}); the hybrid analogues of \cite[Props.~2.1--2.3]{Ba1}, the decay and regularity of $G_j$, $\Delta^{(j)}$, $C^{(j)}$ along the hybrid flow with the paired contours, with \cite[Prop.~5.6]{K2} for the restricted inverses on unions of big blocks, which is \Cref{rot:prop:BD}(a) at every level ($j\ge2$ by its proof, $j\in\{0,1\}$ by the same proofs at the tile and first $Y$-frame levels); the cumulant lower bound for the step Gaussian, \Cref{rot:thm:R1}, whose extra lower-bound term $4n|\Lambda|e^{-\gamma_0\gamma_K^8b^2/(2n)}$ at the threshold $b_j=p_1(t_j)=b_1(1+\ln t_j^{-1})^{p_1}$, $t_j=5^j\eps$, $p_1\ge1$ \cite[(3.43), (3.50)]{Ba1}, on $\Lambda=T^{(j)}_1$ with $|T^{(j)}_1|=|T|t_j^{-3}$, is $\le C_\kappa t_j^\kappa|T^{(j)}_1|$ for every $\kappa$, of the form of Ba{\l}aban's own step error, and summed over $j<r$ is $\le C'_\sigma\delta_r^\sigma|T|$, uniformly in $T$, $k$, $r$; and Ba{\l}aban's lower-bound argument \cite[\S\S3--4]{Ba1} at block factor $5$ (\Cref{thm:a06:BALA}) on the torus on which the steps are performed. The hybrid flow runs step $1$ in $X$'s frame and steps $2,\dots,r$ in $Y$'s frame. On $T_1$, aligned for $X$ and skew for $Y$, steps $2,\dots,r$ are steps on a torus skew in the frame of their blocks, and the lower-bound argument is the hybrid instance of \Cref{rot:prop:SK}, with the constant $C(1+k)^p$ of \Cref{rot:cor:OF} for the hybrid flow (Ba{\l}aban's Neumann cube step being replaced by the constant-background reduction of \Cref{rot:lem:CB}). On $T_2$, skew for $X$ and aligned for $Y$, only step $1$ is skew: the fine axis lattice $\Xk$ on the torus with period lattice $R^{-1}P_J$; its blocks tile $T_2$ (\Cref{rot:lem:adm}(ii), \Cref{rot:lem:tile}(i)), its fluctuation covariance has mass of order $\eps^{-1}$ in physical units, so the torus images of its kernels are $O(e^{-c\,50MJ/\eps})$, and the step lands on $Y_{k-1}\subset T_2$, aligned for every later step. Both instances pass through the bridge \Cref{rot:prop:CT} (King adopts the free-boundary counterterm on \cite[p.~330]{K2} and appeals to \cite{Ba1} for (4.5)) and are stated for the centred-block hybrid action (\Cref{rot:conv:C}).
\end{proof}

\begin{proposition}[the polynomial lower bound for the rotated lattice]\label{rot:prop:H5}
For $\Yk$ on $T$ and every source $s=(zG,h_0)$ with $h_0=-n\psi$, $n\ge0$: $\Zun_{\Yk,T}(s)\ge e^{-C_1|T|}\eps^{\,q(s)}$ for some $q(s)<\infty$.
\end{proposition}
\begin{proof}
At $h=0$ this is the lower half of \cite[(4.2)]{K2} (``The results of [Ba 1--4] apply to both these models''), Ba{\l}aban's lower bound \cite[(1.14)]{Ba1}: on $T_2$, aligned in $Y$'s frame, by \Cref{thm:a06:T1,thm:a06:BALA} after the passive rotation, with the bridge \Cref{rot:prop:CT}; on $T_1$ by the skew instance \Cref{rot:prop:SK}. The source $zG$ is removed by Jensen's inequality: $\Zun_{\Yk,T}(zG,h)=\Zun_{\Yk,T}(0,h)\langle e^{zF(G)}\rangle_{(T,h),\Yk}\ge\Zun_{\Yk,T}(0,h)\,e^{z\langle F(G)\rangle}=\Zun_{\Yk,T}(0,h)$ for real $z$, since $\langle F(G)\rangle=0$ under the charge conjugation of \Cref{rot:def:tori}. At the mass-increasing wall $h_0=-n\psi$, $n>0$, the two-sided $\eps$-uniform bound \cite[(4.2)]{K2} is false (\Cref{rot:rem:log}) and only the polynomial form is asked for; it is \Cref{rot:thm:FS}, at every King point, or alternatively \Cref{rot:prop:J} from the source-free instance and an upper bound on the one-point Wick function.
\end{proof}

\begin{proposition}[Jensen: the wall instance from the source-free instance]\label{rot:prop:J}
Let $h_0=-n\psi\le0$ be a wall. For every real $t$, $\Lambda\in\{\Xk,\Yk\}$ and $T\in\{T_1,T_2\}$,
\[
\Zun_{\Lambda,T}(tG,h_0)\ \ge\ \Zun_{\Lambda,T}(0,0)\,\exp\Big[-n\sum_x\eps^3\psi(x)\big\langle|\phi(x)|^2-2\sW_\eps\big\rangle_{T,0,\Lambda}\Big],
\]
the expectation in the source-free periodic state. On $\Yk$ the one-point function $w_k(T)\defeq\langle|\phi(x)|^2-2\sW_\eps\rangle_{T,0,\Yk}$ is independent of $x$ ($P_J\subset\Yk$, \Cref{rot:lem:adm}), so the right side is $\Zun_{\Yk,T}(0,0)e^{-n\sS_{\Yk}(\psi)w_k(T)}$, and an upper bound $w_k(T)\le C(1+k)$ gives the polynomial form of \Cref{rot:prop:H5} at the wall with $q(tG,h_0)=q(0,0)+nC\sS(\psi)/\ln5$.
\end{proposition}
\begin{proof}
$W_\Lambda(h_0)=\sum_x\eps^3h_0(x)|\phi(x)|^2-2\sW_\eps\sS_\Lambda(h_0)$ gives $\Zun(tG,h_0)=e^{-2\sW_\eps\sS(h_0)}\Zun(0,0)\langle e^{tF(G)+\sum\eps^3h_0|\phi|^2}\rangle_{T,0,\Lambda}$. Put $X\defeq-tF(G)+n\sum_x\eps^3\psi(x)|\phi(x)|^2$; then $\langle X\rangle=n\sum_x\eps^3\psi(x)\langle|\phi(x)|^2\rangle$ because $\langle F(G)\rangle=0$ by charge conjugation. Jensen's inequality $\langle e^{-X}\rangle\ge e^{-\langle X\rangle}$ and $-2\sW_\eps\sS(h_0)-\langle X\rangle=-n\sum_x\eps^3\psi(x)(\langle|\phi(x)|^2\rangle-2\sW_\eps)$ give the bound. Only an upper bound on the one-point function is used, only on $\Yk$ and only without sources; the mass-decreasing sign would need a lower bound, which Jensen's inequality does not give.
\end{proof}

\begin{theorem}[floor shift: the polynomial lower bound at the wall at every King point]\label{rot:thm:FS}
Let $m^2$ be any King point and $h$ a real $P_J$-periodic function (it may depend on $k$) with $\|h\|_\infty\le H$ uniformly in $k$. Write $a_s\defeq s+\dm(s;\eps)$ for the bare scalar coefficient at the King point $s$ and let $M^2$ be a King point with $\Delta_0(M^2)\defeq\liminf_{\eps\to0}[a_M-a_m]>2H$; such an $M^2$ exists for every coupling and every $H$, since $\Delta_0(M^2)\to+\infty$ as $M^2\to\infty$ (the coefficients of $\partial_s\dm$ are non-negative at order $\lambda$, $O(s^{-1/2})$ at order $e^2$ and $O(s^{-1})$ at the two-loop orders, \cref{app:a05:B}). Then on $\Lambda\in\{\Yk,\Xk\}$, $T\in\{T_1,T_2\}$,
\[
\ln\Zun^{m^2}_{\Lambda,T}(0,h)\ge-C|T|-C(1+k)|T|-qk\ln5,\qquad\text{hence}\qquad\Zun^{m^2}_{\Lambda,T}(tG,h)\ge e^{-C_1|T|}\eps^{\,q'}\quad\text{for every real }t .
\]
This is the polynomial form of \Cref{rot:prop:H5} at the wall $h_0=-n\psi$ for every $n>0$; the sign of $h$ is used nowhere.
\end{theorem}
\begin{proof}
The inputs are: (i) the source-free lower bound \cite[(1.14)]{Ba1} at the one fixed King point $M^2$ (\Cref{thm:a06:T1}); (ii) the upper bound \cite[(1.3)]{Ba2} at one fixed $M_+^2$ with $\liminf_{\eps\to0}[a_{M_+}-a_M]>0$ (which exists by the same growth); on $T_1$ the skew instances of \Cref{rot:prop:SK}; (iii) the bounds on the coefficients of the one-point Wick function through order $13$ in the couplings on a compact mass interval, from Reisz's power-counting theorem \cite[Thm.~1]{Rei88} with the counterterm conditions \cite[(3.33)]{K1} (the order-$e^2$ coefficient is computed in \cref{app:a05:A}), and the boundedness of $Y\defeq1+\partial_s\dm$ on $[m^2,M_+^2]$ (\cref{app:a05:B}); (iv) exact identities: the bare-mass form $\Zun^{s}(T;0,u\mathbf 1)=e^{-P_s-2c_su|T|}\hat Z(a_s-2u)$ with $\hat Z(\gamma)$ the convention-free integral at bare coefficient $\gamma$ and $P_s=E_0(s)+E_1(s)$, the integrated counterterm identity for $P_M-P_m$, Jensen's inequality, the monotonicity of $\langle\sum_x\eps^3|\phi|^2\rangle_\gamma$ in $\gamma$, translation invariance on $\Lambda$ and charge conjugation.
\emph{Step 1 (floor shift).} $e^{\sum\eps^3h|\phi|^2}=e^{h_{\min}\sum\eps^3|\phi|^2}e^{\sum\eps^3(h-h_{\min})|\phi|^2}$ moves the wall into the bare coefficient, $\gamma'\defeq a_m-2h_{\min}$, and Jensen's inequality bounds the remaining factor below by $\exp[\sum\eps^3(h-h_{\min})\langle|\phi|^2\rangle_{\gamma'}]\ge1$.
\emph{Step 2 (the higher King point).} $u\defeq\frac12(a_M-\gamma')=\frac12(a_M-a_m+2h_{\min})$ lies in $[0,C]$ for $\eps\le\eps_0$, by $h_{\min}\ge-H$, $\Delta_0(M^2)>2H$ and the finite upper limit of $a_M-a_m$; since $-\partial_\gamma\ln\hat Z=\frac12\langle\sum\eps^3|\phi|^2\rangle_\gamma$ is nonincreasing in $\gamma$, $\ln\hat Z(\gamma')\ge\ln\hat Z(a_M)+u|T|(2c_M+w_k(M))$, where input (i) bounds $\ln\hat Z(a_M)=\ln\Zun^{M^2}(0,0)+P_M$ below by $-C|T|$ plus the $c$-numbers.
\emph{Step 3 ($c$-numbers).} By the integrated counterterm identity the four $\eps^{-1}$-divergent $c$-numbers collapse exactly to $2u(c_M-c_m)|T|=O(|T|)$; the remainder $|T|\int_{m^2}^{M^2}(c_m-c_s)Y\,ds$ is bounded below by $-C|T|$ from $0\le c_m-c_s\le C$ and $|Y|\le C$ (its sign is not needed), and the perturbative remainder is $\ge-C(1+k)|T|$ by (iii).
\emph{Step 4.} $w_k(M)\ge-C(1+k)$ is the mirror bound from (i)--(iii) at $M^2$, $M_+^2$: the same floor shift with the roles of the two points exchanged and the upper bound (ii) in place of the lower one.
\emph{Step 5.} Assemble; the source $tG$ is removed as in \Cref{rot:prop:H5}. No stability bound at $m^2$ itself, no bound on $w_k(m^2)$ and no two-sided Wick bound is used. The statement is in the free-boundary convention; the $c$-number mismatch with the torus convention is $\Delta c_T=O(e^{-c\ell(T)})$ (\Cref{rot:prop:CT}(b)).
\end{proof}

% r03_fill_K5b: proof of (K_0), part VIII: the skew-torus instances and the counterterm bridge.

\subsubsection*{The skew-torus instances}

Call a torus $\R^3/\Pi$ \emph{aligned} in the frame $(f_\alpha)=(Oe_\alpha)$ of a lattice $O(s\Z^3)$ if $\Pi=O(\ell_1\Z\times\ell_2\Z\times\ell_3\Z)$ for some $\ell_\mu>0$, that is, if the torus is a rectangular box with periodic boundary conditions whose sides are parallel to the frame vectors, and \emph{skew} in that frame otherwise. \cite{Ba1,Ba2,Ba3} and \cite{K1} are written for a torus aligned in the frame of the blocks (\cite[(2.21)]{K1}). $T_1=\DJ$ is aligned (a cube) in the frame of $\Xk$ and skew in the frame of $\Yk$; $T_2=R^{-1}\DJ$ is aligned for $\Yk$ and skew for $\Xk$. On $T_1$ every block-spin step performed in $Y$'s frame, that is all $r$ steps of the standard flow and steps $2,\dots,r$ of the hybrid flow, is a step on a torus skew in the frame of its blocks; on $T_2$ only the first hybrid step is.

\begin{proposition}[the skew-torus instances]\label{rot:prop:SK}
King's Theorem~3.1 \cite[(3.4)]{K1} (Ba{\l}aban's upper bound with the renormalisation of \cite{Ba3}) and Ba{\l}aban's lower-bound machinery \cite[\S\S2--4, (1.14)]{Ba1} hold for the standard flow of $\Yk$ on $T_1$, the torus $R\DJ$ in $Y$'s frame, which is skew, with the constants of the aligned case times at most $125$ in the linear part; and Ba{\l}aban's lower-bound argument holds for steps $2,\dots,r$ of the hybrid flow on $T_1$ with the constant $C(1+k)^p$ of \Cref{rot:cor:OF}. In both, the torus enters \cite{Ba1,Ba2,Ba3} only through the embeddings of the block lattices, the external fields and the image terms of the periodised kernels; no torus symmetry and no torus Fourier analysis is used in the proof of \cite[Prop.~1]{Ba3}. The statements hold under the provisos $J\in5^{n_*}\N$ and $J\ge J_0(\delta_r)$ of \Cref{rot:thm:CT}.
\end{proposition}

The proof consists of three transfer lemmas. Let $T'\defeq R\DJ$ be the skew torus in $Y$'s frame with frame period lattice $\Gamma'\defeq R^{\mathrm T}(50MJ\,\Z^3)$ (read in the frame coordinates), its cubic cover $\tilde T\defeq\R^3/250MJ\,\Z^3$ (a cube aligned in $Y$'s frame, side $250MJ\in M\N$), the covering $\pi:\tilde T\to T'$ with deck group $\cG=\Gamma'/250MJ\,\Z^3$ of order $125$, and the sup-systole $s_\infty(\Gamma')=40MJ$ (Ba{\l}aban's metric is the sup metric, \cite[p.~605]{Ba1}).

\begin{lemma}[folding]\label{rot:lem:F}
$250MJ\,R^{-1}\Z^3=50MJ\,(5R^{-1})\Z^3\subset50MJ\,\Z^3=P_J$ because $5R^{-1}$ is integral, with index $|\det5R^{-1}|=125$; so $\tilde T$ covers $T'$ with a deck group of order $125$. Every block lattice $Y_{k-j}$ contains $P_J$ (\Cref{rot:lem:adm}(iii)), so the deck translations commute with every block average and the Gaussian quadratic forms of the flow are deck-invariant; $\pi$ preserves bonds, blocks, big blocks and the contours, $\dist_{\tilde T}(\tilde x,\pi^{-1}S)=\dist_{T'}(x,S)$ in the quotient sup metric, and lifted operators intertwine: $G_j(\pi^{-1}\Omega,A\circ\pi)(f\circ\pi)=(G_j(\Omega,A)f)\circ\pi$, likewise $D^\eta$, the H\"older quotients, $\delta G_j$, $C^{(j)}$. Hence every kernel of a $Y$-frame flow on $T'$ ($G_j$, $\Delta^{(j)}$, $C^{(j)}$, the slice kernels) is the $125$-fold periodisation of the corresponding kernel on the aligned torus $\tilde T$, and \cite[Props.~2.1--2.3]{Ba1} ($=$\cite{Ba4}) and \cite[Thm.~3.3]{K1} hold on $T'$ with the $\tilde T$ constants times at most $125$ (fibre points are at least $s_\infty$ apart, and the finite sum of exponentially decaying terms is absorbed).
\end{lemma}
\begin{proof}
The covering step of the aligned case with $(\Z/3)^3$ replaced by $\cG$: a non-wrapping parallelepiped lifts to disjoint parallelepipeds \cite[p.~611]{Ba1}, and positivity on $T'$ is positivity on $\cG$-invariant functions. The deck group commutes with the tile average as well ($P_J\subset Y_{k-1}$), so the lemma covers the hybrid steps.
\end{proof}

\begin{lemma}[restriction principle]\label{rot:lem:D}
Call a functional $\Phi$ liftable of degree $a\in\{0,1\}$ if $\Phi_{\tilde T}(\xi\circ\pi)=125^a\Phi_{T'}(\xi)$ (local sums $a=1$, sup-type $a=0$, quadratic forms in Ba{\l}aban's operators on unions of big blocks $a=1$, background fields $a=0$). If an inequality $\Phi\le\Psi$ between liftable functionals of the same degree holds on $\tilde T$ for all configurations obeying block-local restrictions $\chi$ and all admissible regions, it holds on $T'$ with the same constants. Moreover $\ell^2(T')$ is the $\cG$-invariant subspace of $\ell^2(\tilde T)$ (inner product times $125$) and lifted operators commute with $\cG$, so $\|X_{T'}\|\le\|\tilde X\|$ and $\tilde X\ge0\Rightarrow X_{T'}\ge0$. The lemma does not carry integrals over field configurations (the $\cG$-invariant model on $\tilde T$ is a different measure); those are handled by \Cref{rot:lem:TC}.
\end{lemma}
\begin{proof}
Lifted configurations obey the lifted restrictions; lifted regions $\pi^{-1}\Lambda$ are unions of big blocks with the same distance and corridor data (\Cref{rot:lem:F}); quadratic forms scale by $125$ by the intertwining; backgrounds are images of intertwined operators (the model is abelian). Evaluate the $\tilde T$ inequality at $\xi\circ\pi$ and divide by $125^a$.
\end{proof}

\begin{lemma}[type closure]\label{rot:lem:TC}
With $\sigma_j\defeq s_\infty(\Gamma')/(5^j\eps)$, every step of \cite[\S\S2--4]{Ba1}, \cite[\S\S2--3]{Ba2}, \cite[\S\S1--3]{Ba3} and \cite[\S3]{K1} is of one of six types and holds verbatim on $T'$ with the aligned constants: \textup{(I)} exact identities on a finite graph (Gaussian integration, composition of the block-spin maps, block normalisations, set decompositions); \textup{(II)} pointwise bounds on the integrand (dropping non-negative terms, Neumann decoupling, indicator $\le$ exponential, rescalings); \textup{(III)} operator inequalities valid on any finite graph (the diamagnetic inequality, $\log\det(1+X)\le\Tr X$, $\Tr(X\mathbf 1_\Lambda)\le\|X\|\,|\Lambda|$, $-\Delta_A\ge0$, also with Neumann restriction); \textup{(IV)} decay and norm bounds for $G_j$, $C_j$, $\Delta^{(j)}$, $H_j$, $\delta G_j$ on unions of large blocks, wrapping allowed (\Cref{rot:lem:F,rot:lem:D}); \textup{(V)} local arguments on a set of sup-diameter $D$ (gauge-outs, Stokes, local expansions, the contour conventions inside a block, ``sets of sup-diameter $<\sigma/2$ do not wrap''), verbatim when $D<\sigma_j$: the sets are blocks, big blocks and corridors whose width in scale-$j$ units grows only logarithmically in $t_j^{-1}$, and since the product of the spacing with the corridor width increases along the flow the binding case is the coarsest scale $\delta_r$, which is the condition $40MJ\ge\delta_r\max\{6\sqrt3\,r(\delta_r)+2,\,2/\delta_0\}$, i.e.\ $J\ge J_0(\delta_r)$; \textup{(VI)} counts ($|T|$, $|\Lambda|$, $|\partial\Lambda|$, sums over scales, and the estimate \cite[(3.42)]{Ba2} in the form of \Cref{thm:a06:BALA}(d), whose box partition descends to $T'$ under the same non-wrapping condition; the printed cube count \cite[(3.45)]{Ba2} and the first inequality of (3.46) are false as printed, at every block factor, and are not used): they depend on the torus only through the block count. The large-field analysis (\cite[\S3]{Ba2}, \cite[\S4]{Ba1}, \cite[\S3.2]{K1}) uses (II)--(III) on large-field blocks, (I), (III), (IV) for the fluctuation integral with Neumann decoupling, (VI) for the sum over large-field configurations, $\prod_j\sum_\Lambda e^{-p(t_j)|\Lambda|}\le e^{C|T|}$, and (I)/(IV)/(V) with the perturbative estimates for the collar terms; the lower bound \cite[\S4]{Ba1} restricts to the whole-torus small-field region by Jensen's inequality and uses (I)--(IV) only. The torus enters the prints only through admissibility, the systole, the kernels and their images, counts, counterterms and Wick images, external-field periodisation and symmetries, and once through a coordinate-parity split of the bonds inside \cite[Prop.~3.1$'$]{Ba4} ($=$\cite[(3.29)]{Ba2}), a quadratic-form statement that folds from the cover (\Cref{rot:lem:D}; on $T'$ the periods are even at every scale, so it also holds verbatim); the one momentum-space formula of \cite{Ba3}, (3.28) on p.~441, uses the cubic symmetry with an image-term defect $e^{-c\sigma}$.
\end{lemma}
\begin{proof}
Each display of \cite[Ch.~1--4]{Ba1}, \cite[pp.~556--594]{Ba2}, \cite[pp.~411--445]{Ba3} and \cite[App.]{K1}, read from the page images, is assigned to a type or to one of the entry points listed; types (I)--(III) are torus-free, (IV) is \Cref{rot:lem:F,rot:lem:D}, (V) is the non-wrapping condition, and (VI) is the block count. The one external lemma, the cumulant bound, is \Cref{rot:thm:R1} on the skew tori.
\end{proof}

\begin{proof}[Proof of \Cref{rot:prop:SK}]
The linear part is \Cref{rot:lem:F}; the nonlinear part (the small- and large-field analysis) is \Cref{rot:lem:TC} relative to the same bounds on aligned tori at block factor $5$, which are \Cref{thm:a06:BALA}. The hybrid steps $2,\dots,r$ act on composite blocks with the hybrid kernels of \Cref{rot:prop:BD}(a); Ba{\l}aban's Neumann cube step is replaced by the constant-background reduction of \Cref{rot:lem:CB} and the whole-torus form of \cite[Prop.~1]{Ba3} for the hybrid flow is \Cref{rot:cor:OF}, whence the constant $C(1+k)^p$. The provisos: $J\in5^{n_*}\N$ because the frame period lattice $\Gamma'$ must lie on the $5$-adic grid $M_b5^{n_*}\Z^3$, $n_*$ the number of large-field scales of the $Y$ flow, fixed by the couplings; and the non-wrapping cap $J\ge J_0(\delta_r)$ of type (V). Both are harmless and cofinal: in the flow of this subsection the coarsest spacing is $\delta_r\to0$ (\Cref{rot:lem:r}), so no condition on $J$ remains for large $k$.
\end{proof}

\subsubsection*{The counterterm bridge}

\begin{proposition}[the counterterm and Wick-constant conventions]\label{rot:prop:CT}
\cite{K1} and \cite{Ba1,Ba2,Ba3} are proved for the model in which the mass counterterm is defined by \cite[(3.33)]{K1} with the sum over the torus and torus propagators in the graphs, the Wick subtraction by \cite[(2.8)]{K1} with the torus propagator, the vacuum counterterm $E_1$ by \cite[(3.35)]{K1} as vacuum graphs on the torus, and the source vertex by \cite[(3.21)]{K1} with torus propagators. \Cref{rot:def:tori} uses the free-boundary conventions of \cite[p.~330]{K2}. For a torus $T$ with shortest period $\ell(T)$ ($=50MJ$ for $T_1$, $T_2$) and spacing $s$, write $\dm_T(s)$, $E_{1,T}(s)$, $c_T(s)$ for the torus-convention quantities and
\[
\Delta_T(s)\defeq\dm(s)-\dm_T(s),\qquad\Delta E_T(s)\defeq E_1(s)-E_{1,T}(s),\qquad\Delta c_T(s)\defeq\sW_s-c_T(s)=-\sum_{\gamma\in P\setminus0}C_s(\gamma).
\]
\begin{enumerate}[label=\textup{(\alph*)},leftmargin=*]
\item With $N=2$,
$\Zun_{\Lambda,T}(G,h)=e^{-\Delta E_T|T|-N\Delta c_T\sS_\Lambda(h)}\int\prod dA\prod d\phi\ e^{-S^{\rm tor}_\Lambda-\frac12\Delta_T\sum_xs^3|\phi(x)|^2+F_\Lambda(G)+W^{\rm tor}_\Lambda(h)}$:
the free-boundary model is the torus-convention model with the bare mass shifted by $\Delta_T$ and its counterterms frozen at $m^2$, times a $c$-number and a shift of the Wick constant. Both block-spin flows carry these $c$-numbers through unchanged, so the same identity holds between the effective actions at every step.
\item $\Delta_T(s)$, $\Delta c_T(s)$ and $\Delta E_T(s)$ are bounded uniformly in $s\le s_0$, and $|\Delta_T(s)-\Delta_T(0)|\le Cs$. Put $\Delta_*\defeq\sup_{s\le s_0}|\Delta_T(s)|$ and $u(s)\defeq-\Delta_T(s)/2$. The free-boundary effective polynomial of \Cref{rot:def:eff} equals King's torus-convention polynomial at the constant-shifted source $h+u(s)\mathbf 1$, minus a $c$-number, up to graphs that carry at least one $u$-vertex and have order $\ge14$ in the couplings, whose sum is bounded by $C\delta_r^\sigma|T|$ on $\supp\chi_r$; and \Cref{rot:prop:Kforward}, the only consumer of \Cref{rot:prop:H3,rot:prop:H4,rot:prop:H5}, is $c$-number blind: multiplying each lattice's $\Zun$ and small-field integral by an arbitrary positive constant leaves every step and the conclusion unchanged. Consequently the extensive bounds \Cref{rot:prop:H3,rot:prop:H4,rot:prop:H5}, proved for the torus-convention model at the sources $(zG,h+u(s)\mathbf 1)$, hold for the model of \Cref{rot:def:tori} as consumed; for \Cref{rot:prop:H3,rot:prop:H4} the constant-source shift needs King's and Ba{\l}aban's bounds with constants uniform over the family $\{h+u\mathbf 1:|u|\le\Delta_*/2\}$, which is the uniformity of \cite[Thm.~3.1]{K1} in the scalar mass over a compact mass set, King's own clause ``$C$ depends on $g$, $h$ and $|T|$''; for \Cref{rot:prop:H5} the shift is removable by Jensen's inequality.
\item The ratios are convention-free up to a shift of the source by a constant: for every $\Lambda$, $G$, $h$,
\[
\Zun^{\rm free}_{\Lambda,T}(G,h)=\exp\big[-(\Delta E_T(s)+\Delta_T(s)c_T(s))|T|-2\Delta c_T(s)\sS_\Lambda(h)\big]\,\Zun^{\rm tor}_{\Lambda,T}\big(G,h-\tfrac12\Delta_T(s)\mathbf 1\big),
\]
with a prefactor independent of $G$; so the walled curvature ratio $\langle e^{tF(G)}\rangle_{\fc,\Lambda}$ is the same number in both conventions at the shifted wall $h_0-\Delta_T(\eps_k)/2$.
\end{enumerate}
\end{proposition}
\begin{proof}
(a) and (c) are exact algebra: the two actions differ by the counterterms $\frac12\Delta_T\sum_xs^3|\phi|^2+\Delta E_T|T|$ and the Wick squares by the $c$-number $-2\Delta c_T\sS_\Lambda(h)$; a field-independent constant passes through every block-spin step (\Cref{lem:a06:CF}). (b) For $\Delta c_T=-\sum_{\gamma\ne0}C_s(\gamma)$ the bound is the exponential decay of the massive lattice propagator, $|C_s(\gamma)|\le Ce^{-c|\gamma|}$ uniformly in $s$, summed over the non-zero periods. For $\Delta_T$, in each of the finitely many graphs of \cite[(3.33)]{K1} (orders $a+2b\le4$) the torus-minus-free difference replaces one propagator by its image sum, which is smooth and removes the coincident-point singularity, so that the remaining singular factors are summable in three dimensions uniformly in $s$; this gives the boundedness, and the rate $|\Delta_T(s)-\Delta_T(0)|\le Cs$ follows at one loop from the explicit integrals and at the two-loop orders $(0,2)$, $(2,1)$, $(4,0)$ from a graph-by-graph computation of the thirteen two-loop topologies at order $(4,0)$, the degree-zero loops evaluated inner-first and symmetrised, using that the $e^2\lambda$ vertex-type loop vanishes by charge conjugation and the exchange composed with the photon self-energy vanishes by lattice transversality. For $\Delta E_T$ the same image-sum argument applies to the vacuum graphs of \cite[(3.35)]{K1}, with one vertex fixed. The statement on the effective polynomials: a constant shift of the scalar source enters King's expansion as a vertex $u\sum\eta^3|\phi|^2$ with a bounded coefficient, and the graphs with a $u$-vertex and order $\le13$ in the couplings are exactly the terms by which the two polynomials differ beyond the $c$-number, so the remainder is a finite set of graphs of order $\ge14$ with running power $\ge7>3$, bounded by $C\delta_r^\sigma|T|$. The $c$-number blindness of \Cref{rot:prop:Kforward} is read off its proof.
\end{proof}

% r03_fill_K6: proof of (K_0), part IX: the existence of the axis-lattice limits, the forward bound, and the assembly.

\subsubsection*{The existence of the walled axis-lattice limits}

\begin{theorem}[the walled torus limit on the axis lattice]\label{rot:thm:H6W}
For every walled configuration $\fc=(J,\psi,n)$ of the class of \Cref{rot:thm:CT} and every real smooth $P_J$-periodic $2$-form $G$ the limit
\[
\Wv_\fc(tG)\defeq\lim_{k\to\infty}\frac{\Zun_{\Xk,T_1}(tG,h_0)}{\Zun_{\Xk,T_1}(0,h_0)}=\lim_{k\to\infty}\langle e^{tF(G)}\rangle_{\fc,\Xk}
\]
exists for every real $t$; it is the same number in the free-boundary convention of \Cref{rot:def:tori} and in King's torus convention at the shifted constant source $h_0-\Delta_T(\eps_k)/2$ (\Cref{rot:prop:CT}(c)), and it equals the limit of the torus-convention ratios $\Zun^{\rm tor}_{\Xk,T_1}(tG,h_\infty)/\Zun^{\rm tor}_{\Xk,T_1}(0,h_\infty)$ at the frozen wall $h_\infty\defeq h_0-\Delta_T(0)/2$. The convergence is locally uniform in $t\in\C$, the limit is entire, and all moments of $F(G)$ converge.
\end{theorem}

This is King's torus theorem \cite[Thm.~2.1]{K1} with the wall as his scalar source, read on the normalised ratio, in the only form in which it is true at a wall: the convergence of King's effective action \cite[Thm.~3.4]{K1} modulo a field-independent constant of its scalar-source sector, which cancels in the ratio, with the drift-free stability bound \cite[(2.23)]{K1} replaced by a polynomial envelope. The proof is given in the following three groups of statements; it compares the walled ratio, which is King's torus-convention ratio at a wall shifted by a bounded constant that moves with $k$ by $O(\eps_k)$, with the torus-convention ratio at a frozen wall on the same lattice, proves the frozen family Cauchy by telescoping one step at a time, and bounds the effect of the constant shift of the source by a per-site bound with its volume factor against the rate of \Cref{rot:prop:CT}(b).

% r03_h6w_a: proof of Theorem rot:thm:H6W, part 1: conventions and the bridge, the a priori bounds, King's flow on the axis lattice, the two one-lattice stability inputs, and the reduction to the small-field ratios.
\providecommand{\Ztor}{\ensuremath{\Zun^{\rm tor}}}
\providecommand{\hinf}{\ensuremath{h_\infty}}

The proof of \Cref{rot:thm:H6W} occupies the rest of this paragraph, up to the subsubsection on the forward bound. It is a one-lattice statement: only the axis lattice $\Xk=\Lk$ on the torus $T\defeq T_1=\DJ$ occurs, with King's standard block-spin flow at block factor $5$ and King's own corner-rooted blocks (the centred blocks of \Cref{rot:conv:C} play no role here), and neither the rotated lattice, nor the hybrid flow, nor \Cref{rot:thm:sandwich} is used. Throughout, $\fc=(J,\psi,n)$ is in the class of \Cref{rot:thm:CT}, $h_0\defeq-n\psi$, $G$ is a real smooth $P_J$-periodic $2$-form sampled in the frame of $\Xk$, $t\in\R$, and $r=r(k)$, $\delta_r=5^r\eps_k=5^{-\lfloor\vr k\rfloor}$ are the depth and the coarse spacing of \Cref{rot:lem:r}. The argument compares two ratios on the same lattice: the walled curvature ratio, which by the bridge of \Cref{rot:prop:CT} is King's torus-convention ratio at a wall that moves with $k$ by a constant of order $\eps_k$, and the torus-convention ratio at a frozen wall; the frozen family is Cauchy by King's renormalisation theorem read modulo the field-independent constant of its scalar-source sector, and the moving family stays close to the frozen one because a constant shift of the scalar source changes the effective action by a bounded amount per site, which the rate $|u_k|\le C\eps_k$ of \Cref{rot:prop:CT}(b) beats.

\subsubsection*{Conventions, the bridge and the a priori bounds}

\begin{lemma}[the moving and the frozen wall]\label{rot:lem:H6conv}
Write $\Ztor_k(t;h)\defeq\Ztor_{\Xk,T}(tG,h)$ for the torus-convention functional of \Cref{rot:prop:CT}(c) and
\[
h_k\defeq h_0-\tfrac12\Delta_T(\eps_k)\mathbf 1,\qquad\hinf\defeq h_0-\tfrac12\Delta_T(0)\mathbf 1,\qquad u_k\defeq h_k-\hinf=\tfrac12\big(\Delta_T(0)-\Delta_T(\eps_k)\big),
\]
a constant with $|u_k|\le C\eps_k$. Then for every $k$ and every real $t$
\[
r_k(t)\defeq\langle e^{tF(G)}\rangle_{\fc,\Xk}=\frac{\Zun_{\Xk,T}(tG,h_0)}{\Zun_{\Xk,T}(0,h_0)}=\frac{\Ztor_k(t;h_k)}{\Ztor_k(0;h_k)},\qquad s_k(t)\defeq\frac{\Ztor_k(t;\hinf)}{\Ztor_k(0;\hinf)} .
\]
The family $\{h_k\}_{k}\cup\{\hinf\}$ is bounded in every $C^N(T)$ norm, $\|h\|_{C^N}\le\|h_0\|_{C^N}+\Delta_*/2$; $\Xk/P_J$ is King's lattice on his torus \cite[(2.21)]{K1}, so the right side of the identity for $r_k$ is literally the object of \cite[Thms.~3.1, 3.4]{K1} with the scalar source $h_k$, and King's constants, which depend on the sources only through bounds on finitely many derivatives (\cite[pp.~664--665, 677]{K1}: ``the sources $g,h$ are bounded by $C$'', ``their derivatives are uniformly bounded''), are uniform over the family. The physical mass $m^2$ is the same in both conventions and is never changed.
\end{lemma}
\begin{proof}
The first identity is the walled-ratio display of \Cref{rot:def:Zun} at $g=0$, with $-w=h_0$ (see also \Cref{rot:lem:walledalg}(c)): the wall is the scalar source $h_0$ up to a constant that cancels in the ratio. The second is \Cref{rot:prop:CT}(c) at $\Lambda=\Xk$, whose prefactor does not depend on $G$ and cancels in the ratio. The existence of $\Delta_T(0)=\lim_{s\to0}\Delta_T(s)$ and $|u_k|\le C\eps_k$ are the rate clause of \Cref{rot:prop:CT}(b), and $|\Delta_T(\eps_k)|,|\Delta_T(0)|\le\Delta_*$ give the $C^N$ bound, the shifts being constants. No sign condition on the scalar source enters King's bounds: $\lambda>0$ makes every single-site measure finite for every real source.
\end{proof}

\begin{lemma}[a priori bounds on the ratios]\label{rot:lem:H6A}
Let $C_G\defeq\sup_kq_{\Xk}(G)/2<\infty$ (\Cref{rot:lem:q}). For every $k$, every real smooth $P_J$-periodic $h$ and every real $t$,
\[
1\le\frac{\Ztor_k(t;h)}{\Ztor_k(0;h)}\le e^{t^2q_{\Xk}(G)/2}\le e^{C_Gt^2},\qquad\Big|\frac{\Ztor_k(z;h)}{\Ztor_k(0;h)}\Big|\le\frac{\Ztor_k(\Re z;h)}{\Ztor_k(0;h)}\le e^{C_G(\Re z)^2}\quad(z\in\C),
\]
and the same in the free-boundary convention. Consequently, if a sequence of such ratios converges for every real $t$, it converges locally uniformly on $\C$ to an entire function, and all moments of $F(G)$ converge to the Taylor coefficients of the limit.
\end{lemma}
\begin{proof}
The torus-convention walled state is a BFS measure of King type in the sense of the paragraph after \Cref{rot:def:Zun}: the torus counterterm and Wick constant only change the Gaussian factor of the single-site measures and a constant. By \Cref{lem:s02:GD} with $F_{\Xk}(G)=\ip A{\ell_{\Xk}}$ (\Cref{rot:lem:CM}), $\langle e^{tF(G)}\rangle\le e^{t^2q_{\Xk}(G)/2}$; by Jensen's inequality and $\langle F(G)\rangle=0$ under the charge conjugation of \Cref{rot:def:tori}, $\langle e^{tF(G)}\rangle\ge1$. For complex $z$, $|e^{zF}|=e^{(\Re z)F}$. The functions $z\mapsto\Ztor_k(z;h)/\Ztor_k(0;h)$ are entire (the integral converges absolutely for every complex $z$ by the same Gaussian bound) and are bounded on every disc $|z|\le\mathsf R$ by $e^{C_G\mathsf R^2}$, uniformly in $k$; convergence on $\R$ then gives locally uniform convergence on $\C$ by Vitali's theorem \cite{Hor90}, and Cauchy's estimates give the moments.
\end{proof}

\subsubsection*{King's flow on the axis lattice}

For a lattice $X_{k'}$ and a depth $r'$ with $5^{r'}\eps_{k'}<1$, King's standard flow \cite[\S3]{K1} on $X_{k'}/P_J$ after $r'$ steps produces the unit lattice $T^{(r')}$ of physical spacing $\delta\defeq5^{r'}\eps_{k'}$, the small-field indicator $\chi_{r'}$ \cite[(3.2)]{K1}, on whose support $|\phi_{r'}|\le p(\delta)(\lambda\delta)^{-1/4}$, and the effective action $S^{(r')}_{k'}(tG,h)$ of \Cref{rot:def:eff}, here in the torus-convention reading for which \cite{K1} is written: its Gaussian parts (quadratic forms, $\ln N_{r'}$, $\ln Z_{r'}(A^{(r')})$) are defined at the mass $m^2$ without sources \cite[(2.13)--(2.16)]{K1}, and the sources enter only through the curvature vertex \cite[(3.20)]{K1} and the $h$-vertex \cite[(3.21)]{K1} of the perturbative part $E'_{r'}$, expanded to King's order $\bar n=13$ with a remainder. Put
\[
I^{(r')}_{k'}(t;h)\defeq\int\chi_{r'}\,e^{-S^{(r')}_{k'}(tG,h)}\,dA_{r'}\,d\phi_{r'},\qquad
\rho^{(r')}_{k'}(t;h)\defeq\frac{I^{(r')}_{k'}(t;h)}{I^{(r')}_{k'}(0;h)},
\]
and $I_k\defeq I^{(r(k))}_k$, $\rho_k\defeq\rho^{(r(k))}_k$ at the lattice's own depth. Two facts are used throughout: the unit lattice has $|T^{(r')}|=|T|\delta^{-3}$ sites (\cite[(2.15)]{K1} in unit-lattice counting); and $5^r\eps_k=5^{r+1}\eps_{k+1}=\delta_r$, so that $\Xk$ at depth $r$ and $X_{k+1}$ at depth $r+1$ live on the same unit lattice $T^{(r)}$ with the same $\chi_r$ and the same field bounds (\cite[p.~656]{K1}: ``since $L^k\eps=L^{k+n}\eps'$, the fields $A_k,\phi_k$ have the same bounds in each integral'').

\subsubsection*{The two stability inputs at the walls}

\begin{proposition}[King's upper bound on the axis lattice, at a wall]\label{rot:prop:H6S1}
For $h$ in the family of \Cref{rot:lem:H6conv}, every $k$ and every admissible depth $r'$, with $\delta=5^{r'}\eps_k$,
\[
e^{-C\delta^\sigma|T|}\,I^{(r')}_k(t;h)\le\Ztor_k(t;h)\le e^{C\delta^\sigma|T|}\,I^{(r')}_k(t;h)+B_k,\qquad B_k\defeq e^{-p(\delta)^2+C|T|}\eps_k^{-q},
\]
with $\sigma>0$, $C=C(t,G,h_0,|T|)$ locally uniform in $t$, and $q=q(t,G,h_0)<\infty$, all independent of $k$, of $r'$ and of $h$ in the family.
\end{proposition}
\begin{proof}
This is \cite[Thm.~3.1, (3.3)--(3.4)]{K1} for King's torus with the sources $(tG,h)$, with the additive large-field term in the form in which King uses it (\cite[(3.11)]{K1}, \cite[(4.7)]{K2}), at block factor $5$ on the torus $\DJ$ aligned with King's blocks: \Cref{thm:a06:BALA}(iii) with \Cref{hyp:k01:BALint}, in the torus convention for which it is printed, so that no bridge is crossed. Two readings are made, the same two as in \Cref{rot:prop:H3}: the uniformity over the bounded source family (\Cref{rot:lem:H6conv}); and the envelope. The factor $e^{C|T|}$ of the large-field term is the ultraviolet stability bound \cite[(2.23)]{K1}, which fails as printed when $\int h\ne0$ (\Cref{rot:rem:log}); what the inductive proof supplies at a wall is a bounded $h$-linear $c$-number per step, $e^{C|T|+C'(1+k)|\int h|}=e^{C|T|}\eps_k^{-q}$. The loss is harmless against $e^{-p(\delta_r)^2}$, since $p(\delta_r)^2$ grows at least quadratically in $k$ (\Cref{rot:lem:r}) while $\ln\eps_k^{-q}=qk\ln5$ is linear. The sources enter the proof through the vertices (3.20)--(3.21) only, with constants depending on $\sup|G|$, $\sup|\nabla G|$, $\sup|h|$, $\sup|\nabla h|$ (\cite[pp.~664--665, 677]{K1}).
\end{proof}

\begin{proposition}[the polynomial lower bound at the wall, torus convention]\label{rot:prop:H6S2}
For $h$ in the family of \Cref{rot:lem:H6conv} and every $k$: $\Ztor_k(0;h)\ge e^{-C_1|T|}\eps_k^{q_1}$, with $C_1$, $q_1$ independent of $k$ and of $h$ in the family.
\end{proposition}
\begin{proof}
This is \Cref{rot:thm:FS} on $\Lambda=\Xk$, $T=T_1$, for the bounded family ($\|h\|_\infty\le H\defeq n\|\psi\|_\infty+\Delta_*/2$; the sign of $h$ is not used there), read in the torus convention: the proof of \Cref{rot:thm:FS} moves the wall into the bare scalar coefficient, compares with a higher King point by Jensen's inequality and the monotonicity of $\langle\sum_x\eps^3|\phi|^2\rangle$ in that coefficient, and uses the source-free stability bounds \cite[(1.14)]{Ba1}, \cite[(1.3)]{Ba2} at the higher points together with the integrated counterterm identity and the perturbative bounds on the one-point Wick function. Each of these steps is available for the torus-convention model with the same constants: \cite[(1.14)]{Ba1} and \cite[(1.3)]{Ba2} are printed for it (\Cref{thm:a06:BALA}(iii)); its bare coefficient, Wick constant and vacuum constant differ from the free-boundary ones by $\Delta_T(s)$, $\Delta c_T(s)$, $\Delta E_T(s)$, which are bounded uniformly in $s$ (\Cref{rot:prop:CT}(b)) and shift the King points by bounded amounts only; and the perturbative coefficients of the torus one-point function differ from the infinite-lattice ones by image sums, exponentially small in $\ell(T)$. The bound reads $\ln\Ztor_k(0;h)\ge-C|T|-C(1+k)|T|-qk\ln5$, which is the claim with $q_1=q+C|T|/\ln5$.

At the mass-increasing wall the same conclusion also follows from Jensen's inequality as in \Cref{rot:prop:J}, in the torus convention: $\Ztor_k(0;h)\ge\Ztor_k(0;0)\exp[\sS_{\Xk}(h)\,w^{\rm tor}_k]$ with $w^{\rm tor}_k\defeq\langle|\phi(x)|^2\rangle^{\rm tor}_{0,0}-2c_T(\eps_k)$, independent of $x$ by translation invariance ($P_J\subset\Xk$), where $\Ztor_k(0;0)\ge e^{-C|T|}$ is \cite[(1.14)]{Ba1} and $\sS_{\Xk}(h)=-n\sS_{\Xk}(\psi)-\frac12\Delta_T|T|<0$ by the wall condition of \Cref{rot:thm:CT}; together with an upper bound $w^{\rm tor}_k\le C(1+k)$ on the source-free one-point Wick function, which is not established here and not needed, this gives the claim with $q_1=C|\sS_{\Xk}(h)|/\ln5$. This is the only place in the proof where the sign of the wall would enter; the route through \Cref{rot:thm:FS} does not use it.
\end{proof}

\subsubsection*{The ratios are small-field ratios up to $o(1)$}

\begin{lemma}[the large-field term is negligible]\label{rot:lem:H6stepA}
Let $h$ be in the family of \Cref{rot:lem:H6conv}, $r'$ an admissible depth for $\Xk$, $\delta=5^{r'}\eps_k$, $x\defeq C\delta^\sigma|T|$ with the constant of \Cref{rot:prop:H6S1}, and $b_k(h)\defeq2B_k/\Ztor_k(0;h)$. Then
\[
b_k(h)\le2e^{-p(\delta)^2+(C+C_1)|T|}\eps_k^{-q-q_1},
\]
and whenever $b_k\le\frac12$ and $2x\le\ln2$:
\begin{equation}\label{rot:eq:H6stepA}
e^{-2x}(1-b_k)\,\rho^{(r')}_k(t;h)\le\frac{\Ztor_k(t;h)}{\Ztor_k(0;h)}\le e^{2x}\rho^{(r')}_k(t;h)+b_k,\qquad
\rho^{(r')}_k(t;h)\le4e^{C_Gt^2},\qquad
\Big|\frac{\Ztor_k(t;h)}{\Ztor_k(0;h)}-\rho^{(r')}_k(t;h)\Big|\le4e^{C_Gt^2}(e^{2x}-1)+(1+4e^{C_Gt^2})\,b_k .
\end{equation}
At the lattice's own depth $r=r(k)$, $x=C\delta_r^\sigma|T|\to0$ and $b_k(h)\to0$ as $k\to\infty$ (\Cref{rot:lem:r}: $p(\delta_r)^2$ grows at least quadratically in $k$, the exponent of $\eps_k^{-q-q_1}$ linearly), so $|r_k(t)-\rho_k(t;h_k)|\to0$ and $|s_k(t)-\rho_k(t;\hinf)|\to0$ for every real $t$, locally uniformly in $t$.
\end{lemma}
\begin{proof}
Write $Z(t)\defeq\Ztor_k(t;h)$, $I(t)\defeq I^{(r')}_k(t;h)$, $\rho\defeq\rho^{(r')}_k(t;h)$, $B\defeq B_k$, $b\defeq b_k(h)$. The bound on $b$ is \Cref{rot:prop:H6S1,rot:prop:H6S2}. By \Cref{rot:prop:H6S1} at $t$ and at $0$: $Z(t)\le e^xI(t)+B$ and $Z(0)\ge e^{-x}I(0)$ give $Z(t)/Z(0)\le e^{2x}\rho+B/Z(0)\le e^{2x}\rho+b$; and $Z(t)\ge e^{-x}I(t)=e^{-x}\rho I(0)$ with $I(0)\ge e^{-x}(Z(0)-B)$ give $Z(t)/Z(0)\ge e^{-2x}\rho(1-B/Z(0))\ge e^{-2x}(1-b)\rho$. Hence, by \Cref{rot:lem:H6A}, $\rho\le e^{2x}(1-b)^{-1}Z(t)/Z(0)\le2e^{2x}e^{C_Gt^2}\le4e^{C_Gt^2}$. For the difference: $Z(t)/Z(0)-\rho\le(e^{2x}-1)\rho+b$ and $\rho-Z(t)/Z(0)\le\rho[(1-e^{-2x})+b]\le\rho[(e^{2x}-1)+b]$; the factor $(1+4e^{C_Gt^2})$ on $b$ covers both sides. The last statement is the case $r'=r(k)$, $h\in\{h_k,\hinf\}$, with the identities of \Cref{rot:lem:H6conv}.
\end{proof}

% r03_h6w_b: proof of Theorem rot:thm:H6W, part 2: graphs and the census of the source subgraphs, the restricted form of King's renormalisation theorem, and the Cauchy estimate at the frozen wall modulo the field-free constant.

\subsubsection*{Graphs, degrees and the census of the source subgraphs}

King's perturbative part $P^{(r)}_k(tG,h)$ is a finite sum of connected renormalised graphs $H$ \cite[(3.29), (3.37)]{K1} with the vertices \cite[(3.18)--(3.26)]{K1}; the number of graphs depends on $\bar n$ only, each graph is multilinear in its source insertions, and its value on $\supp\chi_r$ is $\sum_{\{y\},\{z\}}E^{(r)}(H;\{y\},\{z\})\times(\text{external fields at }\{y\})$, with a kernel that depends on the coarse background $A^{(r)}$ through its lines $G_r(A^{(r)})$ \cite[(3.17)]{K1}, so that a graph without external legs still depends on the background. Degrees are King's \cite[(3.30)--(3.31)]{K1}, as in \Cref{rot:lem:S1}, computed on the point-split graph (the two summation points of a block product as vertices, as in \Cref{rot:cor:OF}), with two conventions. The curvature vertex \cite[(3.20)]{K1} is counted with $D_v=\frac52$, the lattice curl of its photon leg having been summed by parts onto the smooth $G$ (exact on the torus, $|d^*_\eta G|\le C\|\nabla G\|_\infty$); with the curl charged to the vertex, $D_v=\frac32$ and the tree with two curvature vertices would have degree $0$, so the summation by parts is necessary. The $h$-vertex \cite[(3.21)]{K1} has $D_v=2$ (\Cref{rot:lem:S1}). Every connected subgraph containing an $h$-vertex and a curvature vertex therefore has $D\ge2+\frac52-3=\frac32>0$, and the smallest vacuum graph with two curvature vertices and two current vertices has $D=5$.

Each graph value is split at zero coarse fields, $E(H)=E(H)|_0+[E(H)-E(H)|_0]$, where $|_0$ means $(A_r,\phi_r)=(0,0)$, and we put $E_N(H)\defeq E(H)$ if $H$ has an external scalar leg or a curvature vertex and $E_N(H)\defeq E(H)-E(H)|_0$ otherwise: the field- or source-dependent part. The field-free, $G$-free values $E(H)|_0$ of the remaining graphs ($h$-linear or $h$-nonlinear vacuum graphs without curvature vertex, and the $h$-free vacuum energy) are numbers depending on $k$, $r$, $h$ only; nothing is claimed about their convergence, and for the order-$e^2$ $h$-linear family it fails (\Cref{rot:rem:log}: the Wick constant drifts by $c_*\ln5$ per step).

\begin{lemma}[the census of the source subgraphs]\label{rot:lem:H6census}
Expand every graph in the coarse background as in \cite[(3.43)--(3.55)]{K1} (\Cref{rot:lem:CB}(1) for one flow): $E(H)$ becomes a finite sum of graphs $\tilde H$ with lines at zero background, explicit background legs $\tilde A$ at block points, and remainder factors, and $E(H)|_0$ is the core $\tilde H_0$ without background legs. Then:
\begin{enumerate}[label=\textup{(\alph*)},leftmargin=*]
\item Every $\tilde H$ entering $E_N(H)$ is connected and has an external leg (scalar, background or $dG$) or a curvature vertex.
\item For such $\tilde H$, every connected ordered subgraph $\gamma$ (lines with scale indices below $r$) has strictly positive King degree or vanishes identically. In detail: if $\gamma$ contains a block point, $D(\gamma)\ge1$ ($\ge2$ with a source vertex); if $\gamma$ contains an $h$-vertex, $D(\gamma)\ge\frac12$ unless $\gamma$ is a vacuum graph with exactly one $h$-vertex and no curvature vertex (then $\gamma=\tilde H$, excluded by (a)) or $\gamma$ is the $h$-loop through a single current vertex \cite[(3.22)]{K1} with its photon leg external ($D=0$: the current tadpole of \Cref{rot:lem:S1} with the $h$-insertion), which vanishes identically at zero background ($\Tr q=0$); if $\gamma$ contains a curvature vertex and no $h$-vertex, $D(\gamma)\ge\frac12$ unless $\gamma$ is the curvature vertex joined by a photon line to a current tadpole ($D=0$), which vanishes identically at zero background for the same reason; if $\gamma$ is source-free, $D(\gamma)>0$ after Ba{\l}aban's resummed rewriting \cite[(3.5)]{Ba3}, the input of \Cref{rot:cor:OF}.
\item The rewriting of the source-free subdivergences inside a graph that carries source vertices (Ba{\l}aban's Taylor subtraction at the subdivergence with its counterterm partner, \cite[(3.5)]{Ba3}, \cite[(3.83)--(3.85)]{K1}) does not lower the degree of a prefix component that contains a source vertex below the value counted in (b).
\item The only connected subgraphs containing a source vertex with King degree $\le0$ in the raw count are the field-free $h$-linear vacuum graphs of order $\le2$ in the couplings, the family $\cT_h$ of \Cref{rot:lem:S1} (which have no external leg, hence occur only as whole graphs), and the two vanishing exceptions of (b).
\end{enumerate}
Consequently the minimum $\alpha_0$ of the degrees of the non-vanishing ordered subgraphs over the finitely many graphs and orderings is strictly positive.
\end{lemma}
\begin{proof}
(a) The core $\tilde H_0$ is the same in $E(H)$ and $E(H)|_0$ and cancels in $E_N(H)$ when $H$ has no external scalar leg and no curvature vertex; every other $\tilde H$ carries a background leg, an external $dG$ line or a remainder factor (the remainders are bounded separately, \Cref{rot:lem:CB}(2)). (b) The degree arithmetic is that of \Cref{rot:lem:S1}: inserting an $h$-vertex into a scalar line adds $2$; source-free connected subgraphs with an external leg have $D\ge-1$ except the current tadpole ($D=-2$), so a subgraph with an $h$-vertex and an external leg has $D\ge1$ except the $h$-loop through one current vertex ($D=0$); a vacuum subgraph with exactly one $h$-vertex has $D\le0$ only at order $\le2$ in the couplings, the exhaustive enumeration of \Cref{rot:lem:S1}, and it is then the whole graph; with two or more $h$-vertices $D\ge1$. For a curvature vertex, which attaches by a photon leg and contributes $\frac52$, the same enumeration leaves the one exception named. Both exceptions vanish at zero background because the current one-point function of the loop vanishes, $\Tr q=0$, exactly as for the current tadpole in \Cref{rot:lem:S1}. A component containing a block point has degree $\ge1$ by the point-split count of \Cref{rot:cor:OF} (every component containing a block point has Ba{\l}aban's degree $\ge1$, \cite[(2.1)--(2.2)]{Ba3}), and a source vertex adds at least $1$ more. (c) Ba{\l}aban's Taylor subtraction at a source-free subdivergence $S$ raises the degree of $S$ by $\theta_S$ and lowers by $\theta_S$ the degree of the external line of $S$ that receives the derivative; the construction, run on the whole graph with the source vertices counted as vertices, rewrites $S$ only in orderings in which every line of $S$ is finer than every external line of $S$. Let $\gamma$ be a component of an ordered subgraph. If the line that receives the derivative lies in $\gamma$, then every line of $S$, being finer, lies in the same ordered subgraph, and $S$, connected to that line, lies in $\gamma$: the two changes cancel. Otherwise the degree of $\gamma$ can only increase. By induction over the forest of subdivergences, finest first, no component falls below its value in (b); the source vertices, which are never inside a rewritten $S$, keep their contributions $2$ and $\frac52$. (d) is read off the same arithmetic.
\end{proof}

\subsubsection*{The restricted renormalisation theorem}

King's renormalisation theorem \cite[Thm.~3.5]{K1} asserts the per-site bound (3.38) and the convergence (3.39) of the graph kernels for all graphs. As printed it is false on the order-$e^2$ $h$-linear field-free family (\Cref{rot:rem:log}), and the convergence of the field-free, $G$-free constants is not needed here. What is used is the following restricted form.

\begin{proposition}[the restricted renormalisation theorem]\label{rot:prop:H6R}
Let $h$ be in the family of \Cref{rot:lem:H6conv}, $t$ real, and let $\gamma_*,\beta_*,\sigma$ be the exponents of \cite[Thm.~3.4]{K1}.
\begin{enumerate}[label=\textup{(\roman*)},leftmargin=*]
\item \textup{(convergence)} For every graph $H$ and every $k$, with $r=r(k)$, the field- or source-dependent parts obey, after the sum over the positions, for any $\gamma'<\gamma_*$, $\beta'>\frac72$, $\sigma'<\sigma$,
\[
\sum_{\text{positions}}\big|E^{(r+1)}_{N,k+1}(H)-E^{(r)}_{N,k}(H)\big|\times|\text{external fields}|\le C\big[5^{-\gamma'r}\delta_r^{-\beta'}+\delta_r^{\sigma'}\big]|T|\quad\text{on }\supp\chi_r,
\]
the two kernels living on the common unit lattice $T^{(r)}$ of $\Xk$ at depth $r$ and $X_{k+1}$ at depth $r+1$; $C=C(t,G,h_0,|T|)$ is independent of $k$ and locally uniform in $t$.
\item \textup{(per-site bound)} For every graph $H$ with at least one $h$-vertex, other than the degree-$0$ members of the order-$e^2$ $h$-linear field-free family,
\[
|E^{(r)}_k(H;\{y\},\{z\})|\le C(1+k-r)^{p_0}\,\delta_r^{\varrho_H}\,p(\delta_r)^{n_1}\,e^{-\delta\dist(\{y\},\{z\})},\qquad\varrho_H\ge0,
\]
and for the degree-$0$ members of that family the same bound with the extra factor $(1+r)$, which is necessary there. For a graph at non-zero background the bound is read through the expansion of \Cref{rot:lem:H6census}: the core at zero background falls under the second clause if $H$ is in that family and under the first otherwise, and every other term has an external leg and falls under the first. $C$ is linear in $\sup|h|$ for a distinguished $h$-vertex and depends on the other sources through sup norms; $p_0$, $n_1$ depend on $\bar n$ only.
\end{enumerate}
\end{proposition}
\begin{proof}
King's proof of (3.39) \cite[pp.~660--662]{K1} reduces it to the convergence of a single graph with one line replaced by the difference of the two flows' lines (3.56), and that to \cite[Prop.~3.6]{K1}, the one-flow estimate for a renormalised family with one element weakened by a small degree, which King asserts on p.~665 ``by redoing the analysis of Sect.~3.4'' and which is \Cref{rot:cor:OF}. \Cref{rot:cor:OF} applies to a family exactly when every connected ordered subgraph of its non-vanishing graphs has strictly positive degree after Ba{\l}aban's rewriting: a zero-degree ultraviolet subgraph containing the replaced line would give $\sum_j(5^j\eta)^{-\gamma}\sim5^{\gamma r}$ and cancel the gain $5^{-\gamma r}$. By \Cref{rot:lem:H6census} this holds, with the minimum $\alpha_0>0$, for every $\tilde H$ entering $E_N(H)$. The replaced line is a fine line of \cite[(3.17)]{K1} with scale index below $r$; the lines with indices between $r$ and $k$ occur only on the counterterm lines of the $\dm$ graphs, never at a source vertex, so the counterterm-attachment clause \cite[(3.78)]{K1} and the one-vertex attachment are the source-free statements and are unchanged by the presence of the sources. The remaining steps are King's: the kernel bounds \cite[Thm.~3.3, Props.~3.7--3.9]{K1} for the lines and the fine-field reconstructions, printed for the torus convention at block factor $5$ (\Cref{thm:a06:BALA}); the sampling of the sources $x'\to x$ (p.~665), which costs $C\eps_k\|\nabla h\|_\infty+C\eps_k\|\nabla G\|_\infty\le C5^{-\gamma r}$; the tree-level $h$-term $\delta_r^2\sum\eta^3h|\phi^{(r)}|^2$, which has no internal line and is compared directly through the convergence of the fine-field reconstructions $a_rG_rQ_r^*$ \cite[Thm.~3.3, \S3.6]{K1}; the non-polynomial terms of (3.14) by \cite[\S3.6, (3.89)--(3.98), Prop.~3.10]{K1}; the Taylor remainders beyond order $\bar n$, $C\delta_r^\sigma|T|$ per flow (p.~665); and the position sum (3.40)--(3.42), which gives $|T|$ with one free position, the others tied by the exponential tree decay. The factor $(1+k-r)^{p_0}$ of \Cref{rot:cor:OF} is a power of $\log\delta_r^{-1}$ and is absorbed by $\gamma'<\gamma_*$, $\sigma'<\sigma$, and the field norms $(\lambda\delta_r)^{-s/4}p(\delta_r)^{n}$ of the external legs by $\beta'>\frac72$. This proves (i).

(ii) For a graph outside the degree-$0$ members of the order-$e^2$ $h$-linear field-free family, \Cref{rot:lem:H6census} gives strictly positive ordered degrees for every non-vanishing derived graph (the excluded vacuum case of the census is now the family itself), and \Cref{rot:cor:OF} without weakening gives the per-site bound with the factor $(1+k-r)^{p_0}$, King's ``possible addition of some power of $\ln(L^k\eps)^{-1}$'' (p.~666); the sup-norm dependence on the sources is King's (p.~664, 677), the distinguished vertex entering linearly. The two $h$-linear field-free members of order $0$ and $\lambda$ are bounded directly: the subtracted tadpole at zero coarse field is the infrared remainder $-N\sum\eta^3\delta_r^2h\,F_{r}(x,x)$ of the coincident propagator, bounded per site at unit scale, and the figure-eight minus its $\dm_1$ partner leaves a bubble of positive ordered degree times an infrared tadpole line, as in step 3 of the proof of \Cref{rot:cor:OF}. For the degree-$0$ members of the order-$e^2$ family (the exchange and seagull tadpoles with their $\dm_2$ partner) every proper ordered subgraph has $D>0$ (a proper component with the $h$-vertex is not a vacuum graph, the $h$-loop vanishes, and the source-free components are the self-energy and the photon tadpole, renormalised by the $\dm_2$ partner) and only the whole graph has $D=0$: the Taylor identity for the self-energy with its $\dm_2$ partner, Ba{\l}aban's general step at the whole graph with no final renormalisation (there is none to perform), as in King's worked case \cite[(3.81)--(3.87)]{K1}, makes every proper prefix component of degree $\ge\alpha$ and leaves the whole graph at degree $0$, so in the multiscale sum every geometric series keeps its constant except the last, $\sum_{j\le r-1}(5^j\eta)^0=r$: this is the factor $(1+r)$, the zero-degree logarithm of \Cref{rot:cor:OF}. It is necessary: the value of the family at a constant source is $u$ times the drift of \Cref{rot:rem:log}, of size $r$ at fixed $\delta_r$. The infrared counterterm lines are handled as in (i).
\end{proof}

\subsubsection*{The Cauchy estimate at the frozen wall}

\begin{proposition}[King's Cauchy estimate modulo the field-free constant]\label{rot:prop:H6I3}
Let $h=\hinf$ and $t\in\R$, and define the field-free constant $a^{(r)}_k(h)\defeq S^{(r)}_k(0,h)\big|_{A_r=\phi_r=0}$ (the whole value of the effective action at zero fields and zero curvature source) and $d_k\defeq a^{(r)}_k(h)-a^{(r+1)}_{k+1}(h)$, a number independent of $t$ and of $G$. Then on $\supp\chi_r$, with $\gamma',\beta',\sigma'$ as in \Cref{rot:prop:H6R}(i),
\[
\big|S^{(r)}_k(tG,h)-S^{(r+1)}_{k+1}(tG,h)-d_k\big|\le\eta_k|T|,\qquad\eta_k\defeq C\big[5^{-\gamma'r}\delta_r^{-\beta'}+\delta_r^{\sigma'}\big],
\]
$C$ independent of $k$ and locally uniform in $t$.
\end{proposition}
\begin{proof}
Both effective actions live on $T^{(r)}$ with the same $\chi_r$. Their Gaussian parts converge by \cite[Thm.~3.3]{K1} (the step King takes for (3.9) before the perturbative part), and their difference is a quadratic form in the fields plus a constant, the constant being part of $d_k$. In the perturbative parts, split every graph as $E(H)=E(H)|_0+E_N(H)$ when $H$ has no external scalar leg and no curvature vertex: the sum of the values $E(H)|_0$ over these graphs, at $G=0$, is by definition the field-free constant $a^{(r)}_k(h)$ minus the Gaussian constant, and the graphs with a curvature vertex and no external leg are included in $E_N$ and converge by \Cref{rot:prop:H6R}(i). So $S^{(r)}_k(tG,h)-S^{(r+1)}_{k+1}(tG,h)-d_k$ is a sum of differences $E^{(r+1)}_N(H)-E^{(r)}_N(H)$ over the finitely many graphs, each bounded by \Cref{rot:prop:H6R}(i), plus the two Taylor remainders, $\le C\delta_r^\sigma|T|$ each, plus the non-polynomial terms, included in \Cref{rot:prop:H6R}(i). This is King's \cite[(3.9)]{K1} without the field-free $h$-sector constant, which is exactly what is false in the printed form at $\int h\ne0$ (\Cref{rot:rem:log}: the order-$e^2$ $h$-linear constant drifts by $c_*\ln5\int h$ per step, so $d_k\not\to0$); $d_k$ is $t$- and $G$-independent by construction and cancels between numerator and denominator of every ratio below. Only one step, $n=1$, is used.
\end{proof}

% r03_h6w_c: proof of Theorem rot:thm:H6W, part 3: the constant-source shift on one lattice, the rate, the frozen family is Cauchy, and the assembly.

\subsubsection*{The constant-source shift on one lattice}

\begin{lemma}[the constant-source shift of the effective action]\label{rot:lem:H6SD}
Let $h,h'$ be in the family of \Cref{rot:lem:H6conv} with $h-h'=u\mathbf 1$ constant. Then on $\supp\chi_r$, for every real $t$,
\begin{equation}\label{rot:eq:H6SD}
\big|S^{(r)}_k(tG,h)-S^{(r)}_k(tG,h')\big|\le C(1+r)\,|u|\,p(\delta_r)^{n_1}\delta_r^{-3}|T|+2C\delta_r^\sigma|T|,
\end{equation}
with $C=C(t,G,h_0,|T|)$ independent of $k$ and $u$ and locally uniform in $t$, $n_1$ the exponent of \Cref{rot:prop:H6R}(ii) enlarged to absorb $(1+k-r)^{p_0}$, and $\sigma$ the remainder exponent.
\end{lemma}
\begin{proof}
\emph{Where $h$ enters.} By \cite[(3.14)--(3.17)]{K1} the Gaussian parts of $S^{(r)}_k$ are defined at the mass $m^2$ without sources, and the sources sit in $E'_r$ only, through the vertices (3.20), (3.21); the $h$-vertex is linear in $h$ with the explicit coupling $\delta_r^2\le1$. So $S^{(r)}_k(tG,h)-S^{(r)}_k(tG,h')=[P^{(r)}_k(h)-P^{(r)}_k(h')]+[R(h)-R(h')]$, $P$ the Taylor polynomial of order $\bar n=13$ and $R$ the remainder.
\emph{The remainder.} $|R(h)|,|R(h')|\le C\delta_r^\sigma|T|$ on $\supp\chi_r$, uniformly over the source family (\cite[p.~665]{K1}); this is the second term, and no derivative of $R$ in $u$ is taken.
\emph{The polynomial.} $P^{(r)}_k(h)$ is a finite sum of connected renormalised graphs with $m\ge0$ $h$-vertices, each multilinear in its source insertions; by multilinearity $P(h)-P(h')$ is the sum over the graphs with $m\ge1$ of $\sum_{i=1}^m[H$ with the $i$-th $h$-vertex carrying the constant $u$ and the others carrying $h$ or $h']$. The value of each such graph on $\supp\chi_r$ is $\sum_{\{y\},\{z\}}E^{(r)}_k(H;\{y\},\{z\})\times(\text{external fields at }\{y\})$. By \Cref{rot:prop:H6R}(ii), $|E^{(r)}_k(H)|\le C(1+r)(1+k-r)^{p_0}\delta_r^{\varrho_H}p(\delta_r)^{n_1}e^{-\delta\dist}$ with $\varrho_H\ge0$ and $C$ linear in $|u|$ for the distinguished vertex (a constant source has zero derivatives), the factor $(1+r)$ being present only for the degree-$0$ members of the order-$e^2$ $h$-linear field-free family and harmless elsewhere; the external scalar fields obey $|\phi_r(y)|\le p(\delta_r)(\lambda\delta_r)^{-1/4}$ \cite[(3.2)]{K1}, with $\delta_r^{-s/4}$ absorbed by $\varrho_H\ge s/4$ for $s$ external scalar legs, and the external photon legs are bounded likewise by (3.2), so the product of fields and kernel is $\le C|u|(1+r)(1+k-r)^{p_0}p(\delta_r)^{n_1+s}$ per position; this is the whole content of King's remark that the sources ``may be bounded when making a perturbative expansion'' and ``do not affect the leading order positivity of the action'' (p.~677). The graph is connected and its kernel decays exponentially in the tree distance of all its vertices, so the sum over all positions is $\le C|T^{(r)}|=C|T|\delta_r^{-3}$: one free position, the others tied. Collecting the finitely many graphs, with $(1+k-r)^{p_0}\le C(1+\log\delta_r^{-1})^{p_0}$ absorbed by enlarging $n_1$, gives the first term of \eqref{rot:eq:H6SD}.
\end{proof}

The bound \eqref{rot:eq:H6SD} keeps no smallness from the explicit $\delta_r^2$ of the $h$-vertex, from the running couplings or from the positive degrees: it is a per-site bound $O((1+r)p^{n_1})$ times the site count $|T|\delta_r^{-3}$. The true size of the worst term, the tree-level term $u\delta_r^2\sum\eta^3|\phi^{(r)}|^2$, is $|u|p^2\lambda^{-1/2}\delta_r^{-3/2}|T|$ on $\supp\chi_r$: a constant scalar source is a bare-mass shift, and its effect is a running-mass term of the same order as the mass term already present, which does not tend to zero. The smallness comes from the rate below, not from the effective action.

\begin{lemma}[the rate and the choice of $\vr$]\label{rot:lem:H6rate}
With $u_k=\frac12(\Delta_T(0)-\Delta_T(\eps_k))$ of \Cref{rot:lem:H6conv}, $|u_k|\le C\eps_k=C5^{-k}$ (\Cref{rot:prop:CT}(b)), and
\[
\eta'_k\defeq C(1+r)|u_k|p(\delta_r)^{n_1}\delta_r^{-3}+2C\delta_r^\sigma\le C(1+k)\,5^{-k+3\lfloor\vr k\rfloor}\big(b_0(1+k\ln5)\big)^{\mathrm pn_1}+2C\,5^{-\sigma\lfloor\vr k\rfloor}\longrightarrow0,
\]
since $\vr<\frac19<\frac13$ (\Cref{rot:lem:r}): the polynomial factor $(1+k)^{1+\mathrm pn_1}$ is beaten by the geometric $5^{-(1-3\vr)k}$. Hence, on $\supp\chi_r$, $|S^{(r)}_k(tG,h_k)-S^{(r)}_k(tG,\hinf)|\le\eta'_k|T|$, and, exponentiating under the common $\chi_r$ (same lattice, same depth, same fields),
\[
I_k(t;h_k)=e^{\theta_k(t)|T|}I_k(t;\hinf),\quad|\theta_k(t)|\le\eta'_k,\qquad\text{so}\qquad
\frac{\rho_k(t;h_k)}{\rho_k(t;\hinf)}\in[e^{-2\eta'_k|T|},e^{2\eta'_k|T|}],\qquad
|\rho_k(t;h_k)-\rho_k(t;\hinf)|\le16e^{C_Gt^2}\eta'_k|T|\to0
\]
for $k$ so large that $2\eta'_k|T|\le1$, locally uniformly in $t$.
\end{lemma}
\begin{proof}
The first display is \eqref{rot:eq:H6SD} with $u=u_k$, $\delta_r=5^{-\lfloor\vr k\rfloor}$, $1+r\le1+k$ and $p(\delta_r)\le b_0(1+k\ln5)^{\mathrm p}$; both terms tend to $0$ exactly when $0<\vr<\frac13$, and every condition of \Cref{rot:lem:r} on $\vr$ is an upper bound, so its window $\vr<\frac19$ needs no change. The second: $I_k(t;h_k)=\int\chi_re^{-S(tG,\hinf)}e^{-[S(tG,h_k)-S(tG,\hinf)]}$ lies in $[e^{-\eta'_k|T|},e^{\eta'_k|T|}]I_k(t;\hinf)$, and the same at $t=0$; divide, and use $\rho_k(t;\hinf)\le4e^{C_Gt^2}$ (\Cref{rot:lem:H6stepA}) with $e^{x}-1\le2x$ for $0\le x\le1$.
\end{proof}

\subsubsection*{The frozen family is Cauchy}

\begin{lemma}[one step at the frozen wall]\label{rot:lem:H6B}
Fix $t\in\R$ and write $r=r(k)$. For $k\ge k_0$ (so that $\delta_r<1$, $b_k,\tilde b_{k+1}\le\frac12$, $2C\delta_r^\sigma|T|\le\ln2$ and $2\eta_k|T|\le1$), with $\tilde\rho_{k+1}(t)\defeq\rho^{(r+1)}_{k+1}(t;\hinf)$ the small-field ratio of $X_{k+1}$ at depth $r+1$ (which is $r(k+1)$ or $r(k+1)+1$),
\[
\epsilon(k)\defeq|s_k(t)-s_{k+1}(t)|\le16e^{C_Gt^2}|T|\big(2C\delta_r^\sigma+\eta_k\big)+(1+4e^{C_Gt^2})\big(b_k+\tilde b_{k+1}\big),
\]
where $b_k\le2e^{(C+C_1)|T|}e^{-p(\delta_r)^2}5^{(q+q_1)k}$ and $\tilde b_{k+1}\le2e^{(C+C_1)|T|}e^{-p(\delta_r)^2}5^{(q+q_1)(k+1)}$. Consequently $\sum_k\epsilon(k)<\infty$, $(s_k(t))_k$ is Cauchy, and $W^{\rm tor}(t)\defeq\lim_ks_k(t)$ exists, locally uniformly in $t$.
\end{lemma}
\begin{proof}
\emph{(B1) The Cauchy estimate at one step.} By \Cref{rot:prop:H6I3}, on $\supp\chi_r$: $S^{(r)}_k(tG,\hinf)=S^{(r+1)}_{k+1}(tG,\hinf)+d_k+R_k(A_r,\phi_r;t)$ with $|R_k|\le\eta_k|T|$ and $d_k$ independent of the fields and of $t$. Exponentiate and integrate against the common $\chi_r$: $I^{(r)}_k(t;\hinf)=e^{-d_k}I^{(r+1)}_{k+1}(t;\hinf)e^{\theta|T|}$ with $|\theta|\le\eta_k$, and the same at $t=0$ with the same $d_k$. Dividing,
\[
\rho_k(t;\hinf)=\tilde\rho_{k+1}(t)\,e^{\theta'|T|},\quad|\theta'|\le2\eta_k,\qquad\text{so}\qquad|\rho_k(t;\hinf)-\tilde\rho_{k+1}(t)|\le\tilde\rho_{k+1}(t)\big(e^{2\eta_k|T|}-1\big).
\]
\emph{(B2) The large-field terms at the two depths.} \Cref{rot:lem:H6stepA} uses only \Cref{rot:prop:H6S1} at the given lattice and depth and \Cref{rot:prop:H6S2} for that lattice, which is depth-free; both hold for $\Xk$ at depth $r$ and for $X_{k+1}$ at depth $r+1$ ($\delta_{r+1}(X_{k+1})=\delta_r<1$), with the stated bounds on $b_k$ and $\tilde b_{k+1}$. Hence $\rho_k(t;\hinf),\tilde\rho_{k+1}(t)\le4e^{C_Gt^2}$ and, by \eqref{rot:eq:H6stepA} at both lattices,
\[
|s_k(t)-\rho_k(t;\hinf)|\le4e^{C_Gt^2}(e^{2C\delta_r^\sigma|T|}-1)+(1+4e^{C_Gt^2})b_k,\qquad
|s_{k+1}(t)-\tilde\rho_{k+1}(t)|\le4e^{C_Gt^2}(e^{2C\delta_r^\sigma|T|}-1)+(1+4e^{C_Gt^2})\tilde b_{k+1}.
\]
\emph{(B3) One step.} By (B1), (B2) and the triangle inequality,
\[
\epsilon(k)\le4e^{C_Gt^2}\big[2(e^{2C\delta_r^\sigma|T|}-1)+(e^{2\eta_k|T|}-1)\big]+(1+4e^{C_Gt^2})(b_k+\tilde b_{k+1}),
\]
and $e^x-1\le2x$ for $0\le x\le1$ gives the claim.
\emph{(B4) Summability.} With $\delta_r=5^{-\lfloor\vr k\rfloor}\le5\cdot5^{-\vr k}$ and $r=k-\lfloor\vr k\rfloor\ge(1-\vr)k$: $\delta_r^{\sigma}\le5^\sigma5^{-\sigma\vr k}$; $5^{-\gamma'r}\delta_r^{-\beta'}\le5^{\beta'}5^{-[\gamma'-(\gamma'+\beta')\vr]k}$, geometric once $\vr<\gamma'/(\gamma'+\beta')$, which holds for $\gamma'<\gamma_*$ close to $\gamma_*$ and $\beta'\in(\frac72,5)$, because the window $\vr<\gamma/(\gamma+\beta)$ of \Cref{rot:lem:r}, with $\gamma\le\gamma_*$ and $\beta\ge5$, lies strictly inside $\vr<\gamma_*/(\gamma_*+\beta')$; and $b_k+\tilde b_{k+1}\le4e^{(C+C_1)|T|}5^{q+q_1}\exp[-b_0^2(1+\lfloor\vr k\rfloor\ln5)^{2\mathrm p}+(q+q_1)k\ln5]$, a quadratic-in-$k$ negative exponent against a linear positive one ($\mathrm p\ge1$, \Cref{rot:lem:r}), hence $\le Ce^{-ck^2+C'k}$. So $\sum_{k\ge k_0}\epsilon(k)<\infty$, and for $k_0\le k<k'$, $|s_k(t)-s_{k'}(t)|\le\sum_{j=k}^{k'-1}\epsilon(j)\le\sum_{j\ge k}\epsilon(j)\to0$. All constants are uniform for $t$ in compacts. No uniformity in the number of steps is used or claimed: a sequence with summable consecutive differences is Cauchy.
\end{proof}

This lemma alone is King's torus theorem \cite[Thm.~2.1]{K1} with the frozen wall $\hinf$ as his scalar source, read on the normalised ratio: the differences from King's argument \cite[(3.10)--(3.13)]{K1} are that the drift-free stability bound $|\ln Z|\le C|T|$ \cite[(2.23)]{K1}, false at a wall, is replaced by the polynomial envelope of \Cref{rot:prop:H6S1,rot:prop:H6S2}, that the field-free $h$-sector constant, which does not converge, is removed by the ratio, and that the uniformity in the number of steps is replaced by summation over single steps.

\subsubsection*{Assembly}

\begin{proof}[Proof of \Cref{rot:thm:H6W}]
Fix $t\in\R$. By \Cref{rot:lem:H6stepA} at the two walls and \Cref{rot:lem:H6rate},
\[
|r_k(t)-s_k(t)|\le|r_k(t)-\rho_k(t;h_k)|+|\rho_k(t;h_k)-\rho_k(t;\hinf)|+|\rho_k(t;\hinf)-s_k(t)|\longrightarrow0 ,
\]
and by \Cref{rot:lem:H6B} $s_k(t)\to W^{\rm tor}(t)$. Hence $\lim_kr_k(t)=\lim_k\langle e^{tF(G)}\rangle_{\fc,\Xk}$ exists and equals $W^{\rm tor}(t)$, for every real $t$; by \Cref{rot:lem:H6conv} this is the limit of the torus-convention ratios at the moving wall $h_k=h_0-\Delta_T(\eps_k)/2$, which coincide with $r_k(t)$ term by term, and of those at the frozen wall $\hinf=h_0-\Delta_T(0)/2$. By \Cref{rot:lem:H6A} the convergence is locally uniform on $\C$, the limit $\Wv_\fc(zG)$ is entire, and all moments of $F(G)$ converge. The limits at $h_0$ itself in the two conventions are not claimed equal, and nothing below uses that.

\emph{What the proof does not use.} No value of the constants $d_k$; no Wick scheme; no convergence of $\Ztor_k(t;h)$ or $\Ztor_k(0;h)$ separately (both drift like $\eps_k^{\pm q}$); no comparison of two lattices, no rotated lattice, no skew torus, no hybrid flow and no centred blocks; no change of mass; the bridge \Cref{rot:prop:CT} only through the exact identity (c), the boundedness of $\Delta_T$ and the rate $|\Delta_T(s)-\Delta_T(0)|\le Cs$ of (b); and the sign of the wall nowhere, when the lower bound is taken from \Cref{rot:thm:FS}. The statement is at fixed $\fc$, $G$ and $t$, with $\eps_k\to0$ along $\eps_k=5^{-k}$; nothing depends on the later infimum over configurations or on the wall-height index of the shell construction.
\end{proof}


\subsubsection*{The forward bound and the assembly}

\begin{proposition}[the forward bound]\label{rot:prop:Kforward}
For $T\in\{T_1,T_2\}$ and every real source $s=(zG,h)$ with $h$ a wall $-n\psi$ or its transport,
\[
\Zun_{\Yk,T}(s)\le(1+\delta_k)\,e^{-c_k\sS_{\Xk}(h)}\,\Zun_{\Xk,T}(s),\qquad\delta_k=\delta_k(T,s)\to0,
\]
with the constant $c_k$ of \Cref{rot:lem:C}, independent of $T$ and of the source, and $|c_k|\le C\eps_k^{-q}$. No bound on $c_k$ is used.
\end{proposition}
\begin{proof}
Write $I_Y\defeq\int\chi_re^{-S^{(r)}}$ and $I_X\defeq\int\chi_re^{-\hs^{(r)}}$. By \Cref{rot:prop:H3}, $\Zun_Y\le e^{C\delta_r^\sigma|T|}I_Y+B$ with $B\defeq e^{-p(\delta_r)^2+C|T|}\eps^{-q}$. By \Cref{rot:prop:H5} (at the wall, \Cref{rot:thm:FS}), $B\le b\,\Zun_Y$ with $b\defeq e^{-p(\delta_r)^2+(C+C_1)|T|}\eps^{-q-q'}$; since $p(\delta_r)^2$ grows at least quadratically in $k$ (\Cref{rot:lem:r}) while $\ln\eps^{-q-q'}=(q+q')k\ln5$ grows linearly, $b\to0$, so $\Zun_Y\le(1-b)^{-1}e^{C\delta_r^\sigma|T|}I_Y$: the large-field term is absorbed on the left, which is why no bound on $c_k$ and no lower bound for the axis lattice is needed. On $\supp\chi_r$, \Cref{rot:thm:Asrc} gives $e^{-S}\le e^{-\hs-c_k\sS_{\Xk}(h)+\vartheta_k|T|}$, so $I_Y\le e^{-c_k\sS_{\Xk}(h)+\vartheta_k|T|}I_X$. By \Cref{rot:prop:H4}, $I_X\le e^{C(1+k)^p\delta_r^\sigma|T|}\Zun_X$. So $\delta_k=(1-b)^{-1}e^{(2C(1+k)^p\delta_r^\sigma+\vartheta_k)|T|}-1$, which tends to $0$ by \Cref{rot:lem:r}. Multiplying $\Zun$ and the small-field integrals of either lattice by a positive constant changes nothing above, which is the $c$-number blindness used in \Cref{rot:prop:CT}(b).
\end{proof}

\begin{proposition}[mass-increasing walls only]\label{rot:prop:Konly}
Let $\cS^-\defeq\{(tG,h_0):t\in\R\}$ with $h_0=-n\psi$, $n>0$, a family not closed under $h\mapsto-h$, and assume the forward bounds of \Cref{rot:prop:Kforward} on $T_1$ for $s\in\cS^-$ and on $T_2$ for their $S$-images, with the common constant $c_k$. Then the conclusions \textup{(a)} and \textup{(c)} of \Cref{rot:thm:sandwich} hold for the sources of $\cS^-$, with $v=v'=0$, $c=c'=c_k$ and $c_k\to0$ (\Cref{rot:prop:CK}); the sources $(tG,+n\psi)$ do not occur.
\end{proposition}
\begin{proof}
Part (a) of \Cref{rot:thm:sandwich} uses only the two forward bounds, one source at a time: the upper bound is the forward bound on $T_1$, the lower bound the forward bound on $T_2$ for $S\cdot s$ pulled back by \Cref{rot:lem:inv}, with $\sS_{\Lambda}(h)=\int h+O(\eps^N)$ (\Cref{rot:lem:lsum}) and $N>q$. Part (b) is the only place where the pair $(G,\pm h)$ enters, and its conclusion $c+c'\to0$ is here supplied directly by $c=c'=c_k\to0$; the interval of (a) being non-empty gives $2c_k\int h_0\le o(1)$ with $\int h_0<0$, that is $\liminf c_k\ge0$, consistently. Part (c) then follows as in that proof: both ends of the interval in (a) tend to $1$.
\end{proof}

\begin{proof}[Proof of \Cref{rot:thm:CT}]
Fix $\fc=(J,\psi,n)$ in the class and $G$, and let $\cS^-=\{(tG,h_0):t\in\R\}$, $h_0=-n\psi$.
\emph{Step 1 (comparison of the effective actions).} \Cref{rot:thm:Asrc} on $T_1$ for $s\in\cS^-$ and on $T_2$ for $S\cdot s$ (\Cref{rot:lem:walledalg}(b) transports the own-frame sources; the $S$-image of a wall is a wall on $T_2$), with the common constant $c_k$ of \Cref{rot:lem:C}, which does not depend on the torus. Inputs: \Cref{rot:thm:H1S}, \Cref{rot:prop:BD}, \cite[Thm.~3.3]{K1}, \Cref{rot:conv:C}(ii), \Cref{rot:rem:CMrg}.
\emph{Step 2 (forward bounds).} \Cref{rot:prop:Kforward} on $T_1$ and on $T_2$ for these sources. Inputs: \Cref{rot:prop:H3,rot:prop:H4,rot:prop:H5} on both tori, hence \Cref{thm:a06:BALA}, \Cref{rot:prop:CT} and \Cref{rot:conv:C} throughout; on $T_1$ every $Y$-frame step is a skew step (\Cref{rot:prop:SK}); on $T_2$ only the first hybrid step is skew, with image terms $O(e^{-c\,50MJ/\eps})$. The instance of \Cref{rot:prop:H5} at the wall is \Cref{rot:thm:FS}.
\emph{Step 3 (sandwich).} \Cref{rot:prop:Konly} with \Cref{rot:prop:CK}: conclusion (c) of \Cref{rot:thm:sandwich} for the balanced pair $s_1=(tG,h_0)$, $s_0=(0,h_0)$, whose scalar sources have the same integral, gives
\[
u_k(t)\defeq\langle e^{tF(G)}\rangle_{\fc,\Yk}-\langle e^{tF(G)}\rangle_{\fc,\Xk}\to0\qquad\text{for every real }t .
\]
\emph{Step 4 (complex coefficient, moments).} By \Cref{lem:s02:GD} and \Cref{rot:lem:CM}, $1\le\langle e^{tF}\rangle\le e^{t^2q_\Lambda/2}$ with $q_\Lambda\to q$ on both lattices (\Cref{rot:lem:q}), and $|\langle e^{zF}\rangle|\le\langle e^{(\Re z)F}\rangle$; so $u_k$ is bounded on strips uniformly in $k$ and tends to $0$ on $\R$; Vitali's theorem \cite{Hor90} gives (i), and Cauchy's estimates the moments.
\emph{Step 5 (existence).} By \Cref{rot:thm:H6W} the axis family converges for real $t$; by the same argument it converges locally uniformly on $\C$ to an entire function $\Wv_\fc(zG)$, and by (i) so does the rotated family, to the same limit. This is (ii).
\emph{Step 6 (other generators).} \Cref{rot:lem:K}(c): the statement for $(R^{-1},T_1)$ and for the other planes follows from the one for $(R,T_1)$ by the passive reflections and permutations, which preserve $\Xk$ and $T_1$ and transport own-frame sources to own-frame sources, and the class of configurations is invariant under them. This is (iii).
No reverse step (``switching around'', \cite[p.~337]{K2}), no comparison of two spacings, no Wick scheme and no value of any $h$-linear constant, no convergence of the rotated family on its own and no infinite-volume limit is used: the two-lattice content is Steps~1--4, and existence is the one-lattice question of Step~5.
\end{proof}

\begin{proof}[Completion of the proof of \Cref{rot:hyp:K}]
By \Cref{rot:lem:K}(c) it suffices to treat $R=R_{12}(\thz)$, and by \Cref{rot:rem:class} the configurations to be treated are those of the class of \Cref{rot:thm:CT}. For such $\fc$ and $G$, Steps~1--2 of the proof of \Cref{rot:thm:CT} give the bound \hypF\ of \Cref{rot:hyp:F} on both tori for the sources $(tG,h_0)$, in total form: $v_k(T_1)=v_k(T_2)=0$ and $c_k(T_1)=c_k(T_2)=c_k$, with $|c_k|\le C\eps_k^{-q}$ (\Cref{rot:lem:C}), so that \hypFz\ holds for these sources, \hypV\ and \hypCc\ hold trivially, and moreover $c_k\to0$ (\Cref{rot:prop:CK}). \Cref{rot:prop:Konly} supplies conclusion (c) of \Cref{rot:thm:sandwich} for the family $\cS^-$ without the sources $(tG,+n\psi)$, and \Cref{rot:cor:walled} then gives $\langle e^{zF(G)}\rangle_{\fc,R^{-1}\Lk}-\langle e^{zF(G)}\rangle_{\fc,\Lk}\to0$ locally uniformly in $z$; with the existence of $\lim_k\langle e^{tF(G)}\rangle_{\fc,\Lk}$ for real $t$ (\Cref{rot:thm:H6W}) its last clause gives \Kz\ for $\fc$ and $G$.
\end{proof}

